% 纲要标题：# 第四卷　复几何与辛几何
% 纲要位置：几何学计划纲要.md，第 3471 行。
% 正文状态：第1—3章初稿；其余正文及书前书后保留框架，尚未正式审核。

% 纲要内容开始
% 
% > 保持复与辛的双入口；两支只在实际需要的相容线性代数、Kähler 和伪全纯理论处连接。曲线、正性、辛拓扑、形变和周期各有深入主线，不要求所有读者按整卷顺读。
% 
% 纲要内容结束

% !TeX TXS-program:compile = cmd /c build-Book4.bat -NoPause
% !TeX TXS-program:view = txs:///view-pdf-internal --embedded "Book4.pdf"

\documentclass[UTF8,zihao=-4,openany,fontset=fandol]{ctexbook}
\newcommand{\BookTitle}{复几何与辛几何}
\newcommand{\BookVolumeName}{第四卷}
\newcommand{\BookNumber}{4}
\newcommand{\BookDraftDate}{2026年10月}
\newcommand{\BookStatus}{写作稿·前三章初稿}
% 《几何学》五卷共用的基础排版与语义宏。
% 各卷主文件须先定义 \BookTitle、\BookVolumeName 和 \BookDraftDate。
\newcommand{\BookAuthor}{R. EverStarine}
% 稿件状态由各卷元数据指定；尚未开写的卷沿用框架状态。
\providecommand{\BookStatus}{框架预览·正文未撰写}

\usepackage[a4paper,inner=28mm,outer=24mm,top=27mm,bottom=28mm,
  headheight=16pt,headsep=8mm,footskip=12mm]{geometry}

\usepackage{amsmath}
\usepackage{mathtools}
\usepackage{amsthm}
\usepackage{aliascnt}
\usepackage{unicode-math}
% 数学字母使用默认数学字体内的对应字形；禁止 mathrm 等默认继承正文族。
% 运算符（sin、log、operatorname 等）随 mathrm 使用数学正体。
\unimathsetup{mathrm=sym,mathit=sym,mathbf=sym,mathsf=sym,mathtt=sym}
\setmainfont{texgyrepagella-regular.otf}[
  BoldFont=texgyrepagella-bold.otf,
  ItalicFont=texgyrepagella-italic.otf,
  BoldItalicFont=texgyrepagella-bolditalic.otf
]
\setsansfont{texgyreheros-regular.otf}[
  BoldFont=texgyreheros-bold.otf,
  ItalicFont=texgyreheros-italic.otf,
  BoldItalicFont=texgyreheros-bolditalic.otf
]
\setmonofont{lmmonolt10-regular.otf}[
  BoldFont=lmmonolt10-bold.otf,
  ItalicFont=lmmonolt10-oblique.otf,
  BoldItalicFont=lmmonolt10-boldoblique.otf
]
% 人物注释中的古希腊文、斯拉夫文字、日文异体字与希伯来文原名使用有相应字形的字体。
\newfontfamily\GeometryOriginalNameFont{Times New Roman}
\NewDocumentCommand{\OriginalName}{m}{{\GeometryOriginalNameFont #1}}
\newCJKfontfamily\GeometryOriginalHanFont{SimSun}
\NewDocumentCommand{\OriginalHanName}{m}{{\GeometryOriginalHanFont #1}}
\newfontfamily\GeometryHebrewFont{Times New Roman}[Script=Hebrew]
\TeXXeTstate=1
\NewDocumentCommand{\OriginalHebrewName}{m}{{\GeometryHebrewFont\beginR#1\endR}}
% 行内、行间及标题内公式均使用 unicode-math 默认 Latin Modern Math。
% 公式之外的正文西文仍由 setmainfont 使用 Pagella；黑板粗体单独采用 AMS。
\let\mathbb\relax
\DeclareMathAlphabet{\mathbb}{U}{msb}{m}{n}
% unicode-math 的 Latin Modern Math 将 \varnothing 映为瘦长椭圆；空集单独
% 采用传统 AMS MSBM 的圆形斜杠字形，不改变其余数学符号的字体。
\DeclareSymbolFont{GeometryAMSb}{U}{msb}{m}{n}
\DeclareMathSymbol{\GeometryRoundEmptySet}{\mathord}{GeometryAMSb}{"3F}
% 当前 Latin Modern Math 缺少横截符号 U+22D4；仅此关系符使用 AMSa 的现有字形。
\DeclareSymbolFont{GeometryAMSa}{U}{msa}{m}{n}
\DeclareMathSymbol{\GeometryTransverse}{\mathrel}{GeometryAMSa}{"74}
\AtBeginDocument{%
  \let\varnothing\GeometryRoundEmptySet
  \let\emptyset\GeometryRoundEmptySet
  \let\pitchfork\GeometryTransverse
}

\usepackage{microtype}
\usepackage{xcolor}
\usepackage{graphicx}
\usepackage{booktabs}
\usepackage{longtable}
\usepackage{array}
\usepackage{tabularx}
\usepackage{algorithm}
\usepackage{algpseudocode}
\usepackage{enumitem}
\usepackage{fancyhdr}
% 双面排版补齐奇数页时，空白页不残留上一章的页眉或页码。
\usepackage{emptypage}
\usepackage[most]{tcolorbox}
\usepackage{tikz-cd}
\usepackage{csquotes}
\usepackage[
  backend=biber,
  style=numeric-comp,
  sorting=nyt,
  giveninits=true,
  maxbibnames=99,
  doi=true,
  url=true,
  isbn=true
]{biblatex}
\addbibresource{../Shared/References.bib}
% 骨架没有引文时保留书目入口；写入真实条目后由完整构建自动排入书目。
\newif\ifGeometryHasBibliography
\AtDataInput{\global\GeometryHasBibliographytrue}
\NewDocumentCommand{\PrintBookBibliography}{}{%
  \ifGeometryHasBibliography
    \printbibliography[heading=bibintoc,title={参考文献}]%
  \else
    \chapter*{参考文献}%
    \addcontentsline{toc}{chapter}{参考文献}%
    \noindent 待撰写\par
  \fi
}
% 访问日期保留在书目库供内部核验，成品参考文献不显示 visited on。
\renewbibmacro*{url+urldate}{\usebibmacro{url}}
% biblatex 在书目内重绑定格式宏；在其作用域中将强调统一为正体，
% 同时避免书名中的中文及标点仍继承旧的斜体样式。
\AtBeginBibliography{\let\emph\textup\let\textit\textup}
% 中文条目的责任者连接符：不可套用西文的 “and”，按 GB/T 7714—2015 的著录
% 惯例改用逗号分隔。中文条目在 Shared/References.bib 中以 langid = {chinese}
% 标注（含译本，因其著录信息为中文形式），西文条目保持 “and” 不变。
% 按条目切换须在每条书目项处理时判定；\DeclareDelimFormat 的可选参数不是
% 语言选择器（实测无效），故用 \AtEveryBibitem 配合 \iffieldequalstr。
\AtEveryBibitem{%
  \iffieldequalstr{langid}{chinese}%
    {\DeclareDelimFormat{finalnamedelim}{\addcomma\space}%
     \DeclareDelimFormat{multinamedelim}{\addcomma\space}}%
    {\DeclareDelimFormat{finalnamedelim}{\addspace\bibstring{and}\space}%
     \DeclareDelimFormat{multinamedelim}{\addcomma\space}}%
}
\usepackage{etoolbox}
\usepackage{xparse}
\usepackage[noautomatic]{imakeidx}

\definecolor{GeometryBlue}{HTML}{244A64}
\definecolor{GeometryBlueLight}{HTML}{EEF4F7}
\definecolor{GeometryGray}{HTML}{4F5962}
\definecolor{GeometryGrayLight}{HTML}{F4F5F6}

\ctexset{
  part = {
    % 编题页分两行：第一行“第X编”，与封面“第一卷”同款（大号无衬线灰色）；
    % 第二行编名，与封面“几何学”同款（大号无衬线蓝色粗体），字号保持原样。
    % 第一行与行距由 number／aftername 给出，format 只负责编名的字体，
    % 以免字号被重置或被重复输出。
    format = {\centering\fontsize{30}{38}\selectfont\sffamily\bfseries\color{GeometryBlue}},
    name = {,},
    number = {\centering\normalfont\normalsize{\fontsize{24}{30}\selectfont\sffamily\color{GeometryGray}第\chinese{part}编}},
    aftername = {\par\vspace{0.7cm}},
    beforeskip = 0pt,
    afterskip = 2.5\baselineskip,
    pagestyle = plain
  },
  chapter = {
    format = \raggedright\Huge\bfseries\color{GeometryBlue},
    name = {第,章},
    number = \chinese{chapter},
    aftername = \quad,
    beforeskip = 0pt,
    afterskip = 30pt,
    pagestyle = plain
  },
  section = {
    format = \Large\bfseries\color{GeometryBlue},
    beforeskip = 2.4ex plus .6ex minus .2ex,
    afterskip = 1.1ex plus .2ex
  },
  subsection = {
    format = \large\bfseries\color{GeometryGray},
    beforeskip = 2ex plus .4ex minus .2ex,
    afterskip = .8ex plus .2ex
  }
}

\setlength{\parindent}{2em}
\setlength{\parskip}{0pt}
\linespread{1.18}
% 仅用于两个连续的普通正文段落；列表、公式、定理、标题等结构边界不用。
% 段间额外留白统一由这一长度控制，避免改动全局 \parskip 后扰乱各类环境。
\newlength{\GeometryProseParagraphGap}
\AtBeginDocument{%
  \setlength{\GeometryProseParagraphGap}{0.40\baselineskip}%
}
\NewDocumentCommand{\ProseParagraphBreak}{}{%
  \par\addvspace{\GeometryProseParagraphGap}%
}
% 与代数学保持一致：小节编号并进入目录，便于按具体主题定位。
\setcounter{secnumdepth}{3}
\setcounter{tocdepth}{2}
\allowdisplaybreaks[2]
\emergencystretch=2em
\clubpenalty=10000
\widowpenalty=10000
\displaywidowpenalty=10000

% 固定列表内部段距，避免 flushbottom 将默认的弹性 \parsep 拉大到超过正文段距。
\setlist[itemize]{leftmargin=2.2em,itemsep=.35ex,parsep=0pt,topsep=.6ex}
\setlist[enumerate]{leftmargin=2.5em,itemsep=.35ex,parsep=0pt,topsep=.6ex,
  font=\normalfont\upshape}
\setlist[enumerate,1]{label=\arabic*.,ref=\arabic*}
\setlist[enumerate,2]{label=\arabic*),ref=\arabic*}
\setlist[enumerate,3]{label=\alph*),ref=\alph*}
\setlist[description]{leftmargin=2.5em,labelindent=0pt,itemsep=.35ex,parsep=0pt,topsep=.6ex}

% “下列命题等价”专用列表；普通罗列仍使用 enumerate。
\newlist{equivalences}{enumerate}{2}
\setlist[equivalences,1]{leftmargin=2.8em,itemsep=.35ex,parsep=0pt,topsep=.6ex,
  label=\Roman*.,ref=\Roman*,font=\normalfont\upshape}
\setlist[equivalences,2]{leftmargin=2.8em,itemsep=.35ex,parsep=0pt,topsep=.6ex,
  label=\roman*),ref=\roman*,font=\normalfont\upshape}

\pagestyle{fancy}
\fancyhf{}
\fancyhead[LE]{\small\color{GeometryGray}\thepage\quad\leftmark}
% 奇数页采用页前最后有效的节标记，避免页中开始新节时提前更换页眉。
% 新节或习题的标记从后续页生效；章首页使用 plain，不显示上一章的标记。
\fancyhead[RO]{\small\color{GeometryGray}\TopMark{2e-right}\quad\thepage}
\renewcommand{\headrulewidth}{0.35pt}
% 保留数学标题的原始大小写，避免 book 类的 \MakeUppercase
% 与 unicode-math 命令在页眉标记中发生冲突。
% 偶数页页眉显示“章序＋章标题”（如“第一章 代数对象、映射与基本构造”）；
% 章序与章标题之间以 \quad 分隔，间距与奇数页节眉“节号＋节标题”之间一致；
% 附录章的章序由 \CTEXthechapter 按本卷附录字母（如 A）给出。
\renewcommand{\chaptermark}[1]{\markboth{\CTEXthechapter\quad #1}{}}
\renewcommand{\sectionmark}[1]{\markright{\thesection\quad #1}}
% 章末习题不编号，但进入习题页后应覆盖上一节留下的奇数页页眉；
% 同时写入目录项。PDF 书签由 hyperref 对 \section*{习题} 自动生成，
% 不再另行 \pdfbookmark，避免“习题”下再嵌套一个“习题”。
\newcommand{\BookExercises}{%
  \section*{习题}%
  \markright{习题}%
  \addcontentsline{toc}{section}{习题}%
}
\fancypagestyle{plain}{
  \fancyhf{}
  \fancyfoot[C]{\small\color{GeometryGray}\thepage}
  \renewcommand{\headrulewidth}{0pt}
}

\usepackage[unicode,colorlinks=true,linkcolor=GeometryBlue,
  citecolor=GeometryBlue,urlcolor=GeometryBlue,pdfborder={0 0 0}]{hyperref}
\usepackage{bookmark}
% PDF 书签沿用目录条目中的编号：正文章为“第几章”，附录章为“附录 A”，
% 节与小节分别保留 x.x、x.x.x；无编号的导读、习题和卷末单元仍不加号。
\bookmarksetup{numbered}

% unicode-math 会调整原始 Unicode 希腊字母的类别码。若把这些字母直接写进
% \texorpdfstring 的书签分支，辅助文件回读时类别码不同，LaTeX 会在每一轮都
% 误报标签发生变化。书签分支统一写成下列可展开宏，既保留希腊字母显示，
% 又使写入和回读的标签文本稳定。
\newcommand{\GeometryPdfAlpha}{α}
\newcommand{\GeometryPdfGamma}{Γ}
\newcommand{\GeometryPdfLambda}{λ}
\newcommand{\GeometryPdfMu}{μ}
\newcommand{\GeometryPdfPi}{π}
\newcommand{\GeometryPdfSigma}{σ}
\newcommand{\GeometryPdfTheta}{θ}
\newcommand{\GeometryPdfXi}{ξ}
\newcommand{\GeometryPdfZeta}{ζ}
\newcommand{\GeometryPdfCapitalLambda}{Λ}
\newcommand{\GeometryPdfOmega}{Ω}
\newcommand{\GeometryPdfCapitalSigma}{Σ}
\newcommand{\GeometryPdfDelta}{δ}
\newcommand{\GeometryPdfEpsilon}{ε}
\newcommand{\GeometryPdfIota}{ι}
\newcommand{\GeometryPdfSmallOmega}{ω}
\newcommand{\GeometryPdfEll}{ℓ}

\usepackage[nameinlink,noabbrev]{cleveref}
\newcommand{\crefpairconjunction}{和}
\newcommand{\crefmiddleconjunction}{、}
\newcommand{\creflastconjunction}{和}

% CTeX 2.5.10 的标题前缀宏带一个参数，而 TeX Live 2025 中 cleveref
% 仍按普通前缀宏展开它们，会把节号写成含 \endcsname 的损坏元数据。
% 改为恒等包装后，CTeX 传入的标题编号保持不变，\ref 与 \cref 均可正常工作。
\makeatletter
\def\p@part#1{#1}
\def\p@chapter#1{#1}
\def\p@section#1{#1}
\def\p@subsection#1{#1}
\def\p@subsubsection#1{#1}
\def\p@paragraph#1{#1}
\def\p@subparagraph#1{#1}
\makeatother
% 非数学西文字母统一正体；不改变括号类别、标点或数学字体。
% 2026-09-11：非数学模式的西文字母统一使用 Pagella 正体，保留粗细区别。
% 保留字母级边界，不重分类数字、括号、标点及中文。
% 此设置也覆盖列表序号字母和术语原文；数学模式保持原样。
\newfontfamily\GeometryLetterFont{texgyrepagella-regular.otf}[
  BoldFont=texgyrepagella-bold.otf,
  ItalicFont=texgyrepagella-italic.otf,
  BoldItalicFont=texgyrepagella-bolditalic.otf
]
\ExplSyntaxOn
\bool_new:N \l_geometry_text_letters_bool
\bool_set_true:N \l_geometry_text_letters_bool
\cs_new_protected:Npn \GeometryPreserveMathText
  { \bool_set_false:N \l_geometry_text_letters_bool }
\cs_new_protected:Npn \geometry_letter_group_begin:
  {
    \group_begin:
    \bool_if:NT \l_geometry_text_letters_bool
      { \GeometryLetterFont \upshape }
  }
% 继承 Default 的全部 xeCJK 转移，保留中西文间距及标点处理；
% 不能用只有开关分组的转移覆盖 Boundary 或 CJK 的原处理。
\xeCJK_new_class:n { GeometryLatinLetter }
\seq_map_inline:Nn \g__xeCJK_class_seq
  {
    \str_if_eq:nnF {#1} { GeometryLatinLetter }
      {
        \xeCJK_copy_inter_class_toks:nnnn
          {#1} {GeometryLatinLetter} {#1} {Default}
        \xeCJK_copy_inter_class_toks:nnnn
          {GeometryLatinLetter} {#1} {Default} {#1}
        \xeCJK_app_inter_class_toks:nnn {#1} {GeometryLatinLetter}
          { \geometry_letter_group_begin: }
        \xeCJK_pre_inter_class_toks:nnn {GeometryLatinLetter} {#1}
          { \group_end: }
      }
  }
% ASCII 字母及常用拉丁扩展字母；明确跳过乘号、除号，不重分类标点。
\xeCJKDeclareCharClass { GeometryLatinLetter }
  { "0041->"005A, "0061->"007A, "00C0->"00D6, "00D8->"00F6,
    "00F8->"024F, "1E00->"1EFF }
\ExplSyntaxOff
% 数学组内局部关闭字母切换，连公式中 \text{...} 的字体也保持原样。
% 只追加，不覆盖 amsmath、unicode-math、xeCJK 已有的数学初始化。
\AtBeginDocument{%
  \everymath\expandafter{\the\everymath\GeometryPreserveMathText}%
  \everydisplay\expandafter{\the\everydisplay\GeometryPreserveMathText}%
}


% 定理类环境（定理、定义、命题、引理、推论、断言、猜想、开放问题、注记、
% 习题）的正文：中文用楷体，西文保持正体，数学字体不受影响。
% 注意 \normalfont 只重置西文字族，中文字族仍随默认（宋体），故须显式 \kaishu；
% 不能退回 \itshape——那会把西文一并变斜。
\newtheoremstyle{geometryplain}
  {1.2ex}{1.2ex}{\normalfont\upshape\kaishu}{}
  {\color{GeometryBlue}\bfseries}{.}{0.6em}{}
\newtheoremstyle{geometrydefinition}
  {1.2ex}{1.2ex}{\normalfont\kaishu}{}
  {\color{GeometryBlue}\bfseries}{.}{0.6em}{}
\newtheoremstyle{geometryremark}
  {1.0ex}{1.0ex}{\normalfont\kaishu}{}
  {\color{GeometryGray}\bfseries}{.}{0.6em}{}

\theoremstyle{geometryplain}
\newtheorem{theorem}{定理}[section]
% 正文默认按“章.节.流水号”编号；未分节的附录省略中间的 0。
\renewcommand{\thetheorem}{%
  \ifnum\value{section}=0
    \thechapter.\arabic{theorem}%
  \else
    \thesection.\arabic{theorem}%
  \fi
}
\newaliascnt{lemma}{theorem}
\newtheorem{lemma}[lemma]{引理}
\aliascntresetthe{lemma}
\newaliascnt{proposition}{theorem}
\newtheorem{proposition}[proposition]{命题}
\aliascntresetthe{proposition}
\newaliascnt{corollary}{theorem}
\newtheorem{corollary}[corollary]{推论}
\aliascntresetthe{corollary}
\newaliascnt{claim}{theorem}
\newtheorem{claim}[claim]{断言}
\aliascntresetthe{claim}
\newaliascnt{conjecture}{theorem}
\newtheorem{conjecture}[conjecture]{猜想}
\aliascntresetthe{conjecture}
\newaliascnt{openproblem}{theorem}
\newtheorem{openproblem}[openproblem]{开放问题}
\aliascntresetthe{openproblem}

\theoremstyle{geometrydefinition}
\newaliascnt{definition}{theorem}
\newtheorem{definition}[definition]{定义}
\aliascntresetthe{definition}
\newaliascnt{example}{theorem}
\newtheorem{example}[example]{例}
\aliascntresetthe{example}
% 习题按章单独编号，不占用定理类环境的流水号。
\newtheorem{exercise}{习题}[chapter]

\theoremstyle{geometryremark}
\newaliascnt{remark}{theorem}
\newtheorem{remark}[remark]{注记}
\aliascntresetthe{remark}
\renewcommand{\proofname}{证明}
% 仅在证明紧接定理类陈述时收紧间距，并按陈述结尾是否为行间公式调整。
% 证明标题（证明／解／证明概要）的中文用楷体，
% 西文保持正体；正文不施加楷体，仍为默认宋体。不能退回 \itshape。
\newif\ifGeometryDirectProof
\newif\ifGeometryCoupledProof
\newif\ifGeometryStatementEndsInDisplay
\newskip\GeometryProofTrailingSkip
% 行间公式结束后，TeX 留下空的水平列表；在定理环境结束前辨认这一状态。
% 证明跟在公式之后时不再减去定理尾距，否则公式下方会比上方更紧。
\newcommand{\GeometryRecordStatementTail}{%
  \ifhmode
    \ifnum\lastnodetype=-1\relax
      \global\GeometryStatementEndsInDisplaytrue
    \else
      \global\GeometryStatementEndsInDisplayfalse
    \fi
  \else
    \global\GeometryStatementEndsInDisplayfalse
  \fi
}
\AddToHook{para/begin}{\global\GeometryDirectProoffalse\global\GeometryStatementEndsInDisplayfalse}
\AddToHook{env/theorem/end}{\GeometryRecordStatementTail}
\AddToHook{env/lemma/end}{\GeometryRecordStatementTail}
\AddToHook{env/proposition/end}{\GeometryRecordStatementTail}
\AddToHook{env/corollary/end}{\GeometryRecordStatementTail}
\AddToHook{env/claim/end}{\GeometryRecordStatementTail}
\AddToHook{env/exercise/end}{\GeometryRecordStatementTail}
\AddToHook{env/theorem/after}{\global\GeometryDirectProoftrue}
\AddToHook{env/lemma/after}{\global\GeometryDirectProoftrue}
\AddToHook{env/proposition/after}{\global\GeometryDirectProoftrue}
\AddToHook{env/corollary/after}{\global\GeometryDirectProoftrue}
\AddToHook{env/claim/after}{\global\GeometryDirectProoftrue}
\AddToHook{env/exercise/after}{\global\GeometryDirectProoftrue}
\AddToHook{env/figure/before}{\global\GeometryDirectProoffalse}
\AddToHook{env/table/before}{\global\GeometryDirectProoffalse}
\makeatletter
\renewenvironment{proof}[1][\proofname]{\par
  \pushQED{\qed}%
  \normalfont \topsep6\p@\@plus6\p@\relax
  \GeometryCoupledProoffalse
  \ifGeometryDirectProof
    \GeometryCoupledProoftrue
    \GeometryProofTrailingSkip=\topsep
    \advance\GeometryProofTrailingSkip by \partopsep
    \ifGeometryStatementEndsInDisplay\else
      \kern-5\p@% 只在陈述以普通文字结尾时抵消部分尾距。
    \fi
    \topsep2\p@\partopsep\z@\relax
  \fi
  \global\GeometryDirectProoffalse
  \global\GeometryStatementEndsInDisplayfalse
  \trivlist
  \item[\hskip\labelsep\upshape\kaishu #1\@addpunct{.}]\ignorespaces
}{%
  \popQED
  \ifGeometryCoupledProof\@topsepadd\GeometryProofTrailingSkip\fi
  \endtrivlist\@endpefalse
}
\makeatother
\newenvironment{solution}
  {\begin{proof}[解]}
  {\end{proof}}
\newenvironment{proofsketch}
  {\begin{proof}[证明概要]}
  {\end{proof}}

% 正文选读内容的星号属于编号而不属于标题；第二参数统一接收标签。
\newcommand{\GeometryOptionalMark}{\texorpdfstring{\textsuperscript{*}}{*}}
% 选做习题沿用同一编号星号。可选参数仍只表示习题自身的简短题名，
% 不得再写“选读”或“选做”；星号由本环境自动附在题号右侧。
\NewDocumentEnvironment{optionalexercise}{o}{%
  \begingroup
  \let\GeometrySavedTheExercise\theexercise
  \renewcommand{\theexercise}{\GeometrySavedTheExercise\GeometryOptionalMark}%
  \IfNoValueTF{#1}{\begin{exercise}}{\begin{exercise}[#1]}%
}{%
  \end{exercise}%
  \endgroup
}
% 可选参数只供目录与 PDF 书签使用；正文题名可借此按语义手动换行，
% 又不把换行命令写入目录或书签。
\NewDocumentCommand{\OptionalChapter}{O{}mm}{%
  \begingroup
  \let\GeometrySavedTheChapter\thechapter
  \renewcommand{\thechapter}{\GeometrySavedTheChapter\GeometryOptionalMark}%
  \ctexset{chapter={name={第,章\GeometryOptionalMark},
    number=\chinese{chapter},aftername=\quad}}%
  \ifstrempty{#1}{\chapter{#2}}{\chapter[#1]{#2}}\label{#3}%
  \endgroup
}
\NewDocumentCommand{\OptionalSection}{mm}{%
  \begingroup
  \let\GeometrySavedTheSection\thesection
  \renewcommand{\thesection}{\GeometrySavedTheSection\GeometryOptionalMark}%
  \ctexset{section={number=\thesection}}%
  \section{#1}\label{#2}%
  \endgroup
}
\NewDocumentCommand{\OptionalSubsection}{mm}{%
  \begingroup
  \let\GeometrySavedTheSubsection\thesubsection
  \renewcommand{\thesubsection}{\GeometrySavedTheSubsection\GeometryOptionalMark}%
  \ctexset{subsection={number=\thesubsection}}%
  \subsection{#1}\label{#2}%
  \endgroup
}

% 编题页必须从奇数页开始，题页仅在页脚显示页码，其后保留一张空白页。
\makeatletter
\NewDocumentCommand{\BookPart}{mm}{%
  \cleardoublepage
  \begingroup
  \def\GeometryPartLabel{#2}%
  % \part 会在返回前翻页，故须在 \@endpart 开始处记录题页标签。
  \pretocmd{\@endpart}{\label{\GeometryPartLabel}}{}%
    {\PackageError{geometry}{无法为编题页写入标签}{请检查当前文档类的 \string\@endpart 定义。}}%
  \part{#1}%
  \endgroup
  \thispagestyle{empty}%
  \null
  \clearpage
}
% 各编的“本编导读”是无编号章，在总目录和 PDF 书签中与所属编的各章并列。
% addcontentsline 同时建立目录项和书签，不再另加同名 pdfbookmark。
\NewDocumentCommand{\BookPartGuide}{O{本编导读}}{%
  \chapter*{#1}%
  \addcontentsline{toc}{chapter}{#1}%
}

% “附录”单独排无编号编题页，使用与正文编名相同的字号、字体和颜色。
% 题页从奇数页开始，其后保留一张空白页，再进入第一个附录章。
% startatroot 使根级书签与书后顶层节点分离；phantomsection 须先于
% addcontentsline，目录中的“附录”行才能获得独立锚点而非沿用前一节的
% section*.N。根级“附录”书签连同其子章由 addcontentsline 的 part 行自动
% 建立，不得再手写 \pdfbookmark（实测会产生重复的“附录”节点）。
\NewDocumentCommand{\BookAppendixPart}{m}{%
  \cleardoublepage
  \begingroup
  \def\GeometryPartLabel{#1}%
  % 在 part* 的题页尚未翻走时建立唯一的目录与书签目标。
  \pretocmd{\@endpart}{%
    \bookmarksetup{startatroot}%
    \phantomsection
    \addcontentsline{toc}{part}{附录}%
    \label{\GeometryPartLabel}%
  }{}{\PackageError{geometry}{无法为附录编题页写入目录与标签}{请检查当前文档类的 \string\@endpart 定义。}}%
  \part*{附录}%
  \endgroup
  \thispagestyle{empty}%
  \null
  \clearpage
}
\makeatother

% 列表项的括号与编号整体承担链接，避免只让中间的编号可点击。
\NewDocumentCommand{\itemref}{m}{\hyperref[#1]{(\ref*{#1})}}

% 章、节、小节的引用：整个“第 X 章／第 X 节／第 X 小节”都是跳转链接。
% 链接必须走 hyperref 的标签解析机制（目标名 chapter.N／section.N.M）；
% 早先用 \hyperlink{裸标签名} 指向不存在的目标，链接全部失效。
%
% 章号用汉字（“第五章”）且必须取“被引用那一章”的编号；\thechapter 是当前
% 计数器，在编导读、前言或跨章引用时会给出错误章号（表现为“第零章”）。
% hyperref 下 \ref 受保护、\getrefnumber 在 \edef 中又失真，故改由 zref 记录：
% 各章在 \chapter 之后调用 \BookChapterNumber，把汉字号存为 zref 属性；
% \zref@extract 可展开，能在 \edef 中安全取出。
\usepackage[user,hyperref]{zref}
\makeatletter
\zref@newprop{geometrychapternum}{}
% 在 \chapter 之后调用，参数即该章标签；把本章汉字号存为 zref 属性
\NewDocumentCommand{\BookChapterNumber}{m}{%
  \zref@setcurrent{geometrychapternum}{\chinese{chapter}}%
  \zref@labelbyprops{zch@#1}{geometrychapternum}%
}
% 取出某章标签对应的汉字号（\zref@extract 可展开，能在 \edef 中安全使用）
\NewDocumentCommand{\GeometryChineseChapterOf}{m}{%
  \zref@extract{zch@#1}{geometrychapternum}}
\makeatother
% \ref* 的分组会遮住编号末端的西文字符类别；编号前的间距可由 xeCJK
% 自动产生，编号后的同类间距须在完整链接内部显式补回。
\NewDocumentCommand{\secref}{m}{%
  \hyperref[#1]{第\ref*{#1}\CJKecglue 节}}
\NewDocumentCommand{\subsecref}{m}{%
  \hyperref[#1]{第\ref*{#1}\CJKecglue 小节}}
\NewDocumentCommand{\chapref}{m}{%
  \begingroup
  \xeCJKsetup{CJKecglue={}}%
  \edef\GeometryChapterDisp{第\GeometryChineseChapterOf{#1}章}%
  \hyperref[#1]{\GeometryChapterDisp}%
  \endgroup}

% 算法采用带编号的步骤，并明确标出输入与输出。
\floatname{algorithm}{算法}
\renewcommand{\listalgorithmname}{算法目录}
\algrenewcommand\algorithmicrequire{\textbf{输入：}}
\algrenewcommand\algorithmicensure{\textbf{输出：}}
\numberwithin{algorithm}{chapter}
% 完整中文编号进入题注、各自目录和引用；附录自动使用字母章号。
% 内部 \nobreakspace 与 cleveref 默认的定理引用一致；引用整串置于同一链接，
% 链接后的中文间距由下方 \cref 包装补齐。
\renewcommand{\thealgorithm}{算法\nobreakspace\thechapter-\arabic{algorithm}}
\renewcommand{\thefigure}{图\nobreakspace\thechapter-\arabic{figure}}
\renewcommand{\thetable}{表\nobreakspace\thechapter-\arabic{table}}
\renewcommand{\figurename}{}
\renewcommand{\tablename}{}
\makeatletter
\renewcommand{\fnum@figure}{\thefigure}
\renewcommand{\fnum@table}{\thetable}
\renewcommand{\fnum@algorithm}{\thealgorithm}
\makeatother
\NewDocumentEnvironment{algorithmsteps}{mm}
  {\begin{algorithmic}[1]\Require #1\Ensure #2}
  {\end{algorithmic}}

\newtcolorbox{learninggoals}{
  enhanced,breakable,colback=GeometryBlueLight,colframe=GeometryBlue,
  boxrule=.45pt,arc=1mm,left=1.2mm,right=1.2mm,top=1mm,bottom=1mm,
  title=本章目标,fonttitle=\bfseries
}
\newtcolorbox{partguide}{
  enhanced,breakable,colback=GeometryGrayLight,colframe=GeometryGray,
  boxrule=.35pt,arc=1mm,left=1.2mm,right=1.2mm,top=1mm,bottom=1mm,
  title=本编导读,fonttitle=\bfseries
}
\newtcolorbox{chaptersummary}{
  enhanced,breakable,colback=GeometryBlueLight,colframe=GeometryBlue,
  boxrule=.45pt,arc=1mm,left=1.2mm,right=1.2mm,top=1mm,bottom=1mm,
  title=本章小结,fonttitle=\bfseries,
  % 与阅读提示框一致：仅本框内容为多段散文，须恢复首行缩进。
  before upper={\parindent=2em\parskip=0pt}
}
% 仅用于尚待高阶数学复核的核心断言；成稿前必须处理或移除。
\newtcolorbox{coreverification}{
  enhanced,breakable,colback=yellow!10,colframe=orange!85!black,
  boxrule=.55pt,arc=1mm,left=1.2mm,right=1.2mm,top=1mm,bottom=1mm,
  title=待数学复核,fonttitle=\bfseries
}
% 兼容旧稿的占位命令：只输出普通文字，不另加方框。
\NewDocumentCommand{\GeometryPlaceholder}{m}{待撰写}
\NewDocumentEnvironment{chapterabstract}{}
  {\small\list{}{\rightmargin\leftmargin
    \listparindent=2\ccwd\itemindent=\listparindent}%
    \item\relax\color{GeometryGray}}
  {\endlist}
% 引诗与短引文统一居中，中文使用定理类环境的楷体；章摘要已在上方独立定义，
% 不继承这里的字形与对齐方式。
\RenewDocumentEnvironment{quote}{}
  {\list{}{\rightmargin\leftmargin}\item\relax\centering\kaishu}
  {\endlist}
% 引诗环境保留为语义清楚的兼容入口，版式与 quote 完全一致。
\NewDocumentEnvironment{poemquote}{}
  {\begin{quote}}
  {\end{quote}}
% 引诗末尾的作者或出处在引诗区域内右对齐。
\NewDocumentCommand{\quoteattribution}{m}{%
  \par\smallskip
  {\raggedleft #1\par}%
}

\makeindex[name=chinese,title=中文索引,columns=2,intoc,options=-q]
\makeindex[name=foreign,title=外文索引,columns=2,intoc,options=-q]
\makeindex[name=symbols,title=符号索引,columns=2,intoc,options=-q]
% 索引采用稍小字号与自然栏底，字母组间仅留少量空白；末页平衡不得越过正文高度。
\indexsetup{othercode={\small\raggedcolumns\maxbalancingoverflow=0pt
  \renewcommand{\indexspace}{\par\addvspace{0.2\baselineskip}}}}

% 有条目时读入 makeindex 生成的真实索引；尚无条目时仍保留书后结构与目录入口。
% 一键脚本会在 .idx 为空时删除旧 .ind，故这里不会误排上一次构建的陈旧索引。
\NewDocumentCommand{\PrintBookIndex}{m m}{%
  \IfFileExists{#1.ind}{%
    % 双栏输出可能丢失页眉标记；在索引范围内直接绑定标题并完成分页。
    \cleardoublepage
    \begingroup
      \fancyhead[LE]{\small\color{GeometryGray}\thepage\quad #2}%
      \fancyhead[RO]{\small\color{GeometryGray}#2\quad\thepage}%
      \printindex[#1]%
      \clearpage
    \endgroup
  }{%
    \chapter*{#2}%
    \addcontentsline{toc}{chapter}{#2}%
    \noindent 待撰写\par
  }%
}

% 所有术语在每卷首次正式定义或引入时都显示完整原文，并登记中外文索引。
% 星号形式与缺失原文均视为错误，防止旧的隐藏原文规则被重新带回源码。
% 定理类正文处于楷体；Fandol 楷体没有粗体字形，直接 \textbf 会退回普通楷体。
% 术语主称因此局部恢复正文的宋体字族，再使用与普通正文一致的粗体。
\NewDocumentCommand{\term}{s O{} m o}{%
  \IfBooleanT{#1}{%
    \PackageError{geometry}{\string\term* 已取消}{请改用 \string\term，并保留末尾的原文注解。}%
  }%
  {\songti\textbf{\textup{#3}}}%
  \IfNoValueTF{#4}{%
    \PackageError{geometry}{术语缺少原文注解}{每个术语首次正式定义或引入时都必须写成 \string\term[排序键]{中文主称}[原文]。}%
  }{%
    \textnormal{(\textup{#4})}%
  }%
  \ifstrempty{#2}{\index[chinese]{#3}}{\index[chinese]{#2@#3}}%
  \IfNoValueF{#4}{\index[foreign]{#4}}%
}
\NewDocumentCommand{\termidx}{O{}m}{%
  \ifstrempty{#1}{\index[chinese]{#2}}{\index[chinese]{#1@#2}}%
}
% 别称与主称一样用 \termidx 登记实际页码，不生成“见某某”的转引条目。
\NewDocumentCommand{\foreignidx}{O{}m}{%
  \ifstrempty{#1}{\index[foreign]{#2}}{\index[foreign]{#1@#2}}%
}
\newcounter{geometrySymbolIndexOrder}
\NewDocumentCommand{\symidx}{m}{%
  \stepcounter{geometrySymbolIndexOrder}%
  \edef\GeometrySymbolSortKey{%
    \number\numexpr10000+\value{geometrySymbolIndexOrder}\relax}%
  \index[symbols]{\GeometrySymbolSortKey@#1}%
}

% 一般强调与术语粗体分离。
\NewDocumentCommand{\keypoint}{m}{\CJKunderdot{#1}}
\NewDocumentCommand{\foreignkeypoint}{m}{\underline{\textup{#1}}}


\crefname{chapter}{章}{章}
\crefname{part}{编}{编}
\crefname{section}{节}{节}
\crefname{subsection}{小节}{小节}
\crefname{theorem}{定理}{定理}
\crefname{lemma}{引理}{引理}
\crefname{proposition}{命题}{命题}
\crefname{corollary}{推论}{推论}
\crefname{claim}{断言}{断言}
\crefname{conjecture}{猜想}{猜想}
\crefname{openproblem}{开放问题}{开放问题}
\crefname{definition}{定义}{定义}
\crefname{example}{例}{例}
\crefname{remark}{注记}{注记}
\crefname{exercise}{习题}{习题}
\crefname{equation}{式}{式}
\crefname{figure}{}{}
\crefname{table}{}{}
\crefname{algorithm}{}{}
\crefformat{figure}{#2#1#3}
\Crefformat{figure}{#2#1#3}
\crefformat{table}{#2#1#3}
\Crefformat{table}{#2#1#3}
\crefformat{algorithm}{#2#1#3}
\Crefformat{algorithm}{#2#1#3}

% cleveref 会在自己的分组中排出“定义 3.1.5”一类完整链接，类别名与编号之间
% 使用 \nobreakspace。分组结束后，xeCJK 无法再从末端识别出西文数字，
% 因而不会在紧随其后的汉字前插入 \CJKecglue。这里只在下一个源码记号确为
% CJK 正文字时补回同一种间距；空格、控制序列以及中英文标点均不处理，
% 从而不在“定义3.1.5，”或“定义3.1.5.”前制造额外空白。
\NewCommandCopy{\GeometryOriginalCref}{\cref}
\NewCommandCopy{\GeometryOriginalCrefCapital}{\Cref}
\ExplSyntaxOn
\cs_new_protected:Npn \geometry_cref_maybe_cjk_glue:
  { \peek_after:Nw \geometry_cref_maybe_cjk_glue_aux: }
\cs_new_protected:Npn \geometry_cref_maybe_cjk_glue_aux:
  {
    \token_if_space:NF \l_peek_token
      {
        \token_if_cs:NF \l_peek_token
          {
            \xeCJK_if_same_class:NNT \l_peek_token 中
              { \CJKecglue }
          }
      }
  }
\RenewDocumentCommand{\cref}{s m}
  {
    \IfBooleanTF{#1}
      { \GeometryOriginalCref*{#2} }
      { \GeometryOriginalCref{#2} }
    \geometry_cref_maybe_cjk_glue:
  }
\RenewDocumentCommand{\Cref}{s m}
  {
    \IfBooleanTF{#1}
      { \GeometryOriginalCrefCapital*{#2} }
      { \GeometryOriginalCrefCapital{#2} }
    \geometry_cref_maybe_cjk_glue:
  }
% 附录标签沿用字母编号；名称与编号整体进入同一链接，并与 \cref
% 采用相同的链接后中西文间距规则。
\NewDocumentCommand{\appref}{m}
  {
    \hyperref[#1]{附录\CJKecglue\ref*{#1}}
    \geometry_cref_maybe_cjk_glue:
  }
\ExplSyntaxOff

\newcommand{\MakeBookTitle}[1][封面]{%
  % 书名页在 PDF 侧栏中作为顶级入口；卷级文件可传入更具辨识性的书签名。
  \bookmarksetup{startatroot}%
  \pdfbookmark[0]{#1}{GeometryCover}%
  \begin{titlepage}
    \centering
    \vspace*{2.2cm}
    {\fontsize{36}{46}\selectfont\sffamily 几何学\par}
    \vspace{0.7cm}
    {\Large\sffamily\color{GeometryGray}\BookVolumeName\par}
    \vspace{1.0cm}
    {\fontsize{24}{32}\selectfont\bfseries\color{GeometryBlue}\BookTitle\par}
    \vfill
    {\large \BookAuthor\ 编写\par}
    \vspace{0.5cm}
    {\normalsize\BookDraftDate\par}
    \vspace{0.7cm}
    {\large\sffamily\bfseries\color{GeometryGray}\BookStatus\par}
    \vspace*{1.8cm}
  \end{titlepage}
}

% 目录也应有独立的顶级 PDF 书签。使用公共命令避免五卷分别遗漏。
\newcommand{\MakeBookContents}{%
  \cleardoublepage
  \bookmarksetup{startatroot}%
  \pdfbookmark[0]{目录}{GeometryContents}%
  \tableofcontents
}

% 图、表、算法目录不进入总目录，也不建立 PDF 书签；各自页面使用原有章级标题。
% 只另起一页而不补单双页。
\makeatletter
\newcommand{\GeometryListStart}{%
  \clearpage
}
% 各列表命令自带 \chapter*，局部消除其中的再次换页，保留原有标题与列表排版。
\newcommand{\GeometryListOnCurrentPage}[1]{%
  \begingroup
    \let\clearpage\relax
    \let\cleardoublepage\relax
    #1%
  \endgroup
}
\newcommand{\MakeBookListOfFigures}{%
  \GeometryListStart
  \GeometryListOnCurrentPage{\listoffigures}%
}
\newcommand{\MakeBookListOfTables}{%
  \begingroup
    \renewcommand{\listtablename}{表目录}%
    \GeometryListStart
    \GeometryListOnCurrentPage{\listoftables}%
  \endgroup
}
\newcommand{\MakeBookListOfAlgorithms}{%
  \begingroup
    \GeometryListStart
    % float.sty 在 listof 内重设编号栏；在标题回调中再加宽它。
    \let\GeometryOriginalFloatListHead\float@listhead
    \renewcommand{\float@listhead}[1]{%
      \renewcommand*\l@algorithm{\@dottedtocline{1}{1.5em}{5em}}%
      \GeometryOriginalFloatListHead{##1}%
    }%
    \GeometryListOnCurrentPage{\listofalgorithms}%
  \endgroup
}
\makeatother

% 无编号章（序言、目录、编导读、后记、参考文献、索引等）同样写入页眉标记，
% 使后续偶数页页眉显示该单元标题。置于 hyperref 之后，保证保留 hyperref 的章锚点。
\makeatletter
\let\Geometry@orig@schapter\@schapter
\def\@schapter#1{%
  \markboth{#1}{#1}%
  \Geometry@orig@schapter{#1}%
}
% 目录按层级着色：编为黑色、小节为深灰；章与节保持默认链接蓝。
\pretocmd{\l@part}{\hypersetup{linkcolor=black}}{}{%
  \PackageWarning{geometry}{无法为目录编条目着色}}
\pretocmd{\l@chapter}{\hypersetup{linkcolor=GeometryBlue}}{}{%
  \PackageWarning{geometry}{无法为目录章条目着色}}
\pretocmd{\l@section}{\hypersetup{linkcolor=GeometryBlue}}{}{%
  \PackageWarning{geometry}{无法为目录节条目着色}}
\pretocmd{\l@subsection}{\hypersetup{linkcolor=GeometryGray}}{}{%
  \PackageWarning{geometry}{无法为目录小节条目着色}}
% 编的目录条目与 PDF 书签须用纯文本“第X编”。ctex 的 \CTEX@part@tocline
% 取 \CTEXthepart，而该宏带上编题页的颜色与字体设置，会把格式代码写进目录，
% 并使书签出现 “=GeometryGray第一编” 之类乱码。此覆盖须置于 \ctexset 之后，
% 否则会被 ctex 重新定义；\CTEXifname 是 ctex 的内部开关，此处沿用。
\def\CTEX@part@tocline#1#2{%
  \CTEXifname
    {第\chinese{part}编\hspace{1em}}%
    {}%
  #2}
\makeatother

\hypersetup{pdftitle={几何学·\BookVolumeName：\BookTitle},pdfauthor={\BookAuthor},pdfsubject={第四卷写作稿：前三章初稿；其余部分保留框架，尚未正式审核}}

\begin{document}
\frontmatter
% !TeX root = ../Book4.tex
\MakeBookTitle[\BookTitle]
\clearpage
\thispagestyle{empty}
\null
\clearpage
% !TeX root = ../Book4.tex
\chapter*{序言}
\addcontentsline{toc}{chapter}{序言}
本卷为《几何学》\BookVolumeName《\BookTitle》的框架预览。目录依据《几何学计划纲要》建立，选读标记与原纲要一致。

待撰写

\MakeBookContents
% 有实际图、表或算法后再启用相应目录。
% \MakeBookListOfFigures
% \MakeBookListOfTables
% \MakeBookListOfAlgorithms


\mainmatter
% !TeX root = ../Book4.tex
% 纲要标题：# 第一编　黎曼曲面、整体解析与模参数
% 纲要位置：几何学计划纲要.md，第 3258 行。
% 正文状态：未撰写；本文件目前是框架。

% 纲要内容开始
% 
% 纲要内容结束

\BookPart{黎曼曲面、整体解析与模参数}{g4:part:01}
\BookPartGuide
待撰写

% !TeX root = ../../Book4.tex
% 纲要标题：## 第1章　复分析的几何方法
% 纲要位置：几何学计划纲要.md，第 3477 行。
% 正文状态：初稿；尚未正式审核。

% 纲要内容开始
% 
% > 本章目标与证明范围：面向研究生重建几何所需的一元复分析，重心是局部映射、解析延拓、紧性和极值构造，而不是把本科章节依次缩写。
% 
% 纲要内容结束

\chapter{复分析的几何方法}
\label{g4:ch:01}
\BookChapterNumber{g4:ch:01}

\begin{chapterabstract}
	一个复导数的存在会同时限制函数的局部形状、零点分布和整体延拓。本章从复可微与积分公式建立解析性，证明非恒定全纯函数的局部幂模型；随后以留数、绕数和辐角原理计数原像，用函数芽与单值定理描述沿路径的延拓。极值原理及圆盘上的收缩性质给出刚性和归一化唯一性，正规族则把有界函数列转化为可用的全纯极限。最后通过导数极值构造，证明每个单连通真平面域都与单位圆盘双全纯等价，并说明域内等价与边界对应各自包含什么信息。
\end{chapterabstract}

\begin{learninggoals}
	\begin{itemize}
		\item 从复差商得到 Cauchy--Riemann 方程，用三角形积分和 Cauchy 公式证明全纯函数的幂级数表示。
		\item 计算零点阶数，把非恒定全纯函数在一点附近化为幂映射，并判断逆分支的存在范围。
		\item 用主部、留数、辐角原理和 Rouché 定理计算零极点与原像重数，辨明边界和内部假设的作用。
		\item 用有限延拓链追踪函数芽，证明同伦不变性，并由单值定理构造全纯对数和根。
		\item 推导 Schwarz--Pick 的两点及微分不等式，解释等号如何确定圆盘自同构。
		\item 以局部有界性取得全纯极限，完成 Riemann 映照定理的极值证明，并判断哪些边界结论需要额外条件。
	\end{itemize}
\end{learninggoals}

方程 \(w^2=z\) 在每个非零点附近都能解出两个全纯分支，却不能在整个穿孔平面上选择一个连续的平方根。另一方面，一块边界形状很不规则的单连通平面区域，仍可能拥有一个把它整体变成圆盘的复坐标。前一个问题要求追踪局部逆沿环路怎样返回，后一个问题要求从一族候选坐标中找出真正覆盖圆盘的映射。一元复分析提供了处理这两种问题的共同工具：积分把局部微分条件传到边界，整数重数在扰动中保持，紧性使一列越来越好的候选映射产生极限。

\ProseParagraphBreak

这些问题也解释了为何复曲线的理论要先建立平面函数的几何。将来在曲面上更换全纯坐标时，导数非零保证坐标之间可逆；允许分歧时，局部坐标中的幂指数给出重数；在不同图中继续一个函数时，还要核对它绕行之后是否回到同一个函数芽。本章在平面区域中逐项完成这些论证，为下一章的复图册、覆盖与分歧准备可直接使用的结果。

\ProseParagraphBreak

本章需要复数的代数、实多元微积分、幂级数的一致收敛，以及开集、紧集和连通性的基本性质。全纯逆函数、同伦下的解析延拓和正规族紧性均在本章证明。带孔区域的积分使用平面 Green 公式，调用时恢复其假设、方向及第一卷中的证明位置；读者掌握该积分公式即可从这里开始复几何路线。第三卷的函数空间与椭圆理论在本章尚无必要。复积分与局部映射可对照\cite[第4章第2--3节]{Ahlfors1979Complex}，正规族可查\cite[第5章第5.3--5.4节]{Ahlfors1979Complex}；映照构造与解析延拓分别见\cite[第6章第1.1节; 第8章第1.4--1.6节]{Ahlfors1979Complex}。

% !TeX root = ../../Book4.tex
% 纲要标题：### 1.1 全纯性和局部映射
% 纲要位置：几何学计划纲要.md，第 3481 行。
% 正文状态：初稿；尚未正式审核。

% 纲要内容开始
% 
% * Cauchy–Riemann 方程与幂级数
% * Cauchy 积分及解析性的证明
% * 零点阶数与局部幂模型
% * 全纯逆函数及分支
% 
% 纲要内容结束

\section{全纯性和局部映射}
\label{g4:sec:0101}

映射 \(z\mapsto z^2\) 把辐角加倍，却在每个非零点附近保持相交曲线之间的有向角。这两个事实分别涉及有限范围内的绕转和一点处的线性近似。全纯函数的几何从这种区别开始：复导数约束无穷小形状，幂级数则把临界点附近的整个映射化为一个幂映射。我们先从复可微出发证明幂级数的存在，再建立这一局部模型。

\ProseParagraphBreak

本章用 \(D(a,r)=\{z\in\mathbb C:|z-a|<r\}\) 表示开圆盘，\(\mathbb D=D(0,1)\) 表示单位圆盘。\term[quyu]{区域}[domain]指非空连通开集；开集未必连通。圆周的正向为逆时针，复平面的定向由有序实坐标 \((x,y)\) 给定。

\subsection{Cauchy–Riemann 方程与幂级数}
\label{g4:subsec:0101:01}
把 \(f=u+\mathrm i v\) 看成实平面的映射时，它的一阶变化通常由一个任意实矩阵给出。复可微要求这个矩阵等于乘以某个复数；非零复数的乘法是一次旋转与正比例伸缩的复合，因而把两个方向之间的角度保留下来。

\begin{definition}
	\label{g4:def:0101:holomorphic}
	设 \(U\subset\mathbb C\) 开，\(a\in U\)。若极限
	\[
		f'(a)=\lim_{h\to0,\ h\ne0}\frac{f(a+h)-f(a)}{h}
	\]
	存在，则称 \(f:U\to\mathbb C\) 在 \(a\) 处复可微；这里 \(h\) 沿复平面中的任意方向趋于零。若 \(f\) 在 \(U\) 每点复可微，则称它在 \(U\) 上\term[quanchun]{全纯}[holomorphic]，全体这样的函数记为 \(\mathcal O(U)\)。在整个 \(\mathbb C\) 上全纯的函数称为\term[zhenghanshu]{整函数}[entire function]。
	\symidx{\(\mathcal O(U)\)：开集 \(U\) 上的全纯函数}
\end{definition}

复导数存在等价于 \(f(a+h)=f(a)+Ah+o(|h|)\)，其中 \(A=f'(a)\)。因此复可微已经蕴含实可微。反过来，若 \(f=u+\mathrm i v\) 在 \(a\) 实可微，则
\[
	\mathrm{D}f_a=
	\begin{pmatrix}u_x&u_y\\v_x&v_y\end{pmatrix}_a
\]
表示复数乘法，当且仅当 \(u_x=v_y\)、\(u_y=-v_x\)。这就是\term[cauchyriemannfangcheng]{Cauchy--Riemann 方程}[Cauchy--Riemann equations]。\footnote{Cauchy，Augustin-Louis Cauchy（1789--1857），法国数学家，研究实分析、复分析与置换群。Riemann，Georg Friedrich Bernhard Riemann（1826--1866），德国数学家，建立复函数的曲面观点，并深刻影响了微分几何与积分理论。}在此情形
\begin{equation}
	\label{g4:eq:0101:complex-linear}
	f'=u_x+\mathrm i v_x,\qquad
	\mathrm{D}f=
	\begin{pmatrix}\operatorname{Re}f'&-\operatorname{Im}f'\\
	\operatorname{Im}f'&\operatorname{Re}f'\end{pmatrix},\qquad
	\det\mathrm{D}f=|f'|^2.
\end{equation}
尤其对 \(C^1\) 函数，满足这些方程与全纯等价。后面还要证明，全纯函数自动具有任意阶连续导数，因而定义中无需另加 \(C^1\) 假设。

\ProseParagraphBreak

只检查一点处的偏导关系并不足够。例如令 \(f(0)=0\)，在 \((x,y)\ne(0,0)\) 时令
\[
	f(x+\mathrm i y)=\frac{x^3y}{x^4+y^2}.
\]
它在原点连续，因为 \(|x^3y|/(x^4+y^2)\le |x|/2\)；沿两条坐标轴为零，所以四个实偏导在原点满足上述方程。但沿 \(y=x^2\)，有 \(f(x+\mathrm i x^2)=x/2\)，故复差商趋于 \(1/2\)，而沿实轴差商为零。这里缺少的正是实可微性。

\begin{definition}
	\label{g4:def:0101:analytic}
	若函数在每点 \(a\) 的某圆盘内都可表示为收敛幂级数 \(\sum_{n\ge0}c_n(z-a)^n\)，则称它\term[jiexihanshu]{解析}[analytic]。
\end{definition}

收敛半径为 \(R>0\) 的幂级数在 \(|z-a|\le r<R\) 上一致绝对收敛，且可逐项求导。具体说，选 \(r<\rho<R\)，则 \(|c_n|\rho^n\) 有界，而 \(\sum n(r/\rho)^{n-1}\) 收敛，控制了导数级数的一致收敛；积分基本定理随后说明其和确为原函数的导数。逐次使用这一估计，得到所有阶导数。因此解析蕴含全纯，且 \(c_n=f^{(n)}(a)/n!\)。困难的反向结论是：仅有一次复导数，为何就能恢复整个幂级数？

\ProseParagraphBreak

全纯函数的和、积、复合仍全纯；分母非零处的商也全纯。这些性质由差商法则直接得到。特别地，\(\exp z=\sum z^n/n!\) 全纯，满足 \((\exp z)'=\exp z\) 和 \(\exp(z+w)=\exp z\exp w\)；后者可由绝对收敛级数的乘积逐项合并得到。取 \(w=-z\) 可知指数从不为零。它将成为构造局部对数和根的工具。

\subsection{Cauchy 积分及解析性的证明}
\label{g4:subsec:0101:02}
边界上的函数值是否足以决定内部？实可微函数可以在圆盘内部任意作紧支扰动；全纯函数的答案却完全不同。关键是先证明小三角形边界上的积分为零，再把这个消去关系用于带一个极点的积分核。

\ProseParagraphBreak

对分片 \(C^1\) 曲线 \(\gamma:[\alpha,\beta]\to U\) 及连续函数 \(f\)，定义
\[
	\int_\gamma f(z)\,\mathrm{d}z
	=\int_\alpha^\beta f(\gamma(t))\gamma'(t)\,\mathrm{d}t.
\]
积分按分段相加，反向改变符号，且绝对值不超过曲线长度乘以曲线上 \(|f|\) 的上界。若 \(F'=f\)，链式法则给出 \(\int_\gamma f\,\mathrm{d}z=F(\gamma(\beta))-F(\gamma(\alpha))\)，所以原函数的存在会使闭曲线积分消失。

\begin{theorem}[Cauchy--Goursat 定理]
	\label{g4:thm:0101:goursat}
	设 \(U\subset\mathbb C\) 开，\(f\in\mathcal O(U)\)。若闭三角形 \(T\) 包含于 \(U\)，则
	\[
		\int_{\partial T}f(z)\,\mathrm{d}z=0.
	\]
\end{theorem}
这称为\term[cauchygoursatdingli]{Cauchy--Goursat 定理}[Cauchy--Goursat theorem]。\footnote{Goursat，Édouard Jean-Baptiste Goursat（1858--1936），法国数学家，研究复分析与微分方程；其积分定理去掉了对全纯函数导数连续性的预设。}
\begin{proof}
	记原积分为 \(I\)。连接三条边的中点，得到四个相似小三角形；内部边两两抵消，故至少一个小三角形的积分绝对值不小于 \(|I|/4\)。重复选取，得到套叠闭三角形 \(T_n\)，直径和周长分别为 \(2^{-n}\operatorname{diam}T\)、\(2^{-n}\ell(\partial T)\)，且其积分绝对值不小于 \(4^{-n}|I|\)。紧性及直径趋零保证交集为一点 \(p\in U\)。

	在 \(p\) 复可微给出
	\[
		f(z)=f(p)+f'(p)(z-p)+(z-p)\varepsilon(z),\qquad \varepsilon(z)\longrightarrow0.
	\]
	常数项和一次项有显式原函数，闭边界积分为零。因此
	\[
		4^{-n}|I|\le 4^{-n}\ell(\partial T)\operatorname{diam}T
		\sup_{z\in T_n}|\varepsilon(z)|.
	\]
	约去 \(4^{-n}\) 并令 \(n\to\infty\)，得到 \(I=0\)。
\end{proof}

\begin{lemma}
	\label{g4:lem:0101:primitive}
	设 \(D\subset\mathbb C\) 是开圆盘，\(g:D\to\mathbb C\) 连续。若 \(g\) 在每个闭三角形 \(T\subset D\) 的边界上积分为零，则 \(g\) 在 \(D\) 上有全纯原函数。特别地，全纯函数在圆盘上有原函数。
\end{lemma}
\begin{proof}
	固定 \(b\in D\)，沿直线段定义 \(G(z)=\int_{[b,z]}g(\zeta)\,\mathrm{d}\zeta\)。圆盘凸，三角形的零积分给出
	\[
		G(z+h)-G(z)=\int_{[z,z+h]}g(\zeta)\,\mathrm{d}\zeta
		=h\int_0^1g(z+th)\,\mathrm{d}t.
	\]
	连续性使右端除以 \(h\) 后趋于 \(g(z)\)。于是 \(G'=g\)，原函数在整个圆盘上定义。
\end{proof}

下一步会用到一个轻微推广：若 \(g\) 在圆盘上连续，仅在一点 \(p\) 可能不全纯，上述三角形积分仍为零。若三角形遇到 \(p\)，先把它分成以 \(p\) 为顶点的至多三个三角形，再在该顶点切去直径 \(\varepsilon\) 的小三角形。剩余多边形可剖分成不经过 \(p\) 的三角形，对它们用\cref{g4:thm:0101:goursat}。原积分与剩余积分的差由总长度 \(O(\varepsilon)\) 的线段积分组成；\(g\) 在附近有界，故差趋零。点落在边上时用同样切割。这样，\cref{g4:lem:0101:primitive}也适用于这个连续延拓。

\begin{theorem}[Cauchy 积分公式]
	\label{g4:thm:0101:cauchy-formula}
	设 \(U\subset\mathbb C\) 开，\(f\in\mathcal O(U)\)，且 \(\overline{D(a,r)}\subset U\)。对每个 \(z\in D(a,r)\)，有
	\begin{equation}
		\label{g4:eq:0101:cauchy-formula}
		f(z)=\frac{1}{2\pi\mathrm i}\int_{|\zeta-a|=r}
		\frac{f(\zeta)}{\zeta-z}\,\mathrm{d}\zeta.
	\end{equation}
	进而 \(f\) 在该圆盘上有收敛幂级数
	\begin{equation}
		\label{g4:eq:0101:power-series}
		f(z)=\sum_{n=0}^\infty c_n(z-a)^n,\qquad
		c_n=\frac{1}{2\pi\mathrm i}\int_{|\zeta-a|=r}
		\frac{f(\zeta)}{(\zeta-a)^{n+1}}\,\mathrm{d}\zeta
		=\frac{f^{(n)}(a)}{n!}.
	\end{equation}
	级数在每个较小闭圆盘上一致绝对收敛。因此全纯与解析等价，全纯函数为 \(C^\infty\)。
\end{theorem}
公式\eqref{g4:eq:0101:cauchy-formula}称为\term[cauchyjifengongshi]{Cauchy 积分公式}[Cauchy integral formula]。
\begin{proof}
	因闭圆盘紧含于 \(U\)，可把半径略增，仍留在 \(U\)。对固定的内部点 \(z\)，在稍大的圆盘上令
	\[
		g(\zeta)=\frac{f(\zeta)-f(z)}{\zeta-z}\quad(\zeta\ne z),\qquad g(z)=f'(z).
	\]
	它在 \(z\) 连续，在其余点全纯。前面的推广及原函数引理说明其圆周积分为零。另一方面，因 \(|z-a|<r\)，在积分圆周上可以一致展开
	\[
		\frac1{\zeta-z}
		=\sum_{n=0}^\infty\frac{(z-a)^n}{(\zeta-a)^{n+1}}.
	\]
	参数化 \(\zeta=a+r e^{\mathrm i t}\) 可见，右端各项积分除 \(n=0\) 给出 \(2\pi\mathrm i\) 外，其余均为零。因此 \(\int(\zeta-z)^{-1}\,\mathrm{d}\zeta=2\pi\mathrm i\)。把 \(g\) 的零积分展开，便得积分公式。

	将同一几何级数代入积分公式。在 \(|z-a|\le s<r\) 上，\(|z-a|/r\le s/r<1\)，所以可一致逐项积分，得到所列 \(c_n\)。若圆周上的最大模为 \(M_r\)，则 \(|c_n|\le M_r/r^n\)，再次保证较小闭圆盘上的一致绝对收敛。幂级数逐项求导已在前面证明，故 \(c_n=f^{(n)}(a)/n!\)，并得到所有阶导数。
\end{proof}

\begin{proposition}[Cauchy 估计]
	\label{g4:prop:0101:cauchy-estimate}
	设 \(f\) 在闭圆盘 \(\overline{D(a,r)}\) 的邻域全纯，\(M_r=\max_{|z-a|=r}|f(z)|\)。则对所有整数 \(n\ge0\)，
	\[
		|f^{(n)}(a)|\le\frac{n!M_r}{r^n}.
	\]
	更一般地，设 \(U\subset\mathbb C\) 开，\(f\in\mathcal O(U)\)，\(K\subset U\) 紧。若 \(\delta>0\)，所有 \(D(z,\delta)\)（\(z\in K\)）的闭包都包含于 \(U\)，且 \(|f|\le M\) 在这些圆盘上成立，则 \(\sup_K|f^{(n)}|\le n!M/\delta^n\)。
\end{proposition}
\begin{proof}
	对\eqref{g4:eq:0101:power-series}中的系数积分取绝对值，使用圆周长度 \(2\pi r\) 即得第一式。对每个 \(z\in K\) 以它为圆心重复同一估计，便得第二式。
\end{proof}

这个估计把函数值的控制转成所有阶导数的控制。例如有界整函数 \(f\) 若满足 \(|f|\le M\)，则对任意 \(a\) 和任意 \(r>0\)，有 \(|f'(a)|\le M/r\)。令 \(r\to\infty\)，得到 \(f'=0\)，再沿线段积分可知 \(f\) 为常数。这里半径能够无限增大，正是定义域为整个平面的作用。

\ProseParagraphBreak

有了光滑性，就能把圆盘公式扩展到带孔区域。所需实变量公式为
\begin{equation}
	\label{g4:eq:0101:green}
	\int_{\partial G}(P\,\mathrm{d}x+Q\,\mathrm{d}y)
	=\int_G(Q_x-P_y)\,\mathrm{d}x\,\mathrm{d}y.
\end{equation}
这里 \(G\) 有界，边界由有限条互不相交的分片 \(C^2\) 正则简单闭曲线组成，角点内角在 \((0,2\pi)\) 内，\(P,Q\) 在 \(\overline G\) 的邻域为 \(C^1\)；边界按区域位于左侧定向。第一卷第9.1节引理9.1.9 已由积分基本定理、有限局部分解和角点截去极限证明这一 Green 公式。\footnote{Green，George Green（1793--1841），英国数学家，以位势理论、Green 函数和积分公式著称。}复系数情形分别用于实部和虚部即可。本章对一般边界区域均使用上述条件；圆盘和圆环是最常用的情形。

\begin{proposition}
	\label{g4:prop:0101:cauchy-domain}
	设 \(G\subset\mathbb C\) 是有界区域，边界满足公式\eqref{g4:eq:0101:green}所述条件，\(f\) 在 \(\overline G\) 的邻域全纯。则
	\[
		\int_{\partial G} f(\zeta)\,\mathrm{d}\zeta=0,
		\qquad
		f(z)=\frac1{2\pi\mathrm i}\int_{\partial G}
		\frac{f(\zeta)}{\zeta-z}\,\mathrm{d}\zeta\quad(z\in G).
	\]
	有孔时，外边界逆时针，内边界顺时针。
\end{proposition}
\begin{proof}
	写 \(f\,\mathrm{d}z=(u\,\mathrm{d}x-v\,\mathrm{d}y)+\mathrm i(v\,\mathrm{d}x+u\,\mathrm{d}y)\)。由\eqref{g4:eq:0101:green}，两部分的面积被积函数分别为 \(-v_x-u_y\) 与 \(u_x-v_y\)，恰由 Cauchy--Riemann 方程为零。使用该公式合法，因为\cref{g4:thm:0101:cauchy-formula}已经给出 \(f\) 的光滑性。

	对第二式，在 \(z\) 周围删去闭小圆盘 \(\overline{D(z,\varepsilon)}\subset G\)。核 \(f(\zeta)/(\zeta-z)\) 在剩余区域的闭包附近全纯，第一式说明其原边界积分等于小圆的逆时针积分。后者写成
	\[
		\mathrm i\int_0^{2\pi} f(z+\varepsilon e^{\mathrm i t})\,\mathrm{d}t
	\]
	并令 \(\varepsilon\to0\)，即得到 \(2\pi\mathrm i f(z)\)。
\end{proof}

\subsection{零点阶数与局部幂模型}
\label{g4:subsec:0101:03}
幂级数的第一个非零项不仅给出近似，还决定零点的重数。把其后的非零因子吸收到新坐标里，便能把近似变成精确等式。

\begin{proposition}
	\label{g4:prop:0101:zero-order}
	设 \(f\) 在点 \(a\) 的邻域全纯且 \(f(a)=0\)。恰有下列一种情形成立：\(f\) 在 \(a\) 的某邻域恒为零；或者存在唯一整数 \(m\ge1\) 及全纯函数 \(u\)，使
	\[
		f(z)=(z-a)^m u(z),\qquad u(a)\ne0.
	\]
	后一情形的 \(a\) 是孤立零点，整数 \(m\) 称为其\term[lingdiandejieshu]{零点阶数}[order of a zero]或重数，记为 \(\operatorname{ord}_a f\)。
\end{proposition}
\begin{proof}
	在 \(a\) 展开幂级数。若全部系数为零，则函数在收敛圆盘上为零；否则取第一个非零系数的指标 \(m\)，提出 \((z-a)^m\) 即得 \(u\)。其常数项非零，连续性保证缩小圆盘后 \(u\) 处处非零。幂级数系数的唯一性给出 \(m\) 和 \(u\) 的唯一性。
\end{proof}

\begin{theorem}[恒等定理]
	\label{g4:thm:0101:identity}
	设 \(\Omega\subset\mathbb C\) 是区域，\(f,g\in\mathcal O(\Omega)\)。若集合 \(\{z:f(z)=g(z)\}\) 在 \(\Omega\) 内有聚点，则 \(f=g\) 在整个 \(\Omega\) 上成立。特别地，非恒零全纯函数在区域中的零点离散。
\end{theorem}
\begin{proof}
	令 \(h=f-g\)。由\cref{g4:prop:0101:zero-order}，零点的聚点附近必有 \(h=0\)。令 \(E\) 为所有具有恒零邻域的点的集合，则 \(E\) 非空且开。若 \(p\in\overline E\cap\Omega\)，且 \(p\notin E\)，任意靠近 \(p\) 的 \(E\) 中点都是 \(h\) 的零点，因而 \(p\) 是零点聚点；前述局部结论又给出 \(p\in E\)，矛盾。所以 \(E\) 在 \(\Omega\) 中也闭。连通性给出 \(E=\Omega\)。
\end{proof}

连通条件用于把一处的相等传播到整个区域；内聚点条件则用于启动传播。函数 \(\sin(\pi/z)\) 在穿孔圆盘上有零点 \(1/n\) 向缺失的原点聚集，却并非恒零。由幂级数定义的正弦全纯，故这个例子确实只在聚点的位置上不同。

\ProseParagraphBreak

为了把 \((z-a)^m u(z)\) 改写为新坐标的 \(m\) 次幂，需要先说明一个导数非零的候选坐标为何可逆。以下证明也给出逆函数的全纯性。

\begin{theorem}[全纯逆函数定理]
	\label{g4:thm:0101:inverse}
	设 \(f\) 在 \(a\in\mathbb C\) 的邻域全纯，且 \(f'(a)\ne0\)。则存在 \(a\) 和 \(f(a)\) 的开邻域 \(V,W\)，使 \(f:V\to W\) 双射，其逆 \(g\) 全纯，且
	\[
		g'(w)=\frac1{f'(g(w))}.
	\]
\end{theorem}
\begin{proof}
	通过平移和除以 \(f'(a)\)，化为 \(a=f(a)=0\)、\(f'(0)=1\)，并写 \(f(z)=z+h(z)\)。由已证的导数连续性，取 \(r>0\)，使 \(h(0)=0\) 且 \(|h'|\le q<1\) 在 \(\overline{D(0,r)}\) 上成立。沿直线段积分得到
	\[
		|h(z)-h(\zeta)|\le q|z-\zeta|,\qquad
		|f(z)-f(\zeta)|\ge(1-q)|z-\zeta|.
	\]
	故 \(f\) 在此圆盘上单射。对 \(|w|<(1-q)r\)，迭代 \(z_{n+1}=w-h(z_n)\)，从 \(z_0=0\) 开始。映射把闭圆盘映入自身，且相邻差按因子 \(q\) 缩小。几何级数估计说明 \((z_n)\) 是 Cauchy 列，其极限 \(z\) 满足 \(f(z)=w\)；又 \(|z|\le |w|/(1-q)<r\)。同一压缩估计保证解唯一。

	令 \(W=D(0,(1-q)r)\)、\(V=f^{-1}(W)\cap D(0,r)\)，便得到所需双射。上述下界使逆 \(g\) 为 Lipschitz 连续。若 \(w'\to w\)，\(z'=g(w')\)、\(z=g(w)\)，则
	\[
		\frac{g(w')-g(w)}{w'-w}
		=\left(\frac{f(z')-f(z)}{z'-z}\right)^{-1}
		\longrightarrow\frac1{f'(z)}.
	\]
	这里 \(|f'(z)|\ge1-q\)，所以极限合法。还原归一化即得结论。
\end{proof}

\begin{theorem}[局部幂模型]
	\label{g4:thm:0101:local-power}
	设 \(\Omega\subset\mathbb C\) 是区域，\(f\in\mathcal O(\Omega)\) 非恒定，\(a\in\Omega\)。令 \(m=\operatorname{ord}_a(f-f(a))\)。则存在 \(a\) 附近的全纯坐标 \(w=\phi(z)\)，满足 \(\phi(a)=0\)、\(\phi'(a)\ne0\)，并且
	\begin{equation}
		\label{g4:eq:0101:local-power}
		f(z)=f(a)+\phi(z)^m.
	\end{equation}
	在充分小的坐标邻域 \(V=\phi^{-1}(D(0,\rho))\) 内，每个 \(0<|b-f(a)|<\rho^m\) 恰有 \(m\) 个不同原像；\(f(a)\) 的唯一原像为 \(a\)，重数为 \(m\)。
\end{theorem}
\begin{proof}
	恒等定理保证 \(f-f(a)\) 不会在 \(a\) 附近恒零。由零点分解，写成 \((z-a)^m u(z)\)。缩小圆盘使 \(|u(z)/u(a)-1|<1\)。级数
	\[
		L(\xi)=\sum_{n=1}^\infty\frac{(-1)^{n+1}}n\xi^n\qquad(|\xi|<1)
	\]
	满足 \(L'=1/(1+\xi)\)、\(L(0)=0\)。对 \(e^{L(\xi)}/(1+\xi)\) 求导得到零，由初值可知 \(e^{L(\xi)}=1+\xi\)。任取 \(c^m=u(a)\)，令
	\[
		v(z)=c\exp\left(\frac1m L\left(\frac{u(z)}{u(a)}-1\right)\right),\qquad
		\phi(z)=(z-a)v(z).
	\]
	便有 \(v^m=u\)、\(\phi'(a)=c\ne0\)。\cref{g4:thm:0101:inverse}使 \(\phi\) 在缩小邻域上成为全纯可逆坐标，且给出\eqref{g4:eq:0101:local-power}。再把坐标像缩为 \(D(0,\rho)\)。方程 \(w^m=b-f(a)\) 的 \(m\) 个根都在这个圆盘内且互不相同，映回便得原像计数。
\end{proof}

\begin{theorem}[开映射定理]
	\label{g4:thm:0101:open-mapping}
	区域 \(\Omega\subset\mathbb C\) 上的非恒定全纯函数把开集映为开集。它在点 \(a\) 附近单射，当且仅当 \(f'(a)\ne0\)。
\end{theorem}
\begin{proof}
	对任意开集中的点 \(a\)，在其中取\cref{g4:thm:0101:local-power}的坐标邻域。其像包含 \(f(a)\) 的圆盘，故像开。若 \(f'(a)=0\)，局部幂模型中的 \(m\ge2\)，足够靠近 \(f(a)\) 的非中心点在任意足够小邻域内有多个原像，不能局部单射。若 \(f'(a)\ne0\)，逆函数定理给出局部单射。
\end{proof}

对非恒定全纯函数，\(f'(a)=0\) 的点称为\term[linjiedian]{临界点}[critical point]，其像称为临界值；局部幂模型中的整数 \(m\) 记录这里的重数。譬如 \(f(z)=z^3-3z\) 在 \(a=1\) 处满足
\[
	f(z)+2=(z-1)^2(z+2),
\]
故局部重数为 \(2\)；在原点导数为 \(-3\)，局部重数为 \(1\)。临界点附近的二重性与函数作为三次多项式的整体次数是不同的计数。

\subsection{全纯逆函数及分支}
\label{g4:subsec:0101:04}
逆函数定理给出的逆只在一个邻域内有效。当目标点移动时，这些局部逆可能接续，也可能绕一圈后互换；这正是下一节解析延拓所要记录的现象。

\begin{definition}
	\label{g4:def:0101:biholomorphic}
	设 \(U,V\subset\mathbb C\) 为开集。若 \(f:U\to V\) 是全纯双射且逆映射也全纯，则称 \(f\) 为\term[shuangquanchunyingshe]{双全纯映射}[biholomorphic map]。区域上的单射全纯函数称为\term[danyehanshu]{单叶函数}[univalent function]。若 \(f\) 全纯且 \(f'\) 处处非零，则它局部保有向角，称为保向\term[gongxingyingshe]{共形映射}[conformal map]；本章讨论的共形映照在声称区域等价时还要求双射。
\end{definition}

由\cref{g4:thm:0101:open-mapping}，单叶函数的导数处处非零，像为开集；全局集合逆在每点附近都等于\cref{g4:thm:0101:inverse}给出的全纯局部逆，所以单叶函数总是到其像的双全纯映射。保角性则由\eqref{g4:eq:0101:complex-linear}直接得到：若两条正则曲线在 \(a\) 的切向量为 \(v_1,v_2\ne0\)，两条像曲线的切向量为 \(f'(a)v_1,f'(a)v_2\)，其比值仍为 \(v_2/v_1\)，故有向夹角不变。与此同时，长度元乘以 \(|f'(a)|\)，面积元乘以 \(|f'(a)|^2\)。

\begin{example}
	\label{g4:exm:0101:square-branches}
	设 \(S=\mathbb C\setminus(-\infty,0]\)。每个 \(z\in S\) 唯一写成 \(r e^{\mathrm i\theta}\)，其中 \(r>0\)、\(-\pi<\theta<\pi\)。函数
	\[
		s(z)=\sqrt r\,e^{\mathrm i\theta/2}
	\]
	连续，取值于右半平面，且 \(s(z)^2=z\)。在每个 \(z_0\in S\)，平方映射在 \(s(z_0)\ne0\) 处有全纯局部逆；连续性保证 \(s\) 在充分小邻域内选中的正是这个逆，因此 \(s\) 全纯。另一个根为 \(-s\)。它们是平方根的两个\term[fenzhi]{分支}[branch]，即在指定区域上单值全纯地选择根。
\end{example}

这两个分支不能作为平方根延伸到整个 \(\mathbb C\setminus\{0\}\)。若存在连续平方根 \(s\)，沿 \(z=e^{\mathrm i t}\)、\(0\le t\le2\pi\) 考察
\[
	q(t)=s(e^{\mathrm i t})e^{-\mathrm i t/2}.
\]
它连续且平方为 \(1\)，故在区间上恒为 \(1\) 或恒为 \(-1\)。于是 \(s(e^{2\pi\mathrm i})=-s(1)\)，与起终点相同矛盾。删去一条射线使辐角可以单值选取；保留整个穿孔平面，则沿环路的返回变化无法消去。

\ProseParagraphBreak

导数处处非零也不能保证整体单射。例如指数映射满足 \((e^z)'=e^z\ne0\)，但 \(e^{z+2\pi\mathrm i}=e^z\)。局部幂模型在每点都是一次幂，整体上仍有无穷多个逆分支。要把局部信息用于整体构造，接下来需要两种工具：边界积分对原像的计数，以及沿路径接续函数芽时的唯一性。
% 纲要条目排为编号小节标题，并进入目录与书签；说明段仍保留于纲要注释中。

% !TeX root = ../../Book4.tex
% 纲要标题：### 1.2 留数、次数与解析延拓
% 纲要位置：几何学计划纲要.md，第 3488 行。
% 正文状态：初稿；尚未正式审核。

% 纲要内容开始
% 
% * 亚纯函数和主部
% * 留数定理及辐角原理
% * Rouché 与映射次数
% * 函数芽的延拓和路径依赖
% 
% 纲要内容结束

\section{留数、次数与解析延拓}
\label{g4:sec:0102}

围绕一个零点走一周时，函数值可能绕原点走若干周；沿一条闭路延拓局部定义的对数时，回到起点的值又可能增加 \(2\pi\mathrm{i}\)。这两种现象都把局部解析数据变成了沿闭路测得的整体信息。前者由对数导数的积分计数，后者则要求比较不同路径上的函数芽。本节先把孤立奇点写成可积分的级数，再讨论零点计数和解析延拓。

\subsection{亚纯函数和主部}
\label{g4:subsec:0102:01}
在圆盘内，全纯函数由非负整数次幂展开；若中心被挖去，\(1/(z-a)\) 已经表明必须允许负次幂。负次幂记录了函数在缺失点附近无法由普通幂级数描述的部分。先在整个环域中构造这种展开，再令内半径趋于零，便能同时处理极点和更复杂的孤立奇点。

\begin{theorem}[Laurent 展开]\label{g4:thm:0102:laurent}
	设 \(0\leq r<R\leq\infty\)，函数 \(f\) 在环域
	\[
		A=\{z\in\mathbb C:r<|z-a|<R\}
	\]
	上全纯。则存在唯一的一列复数 \((c_n)_{n\in\mathbb Z}\)，使得
	\[
		f(z)=\sum_{n=-\infty}^{\infty}c_n(z-a)^n,
		\qquad
		c_n=\frac{1}{2\pi\mathrm{i}}
		\int_{|\zeta-a|=\rho}\frac{f(\zeta)}{(\zeta-a)^{n+1}}\,\mathrm{d}\zeta
	\]
	对任意 \(r<\rho<R\) 成立，积分圆周取逆时针方向。正次与负次两部分分别在 \(A\) 的每个紧子集上一致绝对收敛。
\end{theorem}
这个展开称为\term[laurentzhankai]{Laurent 展开}[Laurent expansion]。\footnote{Pierre Alphonse Laurent（1813—1854），法国数学家、工程师，研究复函数的级数展开与变分问题。}
\begin{proof}
	关键是把环域的外边界和内边界分开：外圆贡献非负次幂，内圆贡献负次幂。固定 \(z\in A\)，选取
	\[
		r<\rho_1<|z-a|<\rho_2<R.
	\]
	在两圆之间的闭环域上应用\cref{g4:prop:0101:cauchy-domain}，内边界的正向是顺时针方向，故
	\[
		f(z)=\frac{1}{2\pi\mathrm{i}}
		\int_{|\zeta-a|=\rho_2}\frac{f(\zeta)}{\zeta-z}\,\mathrm{d}\zeta
		-\frac{1}{2\pi\mathrm{i}}
		\int_{|\zeta-a|=\rho_1}\frac{f(\zeta)}{\zeta-z}\,\mathrm{d}\zeta.
	\]
	外圆上用几何级数
	\[
		\frac{1}{\zeta-z}
		=\sum_{n=0}^{\infty}\frac{(z-a)^n}{(\zeta-a)^{n+1}},
	\]
	内圆上则用
	\[
		\frac{1}{\zeta-z}
		=-\sum_{k=0}^{\infty}\frac{(\zeta-a)^k}{(z-a)^{k+1}}.
	\]
	在任意较小的闭环域 \(\rho_1<r'\leq|z-a|\leq R'<\rho_2\) 上，两式的公比分别不超过 \(R'/\rho_2\) 和 \(\rho_1/r'\)，均严格小于 \(1\)。因此可以逐项积分，所得两级数在该闭环域上一致绝对收敛。

	对每个整数 \(n\)，被积函数 \(f(\zeta)/(\zeta-a)^{n+1}\) 在 \(A\) 上全纯。在两条半径之间的闭环域上应用边界积分为零的结论，说明其圆周积分与半径无关，故上述系数都可用同一个 \(\rho\) 表示。任意紧子集均包含在一个这样的较小闭环域中，得到所需收敛性。若已有另一展开，在 \(|z-a|=\rho\) 上逐项乘以 \((z-a)^{-m-1}\) 积分，只有第 \(m\) 项的积分非零，因而必恢复同一系数 \(c_m\)，唯一性得证。
\end{proof}

取 \(r=0\) 时，点 \(a\) 称为 \(f\) 的\term[guliqidian]{孤立奇点}[isolated singularity]。展开中的
\[
	\sum_{k=1}^{\infty}c_{-k}(z-a)^{-k}
\]
称为该奇点的\term[zhubu]{主部}[principal part]。若主部为零，非负次幂部分给出穿过 \(a\) 的全纯延拓，称 \(a\) 为\term[kequqidian]{可去奇点}[removable singularity]。若主部恰有有限个非零项，且最负的非零项为 \(c_{-m}(z-a)^{-m}\)，则称 \(a\) 为 \(m\) 阶\term[jidian]{极点}[pole]；等价地，
\[
	f(z)=(z-a)^{-m}u(z),\qquad u\text{ 在 }a\text{ 附近全纯},\quad u(a)\neq0.
\]
主部有无穷多个非零项时，称 \(a\) 为\term[benzhiqidian]{本性奇点}[essential singularity]。这三种情况穷尽孤立奇点，因为它们正是负次幂系数的三种可能。

\begin{proposition}\label{g4:prop:0102:singularity-tests}
	设 \(f\) 在 \(0<|z-a|<R\) 上全纯。则 \(a\) 可去当且仅当 \(f\) 在某个去心邻域上有界；\(a\) 是极点当且仅当 \(|f(z)|\rightarrow\infty\) 当 \(z\rightarrow a\)。
\end{proposition}
\begin{proof}
	若 \(|f|\leq M\)，则对于每个 \(k\geq1\)，Laurent 系数公式给出
	\[
		|c_{-k}|\leq M\rho^k.
	\]
	令 \(\rho\rightarrow0\)，所有负次幂系数都为零。反向由全纯延拓的连续性得到。极点的局部表示直接给出模趋于无穷。反之，若 \(|f(z)|\rightarrow\infty\)，则 \(f\) 在充分小的去心圆盘上无零点，且 \(1/f\rightarrow0\)。可去性判据把 \(1/f\) 延拓为在 \(a\) 取零的全纯函数；它在去心圆盘上不为零，故由\cref{g4:thm:0101:identity}，该零点具有有限阶。取倒数便得到极点表示。
\end{proof}

\begin{definition}\label{g4:def:0102:meromorphic}
	设 \(\Omega\subset\mathbb C\) 为域。若 \(P\subset\Omega\) 在 \(\Omega\) 内无聚点，\(f\) 在 \(\Omega\setminus P\) 上全纯，且每个 \(a\in P\) 都是 \(f\) 的极点，则称 \(f\) 为 \(\Omega\) 上的\term[yachunhanshu]{亚纯函数}[meromorphic function]，在极点处约定 \(f(a)=\infty\)。允许 \(P=\varnothing\)，所以全纯函数也是亚纯函数。若 \(f\) 不恒为零，则在每个 \(a\in\Omega\) 附近可以唯一地写出阶数
	\[
		f(z)=(z-a)^m u(z),\qquad m\in\mathbb Z,\quad u(a)\neq0,
	\]
	其中 \(u\) 全纯。整数 \(m\) 记为 \(\operatorname{ord}_a f\)：零点处为正，极点处为负，其余点为零。\symidx{\(\operatorname{ord}_a f\)：亚纯函数的阶数}
\end{definition}

若 \(g,h\) 在域上全纯且 \(h\) 不恒为零，则 \(g/h\) 在 \(h\neq0\) 处定义，在共同零点处约去相同次数的因子后，所得函数是亚纯的。局部阶数满足
\[
	\operatorname{ord}_a(fg)=\operatorname{ord}_a f+\operatorname{ord}_a g,
	\qquad
	\operatorname{ord}_a(f/g)=\operatorname{ord}_a f-\operatorname{ord}_a g
\]
（各函数均不恒为零）。这些等式把乘除法转成整数的加减，是后续用除子记录零极点的基础。

\begin{example}\label{g4:exm:0102:laurent-domains}
	同一函数在不同环域上的 Laurent 展开可以不同。例如
	\[
		\frac{1}{z(z-1)}=-\frac1z+\frac1{z-1}
		=\begin{cases}
			-\dfrac1z-\displaystyle\sum_{n=0}^{\infty}z^n,&0<|z|<1,\\[6pt]
			\displaystyle\sum_{n=2}^{\infty}z^{-n},&|z|>1.
		\end{cases}
	\]
	第二式不能用于读取 \(0\) 处的主部，因为它的环域并不贴近 \(0\)。另一方面，\(e^{1/z}=\sum_{n=0}^{\infty}z^{-n}/n!\) 在 \(0\) 处具有本性奇点，所以它不是整个复平面上的亚纯函数；在去心平面上它仍是全纯函数。奇点的类型必须连同所讨论的定义域一起判断。
\end{example}

\subsection{留数定理及辐角原理}
\label{g4:subsec:0102:02}
对圆周积分，Laurent 级数中只有 \((z-a)^{-1}\) 一项留下贡献。因此，一个系数便能控制围绕奇点的积分。

\begin{definition}\label{g4:def:0102:residue}
	设 \(f\) 在 \(a\) 的去心邻域上全纯。其 Laurent 展开的系数 \(c_{-1}\) 称为 \(f\) 在 \(a\) 处的\term[liushu]{留数}[residue]，记为 \(\operatorname{Res}_{z=a}f\)。\symidx{\(\operatorname{Res}_{z=a}f\)：留数}对于充分小的正向圆周，
	\[
		\operatorname{Res}_{z=a}f
		=\frac{1}{2\pi\mathrm{i}}\int_{|z-a|=\rho}f(z)\,\mathrm{d}z.
	\]
\end{definition}

若 \(a\) 是至多 \(m\) 阶极点，则把 \((z-a)^m f(z)\) 展开成普通幂级数，可得实用公式
\[
	\operatorname{Res}_{z=a}f
	=\frac{1}{(m-1)!}
	\left.\frac{\mathrm{d}^{m-1}}{\mathrm{d}z^{m-1}}
	\bigl((z-a)^m f(z)\bigr)\right|_{z=a}.
\]
这里先把括号内函数全纯延拓到 \(a\)，再求导。特别地，若 \(g,h\) 在 \(a\) 附近全纯，\(h(a)=0\)、\(h'(a)\neq0\)，则
\[
	\operatorname{Res}_{z=a}\frac{g(z)}{h(z)}=\frac{g(a)}{h'(a)}.
\]

\begin{theorem}[留数定理]\label{g4:thm:0102:residue}
	设 \(G\subset\mathbb C\) 是有界域，边界由有限条互不相交的分片 \(C^2\) 正则简单闭曲线组成，角点内角在 \((0,2\pi)\) 内，并按使 \(G\) 位于左侧的方向定向。设 \(f\) 在 \(\overline G\) 的一个开邻域中，除有限个内点 \(a_1,\ldots,a_N\in G\) 外全纯，各 \(a_j\) 都是孤立奇点。则
	\[
		\int_{\partial G}f(z)\,\mathrm{d}z
		=2\pi\mathrm{i}\sum_{j=1}^N\operatorname{Res}_{z=a_j}f.
	\]
\end{theorem}
这就是\term[liushudingli]{留数定理}[residue theorem]的有限多边界版本；内边界取顺时针方向。这里每条边界弧的正则性指其 \(C^2\) 参数化的一阶导数不为零，角点内角指两条相邻边界弧在区域一侧所夹的角；这些条件与\eqref{g4:eq:0101:green}采用的边界条件一致。
\begin{proof}
	在每个 \(a_j\) 周围挖去一个闭圆盘，令这些圆盘两两不交且包含于 \(G\)。剩余闭区域的一个邻域上 \(f\) 全纯，由\cref{g4:prop:0101:cauchy-domain}的边界积分结论，其总边界积分为零。因此
	\[
		\int_{\partial G}f(z)\,\mathrm{d}z
		=\sum_{j=1}^N\int_{|z-a_j|=\rho_j}f(z)\,\mathrm{d}z,
	\]
	右侧各圆均取逆时针方向。每个圆周积分由 Laurent 系数公式恰等于对应留数的 \(2\pi\mathrm{i}\) 倍。若挖孔后出现多个连通分支，对各分支相加得到同一等式。
\end{proof}

留数也可以内在地赋给亚纯一形式 \(f(z)\,\mathrm{d}z\)。若换用局部全纯坐标 \(z=a+\phi(w)\)，其中 \(\phi(0)=0\)、\(\phi'(0)\neq0\)，则
\[
	\operatorname{Res}_{w=0}\bigl(f(a+\phi(w))\phi'(w)\bigr)
	=\operatorname{Res}_{z=a}f(z).
\]
为看清这一点，把亚纯函数的有限主部逐项处理。除 \((z-a)^{-1}\mathrm{d}z\) 外，各负次幂项都是某个亚纯函数的微分，换坐标后仍是微分，其 Laurent 展开的 \(w^{-1}\mathrm{d}w\) 系数为零；全纯部分也不产生负次幂。剩下的一项满足
\[
	\frac{\phi'(w)}{\phi(w)}\,\mathrm{d}w
	=\frac{\mathrm{d}w}{w}+\frac{v'(w)}{v(w)}\,\mathrm{d}w,
	\qquad \phi(w)=wv(w),\quad v(0)\neq0,
\]
故留数为 \(1\)。这一不变性解释了为什么曲面上的留数应当属于一形式：只改变函数表达式而遗漏微分因子，数值一般会改变。

\ProseParagraphBreak

接下来把积分用于数零点。对一条避开 \(0\) 的分片 \(C^1\) 闭曲线 \(\eta:[0,1]\rightarrow\mathbb C\setminus\{0\}\)，记
\[
	\operatorname{Ind}(\eta,0)
	\coloneqq\frac{1}{2\pi\mathrm{i}}\int_\eta\frac{\mathrm{d}w}{w}.
\]
这个数称为曲线关于 \(0\) 的\term[raoshu]{绕数}[winding number]。它是整数：把参数区间分成有限个小段，使每段的像落在某个可选连续辐角的小圆盘中，在相邻段上调整辐角的 \(2\pi\) 整数倍，使之连续拼接。于是有分片 \(C^1\) 的函数 \(r(t)>0\)、\(\theta(t)\in\mathbb R\)，满足 \(\eta(t)=r(t)e^{\mathrm{i}\theta(t)}\)，且
\[
	\int_\eta\frac{\mathrm{d}w}{w}
	=\log r(1)-\log r(0)+\mathrm{i}\bigl(\theta(1)-\theta(0)\bigr)
	=2\pi\mathrm{i}k,
\]
其中 \(k\in\mathbb Z\)，因为 \(\eta(1)=\eta(0)\)。所以绕数记录的是总辐角增量，而不是曲线图形的某种单调转动次数。

\begin{theorem}[辐角原理]\label{g4:thm:0102:argument}
	设 \(G\subset\mathbb C\) 是有界域，边界由有限条互不相交的分片 \(C^2\) 正则简单闭曲线组成，角点内角在 \((0,2\pi)\) 内，并取正向。设 \(f\) 在包含 \(\overline G\) 的某个域上亚纯且不恒为零，在 \(\partial G\) 上没有零点或极点。则
	\[
		\frac{1}{2\pi\mathrm{i}}\int_{\partial G}\frac{f'(z)}{f(z)}\,\mathrm{d}z
		=\sum_{a\in G}\operatorname{ord}_a f
		=N_G(f)-P_G(f),
	\]
	其中 \(N_G(f)\)、\(P_G(f)\) 分别按重数统计 \(G\) 内的零点数和极点数。左侧也等于各条边界像曲线 \(f\circ\gamma\) 关于 \(0\) 的绕数之和。
\end{theorem}
这里的\term[fujiaoyuanli]{辐角原理}[argument principle]把解析阶数变为整数值的边界数据。
\begin{proof}
	零点和极点在 \(\overline G\) 的邻域中离散，而边界上没有这两类点；紧性保证 \(G\) 内只有有限个。若
	\[
		f(z)=(z-a)^m u(z),\qquad u(a)\neq0,
	\]
	则
	\[
		\frac{f'(z)}{f(z)}=\frac{m}{z-a}+\frac{u'(z)}{u(z)},
	\]
	后项在 \(a\) 附近全纯。因此 \(f'/f\) 在 \(a\) 的留数为 \(m\)，即零点阶数取正、极点阶数取负。对 \(f'/f\) 应用\cref{g4:thm:0102:residue}便得等式。沿每条边界曲线作代换 \(w=f(z)\)，即得到像曲线的绕数表达式。
\end{proof}

\subsection{Rouché 与映射次数}
\label{g4:subsec:0102:03}
零点的位置可能随函数扰动而移动，而绕数只取整数值。只要扰动过程中边界像不经过原点，整数就无法连续改变。这给出了不必求出零点位置的计数方法。

\begin{theorem}[Rouché 定理]\label{g4:thm:0102:rouche}
	设 \(G\subset\mathbb C\) 是有界域，边界由有限条互不相交的分片 \(C^2\) 正则简单闭曲线组成，角点内角在 \((0,2\pi)\) 内。函数 \(f,g\) 在 \(\overline G\) 的一个开邻域上全纯，并且
	\[
		|g(z)-f(z)|<|f(z)|\qquad(z\in\partial G).
	\]
	则 \(f\) 与 \(g\) 在 \(G\) 内的零点总数相同，零点均按重数统计。
\end{theorem}
这一零点比较结论称为\term[rouchedingli]{Rouché 定理}[Rouché's theorem]。\footnote{Eugène Rouché（1832—1910），法国数学家，研究复函数的零点计数，并参与编写几何学教材。}
\begin{proof}
	以线段插值连接两个函数：\(f_t=f+t(g-f)\)，\(0\leq t\leq1\)。边界上有
	\[
		|f_t(z)|\geq|f(z)|-t|g(z)-f(z)|>0.
	\]
	边界紧，且 \((t,z)\mapsto f_t(z)\) 连续，故分母在 \([0,1]\times\partial G\) 上与零保持正距离。于是
	\[
		n(t)=\frac{1}{2\pi\mathrm{i}}\int_{\partial G}\frac{f_t'(z)}{f_t(z)}\,\mathrm{d}z
	\]
	是 \(t\) 的连续函数。由\cref{g4:thm:0102:argument}，它等于 \(f_t\) 的零点数，因而取整数值。连续的整数值函数在区间上为常数，所以 \(n(0)=n(1)\)。
\end{proof}

\begin{example}\label{g4:exm:0102:rouche-polynomial}
	令 \(p(z)=z^5+3z+1\)。在 \(|z|=1\) 上，
	\[
		|p(z)-3z|=|z^5+1|\leq2<3=|3z|,
	\]
	故单位圆盘内恰有一个零点。在 \(|z|=2\) 上，
	\[
		|p(z)-z^5|=|3z+1|\leq7<32=|z^5|,
	\]
	故半径 \(2\) 的圆盘内恰有五个零点。因此环域 \(1<|z|<2\) 内共有四个零点，均按重数计；两个边界圆上都没有零点。比较项由所选边界上的大小决定，同一个多项式在两条圆周上适合比较的主项并不相同。
\end{example}

设 \(G\) 为边界满足\cref{g4:thm:0102:rouche}条件的有界域，\(f\) 在 \(\overline G\) 的邻域上全纯，且在 \(G\) 上不恒定。对任意 \(w\notin f(\partial G)\)，定义
\begin{equation}\label{g4:eq:0102:degree}
	\deg(f,G,w)\coloneqq
	\frac{1}{2\pi\mathrm{i}}\int_{\partial G}\frac{f'(z)}{f(z)-w}\,\mathrm{d}z
	=\sum_{\substack{z\in G\\f(z)=w}}\operatorname{ord}_z(f-w).
\end{equation}
这是 \(f\) 在 \(G\) 中关于目标点 \(w\) 的\term[yingshecishu]{映射次数}[mapping degree]。等式由辐角原理应用于 \(f-w\) 得到，也说明该数是非负整数。由于积分依 \(w\) 连续，它在 \(\mathbb C\setminus f(\partial G)\) 的每个连通分支上恒定。因而，只要目标点不越过边界像，纤维的重数总和就保持不变。

\ProseParagraphBreak

若 \(w\) 的所有原像都是正则点，即 \(f'(z)\neq0\)，\cref{g4:thm:0101:inverse}给出每个原像附近的一张全纯逆分支，\eqref{g4:eq:0102:degree}便是这些原像的个数。在临界点附近，\cref{g4:thm:0101:local-power}把映射化成 \(\zeta\mapsto\zeta^m\)：非零的小目标值有 \(m\) 个互异原像，目标值到达临界值后它们合并成一个 \(m\) 重原像。次数保存的正是合并前后的这个总数。

\begin{proposition}\label{g4:prop:0102:local-count}
	设 \(f\) 在 \(a\in\mathbb C\) 的邻域上全纯，且 \(f-f(a)\) 在 \(a\) 有 \(m\geq1\) 阶零点。则存在 \(r>0\) 和 \(\varepsilon>0\)，使得对每个满足 \(|w-f(a)|<\varepsilon\) 的 \(w\)，方程 \(f(z)=w\) 在 \(|z-a|<r\) 内恰有 \(m\) 个解，按重数统计。
\end{proposition}
\begin{proof}
	取充分小的 \(r\)，使 \(f-f(a)\) 在 \(0<|z-a|\leq r\) 上无零点，并在闭圆盘的邻域上全纯。令
	\[
		\varepsilon=\min_{|z-a|=r}|f(z)-f(a)|>0.
	\]
	在该圆周上，\(|(f-w)-(f-f(a))|=|w-f(a)|<|f-f(a)|\)。由\cref{g4:thm:0102:rouche}，\(f-w\) 与 \(f-f(a)\) 的零点总数同为 \(m\)。
\end{proof}

这个证明还显示，小扰动保持的是总重数，未必保持互异零点数。例如 \(z^2-t\) 在 \(t\neq0\) 时有两个单零点，在 \(t=0\) 时只有一个二重零点。

\subsection{函数芽的延拓和路径依赖}
\label{g4:subsec:0102:04}
幂级数只在某个圆盘中给出函数，但圆盘的边界未必是函数真正的障碍。例如 \(\sum_{n\geq0}z^n\) 的收敛半径为 \(1\)，它的和 \(1/(1-z)\) 却可以穿过单位圆上除 \(1\) 外的所有点。延拓要保留的是邻域中函数的全部局部信息，而不是某个圆盘或某个展开式。

\begin{definition}\label{g4:def:0102:germ-continuation}
	设 \(a\in\Omega\subset\mathbb C\)，其中 \(\Omega\) 为域。两个在 \(a\) 的邻域上定义的全纯函数，若在某个更小邻域中相等，则视为等价；等价类称为 \(a\) 处的\term[hanshuya]{全纯函数芽}[germ of a holomorphic function]，记为 \([f]_a\)。

	设 \(\gamma:[0,1]\rightarrow\Omega\) 为从 \(a\) 出发的连续路径。一个沿 \(\gamma\) 的\term[jiexiyantuo]{解析延拓}[analytic continuation]由有限分割
	\[
		0=t_0<t_1<\cdots<t_N=1
	\]
	以及圆盘 \(U_j\subset\Omega\) 上的全纯函数 \(f_j\) 给出，满足 \(\gamma([t_{j-1},t_j])\subset U_j\)，第一项代表指定的初始函数芽，且
	\[
		[f_j]_{\gamma(t_j)}=[f_{j+1}]_{\gamma(t_j)}
		\qquad(1\leq j<N).
	\]
	末项给出的 \([f_N]_{\gamma(1)}\) 称为延拓的终点函数芽。
\end{definition}

这里的相等要求是在交点的一个邻域中相等；只要求数值相等无法传递导数和幂级数。两个圆盘的交集是连通的，所以由\cref{g4:thm:0101:identity}，相邻代表实际上在整个交集上一致。分割与圆盘只是记录延拓的工具，可以继续细分或缩小而不改变沿路径所得的函数芽。

\begin{lemma}\label{g4:lem:0102:continuation-unique}
	设 \(\Omega\subset\mathbb C\) 为域，\(\gamma:[0,1]\rightarrow\Omega\) 为连续路径。指定 \(\gamma(0)\) 处的全纯函数芽后，若沿 \(\gamma\) 的解析延拓存在，则每个 \(\gamma(t)\) 处随路径传递的函数芽都唯一，特别是终点函数芽唯一。
\end{lemma}
\begin{proof}
	比较两条延拓链，先对它们的分割取公共细分。在每个小参数区间内，两条链分别由两个圆盘上的函数代表，且路径段同时落在这两个圆盘中。开始时两个代表的芽相同；圆盘交集连通，恒等定理使两个代表在交集上相同，于是这一段的终点芽也相同。接着利用各自延拓链在接点处的相容性，进入下一段。有限次归纳便得到所有分割点的芽相同；把任意给定 \(t\) 加入分割，便知该处的芽也相同。
\end{proof}

固定路径上的唯一性仍允许不同路径给出不同结果。在 \(1\) 附近取满足 \(L(1)=0\)、\(e^{L(z)}=z\) 的局部对数，沿单位圆 \(\gamma(t)=e^{2\pi\mathrm{i}t}\) 延拓。连续辐角在这条路径上是 \(2\pi t\)，故沿途的值为 \(2\pi\mathrm{i}t\)，终点对数芽在 \(1\) 的值变成 \(2\pi\mathrm{i}\)。从同一起点走常值路径，终点值仍为 \(0\)。同样，从 \(\sqrt1=1\) 出发的平方根沿此圆周的值为 \(e^{\pi\mathrm{i}t}\)，一周后变成 \(-1\)。

\ProseParagraphBreak

上述值确实来自解析延拓链。对任意 \(b\neq0\)，在 \(|z-b|<|b|\) 内写成
\[
	L_b(z)=\ell_b+\sum_{n=1}^{\infty}\frac{(-1)^{n+1}}n
	\left(\frac{z-b}{b}\right)^n,\qquad e^{\ell_b}=b.
\]
逐项求导得 \(L_b'(z)=1/z\)，且 \(e^{L_b(z)}/z\) 的导数为零，在 \(b\) 取值 \(1\)，所以 \(e^{L_b(z)}=z\)。相邻圆盘上两个这样的对数相差 \(2\pi\mathrm{i}\) 的整数倍；调整常数后即可沿路径拼接。平方根由 \(e^{L_b/2}\) 得到。这也说明导致路径依赖的是绕过缺失点后的常数变化。

\ProseParagraphBreak

回顾这里的拓扑条件：两条同起终点路径\term[dingduandiantonglun]{定端点同伦}[homotopic relative to endpoints]，是指存在连续变形 \(H(t,s)\)，在变形过程中始点和终点始终固定。域内每条闭路都能这样收缩到常值路径时，称该域\term[danliantong]{单连通}[simply connected]。等价地，域内任意两条同起终点路径都定端点同伦；把一条路径接上另一条的反向路径，便得到这两个表述之间的对应。

\begin{theorem}[单值定理]\label{g4:thm:0102:monodromy}
	设 \(\Omega\subset\mathbb C\) 为域，\(a\in\Omega\)，且 \(\mathfrak f\) 是 \(a\) 处的全纯函数芽。假设 \(\mathfrak f\) 可以沿 \(\Omega\) 中每条从 \(a\) 出发的连续路径解析延拓。则沿两条定端点同伦路径得到的终点函数芽相同。特别地，若 \(\Omega\) 单连通，则存在唯一的全纯函数 \(F\) 定义在整个 \(\Omega\) 上，使 \([F]_a=\mathfrak f\)，且所有上述延拓均给出 \(F\) 的相应局部函数芽。
\end{theorem}
这里的\term[danzhidingli]{单值定理}[monodromy theorem]要求沿每条路径都能延拓。单连通是消除路径选择的拓扑条件，延拓存在性则是关于所给函数芽的解析条件。
\begin{proof}
	先证明一条有限延拓链在路径的小变动下仍然有效。固定从 \(a\) 到 \(b\) 的路径 \(\gamma\)，取延拓链 \((U_j,f_j)\) 和相应分割。每个紧集 \(\gamma([t_{j-1},t_j])\) 都包含于开圆盘 \(U_j\)，故存在 \(\delta_j>0\)，使这一紧集的 \(\delta_j\) 邻域仍包含于 \(U_j\)。令 \(\delta\) 为有限个 \(\delta_j\) 的最小值。若另一条路径 \(\widetilde\gamma\) 具有相同端点，并满足
	\[
		\sup_{0\leq t\leq1}|\widetilde\gamma(t)-\gamma(t)|<\delta,
	\]
	则它的每个相应路径段仍位于 \(U_j\)。相邻函数在圆盘交集上相等，所以同一列 \(f_j\) 也给出沿 \(\widetilde\gamma\) 的延拓；起终点固定，两条路径的初始芽与终点芽分别相同。由\cref{g4:lem:0102:continuation-unique}，这就是 \(\widetilde\gamma\) 的唯一终点芽。

	现在设 \(H:[0,1]^2\rightarrow\Omega\) 是定端点同伦，令 \(\gamma_s(t)=H(t,s)\)，其中 \(H(0,s)=a\)、\(H(1,s)=b\)。按假设，每条 \(\gamma_s\) 都有延拓。固定 \(s_0\)，对 \(\gamma_{s_0}\) 使用上一段的有限链；由 \(H\) 在紧方形上一致连续，\(s\) 充分靠近 \(s_0\) 时，\(\gamma_s\) 与 \(\gamma_{s_0}\) 在上述一致距离下小于 \(\delta\)。因此终点芽关于 \(s\) 局部不变。区间 \([0,1]\) 连通，一个局部不变的取值在全区间上必相同：给定一个值，其原像和补集都是开集，且原像非空。所以 \(s=0\) 与 \(s=1\) 的终点芽相同。

	若 \(\Omega\) 单连通，则任意两条具有相同端点的路径定端点同伦。对每个 \(b\in\Omega\)，用任意一条从 \(a\) 到 \(b\) 的路径延拓，并把终点芽的值记为 \(F(b)\)；刚才的结论保证选择无关。为证明 \(F\) 全纯，固定一条到 \(b\) 的路径，令它的末项代表为圆盘 \(U\) 上的函数 \(h\)。在 \(b\) 周围取更小圆盘 \(V\subset U\)。对于 \(c\in V\)，先沿原路径到 \(b\)，再沿线段到 \(c\)，末段使用同一个函数 \(h\)，便得 \(F(c)=h(c)\)。因此 \(F\) 在每点附近都是全纯的，并具有指定初始芽。唯一性由域上的恒等定理得到。
\end{proof}

例如，\(1/(1-z)\) 在 \(0\) 处的函数芽不能沿 \(\mathbb C\) 中终点为 \(1\) 的路径全纯延拓：若终点附近存在全纯代表 \(h\)，在与前一段重合的开集上便有 \((1-z)h(z)=1\)，恒等定理会使此式在 \(1\) 附近成立，在 \(z=1\) 处矛盾。因此，尽管 \(\mathbb C\) 单连通，单值定理的延拓存在性假设仍不满足。

\begin{proposition}\label{g4:prop:0102:logarithm}
	设 \(\Omega\subset\mathbb C\) 是单连通域，\(f\) 在 \(\Omega\) 上全纯且处处非零。则存在全纯函数 \(L\) 使 \(e^L=f\)。指定 \(a\in\Omega\) 及一个满足 \(e^\ell=f(a)\) 的数 \(\ell\) 后，可唯一选定满足 \(L(a)=\ell\) 的 \(L\)。此外，对每个整数 \(m\geq1\)，\(f\) 有全纯 \(m\) 次根；任意两个这样的根相差一个常数 \(m\) 次单位根。
\end{proposition}
\begin{proof}
	每个 \(b\in\Omega\) 附近，都可选圆盘 \(U_b\subset\Omega\)，使 \(|f(z)/f(b)-1|<1\)。选择 \(\ell_b\) 满足 \(e^{\ell_b}=f(b)\)，并令
	\[
		L_b(z)=\ell_b+\sum_{n=1}^{\infty}\frac{(-1)^{n+1}}n
		\left(\frac{f(z)}{f(b)}-1\right)^n.
	\]
	与前面的局部对数计算相同，\(L_b'=f'/f\)、\(e^{L_b}=f\)。若两个局部对数定义在相交圆盘上，其差的指数恒为 \(1\)，故差只取 \(2\pi\mathrm{i}\mathbb Z\) 中的值；连续性和交集连通性说明这个差是常数。

	从 \(a\) 出发任取一条连续路径。用上述圆盘覆盖它的紧像，取参数区间的有限细分，使每个小路径段落在某个圆盘中。先把第一项对数的值调整为指定的 \(\ell\)，再依次加上 \(2\pi\mathrm{i}\) 的整数倍，使相邻对数在交点的芽相同。这样便构成沿该路径的解析延拓。由\cref{g4:thm:0102:monodromy}，这些局部对数粘成整个 \(\Omega\) 上的全纯 \(L\)。恒等式 \(e^L=f\) 在局部成立，故处处成立。两个全纯对数之差在连通域上是常数 \(2\pi\mathrm{i}k\)，指定初值便确定唯一选择。

	令 \(g=e^{L/m}\)，则 \(g^m=f\)。若 \(\widetilde g^m=f\)，因为 \(f\neq0\)，比值 \(\widetilde g/g\) 全纯且只取有限个 \(m\) 次单位根中的值，连通性使之恒定。
\end{proof}

全纯对数的存在还可由闭路积分检测。在任意域 \(\Omega\) 上，无零点全纯函数 \(f\) 有全纯对数，当且仅当
\[
	\int_\gamma\frac{f'(z)}{f(z)}\,\mathrm{d}z=0
	\qquad\text{对每条分片 }C^1\text{ 闭路 }\gamma\subset\Omega.
\]
必要性由 \(f'/f=L'\) 得到。反之，固定 \(a\) 及 \(e^\ell=f(a)\)，定义
\[
	L(z)=\ell+\int_a^z\frac{f'(\zeta)}{f(\zeta)}\,\mathrm{d}\zeta.
\]
域中任意两点可用折线路径连接，闭路积分为零使定义与所选折线路径无关。在小圆盘内沿最后一条线段求导，得 \(L'=f'/f\)，从而 \((fe^{-L})'=0\)，初值使 \(fe^{-L}=1\)。在 \(\mathbb C\setminus\{0\}\) 上取 \(f(z)=z\)，单位圆的积分为 \(2\pi\mathrm{i}\)，所以不存在全域对数。留数、绕数与延拓的路径依赖在这个例子中给出了同一障碍的三种表达。
% 纲要条目排为编号小节标题，并进入目录与书签；说明段仍保留于纲要注释中。

% !TeX root = ../../Book4.tex
% 纲要标题：### 1.3 极值原理和刚性
% 纲要位置：几何学计划纲要.md，第 3495 行。
% 正文状态：初稿；尚未正式审核。

% 纲要内容开始
% 
% * 最大模及调和函数
% * Schwarz 引理与等号
% * Schwarz–Pick 的度量形式
% * 双全纯映射与自同构
% 
% 纲要内容结束

\section{极值原理和刚性}
\label{g4:sec:0103}

一个全纯函数在小圆盘内可以有任意高阶的零点，但它的模不能在内点形成孤立的峰。这个限制把局部解析性转化为整体控制：边界上的估计约束内部，极值处的等号又往往迫使整个函数具有特殊形式。圆盘自映射最清楚地体现了这一点。先固定一个点，把函数与恒等映射比较，再改变所固定的点，就能得到适合所有圆盘自映射的长度尺度。

\subsection{最大模及调和函数}
\label{g4:subsec:0103:01}
先从几何图像理解模的极值。若 \(f\) 非常值，则\cref{g4:thm:0101:open-mapping}说明它把开集映成开集。假如 \(f(a)\ne0\)，一个包含 \(f(a)\) 的小圆盘必有模大于 \(\lvert f(a)\rvert\) 的点；因此像集不可能在 \(f(a)\) 处贴住以原点为中心的圆而完全留在圆内。

\begin{theorem}[最大模原理]
	\label{g4:thm:0103:maximum-modulus}
	设 \(\Omega\subset\mathbb C\) 是区域，\(f\colon\Omega\rightarrow\mathbb C\) 全纯。若 \(\lvert f\rvert\) 在 \(\Omega\) 的某一点取得局部最大值，则 \(f\) 在 \(\Omega\) 上为常值。若进一步假设 \(\Omega\) 有界，且 \(f\) 连续延拓到 \(\overline\Omega\)，则
	\[
		\max_{\overline\Omega}\lvert f\rvert
		=\max_{\partial\Omega}\lvert f\rvert.
	\]
\end{theorem}

\begin{proof}
	设 \(a\in\Omega\)，存在邻域 \(U\subset\Omega\) 使 \(\lvert f(z)\rvert\leq\lvert f(a)\rvert\)。若 \(f\) 非常值，开映射定理使 \(f(U)\) 成为包含 \(f(a)\) 的开集，因而包含模大于 \(\lvert f(a)\rvert\) 的点，矛盾。这也覆盖 \(f(a)=0\) 的情形；此时局部上界本身就迫使 \(f|_U=0\)。

	在有界域的情形，\(\overline\Omega\) 是非空紧集，连续函数 \(\lvert f\rvert\) 在其上取得最大值。若最大值在内点取得，前一结论使 \(f\) 为常值，边界也取得同一值；否则最大值直接在边界取得。有界非空区域的边界非空，所以右端有意义。
\end{proof}

这里的\term[zuidamoyuanli]{最大模原理}[maximum modulus principle]控制的是复函数的模。实部和虚部同样受到极值限制，其原因可以用一个实偏微分方程表达。

\begin{definition}
	\label{g4:def:0103:harmonic}
	设 \(U\subset\mathbb R^2\) 开，\(u\colon U\rightarrow\mathbb R\) 为 \(C^2\) 函数。若
	\[
		\Delta u=-(u_{xx}+u_{yy})=0,
	\]
	则称 \(u\) 为\term[tiaohehanshu]{调和函数}[harmonic function]。这里 \(\Delta\) 沿用全书的非负 Laplace 算子约定。\footnote{Pierre-Simon Laplace（1749--1827），法国数学家，研究天体力学、概率论与位势理论。}若实函数 \(v\) 使 \(u+\mathrm{i}v\) 在 \(U\subset\mathbb C\) 上全纯，则称 \(v\) 为 \(u\) 的\term[tiaohegonge]{调和共轭}[harmonic conjugate]。
\end{definition}

全纯函数 \(f=u+\mathrm{i}v\) 的实部与虚部都是调和函数。事实上，全纯函数由幂级数表示，因而可以连续求任意阶偏导；对 Cauchy–Riemann 方程 \(u_x=v_y\)、\(u_y=-v_x\) 再求导，得到 \(u_{xx}+u_{yy}=v_{yx}-v_{xy}=0\)，对 \(v\) 也有同样结论。反方向的困难在于能否由 \(u\) 恢复 \(v\)：所需导数必须满足
\[
	v_x=-u_y,\qquad v_y=u_x.
\]
调和方程正是这两个导数的相容条件。

\begin{proposition}
	\label{g4:prop:0103:harmonic-conjugate}
	设 \(u\) 是区域 \(\Omega\subset\mathbb C\) 上的实值调和函数。则每一点附近都有调和共轭。在整个 \(\Omega\) 上存在调和共轭，当且仅当对每条分片 \(C^1\) 闭曲线 \(\gamma\subset\Omega\) 都有
	\begin{equation}
		\label{g4:eq:0103:conjugate-period}
		\int_\gamma(-u_y\,\mathrm{d}x+u_x\,\mathrm{d}y)=0.
	\end{equation}
	若全局调和共轭存在，则相差一个实常数。
\end{proposition}

\begin{proof}
	先在包含给定点 \((x_0,y_0)\) 且闭包位于 \(\Omega\) 内的小开矩形上构造。沿水平线段再沿竖直线段积分，令
	\[
		v(x,y)=-\int_{x_0}^{x}u_y(s,y_0)\,\mathrm{d}s
		+\int_{y_0}^{y}u_x(x,t)\,\mathrm{d}t.
	\]
	由微积分基本定理，\(v_y=u_x\)。另一方向的导数满足
	\[
		v_x=-u_y(x,y_0)+\int_{y_0}^{y}u_{xx}(x,t)\,\mathrm{d}t
		=-u_y(x,y_0)-\int_{y_0}^{y}u_{yy}(x,t)\,\mathrm{d}t
		=-u_y(x,y).
	\]
	因此 \(u+\mathrm{i}v\) 满足 Cauchy–Riemann 方程且为 \(C^1\)，从而全纯。这也说明 \(u\) 实际上局部是解析函数的实部，具有任意阶连续偏导。

	若全局的 \(v\) 存在，被积一形式就是 \(\mathrm{d}v\)，其沿闭曲线的积分为零。反过来，假设式\eqref{g4:eq:0103:conjugate-period}成立。固定 \(a\in\Omega\)，对任意 \(z\in\Omega\) 取从 \(a\) 到 \(z\) 的分片 \(C^1\) 道路，定义 \(v(z)\) 为该道路上的积分。开连通的平面区域内任意两点都能用折线连接：从 \(a\) 能以折线到达的点构成非空开闭子集，故就是全域。两条道路的积分之差是拼接闭曲线上的积分，因此 \(v\) 与道路选择无关。把终点在一个小矩形内移动，局部积分公式给出 \(v_x=-u_y\)、\(v_y=u_x\)，所以所得 \(v\) 确为全局调和共轭。两个共轭之差的两个偏导数都为零，沿连接任意两点的折线积分便知其为常数。
\end{proof}

\begin{example}
	\label{g4:exm:0103:log-modulus}
	函数 \(u(z)=\log\lvert z\rvert\) 在 \(\mathbb C^*\) 上调和。写成 \(u(x,y)=\tfrac12\log(x^2+y^2)\)，直接求导得
	\[
		-u_y\,\mathrm{d}x+u_x\,\mathrm{d}y
		=\frac{-y\,\mathrm{d}x+x\,\mathrm{d}y}{x^2+y^2}.
	\]
	在单位圆 \(z=\mathrm{e}^{\mathrm{i}t}\)、\(0\leq t\leq2\pi\) 上，其积分为 \(2\pi\)。因此 \(u\) 没有全局调和共轭。局部共轭是辐角；沿单位圆走一周，辐角增加 \(2\pi\)，正是式\eqref{g4:eq:0103:conjugate-period}检测到的障碍。
\end{example}

调和函数的值还可以从周围的圆周恢复。对 \(u=\operatorname{Re}f\)，这一现象来自 Cauchy 积分公式；以下证明直接从调和方程出发，因而不需要先在整个圆盘上选择共轭。

\begin{proposition}[调和函数的平均值性质]
	\label{g4:prop:0103:harmonic-mean}
	设 \(u\) 是开集 \(U\subset\mathbb C\) 上的实值调和函数，\(\overline{D(a,R)}\subset U\)。则对 \(0<r\leq R\)，有
	\[
		u(a)=\frac1{2\pi}\int_0^{2\pi}u(a+r\mathrm{e}^{\mathrm{i}\theta})\,\mathrm{d}\theta
		=\frac1{\pi r^2}\int_{D(a,r)}u(x,y)\,\mathrm{d}x\,\mathrm{d}y.
	\]
\end{proposition}

\begin{proof}
	令 \(v(r,\theta)=u(a+r\mathrm{e}^{\mathrm{i}\theta})\)，圆周平均为 \(m(r)=(2\pi)^{-1}\int_0^{2\pi}v(r,\theta)\,\mathrm{d}\theta\)。极坐标中的链式法则给出
	\[
		m''(r)+\frac1r m'(r)
		=\frac1{2\pi}\int_0^{2\pi}
		\left((u_{xx}+u_{yy})(a+r\mathrm{e}^{\mathrm{i}\theta})
		-\frac1{r^2}v_{\theta\theta}(r,\theta)\right)
		\,\mathrm{d}\theta=0.
	\]
	这里角向二阶导数的一周积分为零，调和性使 \(u_{xx}+u_{yy}=0\)。因此 \((rm'(r))'=0\)。\(u\) 在中心附近的梯度有界，所以 \(m'(r)\) 有界，\(rm'(r)\rightarrow0\)。于是 \(m'(r)=0\)，再由 \(m(r)\rightarrow u(a)\) 得第一式。用极坐标对圆周平均再乘 \(2r\,\mathrm{d}r\) 积分，得到圆盘平均公式。
\end{proof}

\begin{corollary}[调和函数的极值原理]
	\label{g4:cor:0103:harmonic-maximum}
	设 \(u\) 是区域 \(\Omega\subset\mathbb C\) 上的实值调和函数。若 \(u\) 在内点取得局部最大值或局部最小值，则 \(u\) 为常值。若 \(\Omega\) 有界且 \(u\in C(\overline\Omega)\)，则 \(u\) 的最大值和最小值均在 \(\partial\Omega\) 上取得。特别地，两个在 \(\Omega\) 上调和、在 \(\overline\Omega\) 上连续且边界值相同的实函数必相等。
\end{corollary}

\begin{proof}
	设 \(a\) 为局部最大值点，取足够小的圆盘使其中 \(u\leq u(a)\)。每个以 \(a\) 为中心的小圆周上的平均值都等于 \(u(a)\)，连续非负函数 \(u(a)-u\) 因积分为零而在该圆周恒为零。让半径变化，便知 \(u\) 在一个开圆盘内为常数。

	要把这一结论传到整个连通域，考虑 \(h=u_x-\mathrm{i}u_y\)。调和方程与混合偏导交换分别给出 \(h\) 的两条 Cauchy–Riemann 方程，所以 \(h\) 在 \(\Omega\) 上全纯。它在刚才的圆盘内为零，由\cref{g4:thm:0101:identity}知 \(h\equiv0\)。于是 \(u\) 的梯度在整个区域为零，故 \(u\) 为常数。对 \(-u\) 应用同一论证得到最小值结论。有界域上的边界结论由闭包紧性与内点极值结论得到；最后对两个函数之差应用最大值和最小值结论，便知其恒为零。
\end{proof}

边界条件因此决定了调和延拓的唯一性。存在性需要另行构造：极值原理比较的是已经存在的两个函数。类似地，若全纯函数 \(f\) 在区域内无零点，对 \(1/f\) 应用最大模原理，便得到 \(\lvert f\rvert\) 的最小值原则。无零点这一条件不可省略，例如 \(f(z)=z\) 的模就在原点取得最小值。

\subsection{Schwarz 引理与等号}
\label{g4:subsec:0103:02}
设 \(f\colon\mathbb D\rightarrow\mathbb D\) 全纯并满足 \(f(0)=0\)。由于原点的零点至少有一阶，可以把 \(f(z)\) 除以 \(z\)；对商函数应用最大模原理，就能把值域的界 \(\lvert f\rvert<1\) 改进为随 \(z\) 缩小的界。这是\term[Schwarzyinli]{Schwarz 引理}[Schwarz lemma]的基本想法。\footnote{Hermann Amandus Schwarz（1843--1921），德国数学家，研究共形映射、极小曲面与变分问题。}

\begin{theorem}[Schwarz 引理]
	\label{g4:thm:0103:schwarz}
	若 \(f\colon\mathbb D\rightarrow\mathbb D\) 全纯且 \(f(0)=0\)，则
	\[
		\lvert f(z)\rvert\leq\lvert z\rvert\quad(z\in\mathbb D),
		\qquad \lvert f'(0)\rvert\leq1.
	\]
	若存在非零 \(z_0\in\mathbb D\) 使 \(\lvert f(z_0)\rvert=\lvert z_0\rvert\)，或者 \(\lvert f'(0)\rvert=1\)，则存在 \(\theta\in\mathbb R\) 使 \(f(z)=\mathrm{e}^{\mathrm{i}\theta}z\)。这些旋转也都满足等号。
\end{theorem}

\begin{proof}
	由 \(f\) 在原点的幂级数，函数
	\[
		g(z)=
		\begin{cases}
			f(z)/z,&z\ne0,\\
			f'(0),&z=0
		\end{cases}
	\]
	在整个 \(\mathbb D\) 上全纯。固定 \(0<r<1\)，在 \(\lvert z\rvert=r\) 上有 \(\lvert g(z)\rvert\leq1/r\)。\cref{g4:thm:0103:maximum-modulus}将此界传到闭圆盘 \(\lvert z\rvert\leq r\)。对每个固定 \(z\in\mathbb D\)，令 \(r\uparrow1\)，得到 \(\lvert g(z)\rvert\leq1\)，这同时给出两条不等式。

	任一等号条件都表示 \(\lvert g\rvert\) 在一个内点取得最大值 \(1\)。最大模原理于是使 \(g\) 为模等于 \(1\) 的常数，故 \(g=\mathrm{e}^{\mathrm{i}\theta}\)。反方向直接代入即可。
\end{proof}

等号所给的信息比不等式本身更强：一个非零点的模或原点的导数达到界，就能决定整个函数。若 \(f'(0)=0\)，还可以继续提取零点因子。具体地，若 \(f\) 在原点的零点阶数至少为正整数 \(m\)，对全纯延拓的 \(f(z)/z^m\) 在半径 \(r\) 的圆上估计，再令 \(r\uparrow1\)，得到 \(\lvert f(z)\rvert\leq\lvert z\rvert^m\)。若在非零点取等号，同一极值论证给出 \(f(z)=\mathrm{e}^{\mathrm{i}\theta}z^m\)。例如 \(z^2\) 把靠近原点的点压得更靠近原点，同时保留了一个二阶分歧点。

\ProseParagraphBreak

固定原点的条件使比较特别简洁。若 \(f(0)\ne0\)，\(\lvert f(z)\rvert\leq\lvert z\rvert\) 在 \(z=0\) 便不成立。要对任意全纯圆盘自映射获得同样有效的结论，需要先把任意选定的点及其像同时移到原点。

\subsection{Schwarz–Pick 的度量形式}
\label{g4:subsec:0103:03}
对 \(a\in\mathbb D\)，定义分式线性映射
\begin{equation}
	\label{g4:eq:0103:disk-normalization}
	\phi_a(z)\coloneqq\frac{z-a}{1-\overline a z}.
\end{equation}
分母在 \(\mathbb D\) 内非零，且直接展开分子分母的模平方可得
\begin{equation}
	\label{g4:eq:0103:disk-normalization-identities}
	1-\lvert\phi_a(z)\rvert^2
	=\frac{(1-\lvert a\rvert^2)(1-\lvert z\rvert^2)}{\lvert1-\overline a z\rvert^2},
	\qquad
	\phi_a'(z)=\frac{1-\lvert a\rvert^2}{(1-\overline a z)^2}.
\end{equation}
第一式说明 \(\phi_a\) 把圆盘映入自身；解出 \(z\) 又得到
\[
	\phi_a^{-1}(w)=\frac{w+a}{1+\overline a w}=\phi_{-a}(w).
\]
因此 \(\phi_a\) 是具有全纯逆的圆盘双射，把 \(a\) 移到 \(0\)。

\begin{theorem}[Schwarz--Pick 定理]
	\label{g4:thm:0103:schwarz-pick}
	若 \(f\colon\mathbb D\rightarrow\mathbb D\) 全纯，则对任意 \(z,w\in\mathbb D\)，有
	\begin{equation}
		\label{g4:eq:0103:schwarz-pick-two-points}
		\left\lvert\frac{f(z)-f(w)}{1-\overline{f(w)}f(z)}\right\rvert
		\leq
		\left\lvert\frac{z-w}{1-\overline w z}\right\rvert,
	\end{equation}
	并且对每个 \(z\in\mathbb D\)，有
	\begin{equation}
		\label{g4:eq:0103:schwarz-pick-derivative}
		\frac{\lvert f'(z)\rvert}{1-\lvert f(z)\rvert^2}
		\leq\frac1{1-\lvert z\rvert^2}.
	\end{equation}
	若第一式在某对不同点取等号，或第二式在某一点取等号，则 \(f\) 是圆盘的全纯双射，且逆映射全纯。反过来，每个这样的映射在上述两式中处处取等号。
\end{theorem}

\begin{proof}
	固定 \(a\in\mathbb D\)，把定义域中的 \(a\) 和值域中的 \(f(a)\) 分别移到原点，令
	\[
		g=\phi_{f(a)}\circ f\circ\phi_a^{-1}.
	\]
	这些映射组成交换图
	\[
		\begin{tikzcd}[column sep=large,row sep=large]
			\mathbb D \arrow[r,"f"] \arrow[d,"\phi_a"']
			& \mathbb D \arrow[d,"\phi_{f(a)}"] \\
			\mathbb D \arrow[r,"g"'] & \mathbb D.
		\end{tikzcd}
	\]
	图中 \(a\) 与 \(f(a)\) 均被竖直映射送到 \(0\)，所以 \(g(0)=0\)。由\cref{g4:thm:0103:schwarz}，
	\[
		\lvert\phi_{f(a)}(f(z))\rvert\leq\lvert\phi_a(z)\rvert,
		\qquad \lvert g'(0)\rvert\leq1.
	\]
	第一式正是所需的两点不等式。又由链式法则和式\eqref{g4:eq:0103:disk-normalization-identities}，
	\[
		g'(0)=\frac{1-\lvert a\rvert^2}{1-\lvert f(a)\rvert^2}f'(a),
	\]
	于是得到导数不等式。

	若在不同点 \(z\ne a\) 的两点不等式中取等号，Schwarz 引理的等号条件使 \(g(\zeta)=\mathrm{e}^{\mathrm{i}\theta}\zeta\)；导数等号也给出相同结论。因此
	\[
		f=\phi_{f(a)}^{-1}\circ R_\theta\circ\phi_a,
		\qquad R_\theta(\zeta)=\mathrm{e}^{\mathrm{i}\theta}\zeta.
	\]
	每个因子均有全纯逆，所以 \(f\) 也是这样的双射。反过来，若 \(f\) 及其逆都全纯，分别对二者应用两点不等式，得到互反的两个不等式，故必须取等号。对导数不等式也同样处理，并使用 \((f^{-1})'(f(z))f'(z)=1\)，即得处处等号。
\end{proof}

这个结果称为\term[SchwarzPickdingli]{Schwarz--Pick 定理}[Schwarz--Pick theorem]。\footnote{Georg Alexander Pick（1859--1942），奥地利数学家，研究复变函数、微分方程与几何；格点多边形的面积公式也以其姓氏命名。}例如，当 \(a\ne0\) 时，\(\lvert\phi_a'(a)\rvert=(1-\lvert a\rvert^2)^{-1}>1\)，上述圆盘双射在 \(a\) 附近会放大平面通常的长度。导数不等式中的分母给出了适合圆盘的修正。

\begin{definition}
	\label{g4:def:0103:hyperbolic-distance}
	在单位圆盘上取长度元素
	\[
		\mathrm{d}s_{\mathbb D}=\lambda_{\mathbb D}(z)\lvert\mathrm{d}z\rvert,
		\qquad \lambda_{\mathbb D}(z)\coloneqq\frac2{1-\lvert z\rvert^2}.
	\]
	这称为圆盘的\term[shuangquduliang]{双曲度量}[hyperbolic metric]。分片 \(C^1\) 曲线 \(\gamma\colon[0,1]\rightarrow\mathbb D\) 的长度定义为
	\[
		L_{\mathbb D}(\gamma)=\int_0^1
		\frac{2\lvert\gamma'(t)\rvert}{1-\lvert\gamma(t)\rvert^2}\,\mathrm{d}t.
	\]
	两点间的双曲距离 \(d_{\mathbb D}(z,w)\) 是连接它们的全部这种曲线长度的下确界。
	\symidx{\(\lambda_{\mathbb D}\)：单位圆盘的双曲长度密度}
	\symidx{\(d_{\mathbb D}\)：单位圆盘的双曲距离}
\end{definition}

长度密度处处为正，而且 \(\lambda_{\mathbb D}\geq2\)，所以每条连接 \(z,w\) 的曲线长度至少为 \(2\lvert z-w\rvert\)。圆盘内的直线段给出有限长度的连接。反向参数化保持长度，拼接曲线的长度相加，因而这一构造确实给出有限的距离函数。密度在接近边界时增大，表示靠近单位圆的同样一小段平面长度会被赋予更大的双曲长度。

\begin{proposition}
	\label{g4:prop:0103:hyperbolic-distance-formula}
	对任意 \(z,w\in\mathbb D\)，令
	\[
		q=\left\lvert\frac{z-w}{1-\overline w z}\right\rvert.
	\]
	则 \(0\leq q<1\)，并且
	\begin{equation}
		\label{g4:eq:0103:hyperbolic-distance}
		d_{\mathbb D}(z,w)=\log\frac{1+q}{1-q}
		=2\operatorname{artanh}q.
	\end{equation}
	每个 \(\phi_a\) 和每个旋转都保持这一距离。
\end{proposition}

\begin{proof}
	由式\eqref{g4:eq:0103:disk-normalization-identities}，
	\[
		\lambda_{\mathbb D}(\phi_a(z))\lvert\phi_a'(z)\rvert
		=\lambda_{\mathbb D}(z).
	\]
	链式法则于是说明 \(\phi_a\) 保持每条曲线的双曲长度；它的逆也把连接曲线一一对应起来，所以保持长度的下确界。旋转的长度保持性直接来自模的不变性。取 \(a=w\)，便把待算距离化为原点到某个模为 \(q\) 的点的距离。

	设 \(\gamma\) 从 \(0\) 连到 \(\zeta\)。径向运动给出最短长度的候选；为比较任意曲线，先用光滑近似处理半径在过原点时可能不可微的问题。对 \(\varepsilon>0\)，令
	\[
		r_\varepsilon(t)=\sqrt{\lvert\gamma(t)\rvert^2+\varepsilon^2}-\varepsilon.
	\]
	它分片 \(C^1\)，满足 \(0\leq r_\varepsilon(t)\leq\lvert\gamma(t)\rvert\)，且由求导与复数模的不等式可得
	\[
		\lvert r_\varepsilon'(t)\rvert
		=\frac{\lvert\operatorname{Re}(\overline{\gamma(t)}\gamma'(t))\rvert}
		{\sqrt{\lvert\gamma(t)\rvert^2+\varepsilon^2}}
		\leq\lvert\gamma'(t)\rvert.
	\]
	令
	\[
		A(r)=\log\frac{1+r}{1-r},\qquad A'(r)=\frac2{1-r^2}.
	\]
	曲线像是圆盘内的紧集，所以这里的分母有正下界。由一元微积分的链式法则，
	\[
		L_{\mathbb D}(\gamma)
		\geq\int_0^1\frac{2\lvert r_\varepsilon'(t)\rvert}{1-r_\varepsilon(t)^2}\,\mathrm{d}t
		\geq A(r_\varepsilon(1))-A(r_\varepsilon(0)).
	\]
	因 \(r_\varepsilon(0)=0\)、\(r_\varepsilon(1)\rightarrow\lvert\zeta\rvert\)，令 \(\varepsilon\downarrow0\) 得 \(L_{\mathbb D}(\gamma)\geq A(\lvert\zeta\rvert)\)。径向线段 \(\gamma(t)=t\zeta\) 恰好达到这个下界，因为半径沿曲线单调增加。因此原点到 \(\zeta\) 的距离是 \(A(\lvert\zeta\rvert)\)，从而得到所示公式。
\end{proof}

Schwarz–Pick 的导数不等式现在可以写成
\[
	\lambda_{\mathbb D}(f(z))\lvert f'(z)\rvert
	\leq\lambda_{\mathbb D}(z).
\]
沿任意分片 \(C^1\) 曲线积分，得到 \(L_{\mathbb D}(f\circ\gamma)\leq L_{\mathbb D}(\gamma)\)。在所有连接两点的曲线上取下确界，便有
\begin{equation}
	\label{g4:eq:0103:hyperbolic-contraction}
	d_{\mathbb D}(f(z),f(w))\leq d_{\mathbb D}(z,w).
\end{equation}
因此全纯圆盘自映射都使双曲距离不增加。式\eqref{g4:eq:0103:hyperbolic-distance}中的实函数 \(q\mapsto\log((1+q)/(1-q))\) 严格递增，所以两个不同点间的距离相等，恰好对应\cref{g4:thm:0103:schwarz-pick}的两点等号情形，迫使映射具有全纯逆。

\begin{example}
	\label{g4:exm:0103:square-contraction}
	对 \(f(z)=z^2\)，双曲长度的局部伸缩比为
	\[
		\frac{\lambda_{\mathbb D}(f(z))\lvert f'(z)\rvert}
		{\lambda_{\mathbb D}(z)}
		=\frac{2\lvert z\rvert(1-\lvert z\rvert^2)}{1-\lvert z\rvert^4}
		=\frac{2\lvert z\rvert}{1+\lvert z\rvert^2}<1.
	\]
	最后一步来自 \((1-\lvert z\rvert)^2>0\)。原点处的伸缩比为零，与分歧相符；趋近边界时伸缩比趋于 \(1\)，所以在整个圆盘上不存在统一小于 \(1\) 的伸缩上界。另一方面，原点到实数 \(r\in(0,1)\) 的距离为 \(\log((1+r)/(1-r))\)，当 \(r\uparrow1\) 时趋于无穷。
\end{example}

\subsection{双全纯映射与自同构}
\label{g4:subsec:0103:04}
保持全纯结构的坐标变换要求正向与反向都全纯，这正是\cref{g4:def:0101:biholomorphic}中的双全纯映射。由\cref{g4:thm:0101:open-mapping}及其后的局部逆论证，全纯单射就是到其像的双全纯映射。区域 \(U\subset\mathbb C\) 到自身的双全纯映射称为 \(U\) 的\term[quanchunzitonggou]{全纯自同构}[holomorphic automorphism]，其全体在复合下构成群 \(\operatorname{Aut}(U)\)。对圆盘来说，Schwarz 引理的等号条件足以确定这个群的所有元素。

\begin{theorem}
	\label{g4:thm:0103:disk-aut}
	单位圆盘的全纯自同构恰好是
	\[
		f(z)=\mathrm{e}^{\mathrm{i}\theta}\frac{z-a}{1-\overline a z},
		\qquad a\in\mathbb D,\quad\theta\in\mathbb R.
	\]
	其中 \(a=f^{-1}(0)\) 唯一，\(\theta\) 按模 \(2\pi\) 唯一。特别地，固定原点的全纯自同构都是旋转。
\end{theorem}

\begin{proof}
	所示映射都是式\eqref{g4:eq:0103:disk-normalization}与旋转的复合，已知具有全纯逆。反过来，设 \(f\) 为全纯自同构，取 \(a=f^{-1}(0)\)，令 \(h=f\circ\phi_a^{-1}\)。则 \(h\) 与 \(h^{-1}\) 都固定 \(0\)。对二者分别应用\cref{g4:thm:0103:schwarz}，得到
	\[
		\lvert h(z)\rvert\leq\lvert z\rvert
		=\lvert h^{-1}(h(z))\rvert\leq\lvert h(z)\rvert.
	\]
	于是处处取等号，Schwarz 引理给出 \(h(z)=\mathrm{e}^{\mathrm{i}\theta}z\)，从而 \(f=R_\theta\circ\phi_a\)。\(a\) 由唯一的零点决定；固定 \(a\) 后，旋转因子也由任一非零点的像决定。
\end{proof}

一个圆盘坐标的选择因此有两种自由：先指定哪个点成为原点，再指定该点的方向。固定这两项就能得到唯一性。

\begin{corollary}
	\label{g4:cor:0103:normalized-unique}
	设 \(\Omega\subset\mathbb C\) 是区域，\(a\in\Omega\)。若 \(f,g\colon\Omega\rightarrow\mathbb D\) 均为双全纯映射，且
	\[
		f(a)=g(a)=0,\qquad f'(a)>0,\qquad g'(a)>0,
	\]
	其中导数的不等式表示正实数，则 \(f=g\)。
\end{corollary}

\begin{proof}
	映射 \(h=g\circ f^{-1}\) 是固定 \(0\) 的圆盘自同构，故由\cref{g4:thm:0103:disk-aut}可写成 \(h(z)=\mathrm{e}^{\mathrm{i}\theta}z\)。又
	\[
		\mathrm{e}^{\mathrm{i}\theta}=h'(0)=\frac{g'(a)}{f'(a)}>0.
	\]
	正实数中只有 \(1\) 的模为 \(1\)，所以 \(h\) 是恒等映射，\(f=g\)。
\end{proof}

\begin{example}
	\label{g4:exm:0103:half-plane}
	上半平面 \(\mathbb H=\{w\in\mathbb C:\operatorname{Im}w>0\}\) 与圆盘双全纯等价。具体映射及其逆为
	\[
		\psi(w)=\frac{w-\mathrm{i}}{w+\mathrm{i}},
		\qquad \psi^{-1}(z)=\mathrm{i}\frac{1+z}{1-z}.
	\]
	由
	\[
		1-\lvert\psi(w)\rvert^2
		=\frac{4\operatorname{Im}w}{\lvert w+\mathrm{i}\rvert^2},
		\qquad \lvert\psi'(w)\rvert=\frac2{\lvert w+\mathrm{i}\rvert^2},
	\]
	既可检查 \(\psi(\mathbb H)\subset\mathbb D\)，也可得到圆盘双曲长度的拉回
	\[
		\lambda_{\mathbb D}(\psi(w))\lvert\psi'(w)\rvert\lvert\mathrm{d}w\rvert
		=\frac{\lvert\mathrm{d}w\rvert}{\operatorname{Im}w}.
	\]
	因此上半平面的长度元素可以从圆盘坐标直接恢复。区域在平面中是否有界，不是双全纯等价的不变量。
\end{example}

圆盘自同构的分类解决了圆盘坐标的唯一性问题。对一个给定的平面区域，能否首先找到这样的坐标，需要构造全纯单射并控制其极限；下一节将用正规族把这一存在问题变成一个导数极值问题。
% 纲要条目排为编号小节标题，并进入目录与书签；说明段仍保留于纲要注释中。

% !TeX root = ../../Book4.tex
% 纲要标题：### 1.4 正规族与共形构造
% 纲要位置：几何学计划纲要.md，第 3502 行。
% 正文状态：初稿；尚未正式审核。

% 纲要内容开始
% 
% * 局部有界和 Montel 紧性
% * Hurwitz 定理与单叶极限
% * 平面 Riemann 映照定理的极值证明
% * 边界行为和非单连通域的限制
% 
% ---
% 
% 纲要内容结束

\section{正规族与共形构造}
\label{g4:sec:0104}

把一个平面域变成圆盘，通常找不到可以直接写出的公式。一个可行的办法是先收集所有满足某种归一化条件的单叶映射，再从中选取在指定点放大最多的映射。这立即带来两个问题：趋近最大值的一列映射能否收敛，极限是否还保留单叶性？Cauchy 估计提供前一个问题所需的紧性，零点数的稳定性处理后一个问题。两者合起来便能把一个极值问题变成共形映射的存在证明。

\subsection{局部有界和 Montel 紧性}
\label{g4:subsec:0104:01}

全纯函数在一点附近的变化速度受较大邻域上的函数值控制。因此，对一族全纯函数，统一的局部大小估计还会给出统一的局部连续性估计；这是一般连续函数族没有的性质。

\begin{definition}\label{g4:def:0104:normal-family}
	设 \(\Omega\subset\mathbb C\) 为域。函数列 \(f_j:\Omega\to\mathbb C\) 在每个紧集 \(K\subset\Omega\) 上一致收敛于 \(f\) 时，称其\term[ju bu yi zhi shou lian]{局部一致收敛}[locally uniform convergence]于 \(f\)。对于全纯函数族 \(\mathcal F\subset\mathcal O(\Omega)\)，若每个紧集 \(K\subset\Omega\) 都有常数 \(M_K<\infty\)，使
	\[
		|f(z)|\le M_K\qquad(f\in\mathcal F,\ z\in K),
	\]
	则称 \(\mathcal F\)\term[ju bu you jie]{局部有界}[locally bounded]。若 \(\mathcal F\) 中每个函数列都有一个局部一致收敛于 \(\Omega\) 上全纯函数的子列，则称它在全纯函数空间中为\term[zheng gui zu]{正规族}[normal family]。这里采用有限复值极限；亚纯函数族还可以使用球面距离定义另一种容许无穷远值的正规性。
\end{definition}

局部一致收敛允许靠近边界时失去统一控制。例如 \(f_j(z)=z^j\) 在单位圆盘 \(\mathbb D\) 上局部一致收敛于 \(0\)，因为每个紧集都包含在某个 \(|z|\le r<1\) 中，且 \(|z|^j\le r^j\)。但 \(\sup_{z\in\mathbb D}|z^j|=1\)，所以它们不在整个圆盘上一致收敛。

\begin{proposition}\label{g4:prop:0104:holomorphic-limit}
	设 \(\Omega\subset\mathbb C\) 为域，全纯函数列 \(f_j\) 在 \(\Omega\) 上局部一致收敛于 \(f\)。则 \(f\) 全纯，而且对每个整数 \(k\ge1\)，\(f_j^{(k)}\) 局部一致收敛于 \(f^{(k)}\)。
\end{proposition}
\begin{proof}
	局部一致极限首先是连续函数。取闭圆盘 \(\overline{D(a,R)}\subset\Omega\)。在 \(|z-a|<R\) 上，\cref{g4:thm:0101:cauchy-formula}给出
	\[
		f_j(z)=\frac1{2\pi\mathrm i}\int_{|\zeta-a|=R}
		\frac{f_j(\zeta)}{\zeta-z}\,\mathrm d\zeta.
	\]
	在积分圆周上的一致收敛允许取极限，得到同一公式中的 \(f_j\) 全部换成 \(f\)。对任意 \(0<r<R\)，当 \(|z-a|\le r\) 时，分母的模至少为 \(R-r\)；核函数对 \(z\) 的各阶导数因此都能在积分号下求出。这证明 \(f\) 在圆盘内全纯，并给出
	\[
		\sup_{|z-a|\le r}|f_j^{(k)}(z)-f^{(k)}(z)|
		\le \frac{k!R}{(R-r)^{k+1}}
		\sup_{|\zeta-a|=R}|f_j(\zeta)-f(\zeta)|.
	\]
	右端趋于零。任意紧集可用有限个这样的较小圆盘覆盖，故各阶导数也局部一致收敛。
\end{proof}

要抽取收敛子列，还需要一个紧空间上的基本论证。设 \(K\) 为紧度量空间，连续函数列在 \(K\) 上一致有界且等度连续。取可数稠密集 \(\{x_1,x_2,\ldots\}\subset K\)，依次抽取在 \(x_1,x_2,\ldots\) 处收敛的子列，再取对角子列。它在每个 \(x_\nu\) 处收敛。给定 \(\varepsilon>0\)，等度连续性给出统一的 \(\delta>0\)，使距离小于 \(\delta\) 时函数值相差小于 \(\varepsilon/3\)。用稠密集中的有限多个点组成 \(\delta\)-网；子列在这些点处同时成为 Cauchy 列后，在整个 \(K\) 上的差就小于 \(\varepsilon\)。因此该子列一致收敛。可数稠密集本身可由半径 \(1,1/2,1/3,\ldots\) 的有限网的并得到。

\ProseParagraphBreak

把这个紧性论证与 Cauchy 估计结合，得到\term[Montel ding li]{Montel 定理}[Montel's theorem]。Montel\footnote{Paul Antoine Aristide Montel（1876—1975），法国数学家，发展了全纯函数族与正规族理论。}的结论在这里承担的任务，是把一族函数的大小控制转化为一个极限函数的存在。

\begin{theorem}[Montel]\label{g4:thm:0104:montel}
	设 \(\Omega\subset\mathbb C\) 为域。每个局部有界的全纯函数族 \(\mathcal F\subset\mathcal O(\Omega)\) 都是正规族。
\end{theorem}
\begin{proof}
	先证明任意紧集 \(K\subset\Omega\) 上的等度连续性。取 \(\delta>0\)，使
	\[
		L=\{z\in\mathbb C:\operatorname{dist}(z,K)\le2\delta\}
		\subset\Omega.
	\]
	集合 \(L\) 紧，因而存在对全部 \(f\in\mathcal F\) 通用的上界 \(M_L\)。在距 \(K\) 不超过 \(\delta\) 的点处，半径 \(\delta\) 的圆盘包含在 \(L\) 中，\cref{g4:prop:0101:cauchy-estimate}给出 \(|f'|\le M_L/\delta\)。若 \(z,w\in K\) 且 \(|z-w|<\delta\)，连接二者的线段仍在这个邻域内，故
	\[
		|f(z)-f(w)|
		\le\frac{M_L}{\delta}|z-w|.
	\]
	这就是在 \(K\) 上与 \(f\) 无关的连续性估计。

	现在选择一列紧集穷竭 \(\Omega\)。当 \(\Omega\ne\mathbb C\) 时，可取
	\[
		K_m=\left\{z\in\Omega:|z|\le m,\
		\operatorname{dist}(z,\mathbb C\setminus\Omega)\ge\frac1m\right\},
		\qquad m\ge1;
	\]
	当 \(\Omega=\mathbb C\) 时，取 \(K_m=\{|z|\le m\}\)。这些紧集递增，其内部覆盖 \(\Omega\)，每个 \(\Omega\) 内的紧集都包含在某个 \(K_m\) 中。

	给定 \(\mathcal F\) 中的函数列，利用前述紧空间论证，先抽取在 \(K_1\) 上一致收敛的子列，再从中抽取在 \(K_2\) 上一致收敛的子列，如此继续；空的 \(K_m\) 可略过。取这些嵌套子列的对角子列，它在每个 \(K_m\) 上一致收敛。各个极限在交集上一致，因而定义 \(\Omega\) 上的函数 \(f\)，且所得收敛是局部一致的。由\cref{g4:prop:0104:holomorphic-limit}，\(f\) 全纯。
\end{proof}

局部有界条件必须对函数族统一成立。每个 \(f_j(z)=jz\) 都在紧集上有界，但这列函数在任意含非零点的邻域上没有统一上界，也不可能有局部一致收敛的子列。

\subsection{Hurwitz 定理与单叶极限}
\label{g4:subsec:0104:02}

收敛的全纯函数列仍可能失去某些几何性质。例如 \(f_j(z)=z/j\) 都是单叶函数，极限却是常值。真正需要排除的是这种整体塌缩；一旦极限不恒定，Rouché 定理对零点的控制就会保住单叶性。

\begin{theorem}[Hurwitz]\label{g4:thm:0104:hurwitz}
	设 \(\Omega\subset\mathbb C\) 为域，\(f_j\in\mathcal O(\Omega)\) 局部一致收敛于 \(f\)。若 \(f\) 不恒为零，且 \(a\in\Omega\) 是它的 \(m\) 阶零点，取 \(r>0\)，使 \(\overline{D(a,r)}\subset\Omega\) 且该闭圆盘内没有 \(f\) 的其他零点，则对充分大的 \(j\)，\(f_j\) 在 \(D(a,r)\) 中恰有 \(m\) 个零点，按重数计数。特别地，若每个 \(f_j\) 都无零点，则极限 \(f\) 或者处处非零，或者恒为零。
\end{theorem}
\begin{proof}
	圆周 \(|z-a|=r\) 紧，且其上 \(f\ne0\)，所以
	\[
		\eta=\min_{|z-a|=r}|f(z)|>0.
	\]
	当 \(j\) 足够大时，圆周上有 \(|f_j-f|<\eta\le|f|\)。\cref{g4:thm:0102:rouche}说明 \(f_j\) 与 \(f\) 在圆盘中的零点总数相同，恰为 \(m\)。若所有 \(f_j\) 无零点而非零函数 \(f\) 在一点为零，这个结论就产生矛盾，故得到最后的二择一结论。
\end{proof}

上述结果称为\term[Hurwitz ding li]{Hurwitz 定理}[Hurwitz's theorem]。Hurwitz\footnote{Adolf Hurwitz（1859—1919），德国数学家，研究复函数、代数曲线与数论。}的零点稳定性允许多个零点在极限中合并，但总重数保持不变；例如 \(z^2-j^{-2}\) 的两个简单零点趋向 \(z^2\) 的二重零点。常数列 \(f_j=1/j\) 则说明，处处非零的函数列确实可以趋于恒零函数。

\begin{corollary}\label{g4:cor:0104:univalent-limit}
	设 \(\Omega\subset\mathbb C\) 为域，单叶全纯函数列 \(f_j:\Omega\to\mathbb C\) 局部一致收敛于 \(f\)。则 \(f\) 或为常值函数，或仍为单叶全纯函数。
\end{corollary}
\begin{proof}
	由\cref{g4:prop:0104:holomorphic-limit}，\(f\) 全纯。设 \(f\) 不恒定，固定任意 \(a\in\Omega\)，定义
	\[
		g_j(z)=
		\begin{cases}
			\dfrac{f_j(z)-f_j(a)}{z-a},&z\ne a,\\[6pt]
			f_j'(a),&z=a.
		\end{cases}
	\]
	幂级数说明 \(g_j\) 在 \(a\) 处全纯。由于 \(f_j\) 单叶，\cref{g4:thm:0101:local-power}中局部幂模型的次数只能是 \(1\)，故 \(f_j'(a)\ne0\)；在其他点，单叶性也给出 \(g_j(z)\ne0\)。因此每个 \(g_j\) 在整个 \(\Omega\) 上无零点。

	这些函数局部一致收敛于同样由 \(f\) 定义的差商 \(g\)。在离开 \(a\) 的紧集上，这直接来自 \(f_j\to f\)；在一个以 \(a\) 为中心的较小圆盘上，可用
	\[
		g_j(z)=\int_0^1 f_j'\bigl(a+t(z-a)\bigr)\,\mathrm dt
	\]
	及导数的局部一致收敛验证。由于 \(f\) 不恒定，\(g\) 不恒为零，\cref{g4:thm:0104:hurwitz}于是给出 \(g\) 处处非零。因此 \(z\ne a\) 时 \(f(z)\ne f(a)\)。点 \(a\) 任意，故 \(f\) 单叶。
\end{proof}

在构造映射时，固定一个点处的导数下界就能排除常值极限。这正是下面归一化极值问题的作用。

\subsection{平面 Riemann 映照定理的极值证明}
\label{g4:subsec:0104:03}

圆盘的自同构已经说明：即使一个域与圆盘双全纯，对应也有选择余地。指定一点的像为 \(0\)，再要求该点的导数为正实数，就消除了平移基点和旋转的自由度。以下 \(\mathbb D=\{w\in\mathbb C:|w|<1\}\)，并沿用圆盘自同构
\[
	\phi_b(w)=\frac{w-b}{1-\overline b w},
	\qquad b\in\mathbb D.
\]

\begin{theorem}[平面 Riemann 映照定理]\label{g4:thm:0104:riemann-mapping}
	设 \(\Omega\subsetneq\mathbb C\) 为非空单连通域，\(a\in\Omega\)。则存在唯一的双全纯映射 \(F:\Omega\to\mathbb D\)，满足
	\[
		F(a)=0,\qquad F'(a)>0.
	\]
\end{theorem}

这个结果称为\term[Riemann ying zhao ding li]{Riemann 映照定理}[Riemann mapping theorem]。证明分为三个相互衔接的构造：先找出至少一个从 \(\Omega\) 到圆盘的单叶映射；再在归一化映射中使 \(F'(a)\) 达到上确界；最后证明只要还遗漏圆盘内一点，就能构造导数更大的映射，与极值性矛盾。这一极值方法可与 Ahlfors 的论证对照\cite[第6章第1.1节, 第230--231页]{Ahlfors1979Complex}。

\begin{proof}
	先构造一个有界单叶函数。取 \(b\in\mathbb C\setminus\Omega\)。函数 \(z-b\) 在单连通域 \(\Omega\) 上全纯且处处非零，\cref{g4:prop:0102:logarithm}给出全纯平方根 \(h\)，使
	\[
		h(z)^2=z-b.
	\]
	若 \(h(z_1)=h(z_2)\)，平方后得到 \(z_1=z_2\)，故 \(h\) 单叶。而 \(h(\Omega)\) 与 \(-h(\Omega)\) 不交：若 \(h(z_1)=-h(z_2)\)，平方后仍有 \(z_1=z_2\)，从而 \(h(z_1)=0\)，与 \(z_1\ne b\) 矛盾。

	令 \(c=h(a)\)。由\cref{g4:thm:0101:open-mapping}，存在 \(\rho>0\) 使 \(D(c,\rho)\subset h(\Omega)\)。于是 \(D(-c,\rho)\subset-h(\Omega)\) 完全避开 \(h(\Omega)\)，并且
	\[
		|h(z)+c|\ge\rho\qquad(z\in\Omega).
	\]
	故
	\[
		q(z)=\frac{\rho}{2(h(z)+c)}
	\]
	是模不超过 \(1/2\) 的单叶全纯函数。将 \(q\) 与 \(\phi_{q(a)}\) 复合，再乘一个模为 \(1\) 的常数，就得到一个 \(f:\Omega\to\mathbb D\)，满足 \(f(a)=0\)、\(f'(a)>0\)。这里导数不为零，因为单叶全纯函数的局部幂模型次数为 \(1\)。

	考虑所有满足这些条件的函数组成的集合
	\begin{gather*}
		\mathcal S=\{f:\Omega\to\mathbb D:
		f\text{ 全纯且单叶},\ f(a)=0,\ f'(a)>0\},\\
		M=\sup_{f\in\mathcal S}f'(a).
	\end{gather*}
	刚才的构造说明 \(M>0\)。取 \(r>0\) 使 \(\overline{D(a,r)}\subset\Omega\)，由\cref{g4:prop:0101:cauchy-estimate}有 \(f'(a)\le1/r\)，所以 \(M<\infty\)。选取 \(f_j\in\mathcal S\)，使 \(f_j'(a)\to M\)。函数族一致有界于 \(1\)，\cref{g4:thm:0104:montel}给出局部一致收敛的子列，仍记作 \(f_j\)，其全纯极限记作 \(F\)。

	由\cref{g4:prop:0104:holomorphic-limit}，
	\[
		F(a)=0,\qquad F'(a)=M>0,\qquad |F(z)|\le1.
	\]
	因 \(F\) 不恒定，\cref{g4:thm:0103:maximum-modulus}保证 \(|F(z)|<1\)。再由\cref{g4:cor:0104:univalent-limit}，\(F\) 单叶。因此 \(F\in\mathcal S\)，上确界确实由一个映射达到。

	下面证明 \(F\) 满射。若有 \(\beta\in\mathbb D\setminus F(\Omega)\)，则 \(\beta\ne0\)。将遗漏点移到原点，令
	\[
		H=\phi_\beta\circ F.
	\]
	这是 \(\Omega\) 上处处非零的单叶全纯函数，值域包含于 \(\mathbb D\)。再次使用单连通性和\cref{g4:prop:0102:logarithm}，取全纯函数 \(Q\) 满足 \(Q^2=H\)。于是 \(|Q|<1\)，且 \(Q\) 单叶，因为相等的 \(Q\) 值必给出相等的 \(H\) 值。记 \(\gamma=Q(a)\)，则
	\[
		\gamma^2=-\beta,\qquad 0<|\gamma|<1.
	\]
	映射 \(G=\phi_\gamma\circ Q\) 仍将 \(\Omega\) 单叶地映入 \(\mathbb D\)，且 \(G(a)=0\)。为比较导数，利用
	\[
		|\phi_\beta'(0)|=1-|\beta|^2,
		\qquad |\phi_\gamma'(\gamma)|=\frac1{1-|\gamma|^2},
		\qquad 2\gamma Q'(a)=H'(a),
	\]
	得到
	\begin{align*}
		|G'(a)|
		&=\frac{1-|\beta|^2}{2|\gamma|(1-|\gamma|^2)}F'(a)\\
		&=\frac{1+|\gamma|^2}{2|\gamma|}M
		>M.
	\end{align*}
	最后一个不等式来自 \((1-|\gamma|)^2>0\)。再旋转 \(G\) 的值域，使导数成为正实数，就得到 \(\mathcal S\) 中导数大于 \(M\) 的函数，矛盾。因此 \(F(\Omega)=\mathbb D\)。单叶性使 \(F'\) 处处非零，\cref{g4:thm:0101:inverse}保证其逆映射全纯。

	最后，若 \(F_1,F_2\) 都满足结论，则 \(A=F_2\circ F_1^{-1}\) 是固定 \(0\) 的圆盘自同构。由\cref{g4:thm:0103:disk-aut}，\(A(w)=e^{\mathrm i\theta}w\)。而
	\[
		A'(0)=\frac{F_2'(a)}{F_1'(a)}>0,
	\]
	故 \(e^{\mathrm i\theta}=1\)，从而 \(F_1=F_2\)。
\end{proof}

平方根在证明中出现两次，解决的是两个具体问题。第一次利用域外点，把原域映到一个避开整块圆盘的区域，随后取倒数得到有界函数。第二次利用像域内的遗漏点，在保持单叶性的同时扩大基点处的导数。单连通性保证这两次平方根都能在整个 \(\Omega\) 上选成单值全纯函数；正规族紧性则保证最大化过程不会失去极限。

\subsection{边界行为和非单连通域的限制}
\label{g4:subsec:0104:04}

Riemann 映照定理给出开域之间的双全纯对应。若一个点列在源域中不断逼近边界，它在圆盘中的像也必须离开每个内部紧集。这个结论只需逆映射的连续性，已经能够说明许多边界现象。

\begin{proposition}\label{g4:prop:0104:boundary-escape}
	设 \(\Omega\subset\mathbb C\) 为域，\(F:\Omega\to\mathbb D\) 为双全纯映射，\(z_j\in\Omega\)。则以下两件事等价：点列 \(z_j\) 最终离开 \(\Omega\) 的每个紧集；\(|F(z_j)|\to1\)。特别地，若 \(z_j\to p\in\partial\Omega\)，或 \(|z_j|\to\infty\)，则 \(|F(z_j)|\to1\)。
\end{proposition}
\begin{proof}
	若 \(|F(z_j)|\) 不趋于 \(1\)，则存在 \(r<1\) 和一个子列，其像都位于 \(\overline{D(0,r)}\)。逆映射连续，故
	\[
		F^{-1}\bigl(\overline{D(0,r)}\bigr)
	\]
	是 \(\Omega\) 中的紧集，与点列最终离开每个紧集矛盾。反过来，对任意紧集 \(K\subset\Omega\)，\(F(K)\) 紧且包含在 \(\mathbb D\) 中，故 \(\max_K|F|<1\)。若 \(|F(z_j)|\to1\)，则充分大的 \(j\) 都有 \(z_j\notin K\)。最后两种点列都最终离开 \(\Omega\) 的每个紧集，因而满足结论。
\end{proof}

上式控制的是像点到圆周的径向距离。若要在每个边界点赋予唯一的圆周像，还需要研究从不同方向逼近同一点时的行为。狭缝域给出一个具体区别。

\begin{example}\label{g4:exm:0104:slit-boundary}
	考虑
	\[
		\Omega=\mathbb D\setminus[0,1).
	\]
	每个点有唯一表示 \(z=re^{\mathrm i\theta}\)，其中 \(0<r<1\)、\(0<\theta<2\pi\)。映射 \(z\mapsto(\log r,\theta)\) 把 \(\Omega\) 同胚到凸集 \((-\infty,0)\times(0,2\pi)\)，所以 \(\Omega\) 单连通；\cref{g4:thm:0104:riemann-mapping}给出它到圆盘的双全纯映射。

	然而，这个映射不可能延拓为 Euclid 闭包之间的同胚。因为
	\[
		\overline\Omega=\overline{\mathbb D},
		\qquad
		\overline\Omega\setminus\Omega=\partial\mathbb D\cup[0,1).
	\]
	若存在延拓同胚，它须把 \(\partial\mathbb D\cup[0,1)\) 同胚地映到 \(\partial\mathbb D\)。从源边界删去狭缝上的一点 \(p\in(0,1)\) 后，该集合分成包含 \(0\) 的一段和其余部分；圆周删去一点却仍然连通，故这样的边界同胚不存在。

	局部公式 \(\sqrt z=\sqrt r\,e^{\mathrm i\theta/2}\) 更直观地显示了狭缝的两侧：从上侧逼近 \(p>0\) 时，它趋向 \(\sqrt p\)；从下侧逼近时，它趋向 \(-\sqrt p\)。源域边界中的同一个 Euclid 点携带着两种不同的内部逼近方向。处理边界对应时，必须记录这种区别。
\end{example}

单连通假设也有可计算的作用。设 \(0<r<1\)，圆环
\[
	A_r=\{z\in\mathbb C:r<|z|<1\}
\]
上的坐标函数 \(z\) 没有全纯对数：若 \(e^{L(z)}=z\)，求导得 \(L'(z)=1/z\)，但沿 \(|z|=\rho\)、\(r<\rho<1\) 积分会同时给出
\[
	\int_{|z|=\rho}L'(z)\,\mathrm dz=0,
	\qquad
	\int_{|z|=\rho}\frac{\mathrm dz}{z}=2\pi\mathrm i.
\]
若有双全纯映射 \(F:A_r\to\mathbb D\)，其逆映射 \(G:\mathbb D\to A_r\) 是处处非零的全纯函数。由\cref{g4:prop:0102:logarithm}，\(G\) 有全纯对数 \(H\)，于是 \(H\circ F\) 就是 \(A_r\) 上坐标函数的全纯对数，产生矛盾。因此圆环不能双全纯地变成圆盘。

\ProseParagraphBreak

另一个被排除的域是整个 \(\mathbb C\)。若 \(F:\mathbb C\to\mathbb D\) 全纯，则在任意 \(a\in\mathbb C\) 处，半径 \(R\) 的 Cauchy 估计给出 \(|F'(a)|\le1/R\)。令 \(R\to\infty\) 得 \(F'(a)=0\)，所以 \(F\) 恒定。这说明定理中的“单连通真子域”同时包含两个作用不同的条件：单连通性消除选取全纯根时的环路障碍，域外点则启动有界单叶函数的构造。


\begin{chaptersummary}
	全纯性首先是一个一阶条件。三角形积分的消去和 Cauchy 公式使这个条件产生全部高阶导数与幂级数，从而把非恒定全纯函数在每点附近化为 \(w\mapsto w^m\)。整数 \(m\) 既是零点重数，也是邻近原像的数目；留数和辐角原理把这些局部整数合成边界上的绕数。Rouché 定理说明，只要边界不经过零，这个计数便在适当扰动中保持。

	\ProseParagraphBreak

	解析延拓的对象是带中心的函数芽。局部相等保证接续唯一，沿所有道路可延拓与同伦不变性进一步给出单连通域上的单值函数。极值原理把值域的控制转化为刚性，Schwarz--Pick 将其写成双曲长度与距离不增加的形式。正规族则利用同一类积分估计，把局部有界性转成紧性；零点稳定性使非退化的单叶极限继续单叶。对归一化映射的导数取极值，再以遗漏值构造更优候选，就得到平面 Riemann 映照定理。所得坐标首先定义在域内，边界的逐点对应另受边界几何的约束。
\end{chaptersummary}

\BookExercises

\begin{exercise}
	\label{g4:ex:0101:rigid-images}
	设 \(\Omega\subset\mathbb C\) 是区域，\(f\in\mathcal O(\Omega)\)。分别证明：若 \(f\) 的值都落在同一条实仿射直线上，或者 \(|f|\) 恒定，则 \(f\) 为常数。第一种情形另用 Cauchy--Riemann 方程证明。
\end{exercise}
\begin{solution}
	直线和圆都不包含平面开集，所以\cref{g4:thm:0101:open-mapping}排除了非常值函数。若常模为零，结论由定义直接成立。

	对直线情形，先用非零复数的乘法和一次平移把该直线送到实轴，得到全纯实值函数 \(g=u+\mathrm i0\)。Cauchy--Riemann 方程给出 \(u_x=u_y=0\)。区域内任意两点能以有限条直线段相连，沿这些线段积分得到 \(u\) 为常数，还原仿射变换即可。
\end{solution}

\begin{exercise}
	\label{g4:ex:0101:polynomial-growth}
	设整函数 \(f\) 满足 \(|f(z)|\le C(1+|z|)^m\)，其中 \(C>0\)、\(m\) 为非负整数。证明 \(f\) 是次数至多为 \(m\) 的多项式。再用有界整函数必为常数这一结论，证明每个非常数复多项式至少有一个复根。
\end{exercise}
\begin{solution}
	对 \(n>m\)，以原点为中心用\cref{g4:prop:0101:cauchy-estimate}，得到
	\[
		|f^{(n)}(0)|\le n!C\frac{(1+R)^m}{R^n}\longrightarrow0\qquad(R\to\infty).
	\]
	因此所有 \(n>m\) 的幂级数系数为零。任意有限圆盘上都可用\cref{g4:thm:0101:cauchy-formula}展开，同一组系数便在整个平面表示 \(f\)，所以它是所述多项式。

	若非常数多项式 \(P\) 处处非零，则 \(1/P\) 整。最高次项表明 \(|P(z)|\to\infty\) 当 \(|z|\to\infty\)，故 \(1/P\) 在某大圆盘外有界；在闭圆盘内又因连续且分母非零而有界。因此 \(1/P\) 为常数，矛盾。
\end{solution}

\begin{exercise}
	\label{g4:ex:0101:composition-order}
	设 \(f\) 在 \(a\) 附近全纯，\(g\) 在 \(b=f(a)\) 附近全纯，二者均不局部恒定。记 \(m=\operatorname{ord}_a(f-b)\)、\(n=\operatorname{ord}_b(g-g(b))\)。证明复合映射在 \(a\) 的局部重数为 \(mn\)。并证明重数不随源、靶的双全纯局部坐标改变。
\end{exercise}
\begin{solution}
	写 \(f(z)-b=(z-a)^m u(z)\)、\(g(w)-g(b)=(w-b)^n v(w)\)，其中 \(u(a)v(b)\ne0\)。代入得
	\[
		g(f(z))-g(b)=(z-a)^{mn}u(z)^n v(f(z)).
	\]
	最后一个因子在 \(a\) 非零，故阶数恰为 \(mn\)。双全纯局部坐标及其逆的导数非零，局部重数都是 \(1\)；前后复合只将重数乘以 \(1\)，于是得到坐标无关性。
\end{solution}

\begin{exercise}
	\label{g4:ex:0101:local-root-radius}
	在原点附近取满足 \(s(0)=1\)、\(s(z)^2=1+z\) 的全纯分支。求其幂级数的前三个非恒定项，并证明以原点为中心的收敛半径恰为 \(1\)。
\end{exercise}
\begin{solution}
	平方映射在 \(1\) 的导数为 \(2\)，逆函数定理保证该分支存在且由初值唯一。将 \(s=1+a_1z+a_2z^2+a_3z^3+\cdots\) 平方比较系数，得到 \(2a_1=1\)、\(2a_2+a_1^2=0\)、\(2a_3+2a_1a_2=0\)，所以
	\[
		s(z)=1+\frac z2-\frac{z^2}8+\frac{z^3}{16}+\cdots.
	\]
	在 \(\mathbb D\) 上，函数 \(1+z\) 不为零，由\cref{g4:prop:0102:logarithm}可取与该初值一致的全纯平方根。故幂级数在每个半径小于 \(1\) 的圆盘收敛，收敛半径至少为 \(1\)。若大于 \(1\)，级数和就在 \(-1\) 附近全纯。恒等定理使 \(s^2=1+z\) 仍成立，于是 \(s(-1)=0\)；对等式求导又要求 \(2s(-1)s'(-1)=1\)，矛盾。
\end{solution}

\begin{exercise}
	\label{g4:ex:0101:residue-order}
	设 \(a\in\mathbb C\setminus\{0\}\)，\(F(z)=(z^2-a^2)/z^3\)。计算 \(F'/F\) 的全部留数，并求当 \(r>0\) 且 \(r\ne|a|\) 时，闭曲线 \(F(r e^{\mathrm i t})\)、\(0\le t\le2\pi\)，绕原点的绕数。解释为何本例可以出现负绕数。
\end{exercise}
\begin{solution}
	对数导数为
	\[
		\frac{F'}F=\frac1{z-a}+\frac1{z+a}-\frac3z.
	\]
	在 \(a,-a,0\) 的留数依次为 \(1,1,-3\)，没有其他极点。由\cref{g4:thm:0102:argument}，\(0<r<|a|\) 时绕数为 \(-3\)，\(r>|a|\) 时为 \(1+1-3=-1\)。当 \(r=|a|\) 时像曲线经过零，绕数不定义。负号来自圆盘内部的三阶极点；对于在整个闭圆盘邻域全纯的函数，辐角原理只有零点贡献，不会产生这种负计数。
\end{solution}

\begin{exercise}
	\label{g4:ex:0101:rouche-polynomial}
	不求出各根的显式表达式，证明 \(P(z)=z^6+4z^2+1\) 在 \(|z|<1\) 内恰有两个根，在 \(1<|z|<2\) 内恰有四个根，均按重数计算。并证明全部根都是单根。
\end{exercise}
\begin{solution}
	在 \(|z|=1\) 上，\(|z^6+1|\le2<4=|4z^2|\)。Rouché 定理给出与 \(4z^2\) 相同的两个零点，且边界无零点。在 \(|z|=2\) 上，\(|4z^2+1|\le17<64=|z^6|\)，所以大圆盘内有六个零点、边界仍无零点。相减得到环域中的四个。

	若 \(z\) 是重根，则 \(P'(z)=2z(3z^4+4)=0\)。因 \(P(0)=1\)，必有 \(z\ne0\)、\(z^4=-4/3\)。代回 \(P(z)=0\) 得 \(8z^2/3+1=0\)，即 \(z^2=-3/8\)。于是 \(z^4=9/64\)，与 \(z^4=-4/3\) 矛盾，所以不存在重根。
\end{solution}

\begin{exercise}
	\label{g4:ex:0101:annular-roots}
	设 \(A=\{r<|z|<R\}\)，其中 \(0<r<R<\infty\)，\(k\in\mathbb Z\)，\(n\ge2\)。判断 \(z^k\) 在 \(A\) 上何时有全纯对数，何时有全纯 \(n\) 次根，并给出充分性构造。
\end{exercise}
\begin{solution}
	取 \(r<\rho<R\)，沿 \(|z|=\rho\) 逆时针积分。若 \(e^L=z^k\)，则 \(L'=k/z\)，所以
	\[
		0=\int_{|z|=\rho}L'(z)\,\mathrm{d}z=2\pi\mathrm i k.
	\]
	必要条件为 \(k=0\)，此时取 \(L=0\) 即可。

	若 \(g^n=z^k\)，则 \(g\) 处处非零，微分得到 \(n g'/g=k/z\)。积分并除以 \(2\pi\mathrm i\)，左端是整数绕数的 \(n\) 倍，所以必须有 \(n\mid k\)。反过来，若 \(k=n\ell\)，取 \(g(z)=z^\ell\) 即可，负的 \(\ell\) 在环域上同样全纯。任意两个根之比取值于有限个 \(n\) 次单位根，连通性使其恒定，故全部根正是这些单位根乘以 \(z^\ell\)。
\end{solution}

\begin{exercise}
	\label{g4:ex:0101:cubic-monodromy}
	从 \(z=1\) 处取值为 \(1\) 的 \(z^{1/3}\) 函数芽出发，沿 \(\gamma(t)=e^{2\pi\mathrm i t}\)、\(0\le t\le1\)，作解析延拓。求终点函数芽，并求绕行三周后的结果。说明为何在局部处处存在逆分支，仍不能在 \(\mathbb C^*\) 上得到单值立方根。
\end{exercise}
\begin{solution}
	沿路径取值 \(w(t)=e^{2\pi\mathrm i t/3}\)。它满足 \(w(t)^3=\gamma(t)\)，且 \(w(t)\ne0\)。立方映射在各 \(w(t)\) 处导数 \(3w(t)^2\) 非零，给出相应的全纯局部逆；紧参数区间可取有限个这样的邻域组成延拓链。终点值为 \(e^{2\pi\mathrm i/3}\)，所以终点芽等于起点芽乘以该常数；此结论也由局部根的唯一性得到。三周之后乘子为 \(e^{2\pi\mathrm i}=1\)，回到初始芽。一次环路的返回已经改变芽，故不能来自 \(\mathbb C^*\) 上的单值全纯函数；上一题的 \(k=1,n=3\) 给出同一障碍的积分表达。
\end{solution}

\begin{exercise}
	\label{g4:ex:0101:two-fixed-points}
	设 \(f:\mathbb D\to\mathbb D\) 全纯。证明：若 \(f\) 固定两个不同的内点，则 \(f\) 是恒等映射；若 \(f(a)=a\) 且 \(f'(a)=1\)，也有同一结论。并给出只固定一个内点而不恒等的例子。
\end{exercise}
\begin{solution}
	令 \(h=\phi_a\circ f\circ\phi_a^{-1}\)，其中 \(\phi_a(z)=(z-a)/(1-\overline a z)\)。它固定 \(0\)。若另有不同固定点 \(b\)，则 \(h\) 还固定非零点 \(\phi_a(b)\)。Schwarz 引理的等号情形使 \(h(z)=e^{\mathrm i\theta}z\)，固定该非零点又使 \(e^{\mathrm i\theta}=1\)。若改用导数条件，则链式法则与 \(f(a)=a\) 给出 \(h'(0)=f'(a)=1\)，同样由等号条件知 \(h\) 为恒等。还原共轭得到 \(f\) 为恒等。旋转 \(f(z)=-z\) 固定原点但并非恒等，说明一个固定点本身不够。
\end{solution}

\begin{exercise}
	\label{g4:ex:0101:harmonic-oscillation}
	设 \(0<r<R\)，\(c_0,c_1\in\mathbb R\)。在圆环 \(A=\{r<|z|<R\}\) 上构造一个调和函数，使它连续延拓到闭圆环，在内圆周取常值 \(c_0\)，在外圆周取常值 \(c_1\)。证明解唯一，并判断何时存在全局调和共轭。
\end{exercise}
\begin{solution}
	调和函数 \(\log|z|\) 在两条圆周上分别取 \(\log r\)、\(\log R\)。因此可作仿射组合
	\[
		u(z)=c_0+\frac{c_1-c_0}{\log(R/r)}\log\frac{|z|}{r}.
	\]
	它调和、连续至边界，并满足所给边值。若有另一个解，两解之差在边界为零，\cref{g4:cor:0103:harmonic-maximum}使其在整个圆环为零，故解唯一。

	令 \(A_0=(c_1-c_0)/\log(R/r)\)。局部调和共轭为 \(A_0\arg z\) 加常数，其微分沿任意正向同心圆的积分为 \(2\pi A_0\)。由\cref{g4:prop:0103:harmonic-conjugate}，全局共轭存在必有 \(A_0=0\)，即 \(c_1=c_0\)。反之，此时 \(u\) 为常数，取任意实常数作其共轭即可。
\end{solution}

\begin{exercise}
	\label{g4:ex:0101:degenerate-limits}
	设 \(\Omega\subset\mathbb C\) 是区域，单叶全纯函数列 \(f_n\) 局部一致收敛于非恒定函数 \(f\)。证明：每个紧集 \(K\subset f(\Omega)\) 都从某一项起包含在 \(f_n(\Omega)\) 内，而且逆映射 \(f_n^{-1}\) 在 \(K\) 上一致趋于 \(f^{-1}\)。
\end{exercise}
\begin{solution}
	由\cref{g4:cor:0104:univalent-limit}，\(f\) 单叶，所以 \(g=f^{-1}\) 有定义。固定 \(w_0\in f(\Omega)\) 及 \(\varepsilon>0\)，取 \(0<\rho<\varepsilon/2\)，使 \(\overline{D(g(w_0),\rho)}\subset\Omega\)。在该圆周上 \(f-w_0\) 无零点，令其模的最小值为 \(\eta>0\)。

	对 \(|w-w_0|<\eta/3\)，Rouché 定理说明 \(f-w\) 在此圆盘中仍恰有一个零点，因此 \(w\in f(\Omega)\)，且 \(g(w)\) 落在圆盘内。又因圆周上 \(f_n\to f\) 一致，充分大的 \(n\) 满足 \(|f_n-f|<\eta/3\)。这时
	\[
		|(f_n-w)-(f-w_0)|<2\eta/3<|f-w_0|
	\]
	在圆周上成立，故 \(f_n-w\) 在圆盘内恰有一个零点。单叶性说明它就是 \(f_n^{-1}(w)\)，于是
	\[
		|f_n^{-1}(w)-g(w)|<2\rho<\varepsilon
	\]
	在整个 \(D(w_0,\eta/3)\) 上成立。

	对 \(K\) 中各点作此构造，取有限个目标圆盘覆盖 \(K\)，再取有限个指标阈值的最大值。由此同时得到 \(K\subset f_n(\Omega)\) 以及 \(K\) 上误差小于 \(\varepsilon\)，证明了所需一致收敛。
\end{solution}

\begin{exercise}
	\label{g4:ex:0101:conformal-radius}
	设 \(\Omega\subsetneq\mathbb C\) 为单连通区域，\(a\in\Omega\)，\(F:\Omega\to\mathbb D\) 为满足 \(F(a)=0\)、\(F'(a)>0\) 的双全纯映射。定义 \(r_\Omega(a)=1/F'(a)\)。证明：若 \(\Omega_1\subset\Omega_2\subsetneq\mathbb C\) 都是单连通区域且 \(a\in\Omega_1\)，则 \(r_{\Omega_1}(a)\le r_{\Omega_2}(a)\)，等号当且仅当两域相同。计算 \(\Omega=D(c,R)\) 时的数值。
\end{exercise}
\begin{solution}
	Riemann 映照定理与归一化唯一性保证定义存在且不依赖选择。取两域的归一化映射 \(F_1,F_2\)。函数 \(H=F_2\circ F_1^{-1}:\mathbb D\to\mathbb D\) 全纯且固定 \(0\)，其中使用包含映射将 \(\Omega_1\) 看作 \(\Omega_2\) 的子域。于是
	\[
		0<H'(0)=\frac{F_2'(a)}{F_1'(a)}
		=\frac{r_{\Omega_1}(a)}{r_{\Omega_2}(a)}\le1.
	\]
	等号时 Schwarz 引理给出 \(H\) 为恒等，因此 \(F_2(\Omega_1)=\mathbb D\)。\(F_2\) 在 \(\Omega_2\) 上单射且已满射到圆盘，故 \(\Omega_1=\Omega_2\)。反向由定义立即成立。

	令 \(\alpha=(a-c)/R\)，则 \(F(z)=\phi_\alpha((z-c)/R)\) 满足所需归一化，且
	\[
		F'(a)=\frac1{R(1-|\alpha|^2)},\qquad
		r_{D(c,R)}(a)=\frac{R^2-|a-c|^2}{R}.
	\]
	这也显示该量依赖基点：在圆心为 \(R\)，靠近圆周时趋于零。
\end{solution}

\begin{exercise}
	\label{g4:ex:0101:transport-metric}
	设 \(\Omega\subsetneq\mathbb C\) 为单连通区域，\(F:\Omega\to\mathbb D\) 为任意双全纯映射。证明
	\[
		\lambda_\Omega(z)=\frac{2|F'(z)|}{1-|F(z)|^2}
	\]
	与 \(F\) 的选择无关。若 \(G:\Omega_1\to\Omega_2\) 是两个这样的区域之间的全纯映射，证明 \(\lambda_{\Omega_2}(G(z))|G'(z)|\le\lambda_{\Omega_1}(z)\)，并说明某点等号的含义。
\end{exercise}
\begin{solution}
	另一双全纯坐标 \(\widetilde F\) 与 \(F\) 相差圆盘自同构 \(A=\widetilde F\circ F^{-1}\)。\cref{g4:thm:0103:schwarz-pick}的等号给出
	\[
		\frac{|A'(w)|}{1-|A(w)|^2}=\frac1{1-|w|^2}.
	\]
	令 \(w=F(z)\) 并乘以 \(2|F'(z)|\)，即得选择无关性。

	取两域的圆盘坐标 \(F_1,F_2\)，对 \(F_2\circ G\circ F_1^{-1}\) 用 Schwarz--Pick 导数不等式，再按链式法则展开，得到所求式。若某点取等号，该圆盘自映射便是自同构，所以 \(G\) 是两域之间的双全纯映射。反之双全纯映射使该不等式处处为等号。
\end{solution}

\begin{exercise}
	\label{g4:ex:0101:boundary-escape}
	设 \(\Omega\subsetneq\mathbb C\) 是单连通区域，以前一题的密度 \(\lambda_\Omega\) 定义曲线长度及距离 \(d_\Omega\)。证明任意双全纯映射 \(F:\Omega\to\mathbb D\) 都是这两个长度距离之间的等距映射，并证明 \((\Omega,d_\Omega)\) 完备。对固定 \(a\in\Omega\)，证明 \(d_\Omega(a,z_j)\to\infty\) 当且仅当点列 \(z_j\) 最终离开 \(\Omega\) 的每个紧集。
\end{exercise}
\begin{solution}
	由密度的定义及链式法则，每条分片 \(C^1\) 曲线满足 \(L_\Omega(\gamma)=L_{\mathbb D}(F\circ\gamma)\)。\(F\) 与其逆把连接两点的曲线一一对应，取下确界得
	\[
		d_\Omega(z,w)=d_{\mathbb D}(F(z),F(w)).
	\]
	所以只需在圆盘上证明完备性。设 \((w_j)\) 是双曲距离的 Cauchy 列。由 \(d_{\mathbb D}(z,w)\ge2|z-w|\)，它在平面中也是 Cauchy 列，趋于某个 \(w\in\overline{\mathbb D}\)。另一方面，Cauchy 条件及三角不等式使 \(d_{\mathbb D}(0,w_j)\) 有界。公式\eqref{g4:eq:0103:hyperbolic-distance}给出
	\[
		d_{\mathbb D}(0,w_j)=\log\frac{1+|w_j|}{1-|w_j|},
	\]
	所以 \(|w_j|\) 有一个严格小于 \(1\) 的统一上界，极限 \(w\) 在圆盘内。在 \(w\) 的一个小闭圆盘内，长度密度有界；用连接 \(w_j\) 与 \(w\) 的直线段估计长度，便有 \(d_{\mathbb D}(w_j,w)\to0\)。这证明完备性。

	最后选择满足 \(F(a)=0\) 的圆盘坐标。上式说明 \(d_\Omega(a,z_j)\to\infty\) 等价于 \(|F(z_j)|\to1\)，后者由\cref{g4:prop:0104:boundary-escape}等价于离开源域的每个紧集。通常平面距离有限的边界，在这个长度结构下也位于无穷远处。
\end{solution}

% !TeX root = ../../Book4.tex
% 纲要标题：## 第2章　黎曼曲面、覆盖与分歧
% 纲要位置：几何学计划纲要.md，第 3294 行。
% 正文状态：未撰写；本文件目前是框架。

% 纲要内容开始
% 
% > 本章目标与证明范围：从复图册建立一维复几何，完整证明非恒定紧曲面映射的有限性、局部分歧模型和 Riemann–Hurwitz。亚纯函数、微分和拓扑覆盖的关系同时展开。
% 
% 纲要内容结束

\chapter{黎曼曲面、覆盖与分歧}
\label{g4:ch:02}
\BookChapterNumber{g4:ch:02}

\begin{chapterabstract}
待撰写
\end{chapterabstract}

\begin{learninggoals}
待撰写
\end{learninggoals}

待撰写

% !TeX root = ../../Book4.tex
% 纲要标题：### 2.1 黎曼曲面的构造
% 纲要位置：几何学计划纲要.md，第 3298 行。
% 正文状态：未撰写；本文件目前是框架。

% 纲要内容开始
% 
% * 复图册及全纯映射
% * 球面、复环面和代数曲线模型
% * 解析延拓的曲面
% * 覆盖空间上复结构的提升
% 
% 纲要内容结束

\section{黎曼曲面的构造}
\label{g4:sec:0201}

待撰写

\subsection{复图册及全纯映射}
\label{g4:subsec:0201:01}
待撰写

\subsection{球面、复环面和代数曲线模型}
\label{g4:subsec:0201:02}
待撰写

\subsection{解析延拓的曲面}
\label{g4:subsec:0201:03}
待撰写

\subsection{覆盖空间上复结构的提升}
\label{g4:subsec:0201:04}
待撰写
% 纲要条目排为编号小节标题，并进入目录与书签；说明段仍保留于纲要注释中。

% !TeX root = ../../Book4.tex
% 纲要标题：### 2.2 分歧与有限映射
% 纲要位置：几何学计划纲要.md，第 3522 行。
% 正文状态：未撰写；本文件目前是框架。

% 纲要内容开始
% 
% * 局部幂模型及分歧指数
% * 紧性、开映射与满射
% * 次数和纤维重数
% * 去掉分歧点后的覆盖及单值
% 
% 纲要内容结束

\section{分歧与有限映射}
\label{g4:sec:0202}

待撰写

\subsection{局部幂模型及分歧指数}
\label{g4:subsec:0202:01}
待撰写

\subsection{紧性、开映射与满射}
\label{g4:subsec:0202:02}
待撰写

\subsection{次数和纤维重数}
\label{g4:subsec:0202:03}
待撰写

\subsection{去掉分歧点后的覆盖及单值}
\label{g4:subsec:0202:04}
待撰写
% 纲要条目排为编号小节标题，并进入目录与书签；说明段仍保留于纲要注释中。

% !TeX root = ../../Book4.tex
% 纲要标题：### 2.3 亚纯微分和 Riemann–Hurwitz
% 纲要位置：几何学计划纲要.md，第 3529 行。
% 正文状态：初稿；尚未正式审核。

% 纲要内容开始
% 
% * 微分在坐标变换下的阶数
% * 留数及整体留数定理
% * 典范除子与微分拉回
% * 分歧除子和 Hurwitz 公式的证明
% 
% 纲要内容结束

\section{亚纯微分和 Riemann–Hurwitz}
\label{g4:sec:0203}

局部模型 \(w=z^e\) 同时改变两种数据。一个普通靶点有 \(e\) 个邻近原像，而微分满足 \(\mathrm{d}w=e z^{e-1}\,\mathrm{d}z\)，在源点多出 \(e-1\) 阶零点。前者影响覆盖的拓扑计数，后者影响微分的零极点。这一节把两种计数放在一起：先使微分的阶数和留数成为坐标无关的量，再证明分歧对典范线丛及曲面亏格的修正公式。

\subsection{微分在坐标变换下的阶数}
\label{g4:subsec:0203:01}
在平面上，积分的对象是 \(a(z)\,\mathrm{d}z\)。若换坐标 \(z=z(\zeta)\)，链式法则要求同时改变函数系数和微分因子：
\[
	a(z)\,\mathrm{d}z
	=a(z(\zeta))\frac{\mathrm{d}z}{\mathrm{d}\zeta}\,\mathrm{d}\zeta.
\]
因此，曲面上的微分应由一族相容的局部表达式描述，单独的系数函数通常不能粘成全局函数。

\begin{definition}
	\label{g4:def:0203:meromorphic-differential}
	设 \(X\) 是黎曼曲面。在每张全纯坐标图 \(z:U\rightarrow\mathbb C\) 中给定表达式 \(\omega=a_z(z)\,\mathrm{d}z\)，其中 \(a_z\) 亚纯，并要求交叠图上的表达式满足
	\begin{equation}
		\label{g4:eq:0203:differential-transition}
		a_\zeta(\zeta)=a_z(z(\zeta))\frac{\mathrm{d}z}{\mathrm{d}\zeta}.
	\end{equation}
	这样的一族局部表达式称为 \(X\) 上的\term[yachunweifen]{亚纯微分}[meromorphic differential]，也称亚纯一形式。若每个局部系数均全纯，则称为\term[quanchunweifen]{全纯微分}[holomorphic differential]。
\end{definition}

坐标变换的导数全纯且处处非零，所以一个图中的系数全纯，当且仅当另一个图中的系数全纯。在没有极点处，这些表达式给出实曲面上的复值光滑一形式。这里选择的形式具有 \(a(z)\,\mathrm{d}z\) 的局部形状；一般复值光滑一形式还可以有 \(\mathrm{d}\overline z\) 项。

\begin{proposition}
	\label{g4:prop:0203:order-invariance}
	设 \(X\) 是连通黎曼曲面，\(\omega\) 是不恒为零的亚纯微分，\(p\in X\)。在满足 \(z(p)=0\) 的全纯坐标中，写
	\[
		\omega=z^m u(z)\,\mathrm{d}z,\qquad m\in\mathbb Z,\quad u(0)\ne0,
	\]
	其中 \(u\) 在原点附近全纯。则整数 \(m\) 与坐标选择无关，称为 \(\omega\) 在 \(p\) 的阶数，记为 \(\operatorname{ord}_p\omega\)。阶数为正时，\(p\) 是相应阶数的零点；阶数为负时，\(p\) 是 \(-m\) 阶极点。
\end{proposition}

\begin{proof}
	若一个图中的系数在开集上恒零，式\eqref{g4:eq:0203:differential-transition}与恒等定理会把恒零性沿重叠图传遍连通曲面。因此，\(\omega\not\equiv0\) 保证每个局部系数都不恒零，可使用\cref{g4:def:0102:meromorphic}后的零极点分解。

	换另一中心坐标 \(\zeta\)，则 \(z=\zeta v(\zeta)\)，其中 \(v(0)\ne0\)，而 \(\mathrm{d}z/\mathrm{d}\zeta\) 在原点也非零。于是新系数为
	\[
		\zeta^m v(\zeta)^m u(z(\zeta))
		\frac{\mathrm{d}z}{\mathrm{d}\zeta}.
	\]
	除 \(\zeta^m\) 外，各因子都在原点全纯且非零，包括 \(m<0\) 时的 \(v^m\)。故新坐标中的阶数仍为 \(m\)。
\end{proof}

零点与极点因局部分解而离散，在紧曲面上只有有限个。微分因子在合法的坐标变换下不改变阶数；在分歧映射的拉回中，它却可能产生零点，因为 \(w=z^e\) 在 \(e>1\) 时不是坐标变换。

\begin{example}
	\label{g4:exm:0203:sphere-differential}
	在球面 \(\mathbb P^1\) 的有限坐标 \(z\) 中取 \(\omega=\mathrm{d}z\)。无穷远坐标为 \(\zeta=1/z\)，所以
	\[
		\omega=-\zeta^{-2}\,\mathrm{d}\zeta.
	\]
	它在有限点没有零极点，在无穷远有二阶极点。因此“系数等于 \(1\)”只描述一个仿射图，不能用来判断整个球面上的正则性。更一般地，
	\[
		z^m\,\mathrm{d}z=-\zeta^{-m-2}\,\mathrm{d}\zeta
	\]
	给出任意整数 \(m\) 时的无穷远阶数。
\end{example}

\begin{example}
	\label{g4:exm:0203:torus-differential}
	设 \(X=\mathbb C/\Lambda\) 是复环面，\(\Lambda\) 为格。商映射的局部逆之间相差平移 \(z\mapsto z+\lambda\)，而 \(\mathrm{d}(z+\lambda)=\mathrm{d}z\)。因此平面上的 \(\mathrm{d}z\) 下降为 \(X\) 上处处非零的全纯微分。这里坐标 \(z\) 本身不能在环面上单值选取，其微分却可以。
\end{example}

\subsection{留数及整体留数定理}
\label{g4:subsec:0203:02}
阶数记录最先出现的 Laurent 项，留数则选出积分绕行一周后留下的那一项。对避开极点的分片 \(C^1\) 道路 \(\gamma\)，把参数区间分成有限段，使每段位于一张图内，并逐段规定
\[
	\int_\gamma\omega
	=\sum_j\int_{t_{j-1}}^{t_j}
	a_{z_j}(z_j(\gamma(t)))\,(z_j\circ\gamma)'(t)\,\mathrm{d}t.
\]
细分不改变和式；式\eqref{g4:eq:0203:differential-transition}与一元换元公式保证换图也不改变积分。

\begin{proposition}
	\label{g4:prop:0203:residue-invariance}
	设 \(\omega\) 是黎曼曲面 \(X\) 上的亚纯微分，\(p\in X\)。在中心坐标 \(z(p)=0\) 中展开
	\[
		\omega=\left(\sum_{n=-N}^{\infty}c_nz^n\right)\mathrm{d}z.
	\]
	系数 \(c_{-1}\) 与坐标选择无关，记为 \(\operatorname{Res}_p\omega\)，称为 \(\omega\) 在 \(p\) 的留数。充分小的正向坐标圆周 \(\gamma\) 满足
	\[
		\int_\gamma\omega=2\pi\mathrm{i}\operatorname{Res}_p\omega.
	\]
	全纯点的留数为零。
\end{proposition}

\begin{proof}
	积分公式由\cref{g4:thm:0102:laurent}的系数公式得到。为核对坐标不变性，将有限主部逐项处理。对 \(n\ne-1\)，
	\[
		z^n\,\mathrm{d}z=\mathrm{d}\left(\frac{z^{n+1}}{n+1}\right).
	\]
	亚纯函数的微分在任何坐标中的 Laurent 展开都没有 \(\zeta^{-1}\,\mathrm{d}\zeta\) 项：逐项求导时，能产生指数 \(-1\) 的原项指数只能是 \(0\)，其导数为零。全纯余项换坐标后仍全纯。最后，对中心坐标变换 \(z=\zeta v(\zeta)\)、\(v(0)\ne0\)，有
	\[
		\frac{\mathrm{d}z}{z}
		=\frac{\mathrm{d}\zeta}{\zeta}
		+\frac{v'(\zeta)}{v(\zeta)}\,\mathrm{d}\zeta.
	\]
	后一项全纯，所以 \(z^{-1}\,\mathrm{d}z\) 在新坐标中的留数仍为 \(1\)。由线性性得到结论。
\end{proof}

在有限平面区域中，留数之和等于外边界积分；紧无边界曲面没有留下这样的外边界，全部局部留数因而必须相互抵消。

\begin{theorem}[整体留数定理]
	\label{g4:thm:0203:global-residue}
	设 \(X\) 是紧连通黎曼曲面，\(\omega\) 是 \(X\) 上的亚纯微分。则
	\[
		\sum_{p\in X}\operatorname{Res}_p\omega=0.
	\]
	该和只有有限个可能非零的项。
\end{theorem}

这个结论称为\term[zhengtiliushudingli]{整体留数定理}[global residue theorem]。证明使用第一卷第14.4节定理14.4.3 的 Stokes 公式：对紧有向光滑实曲面 \(M\) 及光滑至边界的一形式 \(\alpha\)，有 \(\int_{\partial M}\alpha=\int_M\mathrm{d}\alpha\)，边界按曲面位于左侧定向。复值形式分别对实部、虚部应用这一公式即可；紧性自动满足该定理的紧支条件。

\begin{proof}
	亚纯微分的极点在紧曲面上离散，因而只有有限个，记为 \(p_1,\ldots,p_s\)。为每个极点取一个中心坐标圆盘 \(D_j\)，使其闭包两两不交，盘内没有其他极点。删去这些开盘，得到紧的带光滑边界曲面
	\[
		M=X\setminus\bigcup_{j=1}^s D_j.
	\]
	复坐标的变换保持定向，给 \(M\) 取由 \(X\) 继承的定向。形式 \(\omega\) 在 \(M\) 的邻域全纯，因此光滑至边界。在任一局部坐标中，
	\[
		\mathrm{d}\omega
		=\mathrm{d}(a(z)\,\mathrm{d}z)
		=a'(z)\,\mathrm{d}z\wedge\mathrm{d}z=0.
	\]
	这里全纯性给出 \(\mathrm{d}a=a'(z)\,\mathrm{d}z\)。Stokes 公式于是给出 \(\int_{\partial M}\omega=0\)。每条孔口边界作为 \(M\) 的边界按顺时针方向走，而计算留数的小圆按逆时针方向走，所以
	\[
		0=-\sum_{j=1}^s\int_{\partial D_j}\omega
		=-2\pi\mathrm{i}\sum_{j=1}^s\operatorname{Res}_{p_j}\omega.
	\]
	其余点的留数为零，得到所需等式。若没有极点，则各项本来就为零。
\end{proof}

在 \(\mathbb P^1\) 上，\(\mathrm{d}z/z\) 在 \(0\) 的留数为 \(1\)，而在 \(\zeta=1/z\) 坐标中等于 \(-\mathrm{d}\zeta/\zeta\)，所以无穷远的留数为 \(-1\)。紧曲面上的亚纯微分不能只有一个简单极点，因为该点留数必非零；只有一个高阶极点则可能成立，\cref{g4:exm:0203:sphere-differential}中的 \(\mathrm{d}z\) 就只有一个二阶极点，其留数为零。紧性也确实参与结论：在非紧曲面 \(\mathbb C\) 上，\(\mathrm{d}z/z\) 只有原点这一个极点，留数之和为 \(1\)。

\subsection{典范除子与微分拉回}
\label{g4:subsec:0203:03}
若把一个微分乘以亚纯函数，它的零极点阶数就逐点相加。用有限整数和记录这些阶数，可以把这类变换写成加法公式。

\begin{definition}
	\label{g4:def:0203:divisor}
	设 \(X\) 是紧连通黎曼曲面。有限形式和
	\[
		D=\sum_{p\in X}n_p[p],\qquad n_p\in\mathbb Z,
	\]
	称为 \(X\) 上的\term[chuzi]{除子}[divisor]。它的次数为 \(\deg D=\sum_p n_p\)；若所有 \(n_p\ge0\)，则称 \(D\) 为\term[youxiaochuzi]{有效除子}[effective divisor]。不恒为零的亚纯函数 \(h\) 确定\term[zhuchuzi]{主除子}[principal divisor]
	\[
		\operatorname{div}(h)=\sum_p\operatorname{ord}_p(h)[p].
	\]
	两个除子之差为主除子时，称它们\term[xianxingdengjia]{线性等价}[linearly equivalent]。
\end{definition}

由阶数的乘除法则，\(\operatorname{div}(h_1h_2)=\operatorname{div}(h_1)+\operatorname{div}(h_2)\)、\(\operatorname{div}(1/h)=-\operatorname{div}(h)\)，所以线性等价确为等价关系。

\begin{proposition}
	\label{g4:prop:0203:principal-degree}
	紧连通黎曼曲面 \(X\) 上的每个主除子次数均为零。因此线性等价的除子具有相同次数。
\end{proposition}

\begin{proof}
	对非零常数 \(h\)，主除子为零。若 \(h\) 非常值，\cref{g4:prop:0201:meromorphic-map}把它视为全纯映射 \(h:X\rightarrow\mathbb P^1\)：在极点处使用靶坐标 \(1/h\)，其零点阶数就是原来的极点阶数。由\cref{g4:thm:0202:degree}，\(0\) 与 \(\infty\) 两条纤维的重数和都等于同一个次数 \(d\)。于是
	\[
		\deg\operatorname{div}(h)
		=\sum_{h(p)=0}e_p(h)-\sum_{h(p)=\infty}e_p(h)
		=d-d=0.
	\]
	最后对两个线性等价除子的差使用这一结论。
\end{proof}

对一条已给定的不恒为零的亚纯微分 \(\omega\)，同样定义
\[
	\operatorname{div}(\omega)=\sum_p\operatorname{ord}_p(\omega)[p],
\]
称为它的\term[dianfanchuzi]{典范除子}[canonical divisor]。若 \(\eta\) 是另一条非零亚纯微分，则局部商 \(a_z/b_z\) 在换坐标时消去同一个导数因子，粘成全局非零亚纯函数 \(h=\omega/\eta\)，并且
\begin{equation}
	\label{g4:eq:0203:canonical-linear-equivalence}
	\operatorname{div}(\omega)-\operatorname{div}(\eta)
	=\operatorname{div}(h).
\end{equation}
因此这些典范除子彼此线性等价，次数也相同。一般紧曲面上非零亚纯微分的存在将在\secref{g4:sec:0304}构造亚纯截面时证明；当前的定义和比较适用于已经给出的微分。

\ProseParagraphBreak

微分的局部标架本身总可以构造。对全纯坐标 \(z\)，取一维复向量空间的局部标架 \(\mathrm{d}z\)，并按
\[
	\mathrm{d}z=\frac{\mathrm{d}z}{\mathrm{d}\zeta}\,\mathrm{d}\zeta
\]
粘合各个局部乘积 \(U\times\mathbb C\)。导数非零保证过渡因子可逆，链式法则保证三重交叠处的相容性。这给出 \(X\) 的\term[dianfanxiancong]{典范线丛}[canonical line bundle]，记为 \(K_X\)。这里全纯线丛指局部与 \(U\times\mathbb C\) 相同、过渡函数为处处非零全纯函数的一维复向量丛；典范线丛的局部全纯或亚纯截面分别就是 \(a(z)\,\mathrm{d}z\) 这样的表达式。一般线丛和除子的对应将在下一章系统建立。
\symidx{\(K_X\)：黎曼曲面的典范线丛}

\begin{proposition}
	\label{g4:prop:0203:pullback-order}
	设 \(F:X\rightarrow Y\) 是连通黎曼曲面之间的非恒定全纯映射，\(\eta\) 是 \(Y\) 上不恒为零的亚纯微分。对 \(p\in X\)，记 \(q=F(p)\)、\(e=e_p(F)\)。则拉回 \(F^*\eta\) 是 \(X\) 上不恒为零的亚纯微分，并且
	\begin{equation}
		\label{g4:eq:0203:pullback-order}
		\operatorname{ord}_p(F^*\eta)
		=e_p(F)\operatorname{ord}_{F(p)}\eta+e_p(F)-1.
	\end{equation}
	留数满足
	\[
		\operatorname{Res}_p(F^*\eta)
		=e_p(F)\operatorname{Res}_{F(p)}\eta.
	\]
\end{proposition}

\begin{proof}
	在局部坐标 \(z,w\) 中，若 \(\eta=a(w)\,\mathrm{d}w\)，则规定
	\[
		F^*\eta=a(w(F(z)))\frac{\mathrm{d}(w\circ F)}{\mathrm{d}z}\,\mathrm{d}z.
	\]
	链式法则使这一表达式与源、靶坐标选择相容。由\cref{g4:thm:0202:local-model}，可取中心坐标使 \(w\circ F=z^e\)。写 \(\eta=w^m u(w)\,\mathrm{d}w\)，其中 \(u(0)\ne0\)，便有
	\[
		F^*\eta=e z^{em+e-1}u(z^e)\,\mathrm{d}z.
	\]
	系数中 \(e u(0)\ne0\)，这证明拉回不局部恒零，并给出阶数公式。对留数，在同样坐标中写 \(\eta=\sum_n c_nw^n\,\mathrm{d}w\)，得到
	\[
		F^*\eta=\sum_n e c_n z^{en+e-1}\,\mathrm{d}z.
	\]
	指数等于 \(-1\) 当且仅当 \(n=-1\)，所以留数为 \(e c_{-1}\)。
\end{proof}

公式中的 \(e-1\) 完全来自微分因子。若 \(\eta\) 在 \(q\) 全纯且非零，则 \(F^*\eta\) 在 \(p\) 恰有 \(e-1\) 阶零点；若 \(\eta\) 在 \(q\) 有简单极点，\(m=-1\)，拉回的阶数仍为 \(-1\)，但留数乘以 \(e\)。例如 \(\mathrm{d}w/w\) 在 \(w=z^e\) 下变成 \(e\,\mathrm{d}z/z\)。

\subsection{分歧除子和 Hurwitz 公式的证明}
\label{g4:subsec:0203:04}
对于紧曲面间的非恒定映射，分歧点只有有限个。因此可以把\cref{g4:prop:0203:pullback-order}中额外出现的 \(e_p-1\) 同时记录为一个除子。以下始终把 \(X,Y\) 取为紧连通黎曼曲面，并令 \(F:X\rightarrow Y\) 非恒定，\(d=\deg F\)。

\begin{definition}
	\label{g4:def:0203:ramification-divisor}
	设 \(X,Y\) 是紧连通黎曼曲面，\(F:X\rightarrow Y\) 是非恒定全纯映射。其\term[fenqichuzi]{分歧除子}[ramification divisor]为
	\[
		R_F\coloneqq\sum_{p\in X}(e_p(F)-1)[p].
	\]
	它是有效除子。对 \(Y\) 上的除子 \(D=\sum_q n_q[q]\)，定义
	\[
		F^*D\coloneqq\sum_q n_q\sum_{F(p)=q}e_p(F)[p].
	\]
	\symidx{\(R_F\)：全纯映射的分歧除子}
\end{definition}

纤维重数公式直接给出
\begin{equation}
	\label{g4:eq:0203:divisor-degree-pullback}
	\deg(F^*D)=d\,\deg D.
\end{equation}
若 \(h\) 是 \(Y\) 上非零亚纯函数，局部幂模型还给出 \(F^*\operatorname{div}(h)=\operatorname{div}(h\circ F)\)，所以除子拉回保持线性等价。

\begin{theorem}[除子形式的 Hurwitz 公式]
	\label{g4:thm:0203:divisor-hurwitz}
	设 \(F:X\rightarrow Y\) 是紧连通黎曼曲面之间的非恒定全纯映射。对 \(Y\) 上每条不恒为零的亚纯微分 \(\eta\)，有
	\begin{equation}
		\label{g4:eq:0203:divisor-hurwitz}
		\operatorname{div}(F^*\eta)
		=F^*\operatorname{div}(\eta)+R_F.
	\end{equation}
	在典范线丛之间，微分给出一个全纯丛映射
	\[
		(\mathrm{d}F)^*:F^*K_Y\longrightarrow K_X,
	\]
	其零点除子恰为 \(R_F\)。相应地有全纯线丛同构
	\begin{equation}
		\label{g4:eq:0203:canonical-bundle-hurwitz}
		K_X\cong F^*K_Y\otimes\mathcal O_X(R_F),
	\end{equation}
	其中 \(\mathcal O_X(R_F)\) 的局部标架在 \(p\) 的中心坐标中为 \(z^{-(e_p(F)-1)}\)。
\end{theorem}

先说明末式涉及的局部对象。在 \(F^*K_Y\) 中，点 \(p\) 处的纤维取为 \(K_Y\) 在 \(F(p)\) 的纤维，靶坐标的 \(\mathrm{d}w\) 给出一个处处非零的抽象标架，记为 \(\widehat{\mathrm{d}w}\)。它经 \((\mathrm{d}F)^*\) 后才成为实际一形式 \(F^*(\mathrm{d}w)\)，后者在分歧点可以为零。对有效除子 \(R=\sum r_p[p]\)，在各支集点附近选两两不交的中心坐标圆盘，并取标架 \(z^{-r_p}\)；在支集之外取标架 \(1\)，便定义了 \(\mathcal O_X(R)\)。交叠处的比值全纯且非零；换中心坐标只使标架乘以一个全纯单位。它的局部全纯截面以亚纯函数表示时，允许在 \(p\) 有至多 \(r_p\) 阶极点。

\begin{proof}
	式\eqref{g4:eq:0203:divisor-hurwitz}逐点比较系数，就是\cref{g4:prop:0203:pullback-order}。

	对线丛陈述，在归一化坐标 \(w\circ F=z^e\) 中，微分的局部表达式为
	\[
		\widehat{\mathrm{d}w}\longmapsto e z^{e-1}\,\mathrm{d}z.
	\]
	换坐标后源、靶标架都乘以全纯单位，链式法则保证这些局部映射粘合。其系数的零点阶数为 \(e-1\)，所以零点除子是 \(R_F\)。在张量积中乘上 \(\mathcal O_X(R_F)\) 的局部标架后，得到
	\[
		\widehat{\mathrm{d}w}\otimes z^{-(e-1)}
		\longmapsto e\,\mathrm{d}z.
	\]
	右边处处非零，因此这给出全纯局部同构。各式在分歧点之外都是原来微分映射乘以相应的亚纯函数，故在交叠处相合；经过分歧点的等式由全纯延拓保持。于是它们粘合成式\eqref{g4:eq:0203:canonical-bundle-hurwitz}。
\end{proof}

微分的局部零点修正还具有一个拓扑解释。每个普通靶点按 \(d\) 张覆盖片计数；一个分支值的纤维中，若某个点的指数为 \(e\)，则 \(e\) 张覆盖片在该点合并，只留下一个顶点，少计 \(e-1\)。把这项损失累加，就得到亏格公式。

\ProseParagraphBreak

所需的曲面拓扑结果来自第一卷第6章。定理6.2.10 给出紧 Hausdorff 第二可数曲面（允许边界）的有限三角剖分；命题6.2.14 证明顶点、边、面数的交错和 \(V-E+F\) 与剖分无关，称为\term[Eulershixingshu]{Euler 示性数}[Euler characteristic]，记为 \(\chi\)。第一卷定理6.5.4 与命题6.5.2 进一步说明，紧连通可定向无边界曲面同胚于球面添 \(g\) 个手柄所得的曲面，其中整数 \(g\ge0\) 唯一，称为\term[kuige]{亏格}[genus]，且 \(\chi=2-2g\)。黎曼曲面由复图册带有自然定向，所以这些结论适用。

\begin{theorem}[Riemann--Hurwitz 公式]
	\label{g4:thm:0203:riemann-hurwitz}
	设 \(X,Y\) 是紧连通黎曼曲面，\(F:X\rightarrow Y\) 为次数 \(d\ge1\) 的非恒定全纯映射。则
	\begin{equation}
		\label{g4:eq:0203:hurwitz-euler}
		\chi(X)=d\,\chi(Y)-\sum_{p\in X}(e_p(F)-1),
	\end{equation}
	等价地，
	\begin{equation}
		\label{g4:eq:0203:riemann-hurwitz}
		2g(X)-2=d\bigl(2g(Y)-2\bigr)+\deg R_F.
	\end{equation}
\end{theorem}

这是\term[RiemannHurwitzgongshi]{Riemann--Hurwitz 公式}[Riemann--Hurwitz formula]。证明先删去分支值周围的小盘，在剩余部分按真正的覆盖计算，再把源、靶的圆盘分别封回。这样，每个分歧点的计数损失都对应一个具体的封盘数。

\begin{proof}
	设分支值集合为 \(B=\{q_1,\ldots,q_s\}\)。由\cref{g4:lem:0202:isolated-fiber}，每个 \(q_j\) 有中心坐标邻域，使其全部逆像分解为纤维中各点的互不相交坐标盘，在这些盘上 \(F\) 写成 \(w=z^{e_p}\)。把靶盘进一步缩小，可取两两不交的闭坐标盘 \(\overline{D_j}\)，使它们各只含一个分支值，且其全部逆像恰为这些源盘内的闭圆盘。若 \(D_j\) 在 \(w\) 坐标中半径为 \(\rho_j\)，则相应源盘在 \(z\) 坐标中的半径为 \(\rho_j^{1/e_p}\)。特别地，每个纤维点贡献一个源盘，包含未分歧的纤维点。

	记 \(r_j=\#F^{-1}(q_j)\)、\(N=\sum_jr_j\)，并设
	\[
		Y_0=Y\setminus\bigcup_jD_j,\qquad
		X_0=X\setminus F^{-1}\left(\bigcup_jD_j\right).
	\]
	它们是带光滑边界的紧曲面；允许先把各连通分支分别计数。所有分歧点已被删去，\cref{g4:thm:0202:unbranched-cover}使 \(F|_{X_0}:X_0\rightarrow Y_0\) 成为 \(d\) 叶覆盖。边界点也满足覆盖条件：沿靶圆周取足够小的半邻域，其各逆分支把这一个半邻域同胚地提升到源圆周旁边。

	恢复第一卷命题6.5.3 所用的计数步骤。先对 \(Y_0\) 取有限三角剖分，使边界为子复形。均匀覆盖邻域给出 \(Y_0\) 的开覆盖。在有限复形的一个度量中，紧性给出 \(\delta>0\)，使每个直径小于 \(\delta\) 的集合都落在某个均匀覆盖邻域中。充分多次重心细分后，每个闭单形的直径均小于 \(\delta\)，因而每个闭单形都落在一个均匀覆盖邻域中。于是每个顶点、开边、开三角形都恰有 \(d\) 个提升；闭三角形的提升与它的各条边、各个顶点相容，组成 \(X_0\) 的有限三角剖分。因此
	\[
		\chi(X_0)=d\,\chi(Y_0).
	\]
	这个计数对有边界或不连通的紧曲面同样成立，因为所有分支的顶点、边、面数相加即可。

	再计算封盘的变化。每个边界圆周先取 \(m\ge3\) 个顶点与 \(m\) 条边，给封入的圆盘添一个中心，连到这 \(m\) 个顶点。新增一个顶点、\(m\) 条边和 \(m\) 个三角形，故交错和增加 \(1-m+m=1\)。所以
	\[
		\chi(Y_0)=\chi(Y)-s,\qquad
		\chi(X_0)=\chi(X)-N.
	\]
	结合覆盖计数得
	\[
		\chi(X)=d\chi(Y)-ds+N.
	\]
	最后由\cref{g4:thm:0202:degree}，对每个 \(q_j\) 有 \(\sum_{F(p)=q_j}e_p=d\)，因而
	\[
		ds-N
		=\sum_{j=1}^s\sum_{F(p)=q_j}(e_p-1)
		=\sum_{p\in X}(e_p-1).
	\]
	所有实际分歧点都出现在这些纤维中，其他点的贡献为零，故得到式\eqref{g4:eq:0203:hurwitz-euler}。若 \(B\) 为空，直接令 \(X_0=X\)、\(Y_0=Y\)、\(s=N=0\)，同一覆盖计数给出结论。代入 \(\chi(X)=2-2g(X)\)、\(\chi(Y)=2-2g(Y)\)，整理即得式\eqref{g4:eq:0203:riemann-hurwitz}。
\end{proof}

\begin{example}
	\label{g4:exm:0203:power-map-hurwitz}
	对球面映射 \(F:\mathbb P^1\rightarrow\mathbb P^1\)、\(F(z)=z^d\)，其中 \(d\ge1\)，原点与无穷远的局部指数都为 \(d\)，其余点为 \(1\)。所以
	\[
		R_F=(d-1)[0]+(d-1)[\infty],\qquad \deg R_F=2d-2.
	\]
	Riemann--Hurwitz 公式成为 \(-2=-2d+(2d-2)\)。另一方面，靶球面的 \(\eta=\mathrm{d}w\) 满足 \(\operatorname{div}(\eta)=-2[\infty]\)，故除子公式给出
	\[
		\operatorname{div}(F^*\mathrm{d}w)
		=-2d[\infty]+R_F
		=(d-1)[0]-(d+1)[\infty].
	\]
	这与直接计算 \(F^*\mathrm{d}w=d z^{d-1}\,\mathrm{d}z\) 完全一致，尤其无穷远的微分因子贡献了额外的二阶极点。
\end{example}

有效性 \(\deg R_F\ge0\) 给出若干直接限制。球面不能非恒定地映到正亏格紧曲面：若 \(g(X)=0\)、\(g(Y)\ge1\)，式\eqref{g4:eq:0203:riemann-hurwitz}左边为 \(-2\)，右边非负，矛盾。两个亏格为 \(1\) 的曲面之间若有非恒定全纯映射，它必无分歧，因为公式迫使 \(\deg R_F=0\)。若 \(F:X\rightarrow\mathbb P^1\) 的次数为 \(2\)，则每个分歧点的指数只能为 \(2\)，每个分支值也只能有这一个重数为 \(2\) 的原像；于是分支值的个数恰为 \(2g(X)+2\)。下一节的显式方程将把这些整数限制实现为具体的曲面。
% 纲要条目排为编号小节标题，并进入目录与书签；说明段仍保留于纲要注释中。

% !TeX root = ../../Book4.tex
% 纲要标题：### 2.4 显式覆盖的几何
% 纲要位置：几何学计划纲要.md，第 3536 行。
% 正文状态：初稿；尚未正式审核。

% 纲要内容开始
% 
% * 超椭圆曲面的分支粘合
% * 三次曲线及亏格
% * 拓扑单值与代数方程的对应
% * 不可约、连通与分支数据
% 
% ---
% 
% 纲要内容结束

\section{显式覆盖的几何}
\label{g4:sec:0204}

方程 \(y^2=P(z)\) 在一个普通的 \(z\) 值上有两个根。沿闭路继续这两个根，返回时可能交换；在 \(P\) 的零点处，两根则合并为一个分歧点。把这些现象组织成一张曲面，需要同时回答三个问题：局部怎样给合并点取坐标，无穷远处应当补几个点，以及所有局部分支究竟连成几张连通曲面。下面先在二次方程中逐项完成构造，再说明一般代数方程中的根置换如何回答连通性问题。

\subsection{超椭圆曲面的分支粘合}
\label{g4:subsec:0204:01}

设
\[
	P(z)=a_d\prod_{j=1}^d(z-a_j),\qquad a_d\ne0,\quad d\ge1,
\]
其中 \(a_1,\ldots,a_d\) 两两不同。先取仿射方程的解集
\[
	X_P^\circ=\{(z,y)\in\mathbb C^2:y^2=P(z)\},
	\qquad \pi(z,y)=z.
\]
在 \(P(z)\ne0\) 处，可以在小圆盘内选取平方根，解集分成两个全纯图像。根相遇的地方要换用 \(y\) 作局部参数；无穷远则必须根据 \(P\) 的次数另取坐标。

\begin{proposition}\label{g4:prop:0204:hyperelliptic-completion}
	设 \(P\in\mathbb C[z]\) 为次数 \(d\ge1\) 的无重根多项式。上述 \(X_P^\circ\) 是连通黎曼曲面，并可补成紧黎曼曲面 \(X_P\)，使 \(\pi\) 延拓为二次全纯映射
	\[
		\pi:X_P\longrightarrow\mathbb P^1.
	\]
	每个有限零点 \(a_j\) 上方有一个分歧指数为 \(2\) 的点。当 \(d\) 为偶数时，无穷远上方有两个无分歧点；当 \(d\) 为奇数时，无穷远上方只有一个点，分歧指数为 \(2\)。其余各点均无分歧。
\end{proposition}
\begin{proof}
	先构造有限处的坐标。若 \(P(z_0)\ne0\)，在 \(z_0\) 的小圆盘上选取全纯平方根 \(s(z)\)，两张图像 \(y=\pm s(z)\) 都以 \(z\) 为坐标，投影没有分歧。若 \(P(a)=0\)，则 \(P'(a)\ne0\)，由一元全纯逆函数定理，\(P\) 在 \(a\) 附近有全纯逆函数 \(\psi\)。因此原方程在 \((a,0)\) 附近写成
	\[
		z=\psi(y^2),
	\]
	并以 \(y\) 为坐标。若进一步写 \(P(z)=(z-a)U(z)\)，在小圆盘上取 \(q(z)^2=U(z)\)，则
	\[
		t=\frac{y}{q(z)},\qquad z-a=t^2
	\]
	是把投影化成二次幂的坐标。这里 \(q(a)\ne0\)，所以 \(t\) 对 \(y\) 的导数不为零。这些图册的过渡由方程中的全纯函数给出。

	令
	\[
		Q(u)=u^dP(1/u)=a_d+a_{d-1}u+\cdots+a_0u^d,
	\]
	于是 \(Q(0)=a_d\ne0\)。取原点附近的全纯平方根 \(q(u)^2=Q(u)\)。

	当 \(d=2m\) 时，在 \(|z|\) 很大处使用
	\[
		u=\frac1z,\qquad v=y\,u^m.
	\]
	方程化为 \(v^2=Q(u)\)。这给出两张图像 \(v=q(u)\) 与 \(v=-q(u)\)，各补一个点
	\[
		p_\infty^+=(0,q(0)),\qquad p_\infty^-=(0,-q(0)).
	\]
	两处都以 \(u\) 为坐标；目标球面在无穷远的坐标同样是 \(u=1/z\)，故投影在两处均无分歧。

	当 \(d=2m+1\) 时，改用
	\[
		u=\frac1z,\qquad v=y\,u^{m+1}.
	\]
	方程化为 \(v^2=uQ(u)\)，在 \(u=0\) 上只有 \((0,0)\)。令 \(t=v/q(u)\)，便得到
	\[
		u=t^2,\qquad v=tq(t^2).
	\]
	这是一张以 \(t\) 为坐标的圆盘，去掉 \(t=0\) 后与原曲面在无穷远附近的部分对应。只补点 \(p_\infty\) 即可，且投影在目标坐标 \(u\) 中为 \(u=t^2\)。在这张坐标中，原函数为
	\begin{equation}\label{g4:eq:0204:odd-infinity}
		z=t^{-2},\qquad y=t^{-(2m+1)}q(t^2).
	\end{equation}

	按照这些坐标与原曲面的公共部分粘合，得到无边界的一维复流形。它是 Hausdorff 的：投影不同的点可由球面中的不交邻域分开；同一有限纤维中的不同点由两张平方根图像分开；偶数次数时新补的两个点有不交的坐标圆盘。原仿射部分第二可数，新增图册只有一张或两张，故整个空间第二可数。

	紧性可以直接检查。对充分大的 \(R\)，集合
	\[
		K_R=\{(z,y)\in X_P^\circ:|z|\le R\}
	\]
	是 \(\mathbb C^2\) 中的闭有界集，因为 \(|y|^2=|P(z)|\) 在 \(|z|\le R\) 上有统一上界，所以 \(K_R\) 紧。无穷远附近的一张或两张闭坐标圆盘覆盖 \(|z|\ge R\) 的部分，并包含所有新补点。因此 \(X_P\) 是有限个紧集的并。

	最后说明连通性。去掉有限分支值后，底空间 \(\mathbb C\setminus\{a_1,\ldots,a_d\}\) 道路连通，投影是二叶覆盖。绕任一 \(a_j\) 的小正向圆周一周，局部公式 \(z-a_j=t^2\) 使两张覆盖片交换。因此基点纤维中的两个点可沿提升道路相连；其他纤维中的点也可沿底空间道路提升连到它们。这个无分歧部分连通，所有分歧点和新补点均在其闭包中，故 \(X_P^\circ\) 和 \(X_P\) 都连通。
\end{proof}

这个构造也可用剪开后跨缝粘合来表示。旋转 \(z\) 平面，使各 \(a_j\) 的实部互不相同，并按实部递增排列。依次用线段连接 \(a_1,a_2\)，连接 \(a_3,a_4\)，如此继续；这些线段的实部区间不交，因而彼此不交。如果 \(d\) 为奇数，再从最后的 \(a_d\) 向实部增加的方向画一条到无穷远的水平射线。删去这些切缝所得的开集 \(U\) 道路连通：连接两点的折线若遇到切缝，可以沿切缝的一侧绕过有限端点，作有限次小的绕行后便避开全部切缝。

\ProseParagraphBreak

在 \(U\) 上可以选取单值全纯平方根。为核对这一点，取 \(U\) 中任一闭路 \(\gamma\)。同一有限切缝的两个端点沿切缝相连而不遇到 \(\gamma\)，所以它们关于 \(\gamma\) 的绕数相等；若有通向无穷远的射线，沿射线绕数恒定，并在足够远处为零，故该端点的绕数为零。于是
\[
	\frac1{2\pi\mathrm{i}}\int_\gamma\frac{P'(z)}{P(z)}\,\mathrm{d}z
	=\sum_{j=1}^d\operatorname{Ind}(\gamma,a_j)
\]
是偶数。固定 \(z_0\in U\) 及 \(c^2=P(z_0)\)，公式
\[
	s(z)=c\exp\left(\frac12\int_{z_0}^z\frac{P'(\zeta)}{P(\zeta)}
	\,\mathrm{d}\zeta\right)
\]
因此与 \(U\) 中的路径选择无关，并满足 \(s^2=P\)。

\ProseParagraphBreak

取两张 \(U\)，分别标记 \(y=s(z)\) 与 \(y=-s(z)\)。跨过一条切缝时，平方根的两种边界值互为相反数：这在端点附近由 \(t^2=z-a_j\) 看出，并沿切缝内部由解析延拓保持。因此粘合规则为
\[
	\begin{array}{c@{\ \longleftrightarrow\ }c}
		\text{第一张的上岸}&\text{第二张的下岸}\\
		\text{第一张的下岸}&\text{第二张的上岸}.
	\end{array}
\]
这里上下岸指沿切缝选定的两侧，粘合保持同一个 \(z\) 值和同一个 \(y\) 值。缝的端点用 \(t\) 圆盘补入，便没有留下边界；无穷远端的补法由次数奇偶性决定。这种粘合描述与方程构造给出同一个带投影的曲面。

\begin{corollary}\label{g4:cor:0204:hyperelliptic-genus}
	设 \(P\in\mathbb C[z]\) 无重根，且 \(\deg P=d\ge1\)。上述紧曲面 \(X_P\) 的亏格为
	\[
		g(X_P)=\left\lfloor\frac{d-1}{2}\right\rfloor.
	\]
\end{corollary}
\begin{proof}
	映射 \(\pi:X_P\to\mathbb P^1\) 的次数为 \(2\)，目标亏格为 \(0\)。每个分歧点的指数为 \(2\)，所以它对分歧除子的次数贡献 \(1\)。由\cref{g4:thm:0203:riemann-hurwitz}，
	\[
		2g(X_P)-2=
		\begin{cases}
			-4+d,&d\text{ 为偶数},\\
			-4+(d+1),&d\text{ 为奇数}.
		\end{cases}
	\]
	分别解出 \(g=(d-2)/2\) 与 \(g=(d-1)/2\) 即得结论。
\end{proof}

亏格至少为 \(2\)，并存在到 \(\mathbb P^1\) 的二次全纯映射的紧黎曼曲面，称为\term[chaotuoyuanqumian]{超椭圆曲面}[hyperelliptic Riemann surface]。上述方程在 \(d\ge5\) 时给出这样的曲面；\(d=1,2\) 时得到亏格 \(0\) 的二叶模型，\(d=3,4\) 时得到亏格 \(1\) 的模型。无重根条件使有限处的合并都恰为 \(t\mapsto t^2\)，并保证仿射方程本身已经光滑。

\subsection{三次曲线及亏格}
\label{g4:subsec:0204:02}

三次方程把上面的奇数次数模型放进普通射影平面，不需要另选加权坐标。取
\[
	P(z)=a_3z^3+a_2z^2+a_1z+a_0,\qquad a_3\ne0,
\]
并假设 \(P\) 无重根。在齐次坐标 \([Z:Y:T]\) 中，考虑射影三次曲线
\begin{equation}\label{g4:eq:0204:projective-cubic}
	Y^2T=a_3Z^3+a_2Z^2T+a_1ZT^2+a_0T^3.
\end{equation}
在 \(T\ne0\) 的仿射图中，\(z=Z/T\)、\(y=Y/T\)，方程恢复为 \(y^2=P(z)\)。

\begin{proposition}\label{g4:prop:0204:cubic}
	设三次多项式 \(P\in\mathbb C[z]\) 无重根。方程\eqref{g4:eq:0204:projective-cubic}定义一条光滑射影三次曲线 \(E\)，它与\cref{g4:prop:0204:hyperelliptic-completion}构造的 \(X_P\) 双全纯。无穷远直线 \(T=0\) 与它只有一个交点 \(O=[0:1:0]\)，且
	\[
		g(E)=1.
	\]
	仿射部分的微分 \(\omega=\mathrm{d}z/y\) 延拓为 \(E\) 上处处非零的全纯微分。
\end{proposition}
\begin{proof}
	在仿射图中，方程函数 \(F(z,y)=y^2-P(z)\) 的两个偏导数为 \(-P'(z)\) 与 \(2y\)。它们不能同时为零，否则 \(P\) 有重根；有限处的全纯坐标已经在前一命题中给出。令 \(T=0\)，齐次方程迫使 \(Z=0\)，于是无穷远只有 \(O=[0:1:0]\)。在齐次方程左减右所得多项式中，对 \(T\) 的偏导数在 \(O\) 为 \(1\)，所以这里也不出现奇点。

	还可直接写出 \(O\) 附近的坐标，核对它就是前面补入的点。令
	\[
		Q(u)=a_3+a_2u+a_1u^2+a_0u^3,\qquad q(u)^2=Q(u),\quad q(0)\ne0.
	\]
	由\eqref{g4:eq:0204:odd-infinity}，仿射部分在无穷远附近表示为
	\[
		z=t^{-2},\qquad y=t^{-3}q(t^2).
	\]
	相应射影点为
	\[
		[Z:Y:T]=[t:q(t^2):t^3].
	\]
	它在 \(t=0\) 取值 \(O\)。在 \(Y\ne0\) 的射影图中，第一坐标 \(Z/Y=t/q(t^2)\) 对 \(t\) 的导数在 \(0\) 为 \(1/q(0)\ne0\)，所以 \(t\) 确为这条射影曲线上的全纯局部坐标。仿射恒等映射因而在补点处也成为双全纯对应。曲面紧性由前一构造得到，亏格由\cref{g4:cor:0204:hyperelliptic-genus}取 \(d=3\) 得到。等价地，投影 \(z:E\to\mathbb P^1\) 在三个有限根和 \(O\) 处各有一次简单分歧，Hurwitz 公式给出 \(2g-2=-4+4=0\)。

	现在逐点检查 \(\omega\)。在 \(y\ne0\) 的有限点处，\(z\) 是坐标，系数 \(1/y\) 全纯且非零。在有限分歧点 \((a,0)\) 附近，\(y\) 是坐标；对方程求微分，得到
	\[
		2y\,\mathrm{d}y=P'(z)\,\mathrm{d}z,
		\qquad
		\omega=\frac{2\,\mathrm{d}y}{P'(z)}.
	\]
	由于 \(P'(a)\ne0\)，这个表达式延拓且不为零。最后，在 \(O\) 的 \(t\) 坐标中，
	\[
		\omega
		=\frac{-2t^{-3}\,\mathrm{d}t}{t^{-3}q(t^2)}
		=-\frac{2}{q(t^2)}\,\mathrm{d}t,
	\]
	仍全纯且非零。
\end{proof}

在一张复坐标图中，\(\mathrm{d}z/y\) 看起来在 \(y=0\) 处有分母为零的问题；换到真正的曲面坐标后，分子同时消去这一阶数。无穷远处也由同一次数计算消去了表面上的极点。因此 \(\operatorname{div}(\omega)=0\)，而 \(\omega\) 给出 \(E\) 的典范线丛的全局非消失截面。

\ProseParagraphBreak

紧黎曼曲面 \(E\) 的亏格为 \(1\)，并选定一点 \(O\in E\) 时，二者组成一条\term[tuoyuanquxian]{椭圆曲线}[elliptic curve]。上述三次模型以唯一的无穷远点作为 \(O\)。这里的方程、分歧点和非消失微分都已由局部计算确定；它们为随后研究积分与周期提供了具体对象。

\subsection{拓扑单值与代数方程的对应}
\label{g4:subsec:0204:03}

局部根的变化可以直接写成纤维上的置换。例如方程
\[
	w^n=z,\qquad n\ge1,
\]
在 \(z\in\mathbb C^\times\) 上定义 \(n\) 叶无分歧覆盖。固定正实基点 \(z_0\)，把它的根标为
\[
	w_j=z_0^{1/n}e^{2\pi\mathrm{i} j/n},\qquad j\in\mathbb Z/n\mathbb Z.
\]
沿 \(z=z_0e^{2\pi\mathrm{i} s}\)、\(0\le s\le1\) 继续时，根 \(w_j\) 变成 \(w_{j+1}\)，所以单值是一个 \(n\)-循环。用 \(w\) 作源坐标，整个紧化映射就是 \(\mathbb P^1\to\mathbb P^1\)、\(w\mapsto w^n\)。在无穷远处令 \(v=1/w\)、\(u=1/z\)，仍有 \(u=v^n\)；\(0\) 与无穷远的分歧指数都为 \(n\)。目标无穷远附近的正向小环路在 \(z\) 平面中沿相反方向绕行，故其单值是前一个循环的逆。

\ProseParagraphBreak

对 \(y^2=P(z)\)，绕每个简单有限零点的单值都是换位 \((12)\)。全部有限零点的环路依次作用，得到 \((12)^d\)，所以绕无穷远的单值在 \(d\) 偶数时为恒等，在 \(d\) 奇数时为换位。这恰好恢复了前面的两个补点与一个二重分歧补点。局部幂坐标、根的解析延拓和拓扑单值因此给出同一组数据。

\ProseParagraphBreak

现在固定球面上的有限非空集合 \(B=\{b_1,\ldots,b_r\}\) 及基点 \(z_0\notin B\)。在 \(U=\mathbb P^1\setminus B\) 中，选取从 \(z_0\) 出发、绕各 \(b_j\) 正向一周再返回的环路 \(\gamma_j\)，使
\[
	\pi_1(U,z_0)=
	\langle\gamma_1,\ldots,\gamma_r:
	\gamma_1*\cdots*\gamma_r=1\rangle.
\]
这个展示可从第一卷命题 5.2.10 直接得到：把 \(b_r\) 送到无穷远后，\(U\) 成为去掉 \(r-1\) 个点的平面，前 \(r-1\) 个绕点环路自由生成基本群；按切开路径的次序，它们的乘积是包围全部有限缺失点的大圆周，而无穷远的正向环路是其反向。这里仍约定 \(\alpha*\beta\) 先走 \(\alpha\) 后走 \(\beta\)。

\begin{theorem}\label{g4:thm:0204:monodromy-construction}
	固定 \(B,z_0,\gamma_1,\ldots,\gamma_r\) 如上，并取整数 \(n\ge1\)。连通紧黎曼曲面的 \(n\) 次全纯映射
	\[
		\pi:X\longrightarrow\mathbb P^1,
		\qquad \text{所有分支值均包含于 }B,
	\]
	在保持目标投影的双全纯等价下，与满足以下条件的置换元组 \((\sigma_1,\ldots,\sigma_r)\in S_n^r\) 在同时共轭下的等价类一一对应：
	\begin{enumerate}
		\item \(\sigma_r\circ\cdots\circ\sigma_1=\operatorname{id}\)；
		\item \(\sigma_1,\ldots,\sigma_r\) 生成的群在 \(\{1,\ldots,n\}\) 上传递。
	\end{enumerate}
	其中 \(\sigma_j\) 是沿 \(\gamma_j\) 提升后作用于基点纤维的终点置换。它的每条长度为 \(e\) 的循环对应 \(b_j\) 上方的一个点，该点的分歧指数为 \(e\)。这里同时共轭指用同一个 \(\tau\in S_n\) 把每个 \(\sigma_j\) 换成 \(\tau\sigma_j\tau^{-1}\)。
\end{theorem}
\begin{proof}
	先从给定映射得到数据。删去 \(B\) 及其原像，得到一个 \(n\) 叶无分歧覆盖。\cref{g4:thm:0202:unbranched-cover}的证明已说明，连通二维流形删去任意有限集后仍道路连通；这里 \(B\) 及其原像都有限，所以这个结论适用。因此基点纤维上的单值传递。环路先走 \(\gamma_1\) 再走 \(\gamma_2\) 时，终点置换为 \(\sigma_2\circ\sigma_1\)，故基本群的唯一关系给出第一项。局部幂模型说明，一个 \(e\) 次分歧点恰好使附近的 \(e\) 张片循环置换。重新编号纤维，正好产生同时共轭。

	反过来，满足第一项的元组给出 \(G=\pi_1(U,z_0)\) 在 \(S=\{1,\ldots,n\}\) 上的右作用，规则为 \(j\cdot\gamma_\nu=\sigma_\nu(j)\)。第一卷第5.5.3小节的覆盖重建定理 5.5.6 适用于 \(U\)：它道路连通，局部同胚于圆盘，所以局部道路连通且半局部单连通。具体地，取从 \(z_0\) 出发的道路类万有覆盖 \(\widetilde U\)，用
	\[
		g\cdot(j,[\alpha])=
		\bigl(j\cdot g^{-1},[\gamma*\alpha]\bigr),
		\qquad g=[\gamma],
	\]
	在 \(S\times\widetilde U\) 上取商，便得到具有指定单值的 \(n\) 叶覆盖 \(E\to U\)。该定理还说明，覆盖的等价恰好对应纤维的等变双射。每个连通分支都遇到基点纤维，而两个纤维点属于同一分支，当且仅当有底空间环路的提升连接它们；第二项因而保证 \(E\) 连通。

	有限叶性使 \(E\) 的覆盖片有可数基；同纤维的不同点由不同覆盖片分离，异纤维点由底空间分离。因此 \(E\) 是 Hausdorff 且第二可数的，\cref{g4:thm:0201:lift-complex-structure}给出唯一使投影局部双全纯的复结构。围绕每个 \(b_j\) 取互不相交的小坐标圆盘 \(D_j\)，把 \(b_j\) 的坐标记为 \(0\)。在穿孔圆盘 \(D_j^*\) 上，\cref{g4:lem:0202:punctured-disk-cover}把覆盖分成与 \(\sigma_j\) 的循环对应的有限个分支；长度为 \(e\) 的分支有坐标 \(t\)，投影写成 \(u=t^e\)。给每个这样的分支补入 \(t=0\)，便得到一张无边界复曲面 \(X\) 和全纯投影 \(\pi\)。原来的 \(E\) 稠密且连通，故 \(X\) 连通。

	所得曲面紧。确实，有限叶覆盖在底空间的每个紧集上有紧原像：用有限个闭小圆盘覆盖这个紧集，且各闭圆盘都包含在均匀覆盖的稍大圆盘内，它们的原像分别是有限个紧副本。球面去掉各 \(b_j\) 的更小圆盘后是紧集，其原像紧；再加上刚补入的有限个闭 \(t\) 圆盘，就覆盖整个 \(X\)。补点引理同时保证各新点有互不混淆的圆盘邻域，因而保留 Hausdorff 性和第二可数性。

	最后，两个单值元组只差同时共轭时，覆盖重建定理给出 \(E\) 之间的等价。它在每张覆盖片上由目标坐标表示为恒等，故全纯；由穿孔圆盘补点引理，它在每个新点处唯一延拓为双全纯映射。反过来，保持 \(\pi\) 的双全纯等价限制到基点纤维后与全部单值置换相容，必给出同时共轭。这也证明补点后没有额外的等价选择。
\end{proof}

定理允许 \(B\) 中有实际上不分歧的点：它们对应 \(\sigma_j=\operatorname{id}\)。如果要求分支值恰为 \(B\)，就要求每个 \(\sigma_j\) 非恒等。给定循环分解之后，补点的个数、每一点的指数以及对 Hurwitz 公式的贡献，都可以逐个读出。

\subsection{不可约、连通与分支数据}
\label{g4:subsec:0204:04}

对一般方程，连通性要在局部根组成的覆盖上判断。考虑关于 \(w\) 首一的多项式
\[
	P(z,w)=w^n+a_1(z)w^{n-1}+\cdots+a_n(z),
	\qquad a_j\in\mathbb C(z),\quad n\ge1.
\]
假设它在 \(\mathbb C(z)[w]\) 中可分，即 \(P\) 与 \(\partial P/\partial w\) 互素。互素性给出 \(A,B\in\mathbb C(z)[w]\) 使 \(AP+BP_w=1\)。清去这个恒等式中有理系数的分母后，说明只有有限多个 \(z\) 可能使专门化方程出现重根。因此可以选有限集合 \(S\subset\mathbb P^1\)，包含无穷远、全部系数极点及所有这样的例外值，使在 \(U=\mathbb P^1\setminus S\) 上，方程始终有 \(n\) 个不同的有限根。

\begin{proposition}\label{g4:prop:0204:root-cover}
	设 \(P(z,w)\in\mathbb C(z)[w]\) 首一、可分，次数为 \(n\ge1\)，并令有限集合 \(S\) 与 \(U\) 满足上述条件。则
	\[
		E_P=\{(z,w)\in U\times\mathbb C:P(z,w)=0\}
		\longrightarrow U,\qquad (z,w)\longmapsto z,
	\]
	是 \(n\) 叶全纯无分歧覆盖。按照穿孔圆盘幂模型补点，得到唯一的紧化 \(X_P\)，它是有限个紧黎曼曲面的不交并；投影延拓为到 \(\mathbb P^1\) 的全纯有限映射。坐标函数 \(z,w\) 均延拓为 \(X_P\) 上的亚纯函数。
\end{proposition}
\begin{proof}
	在 \(z_0\in U\) 上，全部根 \(w_1,\ldots,w_n\) 都是简单根。围绕它们各取一个小圆周 \(C_j\)，使相应闭圆盘两两不交且各只含这一根。由于 \(P\) 的系数在 \(z_0\) 附近全纯，缩小 \(z_0\) 的邻域后，\cref{g4:thm:0102:rouche}使每个 \(C_j\) 内始终恰有一个根。由\cref{g4:thm:0102:residue}，这个根可以写成
	\[
		w_j(z)=\frac1{2\pi\mathrm{i}}\int_{C_j}
		\frac{\xi\,P_w(z,\xi)}{P(z,\xi)}\,\mathrm{d}\xi.
	\]
	积分圆周上的分母不为零，并在缩小的闭参数圆盘上一致远离零，故可对 \(z\) 在积分号下求导；于是 \(w_j(z)\) 全纯。这些根彼此不同，且已给出次数 \(n\) 方程的全部根，所以得到覆盖的局部形式。

	在每个 \(b\in S\) 上方，应用\cref{g4:lem:0202:punctured-disk-cover}给各连通分支补点；与\cref{g4:thm:0204:monodromy-construction}证明中的紧性论证相同，有限叶性保证补成的空间紧。每个连通分支至少遇到一次 \(n\) 点纤维，所以分支数有限。补点的坐标与等价延拓唯一，故这一紧化在保持原覆盖的双全纯等价下唯一。

	剩下的实质问题是 \(w\) 在新点处会不会出现本性奇点。对任意首一复系数多项式
	\[
		w^n+c_1w^{n-1}+\cdots+c_n=0,
	\]
	每个根都满足
	\begin{equation}\label{g4:eq:0204:root-bound}
		|w|\le1+\max_{1\le j\le n}|c_j|.
	\end{equation}
	事实上，记右边的最大值部分为 \(A\)。若 \(|w|>1+A\)，则
	\[
		\left|\sum_{j=1}^n c_jw^{n-j}\right|
		\le A\sum_{j=1}^n|w|^{n-j}
		=A\frac{|w|^n-1}{|w|-1}<|w|^n,
	\]
	与方程矛盾。

	在 \(b\) 附近取目标坐标 \(u\)，使 \(u(b)=0\)；若 \(b=\infty\)，取 \(u=1/z\)。有理系数 \(a_j\) 在 \(u=0\) 至多有极点，因此存在整数 \(N\ge0\) 及常数 \(C\)，使穿孔小圆盘中
	\[
		|a_j(z(u))|\le C|u|^{-N}\qquad(1\le j\le n).
	\]
	由\eqref{g4:eq:0204:root-bound}，每个根满足 \(|w|\le C'|u|^{-N}\)。在某个新点的坐标 \(t\) 中，投影为 \(u=t^e\)，故
	\[
		|w(t)|\le C'|t|^{-eN}.
	\]
	函数 \(t^{eN}w(t)\) 有界且在穿孔圆盘上全纯，由\cref{g4:prop:0102:singularity-tests}延拓到 \(0\)。所以 \(w\) 在新点至多有极点。函数 \(z\) 的延拓更直接：在有限 \(b\) 上有 \(z=b+t^e\)，在无穷远上有 \(z=t^{-e}\)。因此二者均亚纯。
\end{proof}

这个根的增长界还把单值的轨道变成有理系数因子。它是“覆盖连通”与“多项式不可约”之间的关键一步。

\begin{theorem}\label{g4:thm:0204:irreducible-connected}
	设 \(P(z,w)\in\mathbb C(z)[w]\) 是次数 \(n\ge1\) 的首一可分多项式。取包含无穷远和全部系数极点的有限集合 \(S\)，使 \(z\notin S\) 时 \(P(z,w)\) 有 \(n\) 个不同根，令 \(E_P\to U=\mathbb P^1\setminus S\) 为根覆盖，\(X_P\) 为上述补点所得的紧化。以下条件等价：
	\begin{equivalences}
		\item \(P\) 在 \(\mathbb C(z)[w]\) 中不可约；
		\item 根覆盖 \(E_P\) 连通；
		\item 单值作用在任一基点的 \(n\) 个根上传递；
		\item \(X_P\) 连通。
	\end{equivalences}
	更精确地，\(P\) 的各个首一不可约因子，与 \(E_P\) 的连通分支、单值轨道及 \(X_P\) 的连通分支分别一一对应。
\end{theorem}
\begin{proof}
	每个覆盖分支都遇到基点纤维；同一纤维中的两点能被道路连接，恰好表示某条底空间环路的提升把前一点送到后一点。因此连通分支与单值轨道对应。每个新补点都有一个圆盘邻域，其穿孔部分连通，只属于原覆盖的一个分支，所以补点不会把两个分支合在一起。于是后三个条件等价。

	设单值不传递，在基点纤维中选择一个大小为 \(k\) 的轨道 \(O\)，其中 \(1\le k<n\)。在基点附近把根写成全纯函数 \(w_1(z),\ldots,w_n(z)\)，并组成
	\[
		Q_O(z,w)=\prod_{j\in O}(w-w_j(z)).
	\]
	其系数是这一组根的初等对称函数。沿任意路径继续根，可以继续这些系数；沿闭路返回时，轨道 \(O\) 被整体保持，根只在其中置换，所以每个系数返回原函数芽。由此它们在 \(U\) 上粘成单值全纯函数。

	还须证明这些系数是有理函数。在任意 \(b\in S\) 的穿孔坐标邻域中，上一命题的根界给出 \(|w_j|\le C|u|^{-N}\)。\(Q_O\) 的第 \(\ell\) 个系数是 \(\binom{k}{\ell}\) 个 \(\ell\) 重乘积的和，因此
	\[
		|q_\ell(u)|\le\binom{k}{\ell}C^\ell|u|^{-N\ell}.
	\]
	所以 \(u^{N\ell}q_\ell(u)\) 在穿孔圆盘上有界；由\cref{g4:prop:0102:singularity-tests}，它全纯延拓到原点。这说明每个 \(q_\ell\) 在包括无穷远在内的全部例外点至多有极点，即成为球面上的亚纯函数。

	球面上的亚纯函数必为有理函数，理由如下。它的极点局部有限，而球面紧，所以只有有限多个。减去所有有限极点的有限主部后，得到一个整函数 \(h(z)\)。由于无穷远至多为极点，存在整数 \(M\ge0\)，使 \(|h(z)|\le C(1+|z|)^M\)。对任意固定 \(a\)，在半径 \(R\) 的圆周上使用 Cauchy 估计，有
	\[
		|h^{(M+1)}(a)|
		\le (M+1)!\,\frac{C(1+|a|+R)^M}{R^{M+1}}
		\longrightarrow0.
	\]
	故 \(h\) 是次数不超过 \(M\) 的多项式。加回各有限主部，就得到有理函数。应用于 \(q_\ell\)，便知 \(Q_O\in\mathbb C(z)[w]\)。对其余轨道的根作同样构造，得到正次数首一多项式 \(R\)，并且 \(P=Q_OR\)。因此不传递必导致可约。

	反过来，设 \(P=QR\)，其中 \(Q,R\in\mathbb C(z)[w]\) 首一且均有正次数。先核对两因子的系数在 \(U\) 上没有极点。在一个 \(z_0\in U\) 的小闭圆盘内，\(P\) 的系数有界，所以全部根由\eqref{g4:eq:0204:root-bound}统一有界。离开 \(Q,R\) 系数的有限个可能极点时，每个因子的根都是 \(P\) 的一部分根，它们的对称函数因而有同一个局部上界；这些系数的可能极点遂由可去奇点判据消去。于是 \(Q,R\) 的系数都在 \(U\) 全纯。每个纤维中二者都至少有一个根，并且没有公共根，否则 \(P\) 有重根。故
	\[
		E_P=\{Q=0\}\ \sqcup\ \{R=0\}
	\]
	是两个非空闭子集的不交并；每个子集的补集也闭，所以它们都在 \(E_P\) 中开。根覆盖因此不连通。

	最后，对每个轨道作出的 \(Q_O\) 必不可约；若它可约，刚才的论证又把这个轨道分成两个非空不变部分，与轨道的定义矛盾。各轨道的积给出 \(P\)，所以得到所述逐因子对应。
\end{proof}

这里的连通对象是由不同根继续并补成的光滑曲面。原始平面解集可能在奇点处把不同分支碰到一起。例如
\[
	w^2-z^2=(w-z)(w+z)
\]
的仿射零点集是两条相交直线，作为集合连通；但在 \(z\ne0\) 上，根 \(w=z\) 与 \(w=-z\) 各自单值，根覆盖有两个分支。补点时每个分支各补自己的点，得到两张不交的射影直线。相反，\(w^2-z^3\) 在 \(\mathbb C(z)[w]\) 中不可约：若可约，就有有理函数 \(r(z)\) 满足 \(r(z)^2=z^3\)，而左边在 \(0\) 的阶数为偶数，右边为 \(3\)，矛盾。它的光滑参数为
\[
	z=t^2,\qquad w=t^3.
\]
在 \(z\ne0\) 处逆参数为 \(t=w/z\)，所以参数化确实恢复根覆盖；\(t=0\) 补一个点，\(t=\infty\) 再补一个点，得到一张连通曲面。方程的奇点、不可约性和覆盖的分歧分别在这里承担不同的信息。

\ProseParagraphBreak

分支位置与各点的分歧指数给出局部单值的循环类型，但没有给出各个置换如何共同作用于同一纤维。下面的例子把这一区别落实为可核验的数据。

\begin{example}\label{g4:exm:0204:same-branch-profile}
	固定球面上四个不同的点 \(b_1,b_2,b_3,b_4\)、基点及满足乘积关系的标准环路。取三点纤维，考虑两个元组
	\[
		\begin{aligned}
			\boldsymbol\sigma&=\bigl((12),(12),(23),(23)\bigr),\\
			\boldsymbol\tau&=\bigl((12),(23),(23),(12)\bigr).
		\end{aligned}
	\]
	两个元组按规定次序的复合均为恒等；它们都生成由 \((12),(23)\) 生成的传递群。因此\cref{g4:thm:0204:monodromy-construction}分别给出连通三次覆盖。每个 \(b_j\) 上方，两种覆盖都各有一个指数为 \(2\) 的点和一个指数为 \(1\) 的点；分支值恰好是同一组四点，所有局部分歧数据完全相同。

	这两个元组却不同时共轭。在 \(\boldsymbol\sigma\) 中，第一项等于第二项；在 \(\boldsymbol\tau\) 中，第一项不等于第二项，而同时共轭保持这种相等关系。因此两覆盖在固定目标投影下不等价。它们的源曲面亏格却相同，因为 Hurwitz 公式都给出
	\[
		2g-2=3(-2)+4=-2,\qquad g=0.
	\]
\end{example}

完整的单值元组因而同时记录局部循环和这些循环在同一组片上的配合；传递性检测连通性，循环长度给出补点与分歧指数，Hurwitz 公式再求出亏格。对有理系数代数方程，单值轨道还能借根增长界转成有理系数因子，把这一几何分解准确地接回不可约分解。
% 纲要条目排为编号小节标题，并进入目录与书签；说明段仍保留于纲要注释中。


\begin{chaptersummary}
待撰写
\end{chaptersummary}

\BookExercises
待撰写

% !TeX root = ../../Book4.tex
% 纲要标题：## 第3章　曲面上的层、对偶与 Riemann–Roch
% 纲要位置：几何学计划纲要.md，第 3328 行。
% 正文状态：未撰写；本文件目前是框架。

% 纲要内容开始
% 
% > 本章目标与证明范围：把局部主部问题、Dolbeault 解算子、Serre 配对和 Riemann–Roch 连成一条完整解析证明。必要椭圆工具按第三卷的明确模块补齐；不借高维消失定理倒推曲线结果。
% 
% 纲要内容结束

\chapter{曲面上的层、对偶与 Riemann–Roch}
\label{g4:ch:03}
\BookChapterNumber{g4:ch:03}

\begin{chapterabstract}
待撰写
\end{chapterabstract}

\begin{learninggoals}
待撰写
\end{learninggoals}

待撰写

% !TeX root = ../../Book4.tex
% 纲要标题：### 3.1 线丛与除子的几何
% 纲要位置：几何学计划纲要.md，第 3332 行。
% 正文状态：未撰写；本文件目前是框架。

% 纲要内容开始
% 
% * 过渡函数、截面和同构
% * 除子及线性等价
% * 主部层和点处修改
% * 线性系与有理映射
% 
% 纲要内容结束

\section{线丛与除子的几何}
\label{g4:sec:0301}

待撰写

\subsection{过渡函数、截面和同构}
\label{g4:subsec:0301:01}
待撰写

\subsection{除子及线性等价}
\label{g4:subsec:0301:02}
待撰写

\subsection{主部层和点处修改}
\label{g4:subsec:0301:03}
待撰写

\subsection{线性系与有理映射}
\label{g4:subsec:0301:04}
待撰写
% 纲要条目排为编号小节标题，并进入目录与书签；说明段仍保留于纲要注释中。

% !TeX root = ../../Book4.tex
% 纲要标题：### 3.2 层上同调的曲面模型
% 纲要位置：几何学计划纲要.md，第 3339 行。
% 正文状态：未撰写；本文件目前是框架。

% 纲要内容开始
% 
% * 预层、层和茎的最低工具
% * Čech 余圈及细化
% * 局部 ∂̄ 可解性
% * Dolbeault 消解和细层无上同调
% 
% 纲要内容结束

\section{层上同调的曲面模型}
\label{g4:sec:0302}

待撰写

\subsection{预层、层和茎的最低工具}
\label{g4:subsec:0302:01}
待撰写

\subsection{Čech 余圈及细化}
\label{g4:subsec:0302:02}
待撰写

\subsection{\texorpdfstring{局部 \(\bar{\partial}\) 可解性}{局部 \={∂} 可解性}}
\label{g4:subsec:0302:03}
待撰写

\subsection{Dolbeault 消解和细层无上同调}
\label{g4:subsec:0302:04}
待撰写
% 纲要条目排为编号小节标题，并进入目录与书签；说明段仍保留于纲要注释中。

% !TeX root = ../../Book4.tex
% 纲要标题：### 3.3 椭圆工具与 Serre 对偶
% 纲要位置：几何学计划纲要.md，第 3563 行。
% 正文状态：初稿；尚未正式审核。

% 纲要内容开始
% 
% * 曲面上的 ∂̄ 主符号
% * Fredholm 和调和代表的最低证明链
% * 积分配对及形式伴随
% * 配对非退化性
% * 与留数配对的比较
% 
% 纲要内容结束

\section{椭圆工具与 Serre 对偶}
\label{g4:sec:0303}

一个光滑 \(L\) 值 \((0,1)\) 形式 \(\alpha\) 在小圆盘上总能写成 \(\bar\partial u\)，但这些局部解未必能粘成全局解。上一节把粘合障碍表示为上同调类；现在要证明障碍空间有限维，并找出检验某个障碍是否为零的全部线性条件。答案是与 \(K_X\otimes L^{-1}\) 的全纯截面积分配对。证明的分析部分从平面上的一个一阶估计开始，紧性负责把局部控制转为全局分解。

\subsection{\texorpdfstring{曲面上的 \(\bar{\partial}\) 主符号}{曲面上的 \={∂} 主符号}}
\label{g4:subsec:0303:01}
设 \(X\) 是紧黎曼曲面，\(L\to X\) 是全纯线丛。在全纯坐标 \(z=x+\mathrm{i}y\) 和全纯标架 \(e\) 中，
\[
	D\coloneqq\bar\partial_L,\qquad
	D(u e)=\frac{\partial u}{\partial\overline z}\,
	\mathrm{d}\overline z\otimes e,\qquad
	\frac{\partial}{\partial\overline z}
	=\frac12(\partial_x+\mathrm{i}\partial_y).
\]
标架变换全纯，所以这一公式在交叠处相容。对一阶算子，主符号保留最高阶导数，并把 \(\partial_{x_j}\) 替成 \(\mathrm{i}\xi_j\)。因此对实余切向量 \(\xi=\xi_x\,\mathrm{d}x+\xi_y\,\mathrm{d}y\)，有
\begin{equation}
	\label{g4:eq:0303:dbar-symbol}
	\sigma_D(\xi)(a e)
	=\frac{\mathrm{i}\xi_x-\xi_y}{2}\,
	a\,\mathrm{d}\overline z\otimes e.
\end{equation}
当 \(\xi\ne0\) 时，这个一维复向量空间之间的映射可逆；这就是 \(D\) 的\term[tuoyuanxing]{椭圆性}[ellipticity]。下面的估计说明符号可逆在函数层面意味着什么：全部一阶导数可以由 \(Du\) 控制，而不只是某个方向的导数。

\begin{lemma}
	\label{g4:lem:0303:plane-estimate}
	对平面上的每个紧支光滑复函数 \(v\)，有
	\begin{equation}
		\label{g4:eq:0303:plane-estimate}
		\|\partial_xv\|_{L^2(\mathbb R^2)}^2+
		\|\partial_yv\|_{L^2(\mathbb R^2)}^2
		=4\|\partial_{\overline z}v\|_{L^2(\mathbb R^2)}^2.
	\end{equation}
	因而 \(\|v\|_{H^1}^2=\|v\|_2^2+4\|\partial_{\overline z}v\|_2^2\)，其中 \(\|v\|_{H^1}^2=\|v\|_2^2+\|\partial_xv\|_2^2+\|\partial_yv\|_2^2\)。
\end{lemma}
\begin{proof}
	先补出所需的 Fourier 展开依据。\footnote{Jean Baptiste Joseph Fourier（1768--1830），法国数学家、物理学家，研究热传导并发展三角级数方法。}在方形环面 \(T^2=(\mathbb R/2\pi\mathbb Z)^2\) 上使用归一化测度 \((2\pi)^{-2}\mathrm{d}x\,\mathrm{d}y\)。函数 \(e_k(x,y)=e^{\mathrm{i}(k_1x+k_2y)}\)、\(k\in\mathbb Z^2\)，由直接积分构成正交单位族。

	它们的线性张成空间在 \(L^2(T^2)\) 中稠密。为说明这一点，先对连续函数使用一维核
	\[
		F_N(t)=\frac1{N+1}\left|\sum_{j=0}^N e^{\mathrm{i}jt}\right|^2.
	\]
	它非负，归一化积分为 \(1\)，且在距 \(2\pi\mathbb Z\) 至少 \(\delta>0\) 的地方满足 \(F_N(t)\le C_\delta/(N+1)\)。后一估计来自几何级数公式。两个这样的核的乘积在 \(T^2\) 上仍非负、积分为 \(1\)，其原点任意固定邻域之外的积分趋于零。把卷积积分分成这个邻域及其补集，由连续函数的一致连续性可知，卷积一致趋于原函数。每个卷积都是有限三角多项式。这里将 \(L^2\) 视为光滑函数按平方积分范数的完备化；光滑函数本身连续，故上述一致逼近给出所需的 \(L^2\) 稠密性。

	记 \(\widehat v(k)=(v,e_k)_{L^2}\)。对有限正交和应用勾股等式，再用刚证得的稠密性，得到
	\[
		\|v\|_2^2=\sum_{k\in\mathbb Z^2}|\widehat v(k)|^2.
	\]
	光滑函数的导数系数可由分部积分计算：
	\[
		\widehat{\partial_xv}(k)=\mathrm{i}k_1\widehat v(k),\qquad
		\widehat{\partial_yv}(k)=\mathrm{i}k_2\widehat v(k),\qquad
		\widehat{\partial_{\overline z}v}(k)
		=\frac{\mathrm{i}k_1-k_2}{2}\widehat v(k).
	\]
	由于 \(|\mathrm{i}k_1-k_2|^2=k_1^2+k_2^2\)，三条范数等式立即给出式\eqref{g4:eq:0303:plane-estimate}的周期版本。对平面紧支函数，取包含其支集的大方框，在框的边界附近函数恒零，故可光滑周期延拓。将方框缩放到上述标准方框，导数同时乘以相应缩放因子，便得到平面公式。归一化测度的常数在等式两侧相同。
\end{proof}

要把这个估计用于整个曲面，需要固定能对极限操作保持完备的函数空间。先取有限全纯坐标域与标架 \((U_j,z_j,e_j)\)，再取从属光滑单位分解 \(\rho_j\)，使 \(\operatorname{supp}\rho_j\Subset U_j\)。这种有限分解可直接构造：先取仍覆盖 \(X\) 的较小坐标圆盘，在对应较大坐标盘中选取非负光滑函数，令其在小盘上等于 \(1\)，最后除以这些函数之和。

\ProseParagraphBreak

同时在 \(L\) 上选光滑 Hermitian 度量，在 \(X\) 上选与复结构相容的光滑度量。\footnote{Charles Hermite（1822--1901），法国数学家，研究椭圆函数与代数方程，并证明 \(e\) 的超越性。}局部写作
\[
	|e|^2=h>0,\qquad \mathrm{d}s^2=\lambda|\,\mathrm{d}z\,|^2,\qquad
	\mathrm{d}A=\lambda\,\mathrm{d}x\,\mathrm{d}y.
\]
它们都可用上述单位分解对局部正定度量加权得到。内积对第一变量线性，故对局部表达 \(s=u e\)、\(\alpha=a\,\mathrm{d}\overline z\otimes e\)，平方范数分别为
\begin{equation}
	\label{g4:eq:0303:l2-metrics}
	\|s\|_0^2=\int_X h|u|^2\,\mathrm{d}A,\qquad
	\|\alpha\|_1^2=\int_X \frac{2h}{\lambda}|a|^2\,\mathrm{d}A.
\end{equation}
这些写法按局部坐标解释并以单位分解求和；被积密度本身与坐标无关。因 \(|\mathrm{d}\overline z|^2=2/\lambda\)，第二个范数也可在每张图中写成 \(2\int h|a|^2\,\mathrm{d}x\,\mathrm{d}y\)。

\begin{definition}
	\label{g4:def:0303:sobolev}
	在紧黎曼曲面 \(X\) 的全纯线丛 \(L\) 上，将光滑截面按平方积分范数完备化，得到 \(L^2(X,L)\)。若 \(s=u_j e_j\)，定义
	\[
		\|s\|_{H^1_{\mathrm{Sob}}}^2
		=\sum_j\|(\rho_j u_j)\circ z_j^{-1}\|_{H^1(\mathbb R^2)}^2,
	\]
	其中每个局部函数在坐标域外补零。按此范数完备化所得空间记为 \(H^1_{\mathrm{Sob}}(X,L)\)，称为一阶\term[Sobolevkongjian]{Sobolev 空间}[Sobolev space]。\footnote{\OriginalName{Сергей Львович Соболев}（Sergei Lvovich Sobolev，1908--1989），苏联数学家，研究偏微分方程，并发展广义函数与函数空间方法。}对 \(L\) 值 \((0,1)\) 形式也可用光滑局部标架作相同定义。
\end{definition}

不同有限图册、标架和单位分解所得范数彼此等价。在每个紧交叠集上，坐标变换及其逆的一阶导数有界，标架变换函数及其逆、各截断函数的一阶导数也有界；换元公式和乘积法则把一套局部范数控制在另一套的有限和之内。再交换两套选择，得到反向控制。

\ProseParagraphBreak

\(H^1_{\mathrm{Sob}}\) 的元素可视为具有平方可积弱一阶导数的截面。这里弱导数由 \(\int(\partial_xu)\varphi=-\int u\,\partial_x\varphi\) 等测试函数等式定义。光滑逼近列在上述范数中收敛时，函数和各导数分别在 \(L^2\) 中收敛，故测试等式保持；反过来，对紧支弱可微函数作光滑核卷积，函数与各弱导数都在 \(L^2\) 中收敛，因卷积与弱导数交换且平移在 \(L^2\) 中连续。截断后逐图操作便给出反向识别。这也证明完备化到 \(L^2\) 的自然映射是单射；下文把它写成连续包含 \(H^1_{\mathrm{Sob}}(X,L)\subset L^2(X,L)\)。上标的 \(\mathrm{Sob}\) 用来与层上同调 \(H^1(X,\mathcal O(L))\) 区分。

\begin{proposition}
	\label{g4:prop:0303:global-estimate}
	设 \(X\) 是紧黎曼曲面，\(L\to X\) 为全纯线丛，度量与上述 Sobolev 范数已经固定。算子
	\[
		D:H^1_{\mathrm{Sob}}(X,L)\longrightarrow L^2(X,\Lambda^{0,1}\mathrm{T}^*X\otimes L)
	\]
	连续，并存在常数 \(C\)，使
	\begin{equation}
		\label{g4:eq:0303:global-estimate}
		\|s\|_{H^1_{\mathrm{Sob}}}
		\le C\bigl(\|Ds\|_1+\|s\|_0\bigr)
	\end{equation}
	对所有 \(s\in H^1_{\mathrm{Sob}}(X,L)\) 成立。
\end{proposition}
\begin{proof}
	先设 \(s\) 光滑。在全纯标架中，乘积法则给出
	\[
		\partial_{\overline z_j}(\rho_j u_j)
		=\rho_j\partial_{\overline z_j}u_j+
		(\partial_{\overline z_j}\rho_j)u_j.
	\]
	对每个紧支局部函数应用\cref{g4:lem:0303:plane-estimate}。第一项由 \(Ds\) 的局部 \(L^2\) 范数控制；第二项因截断函数的一阶导数有界，由 \(s\) 的 \(L^2\) 范数控制。支集上度量权函数与坐标面积密度都有正的上下界，有限求和得到式\eqref{g4:eq:0303:global-estimate}。反向的连续性估计 \(\|Ds\|_1\le C'\|s\|_{H^1_{\mathrm{Sob}}}\) 由一阶导数表达和有限局部化得到。最后按定义用光滑截面逼近，将算子及估计延拓到完备空间。
\end{proof}

\subsection{Fredholm 和调和代表的最低证明链}
\label{g4:subsec:0303:02}
估计\eqref{g4:eq:0303:global-estimate}右边仍有 \(\|s\|_0\)。这一较弱范数项不能直接删掉，因为全纯截面属于 \(D\) 的核。紧嵌入将表明：去掉有限维的核之后，它可以被 \(\|Ds\|_1\) 吸收。

\begin{lemma}[Rellich 紧嵌入]
	\label{g4:lem:0303:rellich}
	设 \(X\) 是紧、无边界的光滑曲面，\(E\) 是其上的有限秩光滑复向量丛。用有限坐标、标架和截断函数定义一阶 Sobolev 范数，则包含 \(H^1_{\mathrm{Sob}}(X,E)\hookrightarrow L^2(X,E)\) 是紧的：每个 \(H^1_{\mathrm{Sob}}\) 有界序列都有 \(L^2\) 收敛子列。
\end{lemma}
Rellich 紧嵌入把“一阶导数有界”变成对振荡的统一限制。\footnote{Franz Rellich（1906--1955），奥地利—德国数学家，研究线性算子的扰动与数学物理中的分析问题。}下面用有限频率截断给出当前版本的证明。
\begin{proof}
	先考虑 \(T^2\) 上的标量函数列 \(v_n\)，满足 \(\|v_n\|_{H^1}\le M\)。上一小节的正交展开及导数系数公式先对光滑函数成立，再由 \(H^1\) 逼近延拓，给出
	\[
		\sum_{k\in\mathbb Z^2}(1+|k|^2)|\widehat v_n(k)|^2\le M^2.
	\]
	令 \(P_Rv=\sum_{|k|\le R}\widehat v(k)e_k\)，则
	\begin{equation}
		\label{g4:eq:0303:fourier-tail}
		\|v_n-P_Rv_n\|_2^2
		=\sum_{|k|>R}|\widehat v_n(k)|^2\le R^{-2}M^2.
	\end{equation}
	对每个整数 \(R\)，\(P_Rv_n\) 位于一个固定有限维空间的有界集内，因而有收敛子列。依次对 \(R=1,2,\ldots\) 取子列，再取对角子列，使所有有限频率截断均收敛。给定 \(\varepsilon>0\)，先选 \(R\) 使 \(2M/R<\varepsilon/2\)，再使这个有限截断中的两项之差小于 \(\varepsilon/2\)。三角不等式与\eqref{g4:eq:0303:fourier-tail}表明对角子列在 \(L^2\) 中为 Cauchy 列，因此收敛。

	在曲面上，将有界截面列乘以定义范数时选定的各截断函数，并在每个坐标方框内补零、周期化。每个标架分量都满足刚证得的结论。坐标与标架只有有限多个，可逐次取子列，使全部局部截断分量在 \(L^2\) 中收敛。由于 \(\sum_j\rho_j=1\)，截面等于这些局部截断截面之和；坐标范数与丛的 \(L^2\) 范数在紧支集上等价，故所得子列在全局 \(L^2\) 中收敛。
\end{proof}

下文用到一个完备内积空间的基本事实。这样的空间称为\term[Hilbertkongjian]{Hilbert 空间}[Hilbert space]。\footnote{David Hilbert（1862--1943），德国数学家，研究不变量、几何公理化与积分方程。}若 \(V\) 是其中的闭线性子空间，则每个向量 \(x\) 唯一写成 \(x=v+w\)，其中 \(v\in V\)、\(w\perp V\)。证明可从到 \(V\) 的距离下确界出发：取逼近最小距离的 \(v_n\)，平行四边形恒等式使 \(v_n\) 为 Cauchy 列；闭性给出极限 \(v\)，再对 \(v+t a\)、\(a\in V\)，分别取实数及纯虚数 \(t\) 比较平方距离，得到 \(x-v\perp V\)。这也证明正交投影的范数至多为 \(1\)。

\begin{definition}
	\label{g4:def:0303:fredholm}
	两个 Hilbert 空间之间的有界线性算子 \(T:H\to H'\)，若核有限维、像闭，并且余核 \(H'/\operatorname{im}T\) 有限维，则称为\term[Fredholmsuanzi]{Fredholm 算子}[Fredholm operator]。其指标为 \(\operatorname{ind}T=\dim\ker T-\dim\operatorname{coker}T\)。\footnote{Ivar Fredholm（1866--1927），瑞典数学家，以积分方程研究推动了线性算子理论的发展。}
\end{definition}

\begin{theorem}
	\label{g4:thm:0303:fredholm}
	设 \(X\) 是紧黎曼曲面，\(L\to X\) 是全纯线丛，并固定光滑度量。算子
	\[
		D=\bar\partial_L:H^1_{\mathrm{Sob}}(X,L)
		\longrightarrow L^2(X,\Lambda^{0,1}\mathrm{T}^*X\otimes L)
	\]
	为 Fredholm 算子。其核恰为全纯截面空间 \(H^0(X,L)\)。若
	\[
		\mathscr H^{0,1}(X,L)
		\coloneqq(\operatorname{im}D)^\perp,
	\]
	则这个空间有限维，且其中所有元素均光滑。每个平方可积形式有唯一分解
	\begin{equation}
		\label{g4:eq:0303:orthogonal-decomposition}
		\alpha=Du+\beta,\qquad
		u\in H^1_{\mathrm{Sob}}(X,L),\quad
		u\perp H^0(X,L),\quad
		\beta\in\mathscr H^{0,1}(X,L).
	\end{equation}
	若 \(\alpha\) 光滑，则 \(u\) 也光滑。因此每个 \(H^1(X,\mathcal O(L))\) 类具有唯一的 \(\mathscr H^{0,1}(X,L)\) 代表。
\end{theorem}
\begin{proof}
	先说明弱解的正则性。若 \(u\in H^1_{\mathrm{Sob}}\) 且 \(Du=0\)，则在全纯标架中，系数满足分布方程 \(\partial_{\overline z}u=0\)。\cref{g4:lem:0302:weak-weyl}说明它等于全纯函数，故 \(\ker D=H^0(X,L)\)。更一般地，若 \(Du=\alpha\) 而 \(\alpha\) 光滑，由\cref{g4:thm:0302:local-dbar}，每个小坐标盘上有光滑局部解 \(v\)，满足 \(Dv=\alpha\)。差 \(u-v\) 弱全纯，因此全纯，故 \(u\) 光滑。

	令 \(N=\ker D\)。若 \(N\) 无限维，按 \(L^2\) 内积逐次正交化，可取其中的正交单位列 \(s_n\)。由\cref{g4:prop:0303:global-estimate}，\(\|s_n\|_{H^1_{\mathrm{Sob}}}\le C\)。Rellich 紧嵌入给出 \(L^2\) 收敛子列；但任意不同两项的 \(L^2\) 距离为 \(\sqrt2\)，矛盾。因此 \(N\) 有限维。

	现在证明去掉核后的估计
	\begin{equation}
		\label{g4:eq:0303:coercive}
		\|u\|_{H^1_{\mathrm{Sob}}}\le C_0\|Du\|_1,
		\qquad u\perp N.
	\end{equation}
	若不存在这样的 \(C_0\)，归一化后可取 \(u_n\perp N\)，满足 \(\|u_n\|_{H^1_{\mathrm{Sob}}}=1\)、\(\|Du_n\|_1\to0\)。取其 \(L^2\) 收敛子列。把\eqref{g4:eq:0303:global-estimate}用于 \(u_n-u_m\)，右边两项都趋于零，故该子列实际在 \(H^1_{\mathrm{Sob}}\) 中收敛到 \(u\)。连续性给出 \(Du=0\)，而 \(L^2\) 正交条件给出 \(u\perp N\)，所以 \(u=0\)。这与 \(\|u\|_{H^1_{\mathrm{Sob}}}=1\) 矛盾。

	若 \(Du_n\) 在 \(L^2\) 中收敛，先从每个 \(u_n\) 减去其到有限维空间 \(N\) 的 \(L^2\) 正交投影；这不改变 \(Du_n\)。估计\eqref{g4:eq:0303:coercive}使修改后的 \(u_n\) 在 \(H^1_{\mathrm{Sob}}\) 中为 Cauchy 列。它的极限 \(u\) 满足 \(Du=\lim Du_n\)，所以 \(\operatorname{im}D\) 闭。

	还需控制正交补。设局部 \(\beta=b\,\mathrm{d}\overline z\otimes e\in(\operatorname{im}D)^\perp\)。用支集在该图内的任意光滑截面 \(u e\) 测试，得到
	\[
		0=(D(u e),\beta)_1
		=2\int h\,(\partial_{\overline z}u)\overline b\,
		\mathrm{d}x\,\mathrm{d}y.
	\]
	按弱导数定义，这等价于 \(\partial_{\overline z}(h\overline b)=0\)。弱全纯正则性说明 \(q=h\overline b\) 全纯，故 \(\beta\) 光滑。这些函数还给出全局截面 \(q\,\mathrm{d}z\otimes e^{-1}\)。事实上，若 \(w=\phi(z)\)、\(e'=g e\)，则
	\[
		b'=\frac{b}{g\,\overline{\phi'}},\qquad
		h'=|g|^2h,\qquad
		h'\overline{b'}=\frac{g}{\phi'}h\overline b,
	\]
	正好是 \(K_X\otimes L^{-1}\) 的截面变换律。因此 \(\beta\mapsto q\,\mathrm{d}z\otimes e^{-1}\) 给出到 \(H^0(X,K_X\otimes L^{-1})\) 的反线性单射。刚才证明核有限维的论证适用于任意全纯线丛，特别适用于 \(K_X\otimes L^{-1}\)，故这个目标有限维，正交补也有限维。

	像已闭，Hilbert 空间的正交投影给出 \(L^2=\operatorname{im}D\oplus\mathscr H^{0,1}(X,L)\)。每个像中的元素都有原像，再从原像减去其到 \(N\) 的投影，便得到\eqref{g4:eq:0303:orthogonal-decomposition}；若有两种这样的表示，相减后先由两个正交子空间得到 \(D(u-u')=0\)、\(\beta=\beta'\)，再由 \(u-u'\perp N\) 得到 \(u=u'\)。若 \(\alpha\) 光滑，则 \(\beta\) 已光滑，前面的弱解正则性应用于 \(Du=\alpha-\beta\)，给出 \(u\) 光滑。

	最后，\cref{g4:thm:0302:dolbeault}已将 \(H^1(X,\mathcal O(L))\) 识别为光滑 \((0,1)\) 形式模去 \(D\) 的光滑像。每个这样的形式由上述分解与某个 \(\beta\) 同类；若 \(\beta=Dv\) 且 \(\beta\) 与像正交，则 \(\|\beta\|_1^2=0\)。所以代表存在且唯一。
\end{proof}

紧性在这里有两项具体作用：有限局部化使全局估计只有有限个常数，紧嵌入则排除了无限个彼此正交而仍有统一一阶控制的解。式\eqref{g4:eq:0303:coercive}还说明，选取与全纯截面正交的解后，解随右端在 \(L^2\) 范数中连续变化。

\subsection{积分配对及形式伴随}
\label{g4:subsec:0303:03}
正交分解中留下的 \(\beta\) 满足什么方程？在有限维线性代数中，像的正交补是伴随的核。这里也有相同关系，但伴随须由积分分部计算，度量权函数与微分形式的次数都会进入公式。

\begin{proposition}
	\label{g4:prop:0303:formal-adjoint}
	设 \(X\) 是紧黎曼曲面，\(L\) 带光滑 Hermitian 度量，局部记号采用\eqref{g4:eq:0303:l2-metrics}。算子 \(D=\bar\partial_L\) 的形式伴随，即满足
	\[
		(Du,\beta)_1=(u,D^*\beta)_0
	\]
	的光滑微分算子，局部为
	\begin{equation}
		\label{g4:eq:0303:adjoint}
		D^*(b\,\mathrm{d}\overline z\otimes e)
		=-\frac{2}{h\lambda}\partial_z(hb)\,e.
	\end{equation}
	这里等式两边均使用 \(L^2\) 内积。\cref{g4:thm:0303:fredholm}中的空间 \(\mathscr H^{0,1}(X,L)\) 恰为光滑形式中 \(D^*\) 的核。
\end{proposition}
\begin{proof}
	若 \(u e\) 紧支于一张图，分部积分给出
	\[
		(D(u e),\beta)_1
		=2\int h(\partial_{\overline z}u)\overline b\,\mathrm{d}x\,\mathrm{d}y
		=-2\int u\,\partial_{\overline z}(h\overline b)\,
		\mathrm{d}x\,\mathrm{d}y.
	\]
	而式\eqref{g4:eq:0303:adjoint}代入 \((u e,D^*\beta)_0\) 后恰为右端。用单位分解把任意光滑 \(u\) 写成有限个这样的紧支截面之和，便得到全局等式；积分等式也保证局部伴随表达式在交叠处一致。

	若 \(\beta\in\mathscr H^{0,1}(X,L)\)，前一定理已证明它光滑。它与 \(Du\) 的内积对所有光滑 \(u\) 为零，故与 \(D^*\beta\) 的内积恒零，取 \(u=D^*\beta\) 得到 \(D^*\beta=0\)。反之，\(D^*\beta=0\) 使它与所有光滑 \(Du\) 正交；光滑截面在 \(H^1_{\mathrm{Sob}}\) 中稠密，故它与整个像正交。
\end{proof}

曲面上没有 \((0,2)\) 形式，所以 \((0,1)\) 形式上的 \(\bar\partial\) 复形已经到达末端。相应的非负算子为
\[
	\Delta_{\bar\partial,L}\beta=DD^*\beta,\qquad
	(\Delta_{\bar\partial,L}\beta,\beta)_1=\|D^*\beta\|_0^2.
\]
因而光滑形式满足 \(\Delta_{\bar\partial,L}\beta=0\) 当且仅当 \(D^*\beta=0\)。这样的形式称为\term[tiaohuelingyi]{调和 \((0,1)\) 形式}[harmonic \((0,1)\)-form]。\cref{g4:thm:0303:fredholm}因此给出了每个上同调类的唯一调和代表。度量确定这个代表，上同调类本身不依赖度量。

\ProseParagraphBreak

现在构造不含度量的配对。对光滑 \(L\) 值 \((0,1)\) 形式 \(\alpha\) 与全纯截面 \(\eta\in H^0(X,K_X\otimes L^{-1})\)，先将 \(L^{-1}\) 与 \(L\) 收缩，再按 \(\eta\) 在前的次序作楔积，得到普通复值二形式。定义
\begin{equation}
	\label{g4:eq:0303:serre-pairing}
	B([\alpha],\eta)\coloneqq
	\frac1{2\pi\mathrm{i}}\int_X\eta\wedge\alpha.
\end{equation}
局部若 \(\eta=q\,\mathrm{d}z\otimes e^{-1}\)、\(\alpha=a\,\mathrm{d}\overline z\otimes e\)，则 \(\eta\wedge\alpha=qa\,\mathrm{d}z\wedge\mathrm{d}\overline z\)，标架因子相消。

\begin{proposition}
	\label{g4:prop:0303:pairing-adjoint}
	设 \(X\) 是紧黎曼曲面，\(L\to X\) 为全纯线丛。式\eqref{g4:eq:0303:serre-pairing}是 \(H^1(X,\mathcal O(L))\times H^0(X,K_X\otimes L^{-1})\) 上良定义的复双线性配对。固定度量后，局部公式
	\begin{equation}
		\label{g4:eq:0303:antilinear-identification}
		\mathcal J\bigl(b\,\mathrm{d}\overline z\otimes e\bigr)
		=-h\overline b\,\mathrm{d}z\otimes e^{-1}
	\end{equation}
	给出反线性同构
	\[
		\mathcal J:\mathscr H^{0,1}(X,L)
		\xrightarrow{\ \cong\ }H^0(X,K_X\otimes L^{-1}).
	\]
	对每个光滑 \(\alpha\) 和每个调和 \(\beta\)，有
	\begin{equation}
		\label{g4:eq:0303:pairing-inner-product}
		B([\alpha],\mathcal J\beta)=\frac1{2\pi}(\alpha,\beta)_1.
	\end{equation}
\end{proposition}
\begin{proof}
	若将 \(\alpha\) 改成 \(\alpha+Du\)，由 \(\eta\) 全纯及楔积的次序，有
	\[
		\eta\wedge Du=-\mathrm{d}(u\eta).
	\]
	右边为全局一形式的外微分。第一卷第14.4节定理14.4.3 的 Stokes 公式应用于紧无边界的 \(X\)，给出其积分为零。于是配对只依赖 Dolbeault 类，双线性直接来自积分。

	\cref{g4:thm:0303:fredholm}证明中的坐标计算已经说明 \(h\overline b\,\mathrm{d}z\otimes e^{-1}\) 是全局截面；添负号不改变变换律。式\eqref{g4:eq:0303:adjoint}说明
	\[
		D^*\beta=0
		\quad\Longleftrightarrow\quad
		\partial_z(hb)=0
		\quad\Longleftrightarrow\quad
		\partial_{\overline z}(h\overline b)=0.
	\]
	所以调和形式经 \(\mathcal J\) 变成全纯截面。反过来，若 \(\eta=q\,\mathrm{d}z\otimes e^{-1}\) 全纯，令 \(b=-\overline q/h\)。同一变换律使这些系数粘成全局光滑形式，且上式给出 \(D^*\beta=0\)。这证明满射与单射，复共轭则给出反线性。

	最后，复坐标的正定向使 \(\mathrm{d}z\wedge\mathrm{d}\overline z=-2\mathrm{i}\,\mathrm{d}x\wedge\mathrm{d}y\)。因此
	\[
		\frac1{2\pi\mathrm{i}}\int_X(\mathcal J\beta)\wedge\alpha
		=\frac1\pi\int_X h a\overline b\,\mathrm{d}x\,\mathrm{d}y
		=\frac1{2\pi}(\alpha,\beta)_1.
	\]
	各式可逐图求和，所以得到全局恒等式。公式\eqref{g4:eq:0303:antilinear-identification}中的负号正是为配合这一楔积次序。
\end{proof}

\subsection{配对非退化性}
\label{g4:subsec:0303:04}
积分为零可能来自正负抵消，所以只知道配对存在，还不能断定它能检测所有上同调类。调和代表提供了一个针对给定类选取测试截面的方法：将代表 \(\beta\) 送到 \(\mathcal J\beta\)，积分就成为严格正的平方范数。

\begin{theorem}[曲面上的 Serre 对偶]
	\label{g4:thm:0303:serre-duality}
	设 \(X\) 是紧黎曼曲面，\(L\to X\) 是全纯线丛。积分配对
	\[
		B:H^1(X,\mathcal O(L))\times
		H^0(X,K_X\otimes L^{-1})\longrightarrow\mathbb C,
		\qquad
		B([\alpha],\eta)=\frac1{2\pi\mathrm{i}}\int_X\eta\wedge\alpha
	\]
	在两个有限维复向量空间之间非退化。因此它给出自然的复线性同构
	\begin{equation}
		\label{g4:eq:0303:serre-duality}
		H^1(X,\mathcal O(L))
		\cong H^0(X,K_X\otimes L^{-1})^\vee,
	\end{equation}
	其中 \({}^\vee\) 表示复线性对偶空间。
\end{theorem}
\begin{proof}
	有限性已由\cref{g4:thm:0303:fredholm}对 \(L\) 及 \(K_X\otimes L^{-1}\) 给出。设 \([\alpha]\ne0\)，取其唯一调和代表 \(\beta\)。该代表不为零，故
	\[
		B([\alpha],\mathcal J\beta)
		=B([\beta],\mathcal J\beta)
		=\frac1{2\pi}\|\beta\|_1^2>0.
	\]
	所以第一个变量非退化。若 \(\eta\ne0\)，由\cref{g4:prop:0303:pairing-adjoint}可唯一写成 \(\eta=\mathcal J\beta\)，其中 \(\beta\ne0\) 调和。同样的等式说明 \(B([\beta],\eta)>0\)，所以第二个变量也非退化。

	由双线性，映射 \([\alpha]\mapsto(\eta\mapsto B([\alpha],\eta))\) 是复线性的。两侧非退化分别给出两个有限维空间到对方对偶的单射，故维数相等，这个映射为同构。它由全局积分定义，因此与证明中选择的度量、坐标和代表无关。
\end{proof}

这就是本章所用的\term[Serreduiou]{Serre 对偶}[Serre duality]。这里的对偶是复线性对偶；证明中的 \(\mathcal J\) 是依赖度量的反线性识别。两个陈述的线性类型不同，但由\eqref{g4:eq:0303:pairing-inner-product}准确地联系起来。

\begin{corollary}
	\label{g4:cor:0303:solvability}
	设 \(X\) 是紧黎曼曲面，\(L\to X\) 为全纯线丛，\(\alpha\) 是光滑 \(L\) 值 \((0,1)\) 形式。方程 \(\bar\partial_Lu=\alpha\) 有全局光滑解，当且仅当
	\[
		\int_X\eta\wedge\alpha=0
		\quad\text{对所有 }\eta\in H^0(X,K_X\otimes L^{-1})\text{ 成立}.
	\]
	解的全体在任一解上加 \(H^0(X,L)\) 得到。若记 \(h^q(X,L)=\dim_{\mathbb C}H^q(X,\mathcal O(L))\)，则
	\[
		h^1(X,L)=h^0(X,K_X\otimes L^{-1}),\qquad
		\operatorname{ind}\bar\partial_L=h^0(X,L)-h^1(X,L).
	\]
\end{corollary}
\begin{proof}
	方程可解恰好表示 \([\alpha]=0\)，再用 Serre 配对的非退化性即可。两个解之差在 \(\bar\partial_L\) 的核中，故为全纯截面。维数等式由对偶同构得到；最后一个等式由 Fredholm 指标的定义和调和代表同构得到。
\end{proof}

\begin{example}
	\label{g4:exm:0303:flat-torus}
	取方形复环面 \(X=\mathbb C/(2\pi\mathbb Z+2\pi\mathrm{i}\mathbb Z)\)，并令 \(L\) 为平凡线丛，\(h=\lambda=1\)。形式 \(\mathrm{d}z\) 与 \(\mathrm{d}\overline z\) 都在平移商上有定义。任一全纯微分除以处处非零的 \(\mathrm{d}z\)，得到全纯函数；\cref{g4:prop:0201:meromorphic-map}使它为常数。因此 \(H^0(X,K_X)=\mathbb C\,\mathrm{d}z\)，而伴随公式给出 \(\mathscr H^{0,1}(X,\mathcal O_X)=\mathbb C\,\mathrm{d}\overline z\)。

	若 \(\alpha=a(x,y)\,\mathrm{d}\overline z\)，其唯一调和代表为
	\[
		\beta=\left(\frac1{4\pi^2}
		\int_{[0,2\pi]^2}a(x,y)\,\mathrm{d}x\,\mathrm{d}y\right)
		\mathrm{d}\overline z.
	\]
	这正是常数函数方向上的正交投影。配对直接计算为
	\[
		B([\alpha],\mathrm{d}z)
		=-\frac1\pi\int_{[0,2\pi]^2}a(x,y)\,
		\mathrm{d}x\,\mathrm{d}y.
	\]
	所以 \(\bar\partial u=\alpha\) 可解，当且仅当 \(a\) 的平均值为零。若平均值为零，以 \(e^{\mathrm{i}(k_1x+k_2y)}\) 展开，唯一平均值为零的解具有系数
	\[
		\widehat u(0)=0,\qquad
		\widehat u(k)=\frac{2\widehat a(k)}{\mathrm{i}k_1-k_2},
		\quad k\ne0.
	\]
	光滑 \(a\) 的系数经任意次分部积分后按任意多项式阶衰减，分母对 \(k\ne0\) 的模为 \(|k|\)，故解级数及其任意阶导数一致收敛。这也直接验证了光滑解的存在。
\end{example}

\subsection{与留数配对的比较}
\label{g4:subsec:0303:05}
主部问题把局部极点数据作为输入：给定有限多个点附近的亚纯截面，要求找到一个全局亚纯截面，恰有这些主部。截断函数把这一问题转成 \(\bar\partial\) 方程；Serre 配对随后把全局积分重新集中到各极点处的留数。

\ProseParagraphBreak

设 \(S=\{p_1,\ldots,p_m\}\subset X\)。在两两不交的坐标圆盘 \(U_j\) 上选全纯标架 \(e_j\) 和中心坐标 \(z_j(p_j)=0\)，把给定主部写成
\[
	P_j=\left(\sum_{\nu=1}^{N_j}c_{j,\nu}z_j^{-\nu}\right)e_j.
\]
取 \(0\le\chi_j\le1\)，使 \(\chi_j\) 支持在 \(U_j\) 内，且在 \(p_j\) 的一个邻域恒等于 \(1\)。在 \(X\setminus S\) 上令 \(s=\sum_j\chi_jP_j\)，并在这个去点曲面上按普通微分计算 \(\alpha=Ds\)。由于 \(P_j\) 在去心盘上全纯，且 \(\chi_j\) 在更小盘中恒为 \(1\)，\(\alpha\) 在各点的去心邻域为零，故可补零成为 \(X\) 上的光滑 \((0,1)\) 形式。这里的 \(\alpha\) 指这一光滑延拓，不表示含极点截面的全局分布导数。例如\cref{g4:thm:0302:local-dbar}证明中的核恒等式给出 \(\partial_{\overline z}(1/z)=\pi\delta_0\)，而 \(1/z\) 在去心盘上的普通导数为零。

\ProseParagraphBreak

按上一节的约定，Čech 差分为 \((\delta b)_{ij}=b_j-b_i\)，局部亚纯提升 \(s_i\) 给出主部的连接余圈 \(c_{ij}=s_j-s_i\)，而比较映射将它送到 \([D(\sum_i\rho_i s_i)]\)。在当前的圆盘覆盖中，取局部提升为各 \(P_j\) 及 \(X\setminus S\) 上的零，便由\cref{g4:prop:0302:principal-parts-class}得到上述 \([\alpha]\)。记这个连接类为 \(\delta P\in H^1(X,\mathcal O(L))\)。改变截断函数只使 \(s\) 改变一个全局光滑截面；将 \(P_j\) 改变一个局部全纯截面也有同样效果。因此 \(\delta P\) 只依赖主部，不依赖这些选择。

\begin{proposition}
	\label{g4:prop:0303:residue-pairing}
	设 \(X\) 是紧黎曼曲面，\(L\to X\) 为全纯线丛，\(P=(P_1,\ldots,P_m)\) 是有限点集 \(S\) 处给定的 \(L\) 值主部，并用上述约定定义连接类 \(\delta P=[\alpha]\)。对任意 \(\eta\in H^0(X,K_X\otimes L^{-1})\)，有
	\begin{equation}
		\label{g4:eq:0303:residue-pairing}
		B(\delta P,\eta)
		=\frac1{2\pi\mathrm{i}}\int_X\eta\wedge\alpha
		=\sum_{j=1}^m\operatorname{Res}_{p_j}(P_j\eta).
	\end{equation}
	乘积 \(P_j\eta\) 是普通亚纯微分，其留数与坐标、标架及主部代表的选择无关。
\end{proposition}
\begin{proof}
	标架 \(e_j\) 与 \(e_j^{-1}\) 在乘积中相消，故 \(P_j\eta\) 是亚纯一形式。其留数的坐标无关性由\cref{g4:prop:0203:residue-invariance}给出；主部代表增加全纯截面时，乘积只增加全纯微分，不改变留数。

	只须计算一个圆盘内的贡献。取 \(\varepsilon>0\) 足够小，使 \(\chi_j=1\) 在 \(|z_j|\le\varepsilon\) 附近，并取较大的坐标圆周，使 \(\chi_j\) 在该圆周附近为零。令 \(A_j\) 为两圆周之间的闭圆环。其内的 \(s_j=\chi_jP_j\) 光滑，且 \(\eta\) 全纯，所以
	\[
		\eta\wedge Ds_j=-\mathrm{d}(s_j\eta).
	\]
	圆环外边界按逆时针定向，内边界按顺时针定向。Stokes 公式因此给出
	\[
		\int_{U_j}\eta\wedge D(\chi_jP_j)
		=-\int_{\partial A_j}s_j\eta
		=\int_{|z_j|=\varepsilon}^{\text{逆时针}}P_j\eta
		=2\pi\mathrm{i}\operatorname{Res}_{p_j}(P_j\eta).
	\]
	第一处积分按前述光滑延拓理解，因其支集位于圆环中而等于圆环上的积分。中间等式的正号来自两个负号：一个来自外微分公式，一个来自内边界方向。对 \(j\) 求和并除以 \(2\pi\mathrm{i}\)，即得所需公式。
\end{proof}

\begin{corollary}
	\label{g4:cor:0303:principal-parts-criterion}
	设 \(X\) 是紧黎曼曲面，\(L\to X\) 为全纯线丛。在有限点集 \(S\) 给定 \(L\) 值主部 \(P_j\)。存在一个全局亚纯截面，在 \(S\) 以外全纯，且在每个 \(p_j\) 具有指定主部，当且仅当
	\[
		\sum_j\operatorname{Res}_{p_j}(P_j\eta)=0
		\quad\text{对所有 }\eta\in H^0(X,K_X\otimes L^{-1})\text{ 成立}.
	\]
\end{corollary}
\begin{proof}
	\cref{g4:prop:0302:principal-parts-class}已证明：指定主部可由全局亚纯截面实现，当且仅当连接类 \(\delta P\) 为零。由 Serre 配对的非退化性，这又等价于 \(B(\delta P,\eta)=0\) 对所有 \(\eta\) 成立。将\cref{g4:prop:0303:residue-pairing}代入，即得所列留数条件。
\end{proof}

例如，在复环面上取平凡线丛，并用来自平面的局部平移坐标给出简单极点主部 \(c_j/z_j\)。全纯微分空间由 \(\mathrm{d}z\) 张成，故条件恰为 \(\sum_jc_j=0\)。一个非零简单极点主部不能单独实现；在两个不同点指定互为相反数的简单主部，则存在实现它们的亚纯函数。这里的留数是 \((c_j/z_j)\,\mathrm{d}z_j\) 的留数，正是积分配对检测到的局部系数。下一节将通过逐点增加可允许的主部，计算这些有限维空间的维数。
% 纲要条目排为编号小节标题，并进入目录与书签；说明段仍保留于纲要注释中。

% !TeX root = ../../Book4.tex
% 纲要标题：### 3.4 Riemann–Roch 的推导
% 纲要位置：几何学计划纲要.md，第 3571 行。
% 正文状态：未撰写；本文件目前是框架。

% 纲要内容开始
% 
% * 线丛 Euler 示性数的有限性
% * 增加一个点的正合列
% * 亚纯截面的存在而非预设
% * 由除子约化得到次数公式
% * 解析亏格与拓扑亏格
% 
% 纲要内容结束

\section{Riemann–Roch 的推导}
\label{g4:sec:0304}

待撰写

\subsection{线丛 Euler 示性数的有限性}
\label{g4:subsec:0304:01}
待撰写

\subsection{增加一个点的正合列}
\label{g4:subsec:0304:02}
待撰写

\subsection{亚纯截面的存在而非预设}
\label{g4:subsec:0304:03}
待撰写

\subsection{由除子约化得到次数公式}
\label{g4:subsec:0304:04}
待撰写

\subsection{解析亏格与拓扑亏格}
\label{g4:subsec:0304:05}
待撰写
% 纲要条目排为编号小节标题，并进入目录与书签；说明段仍保留于纲要注释中。

% !TeX root = ../../Book4.tex
% 纲要标题：### 3.5 线性系的实际应用
% 纲要位置：几何学计划纲要.md，第 3579 行。
% 正文状态：未撰写；本文件目前是框架。

% 纲要内容开始
% 
% * 高次除子的截面和基点
% * 分离点与切向量的次数条件
% * 典范映射和特殊除子
% * 射影模型及低亏格曲线
% 
% ---
% 
% 纲要内容结束

\section{线性系的实际应用}
\label{g4:sec:0305}

待撰写

\subsection{高次除子的截面和基点}
\label{g4:subsec:0305:01}
待撰写

\subsection{分离点与切向量的次数条件}
\label{g4:subsec:0305:02}
待撰写

\subsection{典范映射和特殊除子}
\label{g4:subsec:0305:03}
待撰写

\subsection{射影模型及低亏格曲线}
\label{g4:subsec:0305:04}
待撰写
% 纲要条目排为编号小节标题，并进入目录与书签；说明段仍保留于纲要注释中。


\begin{chaptersummary}
待撰写
\end{chaptersummary}

\BookExercises
待撰写

% !TeX root = ../../Book4.tex
% 纲要标题：## 第4章　周期、Abel 定理与 Jacobi 簇
% 纲要位置：几何学计划纲要.md，第 3588 行。
% 正文状态：未撰写；本文件目前是框架。

% 纲要内容开始
% 
% > 本章目标与证明范围：由具体周期积分构造 Jacobi 簇，再证明主除子的判据和反演。椭圆函数是完整低亏格模型，不仅作为一条跨学科联系。
% 
% 纲要内容结束

\chapter{周期、Abel 定理与 Jacobi 簇}
\label{g4:ch:04}
\BookChapterNumber{g4:ch:04}

\begin{chapterabstract}
待撰写
\end{chapterabstract}

\begin{learninggoals}
待撰写
\end{learninggoals}

待撰写

% !TeX root = ../../Book4.tex
% 纲要标题：### 4.1 微分的周期
% 纲要位置：几何学计划纲要.md，第 3592 行。
% 正文状态：未撰写；本文件目前是框架。

% 纲要内容开始
% 
% * 同调辛基与周期矩阵
% * Riemann 双线性关系
% * 归一化全纯微分
% * 周期格的离散性与正性
% 
% 纲要内容结束

\section{微分的周期}
\label{g4:sec:0401}

待撰写

\subsection{同调辛基与周期矩阵}
\label{g4:subsec:0401:01}
待撰写

\subsection{Riemann 双线性关系}
\label{g4:subsec:0401:02}
待撰写

\subsection{归一化全纯微分}
\label{g4:subsec:0401:03}
待撰写

\subsection{周期格的离散性与正性}
\label{g4:subsec:0401:04}
待撰写
% 纲要条目排为编号小节标题，并进入目录与书签；说明段仍保留于纲要注释中。

% !TeX root = ../../Book4.tex
% 纲要标题：### 4.2 Jacobi 簇和 Abel–Jacobi 映射
% 纲要位置：几何学计划纲要.md，第 3382 行。
% 正文状态：未撰写；本文件目前是框架。

% 纲要内容开始
% 
% * 复向量空间模周期格
% * 路径积分与基点依赖
% * 对称积上的映射
% * 除子类的几何解释
% 
% 纲要内容结束

\section{Jacobi 簇和 Abel–Jacobi 映射}
\label{g4:sec:0402}

待撰写

\subsection{复向量空间模周期格}
\label{g4:subsec:0402:01}
待撰写

\subsection{路径积分与基点依赖}
\label{g4:subsec:0402:02}
待撰写

\subsection{对称积上的映射}
\label{g4:subsec:0402:03}
待撰写

\subsection{除子类的几何解释}
\label{g4:subsec:0402:04}
待撰写
% 纲要条目排为编号小节标题，并进入目录与书签；说明段仍保留于纲要注释中。

% !TeX root = ../../Book4.tex
% 纲要标题：### 4.3 Abel 定理的证明
% 纲要位置：几何学计划纲要.md，第 3606 行。
% 正文状态：未撰写；本文件目前是框架。

% 纲要内容开始
% 
% * 主除子的必要周期条件
% * 第三类微分的构造
% * 整数周期与指数积分
% * 由周期条件恢复亚纯函数
% 
% 纲要内容结束

\section{Abel 定理的证明}
\label{g4:sec:0403}

待撰写

\subsection{主除子的必要周期条件}
\label{g4:subsec:0403:01}
待撰写

\subsection{第三类微分的构造}
\label{g4:subsec:0403:02}
待撰写

\subsection{整数周期与指数积分}
\label{g4:subsec:0403:03}
待撰写

\subsection{由周期条件恢复亚纯函数}
\label{g4:subsec:0403:04}
待撰写
% 纲要条目排为编号小节标题，并进入目录与书签；说明段仍保留于纲要注释中。

% !TeX root = ../../Book4.tex
% 纲要标题：### 4.4 Jacobi 反演及 theta 几何
% 纲要位置：几何学计划纲要.md，第 3613 行。
% 正文状态：未撰写；本文件目前是框架。

% 纲要内容开始
% 
% * 反演的满射性与 Riemann–Roch
% * 特殊除子及纤维维数
% * theta 除子和主极化
% * 进一步 Torelli 问题的准确位置
% 
% 纲要内容结束

\section{Jacobi 反演及 theta 几何}
\label{g4:sec:0404}

待撰写

\subsection{反演的满射性与 Riemann–Roch}
\label{g4:subsec:0404:01}
待撰写

\subsection{特殊除子及纤维维数}
\label{g4:subsec:0404:02}
待撰写

\subsection{theta 除子和主极化}
\label{g4:subsec:0404:03}
待撰写

\subsection{进一步 Torelli 问题的准确位置}
\label{g4:subsec:0404:04}
待撰写
% 纲要条目排为编号小节标题，并进入目录与书签；说明段仍保留于纲要注释中。

% !TeX root = ../../Book4.tex
% 纲要标题：### 4.5 椭圆曲线的完整模型
% 纲要位置：几何学计划纲要.md，第 3403 行。
% 正文状态：未撰写；本文件目前是框架。

% 纲要内容开始
% 
% * Weierstrass 函数及微分方程
% * 复环面到光滑三次曲线
% * 加法公式与几何群律
% * 格变换、j 参数和退化
% 
% ---
% 
% 纲要内容结束

\section{椭圆曲线的完整模型}
\label{g4:sec:0405}

待撰写

\subsection{Weierstrass 函数及微分方程}
\label{g4:subsec:0405:01}
待撰写

\subsection{复环面到光滑三次曲线}
\label{g4:subsec:0405:02}
待撰写

\subsection{加法公式与几何群律}
\label{g4:subsec:0405:03}
待撰写

\subsection{格变换、j 参数和退化}
\label{g4:subsec:0405:04}
待撰写
% 纲要条目排为编号小节标题，并进入目录与书签；说明段仍保留于纲要注释中。


\begin{chaptersummary}
待撰写
\end{chaptersummary}

\BookExercises
待撰写

% !TeX root = ../../Book4.tex
% 纲要标题：## 第5章　势论、统一化与双曲曲面
% 纲要位置：几何学计划纲要.md，第 3412 行。
% 正文状态：未撰写；本文件目前是框架。

% 纲要内容开始
% 
% > 本章目标与证明范围：为一般统一化保留独立而连续的证明模块：先做带边界曲面上的势论，再用拓扑穷竭及正规族处理单连通非紧曲面。不能只引用平面映照定理后宣称所有曲面均已统一化。
% 
% 纲要内容结束

\chapter{势论、统一化与双曲曲面}
\label{g4:ch:05}
\BookChapterNumber{g4:ch:05}

\begin{chapterabstract}
待撰写
\end{chapterabstract}

\begin{learninggoals}
待撰写
\end{learninggoals}

待撰写

% !TeX root = ../../Book4.tex
% 纲要标题：### 5.1 曲面上的势论工具
% 纲要位置：几何学计划纲要.md，第 3416 行。
% 正文状态：未撰写；本文件目前是框架。

% 纲要内容开始
% 
% * 辅助共形度量和调和性
% * 光滑相对紧区域的 Dirichlet 问题
% * Green 函数及对数奇点
% * 调和共轭、周期和能量
% * 最大值及 Harnack 紧性
% 
% 纲要内容结束

\section{曲面上的势论工具}
\label{g4:sec:0501}

待撰写

\subsection{辅助共形度量和调和性}
\label{g4:subsec:0501:01}
待撰写

\subsection{光滑相对紧区域的 Dirichlet 问题}
\label{g4:subsec:0501:02}
待撰写

\subsection{Green 函数及对数奇点}
\label{g4:subsec:0501:03}
待撰写

\subsection{调和共轭、周期和能量}
\label{g4:subsec:0501:04}
待撰写

\subsection{最大值及 Harnack 紧性}
\label{g4:subsec:0501:05}
待撰写
% 纲要条目排为编号小节标题，并进入目录与书签；说明段仍保留于纲要注释中。

% !TeX root = ../../Book4.tex
% 纲要标题：### 5.2 带边界单连通区域的映照
% 纲要位置：几何学计划纲要.md，第 3424 行。
% 正文状态：未撰写；本文件目前是框架。

% 纲要内容开始
% 
% * 拓扑圆盘和边界正则化
% * 由 Green 函数构造全纯映射
% * 边界模长与映射次数
% * 单叶性和圆盘映照
% 
% 纲要内容结束

\section{带边界单连通区域的映照}
\label{g4:sec:0502}

待撰写

\subsection{拓扑圆盘和边界正则化}
\label{g4:subsec:0502:01}
待撰写

\subsection{由 Green 函数构造全纯映射}
\label{g4:subsec:0502:02}
待撰写

\subsection{边界模长与映射次数}
\label{g4:subsec:0502:03}
待撰写

\subsection{单叶性和圆盘映照}
\label{g4:subsec:0502:04}
待撰写
% 纲要条目排为编号小节标题，并进入目录与书签；说明段仍保留于纲要注释中。

% !TeX root = ../../Book4.tex
% 纲要标题：### 5.3 一般单连通曲面的统一化
% 纲要位置：几何学计划纲要.md，第 3648 行。
% 正文状态：未撰写；本文件目前是框架。

% 纲要内容开始
% 
% * 非紧单连通曲面的适应穷竭引理
% 
%   利用相对紧曲面区域、边界曲线的填洞及单连通性建立圆盘穷竭；本引理在本章证明，不假定已掌握一般非紧曲面分类。
% 
% * 相容区域的归一化映照
% * 正规族极限及非退化性
% 
%   单列归一化导数、单叶函数畸变控制以及逆映射紧性的引理；分别排除常值极限并验证极限像，不以“由正规族”略过核心存在步骤。
% 
% * 双曲型与抛物型的区分
% * 紧单连通情形和黎曼球面
% * 从万有覆盖得到一般曲面的商描述
% 
% 纲要内容结束

\section{一般单连通曲面的统一化}
\label{g4:sec:0503}

待撰写

\subsection{非紧单连通曲面的适应穷竭引理}
\label{g4:subsec:0503:01}
待撰写

\subsection{相容区域的归一化映照}
\label{g4:subsec:0503:02}
待撰写

\subsection{正规族极限及非退化性}
\label{g4:subsec:0503:03}
待撰写

\subsection{双曲型与抛物型的区分}
\label{g4:subsec:0503:04}
待撰写

\subsection{紧单连通情形和黎曼球面}
\label{g4:subsec:0503:05}
待撰写

\subsection{从万有覆盖得到一般曲面的商描述}
\label{g4:subsec:0503:06}
待撰写
% 纲要条目排为编号小节标题，并进入目录与书签；说明段仍保留于纲要注释中。

% !TeX root = ../../Book4.tex
% 纲要标题：### 5.4 共形度量与双曲收缩
% 纲要位置：几何学计划纲要.md，第 3446 行。
% 正文状态：未撰写；本文件目前是框架。

% 纲要内容开始
% 
% * 曲率的共形变换
% * Schwarz–Pick 和 Ahlfors–Schwarz
% * Poincaré 度量及完备性
% * 尖点、环域与漏斗
% 
% 纲要内容结束

\section{共形度量与双曲收缩}
\label{g4:sec:0504}

待撰写

\subsection{曲率的共形变换}
\label{g4:subsec:0504:01}
待撰写

\subsection{Schwarz–Pick 和 Ahlfors–Schwarz}
\label{g4:subsec:0504:02}
待撰写

\subsection{Poincaré 度量及完备性}
\label{g4:subsec:0504:03}
待撰写

\subsection{尖点、环域与漏斗}
\label{g4:subsec:0504:04}
待撰写
% 纲要条目排为编号小节标题，并进入目录与书签；说明段仍保留于纲要注释中。

% !TeX root = ../../Book4.tex
% 纲要标题：### 5.5 极值长度和模量
% 纲要位置：几何学计划纲要.md，第 3453 行。
% 正文状态：未撰写；本文件目前是框架。

% 纲要内容开始
% 
% * 曲线族的极值长度
% * 共形不变性与单调性
% * 环域模量的计算
% * 退化曲面的长颈模型
% 
% ---
% 
% 纲要内容结束

\section{极值长度和模量}
\label{g4:sec:0505}

待撰写

\subsection{曲线族的极值长度}
\label{g4:subsec:0505:01}
待撰写

\subsection{共形不变性与单调性}
\label{g4:subsec:0505:02}
待撰写

\subsection{环域模量的计算}
\label{g4:subsec:0505:03}
待撰写

\subsection{退化曲面的长颈模型}
\label{g4:subsec:0505:04}
待撰写
% 纲要条目排为编号小节标题，并进入目录与书签；说明段仍保留于纲要注释中。


\begin{chaptersummary}
待撰写
\end{chaptersummary}

\BookExercises
待撰写

% !TeX root = ../../Book4.tex
% 纲要标题：## 第6章*　Teichmüller 空间与曲线的形变
% 纲要位置：几何学计划纲要.md，第 3462 行。
% 正文状态：未撰写；本文件目前是框架。

% 纲要内容开始
% 
% > 本章目标与证明范围：将曲线模参数从高维复形变中分离。无穷小空间和实际局部参数的联系给明确证明任务；一般可测 Riemann 映照及全局 Teichmüller 深层结果作为逐项外部输入。
% 
% 纲要内容结束

\OptionalChapter{Teichmüller 空间与曲线的形变}{g4:ch:06}
\BookChapterNumber{g4:ch:06}

\begin{chapterabstract}
待撰写
\end{chapterabstract}

\begin{learninggoals}
待撰写
\end{learninggoals}

待撰写

% !TeX root = ../../Book4.tex
% 纲要标题：### 6.1 标记与等价关系
% 纲要位置：几何学计划纲要.md，第 3683 行。
% 正文状态：未撰写；本文件目前是框架。
% 本次扩充的范围、先修与边界：先修为 I.6 的闭曲面分类、II.2 的一阶同调及闭曲面的同调交叉配对，并使用本节已有的标记、同位和群作用。固定闭连通定向曲面，构造 Dehn 扭转并核对方向和交叉数符号，计算其同调作用，将 Torelli 群定义为作用的核。Nielsen–Thurston 三分法限亏格至少二的精确陈述与代表模型；不承担完整分类证明、映射类群的全部生成关系或 Torelli 群的专门结构理论。

% 纲要内容开始
% 
% > 补充写作范围与外部输入：先修为 I.6 的闭曲面分类、II.2 的一阶同调及闭曲面的同调交叉配对，并使用本节已有的标记、同位和群作用。固定闭连通定向曲面，构造 Dehn 扭转并核对方向和交叉数符号，计算其同调作用，将 Torelli 群定义为作用的核。Nielsen–Thurston 三分法限亏格至少二的精确陈述与代表模型；不承担完整分类证明、映射类群的全部生成关系或 Torelli 群的专门结构理论。
% 
% * 标记曲面、同位和双全纯
% * Teichmüller 空间与模空间的区别
% * 映射类群作用及稳定子
% 
% * Dehn 扭转的环带构造与同位类
% * Dehn 扭转对一阶同调的作用
% * Torelli 群作为同调作用的核
% * Nielsen–Thurston 三分法的陈述与模型
% 
% 
% 纲要内容结束

\section{标记与等价关系}
\label{g4:sec:0601}

待撰写

\subsection{标记曲面、同位和双全纯}
\label{g4:subsec:0601:01}
待撰写

\subsection{Teichmüller 空间与模空间的区别}
\label{g4:subsec:0601:02}
待撰写

\subsection{映射类群作用及稳定子}
\label{g4:subsec:0601:03}
待撰写

\subsection{Dehn 扭转的环带构造与同位类}
\label{g4:subsec:0601:04}
待撰写

\subsection{Dehn 扭转对一阶同调的作用}
\label{g4:subsec:0601:05}
待撰写

\subsection{Torelli 群作为同调作用的核}
\label{g4:subsec:0601:06}
待撰写

\subsection{Nielsen–Thurston 三分法的陈述与模型}
\label{g4:subsec:0601:07}
待撰写

% 纲要条目排为编号小节标题，并进入目录与书签；说明段仍保留于纲要注释中。

% !TeX root = ../../Book4.tex
% 纲要标题：### 6.2 Beltrami 方程与复结构
% 纲要位置：几何学计划纲要.md，第 3472 行。
% 正文状态：未撰写；本文件目前是框架。

% 纲要内容开始
% 
% * Beltrami 微分和局部变形
% * 光滑小参数解的构造
% * 可测 Riemann 映照定理的准确引用
% * 复结构变化与坐标变化
% 
% 纲要内容结束

\section{Beltrami 方程与复结构}
\label{g4:sec:0602}

待撰写

\subsection{Beltrami 微分和局部变形}
\label{g4:subsec:0602:01}
待撰写

\subsection{光滑小参数解的构造}
\label{g4:subsec:0602:02}
待撰写

\subsection{可测 Riemann 映照定理的准确引用}
\label{g4:subsec:0602:03}
待撰写

\subsection{复结构变化与坐标变化}
\label{g4:subsec:0602:04}
待撰写
% 纲要条目排为编号小节标题，并进入目录与书签；说明段仍保留于纲要注释中。

% !TeX root = ../../Book4.tex
% 纲要标题：### 6.3 切空间、余切空间与维数
% 纲要位置：几何学计划纲要.md，第 3704 行。
% 正文状态：未撰写；本文件目前是框架。

% 纲要内容开始
% 
% * 平凡形变和向量场
% * Beltrami 类与上同调
% * 二次微分及 Serre 对偶
% * 稳定亏格情形的维数计算
% 
% 纲要内容结束

\section{切空间、余切空间与维数}
\label{g4:sec:0603}

待撰写

\subsection{平凡形变和向量场}
\label{g4:subsec:0603:01}
待撰写

\subsection{Beltrami 类与上同调}
\label{g4:subsec:0603:02}
待撰写

\subsection{二次微分及 Serre 对偶}
\label{g4:subsec:0603:03}
待撰写

\subsection{稳定亏格情形的维数计算}
\label{g4:subsec:0603:04}
待撰写
% 纲要条目排为编号小节标题，并进入目录与书签；说明段仍保留于纲要注释中。

% !TeX root = ../../Book4.tex
% 纲要标题：### 6.4 实际参数与全局几何
% 纲要位置：几何学计划纲要.md，第 3486 行。
% 正文状态：未撰写；本文件目前是框架。

% 纲要内容开始
% 
% * 局部坐标及自同构的影响
% * 极值长度和 Teichmüller 度量
% * Fenchel–Nielsen 参数的模型
% * 退化、节点和紧化的入口
% 
% 纲要内容结束

\section{实际参数与全局几何}
\label{g4:sec:0604}

待撰写

\subsection{局部坐标及自同构的影响}
\label{g4:subsec:0604:01}
待撰写

\subsection{极值长度和 Teichmüller 度量}
\label{g4:subsec:0604:02}
待撰写

\subsection{Fenchel–Nielsen 参数的模型}
\label{g4:subsec:0604:03}
待撰写

\subsection{退化、节点和紧化的入口}
\label{g4:subsec:0604:04}
待撰写
% 纲要条目排为编号小节标题，并进入目录与书签；说明段仍保留于纲要注释中。

% !TeX root = ../../Book4.tex
% 纲要标题：### 6.5 四点标记球面的参数与退化
% 纲要位置：几何学计划纲要.md，第 3493 行。
% 正文状态：未撰写；本文件目前是框架。

% 纲要内容开始
% 
% 这里完成解析参数和三类退化的识别；跨任意代数基底的通用性证明在 V.34.6 进行。交比参数化的是模空间，不把它等同于其带标记的 Teichmüller 覆盖。
% 
% * 有序四点、Möbius 归一化与交比参数
% * $\mathbb P^1\setminus\{0,1,\infty\}$ 上的光滑四点族
% * 标记碰撞、长颈与结点极限
% * 两两分配标记的三种稳定边界类型
% * 有序标记与置换商、自同构的区别
% * 与 V.34.6 的代数通用族及完整紧化构造比较
% 
% ---
% 
% 纲要内容结束

\section{四点标记球面的参数与退化}
\label{g4:sec:0605}

待撰写

\subsection{有序四点、Möbius 归一化与交比参数}
\label{g4:subsec:0605:01}
待撰写

\subsection{\texorpdfstring{\(\mathbb P^1\setminus\{0,1,\infty\}\) 上的光滑四点族}{P¹∖\{0,1,∞\} 上的光滑四点族}}
\label{g4:subsec:0605:02}
待撰写

\subsection{标记碰撞、长颈与结点极限}
\label{g4:subsec:0605:03}
待撰写

\subsection{两两分配标记的三种稳定边界类型}
\label{g4:subsec:0605:04}
待撰写

\subsection{有序标记与置换商、自同构的区别}
\label{g4:subsec:0605:05}
待撰写

\subsection{与 V.34.6 的代数通用族及完整紧化构造比较}
\label{g4:subsec:0605:06}
待撰写
% 纲要条目排为编号小节标题，并进入目录与书签；说明段仍保留于纲要注释中。


\begin{chaptersummary}
待撰写
\end{chaptersummary}

\BookExercises
待撰写


% !TeX root = ../Book4.tex
% 纲要标题：# 第二编　辛结构、Hamilton 系统与环面几何
% 纲要位置：几何学计划纲要.md，第 3506 行。
% 正文状态：未撰写；本文件目前是框架。

% 纲要内容开始
% 
% 纲要内容结束

\BookPart{辛结构、Hamilton 系统与环面几何}{g4:part:02}
\BookPartGuide
待撰写

% !TeX root = ../../Book4.tex
% 纲要标题：## 第7章　线性辛几何、余切丛与相容结构
% 纲要位置：几何学计划纲要.md，第 3508 行。
% 正文状态：未撰写；本文件目前是框架。

% 纲要内容开始
% 
% > 本章目标与证明范围：在辛基础中建立相容概复结构的存在和可缩性，不要求先经过全纯曲线或 Kähler。闭性、非退化性和可积性始终是不同条件。
% 
% 纲要内容结束

\chapter{线性辛几何、余切丛与相容结构}
\label{g4:ch:07}
\BookChapterNumber{g4:ch:07}

\begin{chapterabstract}
待撰写
\end{chapterabstract}

\begin{learninggoals}
待撰写
\end{learninggoals}

待撰写

% !TeX root = ../../Book4.tex
% 纲要标题：### 7.1 线性辛空间
% 纲要位置：几何学计划纲要.md，第 3745 行。
% 正文状态：未撰写；本文件目前是框架。

% 纲要内容开始
% 
% * 辛基与交替型标准形
% * 各向同性、余各向同性和 Lagrange
% * 辛约化的线性模型
% * 辛群及其基本拓扑
% 
% 纲要内容结束

\section{线性辛空间}
\label{g4:sec:0701}

待撰写

\subsection{辛基与交替型标准形}
\label{g4:subsec:0701:01}
待撰写

\subsection{各向同性、余各向同性和 Lagrange}
\label{g4:subsec:0701:02}
待撰写

\subsection{辛约化的线性模型}
\label{g4:subsec:0701:03}
待撰写

\subsection{辛群及其基本拓扑}
\label{g4:subsec:0701:04}
待撰写
% 纲要条目排为编号小节标题，并进入目录与书签；说明段仍保留于纲要注释中。

% !TeX root = ../../Book4.tex
% 纲要标题：### 7.2 余切丛与典范构造
% 纲要位置：几何学计划纲要.md，第 3752 行。
% 正文状态：未撰写；本文件目前是框架。

% 纲要内容开始
% 
% * 典范一形式和坐标无关性
% * ω=−dθ 的符号
% * 余切提升与辛不变性
% * 闭一形式的 Lagrange 图
% 
% 纲要内容结束

\section{余切丛与典范构造}
\label{g4:sec:0702}

待撰写

\subsection{典范一形式和坐标无关性}
\label{g4:subsec:0702:01}
待撰写

\subsection{\texorpdfstring{\(\omega=-\mathrm{d}\theta\) 的符号}{\GeometryPdfSmallOmega{}=−d\GeometryPdfTheta{} 的符号}}
\label{g4:subsec:0702:02}
待撰写

\subsection{余切提升与辛不变性}
\label{g4:subsec:0702:03}
待撰写

\subsection{闭一形式的 Lagrange 图}
\label{g4:subsec:0702:04}
待撰写
% 纲要条目排为编号小节标题，并进入目录与书签；说明段仍保留于纲要注释中。

% !TeX root = ../../Book4.tex
% 纲要标题：### 7.3 相容三元组
% 纲要位置：几何学计划纲要.md，第 3526 行。
% 正文状态：未撰写；本文件目前是框架。

% 纲要内容开始
% 
% * 由 ω 与辅助内积构造 J
% * 正定配对 g(u,v)=ω(u,Jv)
% * 相容结构空间的可缩性
% * 光滑参数族与流形上的整体构造
% 
% 纲要内容结束

\section{相容三元组}
\label{g4:sec:0703}

待撰写

\subsection{\texorpdfstring{由 \(\omega\) 与辅助内积构造 J}{由 \GeometryPdfSmallOmega{} 与辅助内积构造 J}}
\label{g4:subsec:0703:01}
待撰写

\subsection{\texorpdfstring{正定配对 \(g(u,v)=\omega(u,Jv)\)}{正定配对 g(u,v)=\GeometryPdfSmallOmega{}(u,Jv)}}
\label{g4:subsec:0703:02}
待撰写

\subsection{相容结构空间的可缩性}
\label{g4:subsec:0703:03}
待撰写

\subsection{光滑参数族与流形上的整体构造}
\label{g4:subsec:0703:04}
待撰写
% 纲要条目排为编号小节标题，并进入目录与书签；说明段仍保留于纲要注释中。

% !TeX root = ../../Book4.tex
% 纲要标题：### 7.4 辛流形的整体约束
% 纲要位置：几何学计划纲要.md，第 3533 行。
% 正文状态：未撰写；本文件目前是框架。

% 纲要内容开始
% 
% * 闭辛流形的体积和上同调
% * 乘积、曲面和余切丛
% * 相容几乎复结构与 Chern 类
% * 不预设 J 的可积性
% 
% ---
% 
% 纲要内容结束

\section{辛流形的整体约束}
\label{g4:sec:0704}

待撰写

\subsection{闭辛流形的体积和上同调}
\label{g4:subsec:0704:01}
待撰写

\subsection{乘积、曲面和余切丛}
\label{g4:subsec:0704:02}
待撰写

\subsection{相容几乎复结构与 Chern 类}
\label{g4:subsec:0704:03}
待撰写

\subsection{不预设 J 的可积性}
\label{g4:subsec:0704:04}
待撰写
% 纲要条目排为编号小节标题，并进入目录与书签；说明段仍保留于纲要注释中。


\begin{chaptersummary}
待撰写
\end{chaptersummary}

\BookExercises
待撰写

% !TeX root = ../../Book4.tex
% 纲要标题：## 第8章　Moser 方法、Darboux 与邻域定理
% 纲要位置：几何学计划纲要.md，第 3775 行。
% 正文状态：未撰写；本文件目前是框架。

% 纲要内容开始
% 
% > 本章目标与证明范围：同一变形方法依次处理整体同位、局部标准形与沿子流形的邻域。每个版本均说明原始一形式、非退化区间及向量场完备性。
% 
% 纲要内容结束

\chapter{Moser 方法、Darboux 与邻域定理}
\label{g4:ch:08}
\BookChapterNumber{g4:ch:08}

\begin{chapterabstract}
待撰写
\end{chapterabstract}

\begin{learninggoals}
待撰写
\end{learninggoals}

待撰写

% !TeX root = ../../Book4.tex
% 纲要标题：### 8.1 时间依赖流与 Moser 方程
% 纲要位置：几何学计划纲要.md，第 3779 行。
% 正文状态：未撰写；本文件目前是框架。

% 纲要内容开始
% 
% * Lie 导数和拉回求导
% * 同上同调类的形式族
% * 求解收缩方程
% * 紧性、紧支和相对保持
% 
% 纲要内容结束

\section{时间依赖流与 Moser 方程}
\label{g4:sec:0801}

待撰写

\subsection{Lie 导数和拉回求导}
\label{g4:subsec:0801:01}
待撰写

\subsection{同上同调类的形式族}
\label{g4:subsec:0801:02}
待撰写

\subsection{求解收缩方程}
\label{g4:subsec:0801:03}
待撰写

\subsection{紧性、紧支和相对保持}
\label{g4:subsec:0801:04}
待撰写
% 纲要条目排为编号小节标题，并进入目录与书签；说明段仍保留于纲要注释中。

% !TeX root = ../../Book4.tex
% 纲要标题：### 8.2 Darboux 定理
% 纲要位置：几何学计划纲要.md，第 3553 行。
% 正文状态：未撰写；本文件目前是框架。

% 纲要内容开始
% 
% * 局部线性化和直线路径
% * 非退化性的小邻域控制
% * 原始形式和径向同伦
% * 局部坐标的构造及无局部曲率不变量
% 
% 纲要内容结束

\section{Darboux 定理}
\label{g4:sec:0802}

待撰写

\subsection{局部线性化和直线路径}
\label{g4:subsec:0802:01}
待撰写

\subsection{非退化性的小邻域控制}
\label{g4:subsec:0802:02}
待撰写

\subsection{原始形式和径向同伦}
\label{g4:subsec:0802:03}
待撰写

\subsection{局部坐标的构造及无局部曲率不变量}
\label{g4:subsec:0802:04}
待撰写
% 纲要条目排为编号小节标题，并进入目录与书签；说明段仍保留于纲要注释中。

% !TeX root = ../../Book4.tex
% 纲要标题：### 8.3 沿子流形的相对方法
% 纲要位置：几何学计划纲要.md，第 3560 行。
% 正文状态：未撰写；本文件目前是框架。

% 纲要内容开始
% 
% * 相对 Poincaré 引理
% * 匹配法向线性数据
% * 辛子流形邻域定理
% * Lagrange 邻域定理
% 
% 纲要内容结束

\section{沿子流形的相对方法}
\label{g4:sec:0803}

待撰写

\subsection{相对 Poincaré 引理}
\label{g4:subsec:0803:01}
待撰写

\subsection{匹配法向线性数据}
\label{g4:subsec:0803:02}
待撰写

\subsection{辛子流形邻域定理}
\label{g4:subsec:0803:03}
待撰写

\subsection{Lagrange 邻域定理}
\label{g4:subsec:0803:04}
待撰写
% 纲要条目排为编号小节标题，并进入目录与书签；说明段仍保留于纲要注释中。

% !TeX root = ../../Book4.tex
% 纲要标题：### 8.4 图像与生成函数
% 纲要位置：几何学计划纲要.md，第 3800 行。
% 正文状态：未撰写；本文件目前是框架。

% 纲要内容开始
% 
% * 闭一形式和局部恰当性
% * 辛映射图像的 Lagrange 性
% * 生成函数及不同投影的适用范围
% 
% 纲要内容结束

\section{图像与生成函数}
\label{g4:sec:0804}

待撰写

\subsection{闭一形式和局部恰当性}
\label{g4:subsec:0804:01}
待撰写

\subsection{辛映射图像的 Lagrange 性}
\label{g4:subsec:0804:02}
待撰写

\subsection{生成函数及不同投影的适用范围}
\label{g4:subsec:0804:03}
待撰写
% 纲要条目排为编号小节标题，并进入目录与书签；说明段仍保留于纲要注释中。

% !TeX root = ../../Book4.tex
% 纲要标题：### 8.5 局部自由与整体障碍
% 纲要位置：几何学计划纲要.md，第 3806 行。
% 正文状态：未撰写；本文件目前是框架。

% 纲要内容开始
% 
% 在这里准确提出球—柱嵌入问题并预告非挤压；证明安排在第24章，不作为本章标准形的前提。
% 
% ---
% 
% 纲要内容结束

\section{局部自由与整体障碍}
\label{g4:sec:0805}

待撰写
% 本节沿用纲要的连续叙述安排，不另造小节标题。


\begin{chaptersummary}
待撰写
\end{chaptersummary}

\BookExercises
待撰写

% !TeX root = ../../Book4.tex
% 纲要标题：## 第9章　Hamilton 动力学、通量与 Poisson 结构
% 纲要位置：几何学计划纲要.md，第 3579 行。
% 正文状态：未撰写；本文件目前是框架。

% 纲要内容开始
% 
% > 本章目标与证明范围：由作用变分进入 Hamilton 系统，再用闭与恰当一形式的区别研究辛同位。Poisson 的退化情况是一条独立延伸，不与非退化辛情形混写。
% 
% 纲要内容结束

\chapter{Hamilton 动力学、通量与 Poisson 结构}
\label{g4:ch:09}
\BookChapterNumber{g4:ch:09}

\begin{chapterabstract}
待撰写
\end{chapterabstract}

\begin{learninggoals}
待撰写
\end{learninggoals}

待撰写

% !TeX root = ../../Book4.tex
% 纲要标题：### 9.1 Hamilton 方程与作用
% 纲要位置：几何学计划纲要.md，第 3583 行。
% 正文状态：未撰写；本文件目前是框架。

% 纲要内容开始
% 
% * ι_XH ω=dH 的固定约定
% * 余切坐标中的 Hamilton 方程
% * 作用及端点条件
% * 时间依赖 Hamilton 和守恒量
% 
% 纲要内容结束

\section{Hamilton 方程与作用}
\label{g4:sec:0901}

待撰写

\subsection{\texorpdfstring{\(\iota_{X_H}\omega=\mathrm{d}H\) 的固定约定}{\GeometryPdfIota{}\_XH \GeometryPdfSmallOmega{}=dH 的固定约定}}
\label{g4:subsec:0901:01}
待撰写

\subsection{余切坐标中的 Hamilton 方程}
\label{g4:subsec:0901:02}
待撰写

\subsection{作用及端点条件}
\label{g4:subsec:0901:03}
待撰写

\subsection{时间依赖 Hamilton 和守恒量}
\label{g4:subsec:0901:04}
待撰写
% 纲要条目排为编号小节标题，并进入目录与书签；说明段仍保留于纲要注释中。

% !TeX root = ../../Book4.tex
% 纲要标题：### 9.2 Poisson 括号与无穷小对称
% 纲要位置：几何学计划纲要.md，第 3823 行。
% 正文状态：未撰写；本文件目前是框架。

% 纲要内容开始
% 
% * 括号的定义与 Jacobi 恒等式
% * 向量场括号的符号核验
% * 第一积分和约束
% * 典型动力系统
% 
% 纲要内容结束

\section{Poisson 括号与无穷小对称}
\label{g4:sec:0902}

待撰写

\subsection{括号的定义与 Jacobi 恒等式}
\label{g4:subsec:0902:01}
待撰写

\subsection{向量场括号的符号核验}
\label{g4:subsec:0902:02}
待撰写

\subsection{第一积分和约束}
\label{g4:subsec:0902:03}
待撰写

\subsection{典型动力系统}
\label{g4:subsec:0902:04}
待撰写
% 纲要条目排为编号小节标题，并进入目录与书签；说明段仍保留于纲要注释中。

% !TeX root = ../../Book4.tex
% 纲要标题：### 9.3 辛同位与通量
% 纲要位置：几何学计划纲要.md，第 3830 行。
% 正文状态：未撰写；本文件目前是框架。

% 纲要内容开始
% 
% * 辛向量场对应闭一形式
% * Hamilton 向量场对应恰当一形式
% * 通量的路径积分及同伦变化
% * 环面平移的显式反例
% 
% 纲要内容结束

\section{辛同位与通量}
\label{g4:sec:0903}

待撰写

\subsection{辛向量场对应闭一形式}
\label{g4:subsec:0903:01}
待撰写

\subsection{Hamilton 向量场对应恰当一形式}
\label{g4:subsec:0903:02}
待撰写

\subsection{通量的路径积分及同伦变化}
\label{g4:subsec:0903:03}
待撰写

\subsection{环面平移的显式反例}
\label{g4:subsec:0903:04}
待撰写
% 纲要条目排为编号小节标题，并进入目录与书签；说明段仍保留于纲要注释中。

% !TeX root = ../../Book4.tex
% 纲要标题：### 9.4* 退化 Poisson 几何
% 纲要位置：几何学计划纲要.md，第 3837 行。
% 正文状态：未撰写；本文件目前是框架。

% 纲要内容开始
% 
% * Poisson 张量和 Schouten 条件
% * 特征分布与辛叶
% * 线性 Lie–Poisson 结构
% * 秩跳跃与奇异叶
% 
% ---
% 
% 纲要内容结束

\OptionalSection{退化 Poisson 几何}{g4:sec:0904}

待撰写

\subsection{Poisson 张量和 Schouten 条件}
\label{g4:subsec:0904:01}
待撰写

\subsection{特征分布与辛叶}
\label{g4:subsec:0904:02}
待撰写

\subsection{线性 Lie–Poisson 结构}
\label{g4:subsec:0904:03}
待撰写

\subsection{秩跳跃与奇异叶}
\label{g4:subsec:0904:04}
待撰写
% 纲要条目排为编号小节标题，并进入目录与书签；说明段仍保留于纲要注释中。


\begin{chaptersummary}
待撰写
\end{chaptersummary}

\BookExercises
待撰写

% !TeX root = ../../Book4.tex
% 纲要标题：## 第10章　矩映射与辛约化
% 纲要位置：几何学计划纲要.md，第 3846 行。
% 正文状态：未撰写；本文件目前是框架。

% 纲要内容开始
% 
% > 本章目标与证明范围：完整证明正则值上自由固有作用的约化定理；先处理矩映射的存在与等变性障碍，再讨论商空间，避免把每个辛作用都当作 Hamilton 作用。
% 
% 纲要内容结束

\chapter{矩映射与辛约化}
\label{g4:ch:10}
\BookChapterNumber{g4:ch:10}

\begin{chapterabstract}
待撰写
\end{chapterabstract}

\begin{learninggoals}
待撰写
\end{learninggoals}

待撰写

% !TeX root = ../../Book4.tex
% 纲要标题：### 10.1 群作用和矩映射
% 纲要位置：几何学计划纲要.md，第 3850 行。
% 正文状态：未撰写；本文件目前是框架。

% 纲要内容开始
% 
% * 基本向量场及作用侧别
% * Hamilton 作用与矩映射分量
% * 等变性和加常数的自由
% * Lie 代数余循环与存在障碍
% 
% 纲要内容结束

\section{群作用和矩映射}
\label{g4:sec:1001}

待撰写

\subsection{基本向量场及作用侧别}
\label{g4:subsec:1001:01}
待撰写

\subsection{Hamilton 作用与矩映射分量}
\label{g4:subsec:1001:02}
待撰写

\subsection{等变性和加常数的自由}
\label{g4:subsec:1001:03}
待撰写

\subsection{Lie 代数余循环与存在障碍}
\label{g4:subsec:1001:04}
待撰写
% 纲要条目排为编号小节标题，并进入目录与书签；说明段仍保留于纲要注释中。

% !TeX root = ../../Book4.tex
% 纲要标题：### 10.2 约化的线性和微分几何
% 纲要位置：几何学计划纲要.md，第 3624 行。
% 正文状态：未撰写；本文件目前是框架。

% 纲要内容开始
% 
% * 水平集的切空间
% * 辛正交与轨道方向
% * 闭二形式下降
% * 非退化性和商光滑性
% * Marsden–Weinstein 约化的完整证明
% 
% 纲要内容结束

\section{约化的线性和微分几何}
\label{g4:sec:1002}

待撰写

\subsection{水平集的切空间}
\label{g4:subsec:1002:01}
待撰写

\subsection{辛正交与轨道方向}
\label{g4:subsec:1002:02}
待撰写

\subsection{闭二形式下降}
\label{g4:subsec:1002:03}
待撰写

\subsection{非退化性和商光滑性}
\label{g4:subsec:1002:04}
待撰写

\subsection{Marsden–Weinstein 约化的完整证明}
\label{g4:subsec:1002:05}
待撰写
% 纲要条目排为编号小节标题，并进入目录与书签；说明段仍保留于纲要注释中。

% !TeX root = ../../Book4.tex
% 纲要标题：### 10.3 可计算的约化模型
% 纲要位置：几何学计划纲要.md，第 3632 行。
% 正文状态：未撰写；本文件目前是框架。

% 纲要内容开始
% 
% * 余切丛约化
% * 射影空间和 Fubini–Study 形式
% * 余伴随轨道
% * 分阶段约化及相容条件
% 
% 纲要内容结束

\section{可计算的约化模型}
\label{g4:sec:1003}

待撰写

\subsection{余切丛约化}
\label{g4:subsec:1003:01}
待撰写

\subsection{射影空间和 Fubini–Study 形式}
\label{g4:subsec:1003:02}
待撰写

\subsection{余伴随轨道}
\label{g4:subsec:1003:03}
待撰写

\subsection{分阶段约化及相容条件}
\label{g4:subsec:1003:04}
待撰写
% 纲要条目排为编号小节标题，并进入目录与书签；说明段仍保留于纲要注释中。

% !TeX root = ../../Book4.tex
% 纲要标题：### 10.4* 奇异约化与局部结构
% 纲要位置：几何学计划纲要.md，第 3872 行。
% 正文状态：未撰写；本文件目前是框架。

% 纲要内容开始
% 
% * 非自由稳定子和轨道类型
% * 奇异值上的分层模型
% * 光滑层的辛结构
% * 一般奇异约化定理的引用范围
% 
% ---
% 
% 纲要内容结束

\OptionalSection{奇异约化与局部结构}{g4:sec:1004}

待撰写

\subsection{非自由稳定子和轨道类型}
\label{g4:subsec:1004:01}
待撰写

\subsection{奇异值上的分层模型}
\label{g4:subsec:1004:02}
待撰写

\subsection{光滑层的辛结构}
\label{g4:subsec:1004:03}
待撰写

\subsection{一般奇异约化定理的引用范围}
\label{g4:subsec:1004:04}
待撰写
% 纲要条目排为编号小节标题，并进入目录与书签；说明段仍保留于纲要注释中。


\begin{chaptersummary}
待撰写
\end{chaptersummary}

\BookExercises
待撰写

% !TeX root = ../../Book4.tex
% 纲要标题：## 第11章　可积系统与作用角坐标
% 纲要位置：几何学计划纲要.md，第 3881 行。
% 正文状态：未撰写；本文件目前是框架。

% 纲要内容开始
% 
% > 本章目标与证明范围：将 Liouville–Arnold 的局部环面理论与矩映射凸性分类分章。完整证明紧连通正则共同水平集的环面结构和邻域坐标，随后展示整体单值障碍。
% 
% 纲要内容结束

\chapter{可积系统与作用角坐标}
\label{g4:ch:11}
\BookChapterNumber{g4:ch:11}

\begin{chapterabstract}
待撰写
\end{chapterabstract}

\begin{learninggoals}
待撰写
\end{learninggoals}

待撰写

% !TeX root = ../../Book4.tex
% 纲要标题：### 11.1 可积性与交换流
% 纲要位置：几何学计划纲要.md，第 3652 行。
% 正文状态：未撰写；本文件目前是框架。

% 纲要内容开始
% 
% * 独立第一积分及 Poisson 交换
% * 共同水平集的 Lagrange 性
% * 完备交换流与轨道
% * 紧连通正则纤维的环面识别
% 
% 纲要内容结束

\section{可积性与交换流}
\label{g4:sec:1101}

待撰写

\subsection{独立第一积分及 Poisson 交换}
\label{g4:subsec:1101:01}
待撰写

\subsection{共同水平集的 Lagrange 性}
\label{g4:subsec:1101:02}
待撰写

\subsection{完备交换流与轨道}
\label{g4:subsec:1101:03}
待撰写

\subsection{紧连通正则纤维的环面识别}
\label{g4:subsec:1101:04}
待撰写
% 纲要条目排为编号小节标题，并进入目录与书签；说明段仍保留于纲要注释中。

% !TeX root = ../../Book4.tex
% 纲要标题：### 11.2 Liouville–Arnold 定理
% 纲要位置：几何学计划纲要.md，第 3892 行。
% 正文状态：未撰写；本文件目前是框架。

% 纲要内容开始
% 
% * 周期格和局部基底
% * 作用积分的定义
% * 角坐标的构造
% * 辛形式的标准表达
% * 参数及格基变化
% 
% 纲要内容结束

\section{Liouville–Arnold 定理}
\label{g4:sec:1102}

待撰写

\subsection{周期格和局部基底}
\label{g4:subsec:1102:01}
待撰写

\subsection{作用积分的定义}
\label{g4:subsec:1102:02}
待撰写

\subsection{角坐标的构造}
\label{g4:subsec:1102:03}
待撰写

\subsection{辛形式的标准表达}
\label{g4:subsec:1102:04}
待撰写

\subsection{参数及格基变化}
\label{g4:subsec:1102:05}
待撰写
% 纲要条目排为编号小节标题，并进入目录与书签；说明段仍保留于纲要注释中。

% !TeX root = ../../Book4.tex
% 纲要标题：### 11.3 整体障碍
% 纲要位置：几何学计划纲要.md，第 3900 行。
% 正文状态：未撰写；本文件目前是框架。

% 纲要内容开始
% 
% * 周期格局部系统
% * 整仿射结构
% * 单值矩阵和不可整体选取作用角坐标
% * 焦点型奇纤维的代表模型
% 
% 纲要内容结束

\section{整体障碍}
\label{g4:sec:1103}

待撰写

\subsection{周期格局部系统}
\label{g4:subsec:1103:01}
待撰写

\subsection{整仿射结构}
\label{g4:subsec:1103:02}
待撰写

\subsection{单值矩阵和不可整体选取作用角坐标}
\label{g4:subsec:1103:03}
待撰写

\subsection{焦点型奇纤维的代表模型}
\label{g4:subsec:1103:04}
待撰写
% 纲要条目排为编号小节标题，并进入目录与书签；说明段仍保留于纲要注释中。

% !TeX root = ../../Book4.tex
% 纲要标题：### 11.4 典型可积系统
% 纲要位置：几何学计划纲要.md，第 3907 行。
% 正文状态：未撰写；本文件目前是框架。

% 纲要内容开始
% 
% * 振子与环面平移
% * 测地流的对称性
% * 球摆及奇异纤维的入口
% 
% ---
% 
% 纲要内容结束

\section{典型可积系统}
\label{g4:sec:1104}

待撰写

\subsection{振子与环面平移}
\label{g4:subsec:1104:01}
待撰写

\subsection{测地流的对称性}
\label{g4:subsec:1104:02}
待撰写

\subsection{球摆及奇异纤维的入口}
\label{g4:subsec:1104:03}
待撰写
% 纲要条目排为编号小节标题，并进入目录与书签；说明段仍保留于纲要注释中。


\begin{chaptersummary}
待撰写
\end{chaptersummary}

\BookExercises
待撰写

% !TeX root = ../../Book4.tex
% 纲要标题：## 第12章　Hamilton 环面作用、凸性与 Delzant 分类
% 纲要位置：几何学计划纲要.md，第 3915 行。
% 正文状态：未撰写；本文件目前是框架。

% 纲要内容开始
% 
% > 本章目标与证明范围：在独立的环面作用模块中完成凸性及光滑紧辛环面流形的 Delzant 分类。像为多面体、由多面体构造、由像决定同构是三个不同证明任务。 凸性主版本取紧连通辛流形上的 Hamilton 环面作用；Delzant 分类另要求作用有效且环面维数为流形实维数的一半。
% 
% 纲要内容结束

\chapter{Hamilton 环面作用、凸性与 Delzant 分类}
\label{g4:ch:12}
\BookChapterNumber{g4:ch:12}

\begin{chapterabstract}
待撰写
\end{chapterabstract}

\begin{learninggoals}
待撰写
\end{learninggoals}

待撰写

% !TeX root = ../../Book4.tex
% 纲要标题：### 12.1 矩映射的局部结构与 Morse–Bott 工具
% 纲要位置：几何学计划纲要.md，第 3919 行。
% 正文状态：未撰写；本文件目前是框架。

% 纲要内容开始
% 
% 就地证明凸性论证需要的局部标准形和次水平集变化，引入第二卷普通 Morse 工具时明确新增的临界子流形与负法丛。
% 
% * 固定点附近的权模型与相容复法丛
% * 正则轨道与局部辛标准形
% * 临界子流形、法向 Hessian 与 Morse–Bott 引理
% * 负法丛盘丛与跨临界值的次水平集变化
% * 一般圆方向矩映射分量的 Morse–Bott 性
% * 指标和余指标的偶性及其拓扑后果
% * 无指标一和无余指标一情形的连通水平集论证
% 
%   主范围为紧连通 Hamilton 环面作用的矩映射分量；核验极值分量、连通性与偶指标条件，不把连通水平集结论宣称为任意 Morse–Bott 函数都成立。
% 
% 纲要内容结束

\section{矩映射的局部结构与 Morse–Bott 工具}
\label{g4:sec:1201}

待撰写

\subsection{固定点附近的权模型与相容复法丛}
\label{g4:subsec:1201:01}
待撰写

\subsection{正则轨道与局部辛标准形}
\label{g4:subsec:1201:02}
待撰写

\subsection{临界子流形、法向 Hessian 与 Morse–Bott 引理}
\label{g4:subsec:1201:03}
待撰写

\subsection{负法丛盘丛与跨临界值的次水平集变化}
\label{g4:subsec:1201:04}
待撰写

\subsection{一般圆方向矩映射分量的 Morse–Bott 性}
\label{g4:subsec:1201:05}
待撰写

\subsection{指标和余指标的偶性及其拓扑后果}
\label{g4:subsec:1201:06}
待撰写

\subsection{无指标一和无余指标一情形的连通水平集论证}
\label{g4:subsec:1201:07}
待撰写
% 纲要条目排为编号小节标题，并进入目录与书签；说明段仍保留于纲要注释中。

% !TeX root = ../../Book4.tex
% 纲要标题：### 12.2 凸性定理
% 纲要位置：几何学计划纲要.md，第 3933 行。
% 正文状态：未撰写；本文件目前是框架。

% 纲要内容开始
% 
% * 圆方向的极值和连通水平集
% * 环面秩的归纳
% * 矩映射像的凸性
% * 固定点像及凸包描述
% 
% 纲要内容结束

\section{凸性定理}
\label{g4:sec:1202}

待撰写

\subsection{圆方向的极值和连通水平集}
\label{g4:subsec:1202:01}
待撰写

\subsection{环面秩的归纳}
\label{g4:subsec:1202:02}
待撰写

\subsection{矩映射像的凸性}
\label{g4:subsec:1202:03}
待撰写

\subsection{固定点像及凸包描述}
\label{g4:subsec:1202:04}
待撰写
% 纲要条目排为编号小节标题，并进入目录与书签；说明段仍保留于纲要注释中。

% !TeX root = ../../Book4.tex
% 纲要标题：### 12.3 Delzant 多面体与构造
% 纲要位置：几何学计划纲要.md，第 3940 行。
% 正文状态：未撰写；本文件目前是框架。

% 纲要内容开始
% 
% * 简单、有理和光滑格条件
% * 从多面体构造辛约化
% * 自由性和紧性的核验
% * 矩映射像的恢复
% 
% 纲要内容结束

\section{Delzant 多面体与构造}
\label{g4:sec:1203}

待撰写

\subsection{简单、有理和光滑格条件}
\label{g4:subsec:1203:01}
待撰写

\subsection{从多面体构造辛约化}
\label{g4:subsec:1203:02}
待撰写

\subsection{自由性和紧性的核验}
\label{g4:subsec:1203:03}
待撰写

\subsection{矩映射像的恢复}
\label{g4:subsec:1203:04}
待撰写
% 纲要条目排为编号小节标题，并进入目录与书签；说明段仍保留于纲要注释中。

% !TeX root = ../../Book4.tex
% 纲要标题：### 12.4 分类的唯一性
% 纲要位置：几何学计划纲要.md，第 3947 行。
% 正文状态：未撰写；本文件目前是框架。

% 纲要内容开始
% 
% * 局部标准模型的比较
% * 自由轨道部分的粘合
% * 边界轨道的延拓
% * 等变辛同构与平移自由
% 
% 纲要内容结束

\section{分类的唯一性}
\label{g4:sec:1204}

待撰写

\subsection{局部标准模型的比较}
\label{g4:subsec:1204:01}
待撰写

\subsection{自由轨道部分的粘合}
\label{g4:subsec:1204:02}
待撰写

\subsection{边界轨道的延拓}
\label{g4:subsec:1204:03}
待撰写

\subsection{等变辛同构与平移自由}
\label{g4:subsec:1204:04}
待撰写
% 纲要条目排为编号小节标题，并进入目录与书签；说明段仍保留于纲要注释中。

% !TeX root = ../../Book4.tex
% 纲要标题：### 12.5 与代数环面簇的比较
% 纲要位置：几何学计划纲要.md，第 3721 行。
% 正文状态：未撰写；本文件目前是框架。

% 纲要内容开始
% 
% * 法扇与不变除子
% * 极化和整多面体
% * 射影空间、乘积及吹起模型
% * 一般非光滑多面体的新现象
% 
% ---
% 
% 纲要内容结束

\section{与代数环面簇的比较}
\label{g4:sec:1205}

待撰写

\subsection{法扇与不变除子}
\label{g4:subsec:1205:01}
待撰写

\subsection{极化和整多面体}
\label{g4:subsec:1205:02}
待撰写

\subsection{射影空间、乘积及吹起模型}
\label{g4:subsec:1205:03}
待撰写

\subsection{一般非光滑多面体的新现象}
\label{g4:subsec:1205:04}
待撰写
% 纲要条目排为编号小节标题，并进入目录与书签；说明段仍保留于纲要注释中。

% !TeX root = ../../Book4.tex
% 纲要标题：### 12.6*　等变局部化与约化体积
% 纲要位置：几何学计划纲要.md，第 3963 行。
% 正文状态：未撰写；本文件目前是框架。
% 写作范围、先修与边界：先修为 IV.10 的正则辛约化、IV.12.1 的固定点权和局部标准形，以及 II.附录 C 的等变上同调；本节保持选读，不反向成为凸性或 Delzant 分类的前提。固定紧连通辛流形上的 Hamilton 环面作用，以实或有理系数并对非零权局部化；完整计算孤立固定点的球面和复射影空间模型，证明圆作用自由正则水平集上约化类及体积的变化。一般固定子流形的局部化和一般分段多项式结论按精确版本外引，局部自由、轨形与奇异约化不混同于自由商。

% 纲要内容开始
% 
% > 写作范围与外部输入：先修为 IV.10 的正则辛约化、IV.12.1 的固定点权和局部标准形，以及 II.附录 C 的等变上同调；本节保持选读，不反向成为凸性或 Delzant 分类的前提。固定紧连通辛流形上的 Hamilton 环面作用，以实或有理系数并对非零权局部化；完整计算孤立固定点的球面和复射影空间模型，证明圆作用自由正则水平集上约化类及体积的变化。一般固定子流形的局部化和一般分段多项式结论按精确版本外引，局部自由、轨形与奇异约化不混同于自由商。
% 
% * Cartan 模型与等变积分的归一化
% * 固定点权、等变 Euler 类与局部化公式
% * 正则值上的约化形式与上同调类变化
% * Duistermaat–Heckman 测度的分段多项式性
% * 球面与复射影空间的局部化及体积计算
% 
% ---
% 
% 纲要内容结束

\OptionalSection{等变局部化与约化体积}{g4:sec:1206}

待撰写

\subsection{Cartan 模型与等变积分的归一化}
\label{g4:subsec:1206:01}
待撰写

\subsection{固定点权、等变 Euler 类与局部化公式}
\label{g4:subsec:1206:02}
待撰写

\subsection{正则值上的约化形式与上同调类变化}
\label{g4:subsec:1206:03}
待撰写

\subsection{Duistermaat–Heckman 测度的分段多项式性}
\label{g4:subsec:1206:04}
待撰写

\subsection{球面与复射影空间的局部化及体积计算}
\label{g4:subsec:1206:05}
待撰写

% 纲要条目排为编号小节标题，并进入目录与书签；说明段仍保留于纲要注释中。


\begin{chaptersummary}
待撰写
\end{chaptersummary}

\BookExercises
待撰写

% !TeX root = ../../Book4.tex
% 纲要标题：## 第13章　接触结构、Reeb 动力学与辛化
% 纲要位置：几何学计划纲要.md，第 3975 行。
% 正文状态：未撰写；本文件目前是框架。

% 纲要内容开始
% 
% > 本章目标与证明范围：完整处理 Darboux、Gray 及 Legendrian 邻域；接触分布与选择的接触形式分别描述。周期轨道和填充属于整体问题，不由局部模型推出。
% 
% 纲要内容结束

\chapter{接触结构、Reeb 动力学与辛化}
\label{g4:ch:13}
\BookChapterNumber{g4:ch:13}

\begin{chapterabstract}
待撰写
\end{chapterabstract}

\begin{learninggoals}
待撰写
\end{learninggoals}

待撰写

% !TeX root = ../../Book4.tex
% 纲要标题：### 13.1 接触结构的定义与模型
% 纲要位置：几何学计划纲要.md，第 3734 行。
% 正文状态：未撰写；本文件目前是框架。

% 纲要内容开始
% 
% * 最大非可积分布
% * 接触形式的缩放
% * 余切单位球丛和一阶喷射空间
% * Reeb 场及形式依赖
% 
% 纲要内容结束

\section{接触结构的定义与模型}
\label{g4:sec:1301}

待撰写

\subsection{最大非可积分布}
\label{g4:subsec:1301:01}
待撰写

\subsection{接触形式的缩放}
\label{g4:subsec:1301:02}
待撰写

\subsection{余切单位球丛和一阶喷射空间}
\label{g4:subsec:1301:03}
待撰写

\subsection{Reeb 场及形式依赖}
\label{g4:subsec:1301:04}
待撰写
% 纲要条目排为编号小节标题，并进入目录与书签；说明段仍保留于纲要注释中。

% !TeX root = ../../Book4.tex
% 纲要标题：### 13.2 局部与稳定性定理
% 纲要位置：几何学计划纲要.md，第 3986 行。
% 正文状态：未撰写；本文件目前是框架。

% 纲要内容开始
% 
% * 接触 Darboux 的变形证明
% * Gray 稳定性的向量场方程
% * 紧性和相对版本
% * Legendrian 邻域模型
% 
% 纲要内容结束

\section{局部与稳定性定理}
\label{g4:sec:1302}

待撰写

\subsection{接触 Darboux 的变形证明}
\label{g4:subsec:1302:01}
待撰写

\subsection{Gray 稳定性的向量场方程}
\label{g4:subsec:1302:02}
待撰写

\subsection{紧性和相对版本}
\label{g4:subsec:1302:03}
待撰写

\subsection{Legendrian 邻域模型}
\label{g4:subsec:1302:04}
待撰写
% 纲要条目排为编号小节标题，并进入目录与书签；说明段仍保留于纲要注释中。

% !TeX root = ../../Book4.tex
% 纲要标题：### 13.3 Legendrian 与 Reeb 现象
% 纲要位置：几何学计划纲要.md，第 3748 行。
% 正文状态：未撰写；本文件目前是框架。

% 纲要内容开始
% 
% * 前投影和 Lagrange 投影
% * 标准接触三空间中的曲线
% * Reeb 弦和周期轨道
% * 作用及接触型超曲面
% 
% 纲要内容结束

\section{Legendrian 与 Reeb 现象}
\label{g4:sec:1303}

待撰写

\subsection{前投影和 Lagrange 投影}
\label{g4:subsec:1303:01}
待撰写

\subsection{标准接触三空间中的曲线}
\label{g4:subsec:1303:02}
待撰写

\subsection{Reeb 弦和周期轨道}
\label{g4:subsec:1303:03}
待撰写

\subsection{作用及接触型超曲面}
\label{g4:subsec:1303:04}
待撰写
% 纲要条目排为编号小节标题，并进入目录与书签；说明段仍保留于纲要注释中。

% !TeX root = ../../Book4.tex
% 纲要标题：### 13.4* 辛化、填充与开书
% 纲要位置：几何学计划纲要.md，第 4000 行。
% 正文状态：未撰写；本文件目前是框架。

% 纲要内容开始
% 
% * 辛化及其锥结构
% * 强、弱和精确填充的区别
% * 开书的局部模型
% * 周期存在与填充障碍的外部输入
% 
% ---
% 
% 纲要内容结束

\OptionalSection{辛化、填充与开书}{g4:sec:1304}

待撰写

\subsection{辛化及其锥结构}
\label{g4:subsec:1304:01}
待撰写

\subsection{强、弱和精确填充的区别}
\label{g4:subsec:1304:02}
待撰写

\subsection{开书的局部模型}
\label{g4:subsec:1304:03}
待撰写

\subsection{周期存在与填充障碍的外部输入}
\label{g4:subsec:1304:04}
待撰写
% 纲要条目排为编号小节标题，并进入目录与书签；说明段仍保留于纲要注释中。


\begin{chaptersummary}
待撰写
\end{chaptersummary}

\BookExercises
待撰写


% !TeX root = ../Book4.tex
% 纲要标题：# 第三编　局部解析、全纯丛与 Kähler 理论
% 纲要位置：几何学计划纲要.md，第 4009 行。
% 正文状态：未撰写；本文件目前是框架。

% 纲要内容开始
% 
% 纲要内容结束

\BookPart{局部解析、全纯丛与 Kähler 理论}{g4:part:03}
\BookPartGuide
待撰写

% !TeX root = ../../Book4.tex
% 纲要标题：## 第14章　多复变与局部解析几何
% 纲要位置：几何学计划纲要.md，第 4011 行。
% 正文状态：未撰写；本文件目前是框架。

% 纲要内容开始
% 
% > 本章目标与证明范围：必要的解析代数从多变量方程内部建立。Hartogs、Weierstrass 和凝聚性不放在一节中统称为“局部理论”；每项有不同工具和用途。
% 
% 纲要内容结束

\chapter{多复变与局部解析几何}
\label{g4:ch:14}
\BookChapterNumber{g4:ch:14}

\begin{chapterabstract}
待撰写
\end{chapterabstract}

\begin{learninggoals}
待撰写
\end{learninggoals}

待撰写

% !TeX root = ../../Book4.tex
% 纲要标题：### 14.1 多元全纯函数
% 纲要位置：几何学计划纲要.md，第 4015 行。
% 正文状态：未撰写；本文件目前是框架。

% 纲要内容开始
% 
% * 多圆盘 Cauchy 公式
% * 幂级数与正常收敛
% * 解析逆函数和常秩模型
% * Hartogs 可微与延拓现象
% 
% 纲要内容结束

\section{多元全纯函数}
\label{g4:sec:1401}

待撰写

\subsection{多圆盘 Cauchy 公式}
\label{g4:subsec:1401:01}
待撰写

\subsection{幂级数与正常收敛}
\label{g4:subsec:1401:02}
待撰写

\subsection{解析逆函数和常秩模型}
\label{g4:subsec:1401:03}
待撰写

\subsection{Hartogs 可微与延拓现象}
\label{g4:subsec:1401:04}
待撰写
% 纲要条目排为编号小节标题，并进入目录与书签；说明段仍保留于纲要注释中。

% !TeX root = ../../Book4.tex
% 纲要标题：### 14.2 Weierstrass 方法
% 纲要位置：几何学计划纲要.md，第 3777 行。
% 正文状态：未撰写；本文件目前是框架。

% 纲要内容开始
% 
% * 正规变量及多项式化
% * 除法定理的构造
% * 预备定理及单位因子
% * 收敛控制与形式版本的区别
% 
% 纲要内容结束

\section{Weierstrass 方法}
\label{g4:sec:1402}

待撰写

\subsection{正规变量及多项式化}
\label{g4:subsec:1402:01}
待撰写

\subsection{除法定理的构造}
\label{g4:subsec:1402:02}
待撰写

\subsection{预备定理及单位因子}
\label{g4:subsec:1402:03}
待撰写

\subsection{收敛控制与形式版本的区别}
\label{g4:subsec:1402:04}
待撰写
% 纲要条目排为编号小节标题，并进入目录与书签；说明段仍保留于纲要注释中。

% !TeX root = ../../Book4.tex
% 纲要标题：### 14.3 解析集合的局部结构
% 纲要位置：几何学计划纲要.md，第 4029 行。
% 正文状态：未撰写；本文件目前是框架。

% 纲要内容开始
% 
% * 解析局部环的有限性
% * 不可约分支与局部维数
% * 正则点和奇异点
% * 有限投影和重数
% 
% 纲要内容结束

\section{解析集合的局部结构}
\label{g4:sec:1403}

待撰写

\subsection{解析局部环的有限性}
\label{g4:subsec:1403:01}
待撰写

\subsection{不可约分支与局部维数}
\label{g4:subsec:1403:02}
待撰写

\subsection{正则点和奇异点}
\label{g4:subsec:1403:03}
待撰写

\subsection{有限投影和重数}
\label{g4:subsec:1403:04}
待撰写
% 纲要条目排为编号小节标题，并进入目录与书签；说明段仍保留于纲要注释中。

% !TeX root = ../../Book4.tex
% 纲要标题：### 14.4 解析层的凝聚性
% 纲要位置：几何学计划纲要.md，第 3791 行。
% 正文状态：未撰写；本文件目前是框架。

% 纲要内容开始
% 
% * 有限关系及局部呈示
% * Oka 凝聚定理的 Weierstrass 归纳模块
% * 理想层与解析子空间
% * 非约化解析空间的模型
% 
% 纲要内容结束

\section{解析层的凝聚性}
\label{g4:sec:1404}

待撰写

\subsection{有限关系及局部呈示}
\label{g4:subsec:1404:01}
待撰写

\subsection{Oka 凝聚定理的 Weierstrass 归纳模块}
\label{g4:subsec:1404:02}
待撰写

\subsection{理想层与解析子空间}
\label{g4:subsec:1404:03}
待撰写

\subsection{非约化解析空间的模型}
\label{g4:subsec:1404:04}
待撰写
% 纲要条目排为编号小节标题，并进入目录与书签；说明段仍保留于纲要注释中。

% !TeX root = ../../Book4.tex
% 纲要标题：### 14.5 多重次调和与 Levi 几何
% 纲要位置：几何学计划纲要.md，第 4043 行。
% 正文状态：未撰写；本文件目前是框架。
% 本次扩充的范围、先修与边界：先修为本章多元全纯坐标、定义函数与 Levi 形式，以及 IV.13 的接触分布。固定复维数至少二的复欧氏空间中的光滑实超曲面，要求定义函数微分非零，构造复切 CR 分布并说明其形式可积性；完成球面、Heisenberg 超曲面及严格伪凸所诱导接触结构的计算。只处理嵌入超曲面模型，不承担抽象 CR 结构的一般可嵌入性、Kohn 理论或一般共形几何。

% 纲要内容开始
% 
% > 补充写作范围与外部输入：先修为本章多元全纯坐标、定义函数与 Levi 形式，以及 IV.13 的接触分布。固定复维数至少二的复欧氏空间中的光滑实超曲面，要求定义函数微分非零，构造复切 CR 分布并说明其形式可积性；完成球面、Heisenberg 超曲面及严格伪凸所诱导接触结构的计算。只处理嵌入超曲面模型，不承担抽象 CR 结构的一般可嵌入性、Kohn 理论或一般共形几何。
% 
% * Levi 形式及 psh 函数
% * 严格 psh 与局部势
% * 伪凸性和域的整体问题
% * 与后续带权 ∂̄ 估计的连接
% 
% * 实超曲面的 CR 分布与形式可积性
% * 球面与 Heisenberg 超曲面的 Levi 形式
% * 严格伪凸超曲面的接触结构
% 
% 
% ---
% 
% 纲要内容结束

\section{多重次调和与 Levi 几何}
\label{g4:sec:1405}

待撰写

\subsection{Levi 形式及 psh 函数}
\label{g4:subsec:1405:01}
待撰写

\subsection{严格 psh 与局部势}
\label{g4:subsec:1405:02}
待撰写

\subsection{伪凸性和域的整体问题}
\label{g4:subsec:1405:03}
待撰写

\subsection{\texorpdfstring{与后续带权 \(\bar{\partial}\) 估计的连接}{与后续带权 \={∂} 估计的连接}}
\label{g4:subsec:1405:04}
待撰写

\subsection{实超曲面的 CR 分布与形式可积性}
\label{g4:subsec:1405:05}
待撰写

\subsection{球面与 Heisenberg 超曲面的 Levi 形式}
\label{g4:subsec:1405:06}
待撰写

\subsection{严格伪凸超曲面的接触结构}
\label{g4:subsec:1405:07}
待撰写

% 纲要条目排为编号小节标题，并进入目录与书签；说明段仍保留于纲要注释中。


\begin{chaptersummary}
待撰写
\end{chaptersummary}

\BookExercises
待撰写

% !TeX root = ../../Book4.tex
% 纲要标题：## 第15章　复流形、概复结构与可积性
% 纲要位置：几何学计划纲要.md，第 4059 行。
% 正文状态：未撰写；本文件目前是框架。

% 纲要内容开始
% 
% > 本章目标与证明范围：从概复的线性型分解推进到全纯坐标的存在问题。Nijenhuis 消失和坐标可积性之间使用明确的 Newlander–Nirenberg 外部定理；不把这一深层 PDE 步骤伪装成线性代数。
% 
% 纲要内容结束

\chapter{复流形、概复结构与可积性}
\label{g4:ch:15}
\BookChapterNumber{g4:ch:15}

\begin{chapterabstract}
待撰写
\end{chapterabstract}

\begin{learninggoals}
待撰写
\end{learninggoals}

待撰写

% !TeX root = ../../Book4.tex
% 纲要标题：### 15.1 复化切丛与型分解
% 纲要位置：几何学计划纲要.md，第 4063 行。
% 正文状态：未撰写；本文件目前是框架。

% 纲要内容开始
% 
% * T¹⁰、T⁰¹ 及共轭
% * 形式的双次数
% * d 的各型分量
% * 全纯坐标下的 ∂ 和 ∂̄
% 
% 纲要内容结束

\section{复化切丛与型分解}
\label{g4:sec:1501}

待撰写

\subsection{\texorpdfstring{\(\mathrm{T}^{1,0}\)、\(\mathrm{T}^{0,1}\) 及共轭}{T¹⁰、T⁰¹ 及共轭}}
\label{g4:subsec:1501:01}
待撰写

\subsection{形式的双次数}
\label{g4:subsec:1501:02}
待撰写

\subsection{d 的各型分量}
\label{g4:subsec:1501:03}
待撰写

\subsection{\texorpdfstring{全纯坐标下的 \(\partial\) 和 \(\bar{\partial}\)}{全纯坐标下的 ∂ 和 \={∂}}}
\label{g4:subsec:1501:04}
待撰写
% 纲要条目排为编号小节标题，并进入目录与书签；说明段仍保留于纲要注释中。

% !TeX root = ../../Book4.tex
% 纲要标题：### 15.2 Nijenhuis 与可积性
% 纲要位置：几何学计划纲要.md，第 3818 行。
% 正文状态：未撰写；本文件目前是框架。

% 纲要内容开始
% 
% * Nijenhuis 张量的内在定义
% * 对合条件和形式可积性
% * 光滑 Newlander–Nirenberg 定理的精确引用
% * 低正则性不自动继承同一结论
% 
% 纲要内容结束

\section{Nijenhuis 与可积性}
\label{g4:sec:1502}

待撰写

\subsection{Nijenhuis 张量的内在定义}
\label{g4:subsec:1502:01}
待撰写

\subsection{对合条件和形式可积性}
\label{g4:subsec:1502:02}
待撰写

\subsection{光滑 Newlander–Nirenberg 定理的精确引用}
\label{g4:subsec:1502:03}
待撰写

\subsection{低正则性不自动继承同一结论}
\label{g4:subsec:1502:04}
待撰写
% 纲要条目排为编号小节标题，并进入目录与书签；说明段仍保留于纲要注释中。

% !TeX root = ../../Book4.tex
% 纲要标题：### 15.3 复流形及解析子空间
% 纲要位置：几何学计划纲要.md，第 3825 行。
% 正文状态：未撰写；本文件目前是框架。

% 纲要内容开始
% 
% * 全纯映射、浸入与子流形
% * 解析空间和局部环
% * 复吹起的坐标模型
% * 正规与光滑的区别
% 
% 纲要内容结束

\section{复流形及解析子空间}
\label{g4:sec:1503}

待撰写

\subsection{全纯映射、浸入与子流形}
\label{g4:subsec:1503:01}
待撰写

\subsection{解析空间和局部环}
\label{g4:subsec:1503:02}
待撰写

\subsection{复吹起的坐标模型}
\label{g4:subsec:1503:03}
待撰写

\subsection{正规与光滑的区别}
\label{g4:subsec:1503:04}
待撰写
% 纲要条目排为编号小节标题，并进入目录与书签；说明段仍保留于纲要注释中。

% !TeX root = ../../Book4.tex
% 纲要标题：### 15.4 Kähler 之前的实例与障碍
% 纲要位置：几何学计划纲要.md，第 3832 行。
% 正文状态：未撰写；本文件目前是框架。

% 纲要内容开始
% 
% * 复环面和 Hopf 曲面
% * 奇数第一 Betti 数的模型
% * 概复而未指定可积结构的例子
% * 复、辛及 Kähler 条件的独立性
% 
% ---
% 
% 纲要内容结束

\section{Kähler 之前的实例与障碍}
\label{g4:sec:1504}

待撰写

\subsection{复环面和 Hopf 曲面}
\label{g4:subsec:1504:01}
待撰写

\subsection{奇数第一 Betti 数的模型}
\label{g4:subsec:1504:02}
待撰写

\subsection{概复而未指定可积结构的例子}
\label{g4:subsec:1504:03}
待撰写

\subsection{复、辛及 Kähler 条件的独立性}
\label{g4:subsec:1504:04}
待撰写
% 纲要条目排为编号小节标题，并进入目录与书签；说明段仍保留于纲要注释中。


\begin{chaptersummary}
待撰写
\end{chaptersummary}

\BookExercises
待撰写

% !TeX root = ../../Book4.tex
% 纲要标题：## 第16章　全纯向量丛、Dolbeault 与 Chern 联络
% 纲要位置：几何学计划纲要.md，第 3841 行。
% 正文状态：未撰写；本文件目前是框架。

% 纲要内容开始
% 
% > 本章目标与证明范围：以全纯结构作为满足可积条件的 ∂̄ 算子来组织丛论，再连接层上同调和 Hermitian 曲率。局部自由层与一般凝聚层不混同。
% 
% 纲要内容结束

\chapter{全纯向量丛、Dolbeault 与 Chern 联络}
\label{g4:ch:16}
\BookChapterNumber{g4:ch:16}

\begin{chapterabstract}
待撰写
\end{chapterabstract}

\begin{learninggoals}
待撰写
\end{learninggoals}

待撰写

% !TeX root = ../../Book4.tex
% 纲要标题：### 16.1 全纯丛的等价描述
% 纲要位置：几何学计划纲要.md，第 3845 行。
% 正文状态：未撰写；本文件目前是框架。

% 纲要内容开始
% 
% * 全纯过渡函数与局部截面
% * 丛值 ∂̄ 算子和平方为零
% * 可积 ∂̄ 算子的局部平凡化
% 
%   固定光滑复向量丛和光滑 (0,1) 联络，按 Koszul–Malgrange 的局部定理处理 ∂̄²=0；本节安排逐变量求解与规范变换的证明模块，不把形式平方为零直接当作已有全纯标架。
% 
% * 丛的扩张及上同调类
% 
% 纲要内容结束

\section{全纯丛的等价描述}
\label{g4:sec:1601}

待撰写

\subsection{全纯过渡函数与局部截面}
\label{g4:subsec:1601:01}
待撰写

\subsection{\texorpdfstring{丛值 \(\bar{\partial}\) 算子和平方为零}{丛值 \={∂} 算子和平方为零}}
\label{g4:subsec:1601:02}
待撰写

\subsection{\texorpdfstring{可积 \(\bar{\partial}\) 算子的局部平凡化}{可积 \={∂} 算子的局部平凡化}}
\label{g4:subsec:1601:03}
待撰写

\subsection{丛的扩张及上同调类}
\label{g4:subsec:1601:04}
待撰写
% 纲要条目排为编号小节标题，并进入目录与书签；说明段仍保留于纲要注释中。

% !TeX root = ../../Book4.tex
% 纲要标题：### 16.2 Dolbeault 消解与上同调
% 纲要位置：几何学计划纲要.md，第 4107 行。
% 正文状态：未撰写；本文件目前是框架。

% 纲要内容开始
% 
% * 局部 ∂̄ 引理
% * 丛值形式的细层消解
% * 与层上同调的比较
% * 紧复流形上的有限维性
% 
% 纲要内容结束

\section{Dolbeault 消解与上同调}
\label{g4:sec:1602}

待撰写

\subsection{\texorpdfstring{局部 \(\bar{\partial}\) 引理}{局部 \={∂} 引理}}
\label{g4:subsec:1602:01}
待撰写

\subsection{丛值形式的细层消解}
\label{g4:subsec:1602:02}
待撰写

\subsection{与层上同调的比较}
\label{g4:subsec:1602:03}
待撰写

\subsection{紧复流形上的有限维性}
\label{g4:subsec:1602:04}
待撰写
% 纲要条目排为编号小节标题，并进入目录与书签；说明段仍保留于纲要注释中。

% !TeX root = ../../Book4.tex
% 纲要标题：### 16.3 Chern 联络的构造
% 纲要位置：几何学计划纲要.md，第 4114 行。
% 正文状态：未撰写；本文件目前是框架。

% 纲要内容开始
% 
% * Hermitian 度量与联络相容性
% * (0,1) 部分固定的条件
% * 存在唯一性和局部矩阵公式
% * 曲率型及 Bianchi 恒等式
% 
% 纲要内容结束

\section{Chern 联络的构造}
\label{g4:sec:1603}

待撰写

\subsection{Hermitian 度量与联络相容性}
\label{g4:subsec:1603:01}
待撰写

\subsection{(0,1) 部分固定的条件}
\label{g4:subsec:1603:02}
待撰写

\subsection{存在唯一性和局部矩阵公式}
\label{g4:subsec:1603:03}
待撰写

\subsection{曲率型及 Bianchi 恒等式}
\label{g4:subsec:1603:04}
待撰写
% 纲要条目排为编号小节标题，并进入目录与书签；说明段仍保留于纲要注释中。

% !TeX root = ../../Book4.tex
% 纲要标题：### 16.4 曲率与丛的几何
% 纲要位置：几何学计划纲要.md，第 4121 行。
% 正文状态：未撰写；本文件目前是框架。

% 纲要内容开始
% 
% * 第一陈类和除子
% * 子丛、商丛与第二基本形式
% * 曲率的 Gauss–Codazzi 型公式
% * 线丛、切丛及射影模型
% 
% ---
% 
% 纲要内容结束

\section{曲率与丛的几何}
\label{g4:sec:1604}

待撰写

\subsection{第一陈类和除子}
\label{g4:subsec:1604:01}
待撰写

\subsection{子丛、商丛与第二基本形式}
\label{g4:subsec:1604:02}
待撰写

\subsection{曲率的 Gauss–Codazzi 型公式}
\label{g4:subsec:1604:03}
待撰写

\subsection{线丛、切丛及射影模型}
\label{g4:subsec:1604:04}
待撰写
% 纲要条目排为编号小节标题，并进入目录与书签；说明段仍保留于纲要注释中。


\begin{chaptersummary}
待撰写
\end{chaptersummary}

\BookExercises
待撰写

% !TeX root = ../../Book4.tex
% 纲要标题：## 第17章　Kähler 结构、恒等式与 Hodge 理论
% 纲要位置：几何学计划纲要.md，第 4130 行。
% 正文状态：未撰写；本文件目前是框架。

% 纲要内容开始
% 
% > 本章目标与证明范围：不把相容结构回顾单独铺成五节；在实际度量模型之后连续证明 Lefschetz 线性代数、Kähler 恒等式、分解、∂∂̄ 引理及双线性关系。
% 
% 纲要内容结束

\chapter{Kähler 结构、恒等式与 Hodge 理论}
\label{g4:ch:17}
\BookChapterNumber{g4:ch:17}

\begin{chapterabstract}
待撰写
\end{chapterabstract}

\begin{learninggoals}
待撰写
\end{learninggoals}

待撰写

% !TeX root = ../../Book4.tex
% 纲要标题：### 17.1 Kähler 条件及局部势
% 纲要位置：几何学计划纲要.md，第 3882 行。
% 正文状态：未撰写；本文件目前是框架。

% 纲要内容开始
% 
% * 相容三元组与可积性回顾
% * dω=0 和平行 J 的关系
% * 局部 Kähler 势
% * Fubini–Study、复环面及 Kähler 约化
% 
% 纲要内容结束

\section{Kähler 条件及局部势}
\label{g4:sec:1701}

待撰写

\subsection{相容三元组与可积性回顾}
\label{g4:subsec:1701:01}
待撰写

\subsection{\texorpdfstring{\(\mathrm{d}\omega=0\) 和平行 J 的关系}{d\GeometryPdfSmallOmega{}=0 和平行 J 的关系}}
\label{g4:subsec:1701:02}
待撰写

\subsection{局部 Kähler 势}
\label{g4:subsec:1701:03}
待撰写

\subsection{Fubini–Study、复环面及 Kähler 约化}
\label{g4:subsec:1701:04}
待撰写
% 纲要条目排为编号小节标题，并进入目录与书签；说明段仍保留于纲要注释中。

% !TeX root = ../../Book4.tex
% 纲要标题：### 17.2 Lefschetz 线性代数
% 纲要位置：几何学计划纲要.md，第 4141 行。
% 正文状态：未撰写；本文件目前是框架。

% 纲要内容开始
% 
% * L、Λ 和次数算子
% * sl₂ 交换关系的形式证明
% * 原始形式及线性分解
% * Hodge 星和正定配对
% 
% 纲要内容结束

\section{Lefschetz 线性代数}
\label{g4:sec:1702}

待撰写

\subsection{\texorpdfstring{L、\(\Lambda\) 和次数算子}{L、\GeometryPdfCapitalLambda{} 和次数算子}}
\label{g4:subsec:1702:01}
待撰写

\subsection{\texorpdfstring{\(\mathfrak{sl}_2\) 交换关系的形式证明}{sl₂ 交换关系的形式证明}}
\label{g4:subsec:1702:02}
待撰写

\subsection{原始形式及线性分解}
\label{g4:subsec:1702:03}
待撰写

\subsection{Hodge 星和正定配对}
\label{g4:subsec:1702:04}
待撰写
% 纲要条目排为编号小节标题，并进入目录与书签；说明段仍保留于纲要注释中。

% !TeX root = ../../Book4.tex
% 纲要标题：### 17.3 Kähler 恒等式
% 纲要位置：几何学计划纲要.md，第 3896 行。
% 正文状态：未撰写；本文件目前是框架。

% 纲要内容开始
% 
% * 正规全纯标架和局部计算
% * Λ 与 ∂、∂̄ 的交换子
% * 两种 Laplace 算子的关系
% * 丛值推广所需的曲率修正
% 
% 纲要内容结束

\section{Kähler 恒等式}
\label{g4:sec:1703}

待撰写

\subsection{正规全纯标架和局部计算}
\label{g4:subsec:1703:01}
待撰写

\subsection{\texorpdfstring{\(\Lambda\) 与 \(\partial\)、\(\bar{\partial}\) 的交换子}{\GeometryPdfCapitalLambda{} 与 ∂、\={∂} 的交换子}}
\label{g4:subsec:1703:02}
待撰写

\subsection{两种 Laplace 算子的关系}
\label{g4:subsec:1703:03}
待撰写

\subsection{丛值推广所需的曲率修正}
\label{g4:subsec:1703:04}
待撰写
% 纲要条目排为编号小节标题，并进入目录与书签；说明段仍保留于纲要注释中。

% !TeX root = ../../Book4.tex
% 纲要标题：### 17.4 Hodge 分解与 ∂∂̄ 引理
% 纲要位置：几何学计划纲要.md，第 4155 行。
% 正文状态：未撰写；本文件目前是框架。

% 纲要内容开始
% 
% * 调和形式的型分解
% * 共轭和 Serre 对称
% * ∂∂̄ 引理的 Green 算子证明
% * Bott–陈、Aeppli 与 de Rham 的比较范围
% 
% 纲要内容结束

\section{Hodge 分解与 \texorpdfstring{\(\partial\bar{\partial}\)}{∂\={∂}} 引理}
\label{g4:sec:1704}

待撰写

\subsection{调和形式的型分解}
\label{g4:subsec:1704:01}
待撰写

\subsection{共轭和 Serre 对称}
\label{g4:subsec:1704:02}
待撰写

\subsection{\texorpdfstring{\(\partial\bar{\partial}\) 引理的 Green 算子证明}{∂\={∂} 引理的 Green 算子证明}}
\label{g4:subsec:1704:03}
待撰写

\subsection{Bott–陈、Aeppli 与 de Rham 的比较范围}
\label{g4:subsec:1704:04}
待撰写
% 纲要条目排为编号小节标题，并进入目录与书签；说明段仍保留于纲要注释中。

% !TeX root = ../../Book4.tex
% 纲要标题：### 17.5 Hard Lefschetz 定理
% 纲要位置：几何学计划纲要.md，第 4162 行。
% 正文状态：未撰写；本文件目前是框架。

% 纲要内容开始
% 
% * 原始调和形式及 sl₂ 分解
% * Lefschetz 算子与 Laplace 交换
% * 同调上的同构
% * 逆算子和原始分解
% 
% 纲要内容结束

\section{Hard Lefschetz 定理}
\label{g4:sec:1705}

待撰写

\subsection{\texorpdfstring{原始调和形式及 \(\mathfrak{sl}_2\) 分解}{原始调和形式及 sl₂ 分解}}
\label{g4:subsec:1705:01}
待撰写

\subsection{Lefschetz 算子与 Laplace 交换}
\label{g4:subsec:1705:02}
待撰写

\subsection{同调上的同构}
\label{g4:subsec:1705:03}
待撰写

\subsection{逆算子和原始分解}
\label{g4:subsec:1705:04}
待撰写
% 纲要条目排为编号小节标题，并进入目录与书签；说明段仍保留于纲要注释中。

% !TeX root = ../../Book4.tex
% 纲要标题：### 17.6 Hodge–Riemann 双线性关系
% 纲要位置：几何学计划纲要.md，第 4169 行。
% 正文状态：未撰写；本文件目前是框架。

% 纲要内容开始
% 
% * 原始型分量的配对
% * Hodge 星的符号
% * 正定性和非退化性
% * 极化及整格
% * 紧性或 Kähler 条件失败的模型
% 
% ---
% 
% 纲要内容结束

\section{Hodge–Riemann 双线性关系}
\label{g4:sec:1706}

待撰写

\subsection{原始型分量的配对}
\label{g4:subsec:1706:01}
待撰写

\subsection{Hodge 星的符号}
\label{g4:subsec:1706:02}
待撰写

\subsection{正定性和非退化性}
\label{g4:subsec:1706:03}
待撰写

\subsection{极化及整格}
\label{g4:subsec:1706:04}
待撰写

\subsection{紧性或 Kähler 条件失败的模型}
\label{g4:subsec:1706:05}
待撰写
% 纲要条目排为编号小节标题，并进入目录与书签；说明段仍保留于纲要注释中。

% !TeX root = ../../Book4.tex
% 纲要标题：### 17.7*　有理 Hodge 类与线丛
% 纲要位置：几何学计划纲要.md，第 4179 行。
% 正文状态：未撰写；本文件目前是框架。
% 写作范围、先修与边界：先修为 IV.16 的全纯线丛和层上同调、IV.17.4 的 Hodge 分解以及第二卷的整上同调与系数变换；射影代数圈的解释另用 IV.20 的比较，不要求先完成 V.27。先在紧 Kähler 流形上由指数正合列证明 Lefschetz (1,1) 定理，核对整系数、挠元和第一 Chern 类的像；有理 Hodge 猜想只在光滑复射影簇上准确陈述，代数圈类的更完整形式指向 V.27。明确有理与整数版本、除子与高余维的区别，不承担 Hodge 猜想的一般证明。

% 纲要内容开始
% 
% > 写作范围与外部输入：先修为 IV.16 的全纯线丛和层上同调、IV.17.4 的 Hodge 分解以及第二卷的整上同调与系数变换；射影代数圈的解释另用 IV.20 的比较，不要求先完成 V.27。先在紧 Kähler 流形上由指数正合列证明 Lefschetz (1,1) 定理，核对整系数、挠元和第一 Chern 类的像；有理 Hodge 猜想只在光滑复射影簇上准确陈述，代数圈类的更完整形式指向 V.27。明确有理与整数版本、除子与高余维的区别，不承担 Hodge 猜想的一般证明。
% 
% * 有理 Hodge 类与整上同调的区别
% * 指数正合列与线丛的第一 Chern 类
% * 线丛的整 Hodge 类与 Lefschetz 定理
% * 代数圈的上同调类与有理 Hodge 猜想
% * 除子与高余维问题的区别
% 
% ---
% 
% 纲要内容结束

\OptionalSection{有理 Hodge 类与线丛}{g4:sec:1707}

待撰写

\subsection{有理 Hodge 类与整上同调的区别}
\label{g4:subsec:1707:01}
待撰写

\subsection{指数正合列与线丛的第一 Chern 类}
\label{g4:subsec:1707:02}
待撰写

\subsection{线丛的整 Hodge 类与 Lefschetz 定理}
\label{g4:subsec:1707:03}
待撰写

\subsection{代数圈的上同调类与有理 Hodge 猜想}
\label{g4:subsec:1707:04}
待撰写

\subsection{除子与高余维问题的区别}
\label{g4:subsec:1707:05}
待撰写

% 纲要条目排为编号小节标题，并进入目录与书签；说明段仍保留于纲要注释中。


\begin{chaptersummary}
待撰写
\end{chaptersummary}

\BookExercises
待撰写


% !TeX root = ../Book4.tex
% 纲要标题：# 第四编　正性、截面构造与射影化
% 纲要位置：几何学计划纲要.md，第 4191 行。
% 正文状态：未撰写；本文件目前是框架。

% 纲要内容开始
% 
% 纲要内容结束

\BookPart{正性、截面构造与射影化}{g4:part:04}
\BookPartGuide
待撰写

% !TeX root = ../../Book4.tex
% 纲要标题：## 第18章　曲率正性与消失定理
% 纲要位置：几何学计划纲要.md，第 4193 行。
% 正文状态：未撰写；本文件目前是框架。

% 纲要内容开始
% 
% > 本章目标与证明范围：以曲率项的代数正性为起点，完整建立 Bochner–Kodaira 和所选 Kodaira–Nakano 消失。区分线丛正性、Griffiths 和 Nakano，不用同一“正”字跳过不同结论。
% 
% 纲要内容结束

\chapter{曲率正性与消失定理}
\label{g4:ch:18}
\BookChapterNumber{g4:ch:18}

\begin{chapterabstract}
待撰写
\end{chapterabstract}

\begin{learninggoals}
待撰写
\end{learninggoals}

待撰写

% !TeX root = ../../Book4.tex
% 纲要标题：### 18.1 不同正性的比较
% 纲要位置：几何学计划纲要.md，第 4197 行。
% 正文状态：未撰写；本文件目前是框架。

% 纲要内容开始
% 
% * 线丛曲率和局部权函数
% * Griffiths 与 Nakano 二次型
% * 张量、对偶和高次幂
% * 半正性及退化例子
% 
% 纲要内容结束

\section{不同正性的比较}
\label{g4:sec:1801}

待撰写

\subsection{线丛曲率和局部权函数}
\label{g4:subsec:1801:01}
待撰写

\subsection{Griffiths 与 Nakano 二次型}
\label{g4:subsec:1801:02}
待撰写

\subsection{张量、对偶和高次幂}
\label{g4:subsec:1801:03}
待撰写

\subsection{半正性及退化例子}
\label{g4:subsec:1801:04}
待撰写
% 纲要条目排为编号小节标题，并进入目录与书签；说明段仍保留于纲要注释中。

% !TeX root = ../../Book4.tex
% 纲要标题：### 18.2 Bochner–Kodaira 恒等式
% 纲要位置：几何学计划纲要.md，第 3940 行。
% 正文状态：未撰写；本文件目前是框架。

% 纲要内容开始
% 
% * 丛值 Kähler 交换公式
% * 曲率交换子及双次数
% * 积分恒等式
% * 边界和非 Kähler 挠率项的区别
% 
% 纲要内容结束

\section{Bochner–Kodaira 恒等式}
\label{g4:sec:1802}

待撰写

\subsection{丛值 Kähler 交换公式}
\label{g4:subsec:1802:01}
待撰写

\subsection{曲率交换子及双次数}
\label{g4:subsec:1802:02}
待撰写

\subsection{积分恒等式}
\label{g4:subsec:1802:03}
待撰写

\subsection{边界和非 Kähler 挠率项的区别}
\label{g4:subsec:1802:04}
待撰写
% 纲要条目排为编号小节标题，并进入目录与书签；说明段仍保留于纲要注释中。

% !TeX root = ../../Book4.tex
% 纲要标题：### 18.3 从基本估计到消失
% 纲要位置：几何学计划纲要.md，第 4211 行。
% 正文状态：未撰写；本文件目前是框架。

% 纲要内容开始
% 
% * 正曲率给出的强制性
% * 调和代表的消失
% * Kodaira 与所选 Nakano 版本
% * Serre 对偶下的另一形式
% 
% 纲要内容结束

\section{从基本估计到消失}
\label{g4:sec:1803}

待撰写

\subsection{正曲率给出的强制性}
\label{g4:subsec:1803:01}
待撰写

\subsection{调和代表的消失}
\label{g4:subsec:1803:02}
待撰写

\subsection{Kodaira 与所选 Nakano 版本}
\label{g4:subsec:1803:03}
待撰写

\subsection{Serre 对偶下的另一形式}
\label{g4:subsec:1803:04}
待撰写
% 纲要条目排为编号小节标题，并进入目录与书签；说明段仍保留于纲要注释中。

% !TeX root = ../../Book4.tex
% 纲要标题：### 18.4 零点、截面与拓扑限制
% 纲要位置：几何学计划纲要.md，第 3954 行。
% 正文状态：未撰写；本文件目前是框架。

% 纲要内容开始
% 
% * 非零截面和有效除子
% * 曲率积分与第一 Chern 类
% * 半正而不消失的模型
% * 正性与射影性问题的提出
% 
% ---
% 
% 纲要内容结束

\section{零点、截面与拓扑限制}
\label{g4:sec:1804}

待撰写

\subsection{非零截面和有效除子}
\label{g4:subsec:1804:01}
待撰写

\subsection{曲率积分与第一 Chern 类}
\label{g4:subsec:1804:02}
待撰写

\subsection{半正而不消失的模型}
\label{g4:subsec:1804:03}
待撰写

\subsection{正性与射影性问题的提出}
\label{g4:subsec:1804:04}
待撰写
% 纲要条目排为编号小节标题，并进入目录与书签；说明段仍保留于纲要注释中。


\begin{chaptersummary}
待撰写
\end{chaptersummary}

\BookExercises
待撰写

% !TeX root = ../../Book4.tex
% 纲要标题：## 第19章　∂̄ 解算子、峰值截面与 Kodaira 嵌入
% 纲要位置：几何学计划纲要.md，第 4227 行。
% 正文状态：未撰写；本文件目前是框架。

% 纲要内容开始
% 
% > 本章目标与证明范围：完整证明正线丛高次幂的射影嵌入。核心不是在消失定理后添一句“于是分离点”，而是将求解、局部修补、喷射控制和统一高次幂逐步建成。
% 
% 纲要内容结束

\chapter{\texorpdfstring{\(\bar{\partial}\)}{\={∂}} 解算子、峰值截面与 Kodaira 嵌入}
\label{g4:ch:19}
\BookChapterNumber{g4:ch:19}

\begin{chapterabstract}
待撰写
\end{chapterabstract}

\begin{learninggoals}
待撰写
\end{learninggoals}

待撰写

% !TeX root = ../../Book4.tex
% 纲要标题：### 19.1 L² 算子及可解性
% 纲要位置：几何学计划纲要.md，第 3967 行。
% 正文状态：未撰写；本文件目前是框架。

% 纲要内容开始
% 
% * 算子定义域、闭算子和 Hilbert 伴随
% * 基本估计与正交障碍
% * 通过对偶构造最小范数解
% * 弱解的椭圆正则性
% * 正线丛高次幂中的常数控制
% 
% 纲要内容结束

\section{\texorpdfstring{\(L^2\)}{L²} 算子及可解性}
\label{g4:sec:1901}

待撰写

\subsection{算子定义域、闭算子和 Hilbert 伴随}
\label{g4:subsec:1901:01}
待撰写

\subsection{基本估计与正交障碍}
\label{g4:subsec:1901:02}
待撰写

\subsection{通过对偶构造最小范数解}
\label{g4:subsec:1901:03}
待撰写

\subsection{弱解的椭圆正则性}
\label{g4:subsec:1901:04}
待撰写

\subsection{正线丛高次幂中的常数控制}
\label{g4:subsec:1901:05}
待撰写
% 纲要条目排为编号小节标题，并进入目录与书签；说明段仍保留于纲要注释中。

% !TeX root = ../../Book4.tex
% 纲要标题：### 19.2 局部模型与截断误差
% 纲要位置：几何学计划纲要.md，第 4239 行。
% 正文状态：未撰写；本文件目前是框架。

% 纲要内容开始
% 
% * 正曲率的正规局部标架
% * Gaussian 型峰值和缩放
% * 常数及一次喷射的局部截面
% * 截断后 ∂̄ 误差的定量估计
% 
% 纲要内容结束

\section{局部模型与截断误差}
\label{g4:sec:1902}

待撰写

\subsection{正曲率的正规局部标架}
\label{g4:subsec:1902:01}
待撰写

\subsection{Gaussian 型峰值和缩放}
\label{g4:subsec:1902:02}
待撰写

\subsection{常数及一次喷射的局部截面}
\label{g4:subsec:1902:03}
待撰写

\subsection{\texorpdfstring{截断后 \(\bar{\partial}\) 误差的定量估计}{截断后 \={∂} 误差的定量估计}}
\label{g4:subsec:1902:04}
待撰写
% 纲要条目排为编号小节标题，并进入目录与书签；说明段仍保留于纲要注释中。

% !TeX root = ../../Book4.tex
% 纲要标题：### 19.3 整体修补及喷射控制
% 纲要位置：几何学计划纲要.md，第 4246 行。
% 正文状态：未撰写；本文件目前是框架。

% 纲要内容开始
% 
% * 解 ∂̄ 消除误差
% * L² 控制到局部 C¹ 控制
% * 修补后零阶与一阶喷射的线性独立
% * 局部估计对中心点的一致性
% 
% 纲要内容结束

\section{整体修补及喷射控制}
\label{g4:sec:1903}

待撰写

\subsection{\texorpdfstring{解 \(\bar{\partial}\) 消除误差}{解 \={∂} 消除误差}}
\label{g4:subsec:1903:01}
待撰写

\subsection{\texorpdfstring{\(L^2\) 控制到局部 \(C^1\) 控制}{L² 控制到局部 C¹ 控制}}
\label{g4:subsec:1903:02}
待撰写

\subsection{修补后零阶与一阶喷射的线性独立}
\label{g4:subsec:1903:03}
待撰写

\subsection{局部估计对中心点的一致性}
\label{g4:subsec:1903:04}
待撰写
% 纲要条目排为编号小节标题，并进入目录与书签；说明段仍保留于纲要注释中。

% !TeX root = ../../Book4.tex
% 纲要标题：### 19.4 分离点和切向量
% 纲要位置：几何学计划纲要.md，第 4253 行。
% 正文状态：未撰写；本文件目前是框架。

% 纲要内容开始
% 
% * 无基点性的证明
% * 切向量分离和局部浸入
% * 近点与远点分别处理
% * 紧性给出统一幂次
% * 由有限维截面空间得到嵌入
% 
% 纲要内容结束

\section{分离点和切向量}
\label{g4:sec:1904}

待撰写

\subsection{无基点性的证明}
\label{g4:subsec:1904:01}
待撰写

\subsection{切向量分离和局部浸入}
\label{g4:subsec:1904:02}
待撰写

\subsection{近点与远点分别处理}
\label{g4:subsec:1904:03}
待撰写

\subsection{紧性给出统一幂次}
\label{g4:subsec:1904:04}
待撰写

\subsection{由有限维截面空间得到嵌入}
\label{g4:subsec:1904:05}
待撰写
% 纲要条目排为编号小节标题，并进入目录与书签；说明段仍保留于纲要注释中。

% !TeX root = ../../Book4.tex
% 纲要标题：### 19.5 极化与代数化的准备
% 纲要位置：几何学计划纲要.md，第 4261 行。
% 正文状态：未撰写；本文件目前是框架。

% 纲要内容开始
% 
% * 高次幂和极化的不同选择
% * 射影度量及原 Kähler 类
% * 正线丛存在的整性条件
% * 不能由任意 Kähler 度量直接推出射影性
% 
% ---
% 
% 纲要内容结束

\section{极化与代数化的准备}
\label{g4:sec:1905}

待撰写

\subsection{高次幂和极化的不同选择}
\label{g4:subsec:1905:01}
待撰写

\subsection{射影度量及原 Kähler 类}
\label{g4:subsec:1905:02}
待撰写

\subsection{正线丛存在的整性条件}
\label{g4:subsec:1905:03}
待撰写

\subsection{不能由任意 Kähler 度量直接推出射影性}
\label{g4:subsec:1905:04}
待撰写
% 纲要条目排为编号小节标题，并进入目录与书签；说明段仍保留于纲要注释中。


\begin{chaptersummary}
待撰写
\end{chaptersummary}

\BookExercises
待撰写

% !TeX root = ../../Book4.tex
% 纲要标题：## 第20章　Chow、解析化与射影 GAGA
% 纲要位置：几何学计划纲要.md，第 4006 行。
% 正文状态：未撰写；本文件目前是框架。

% 纲要内容开始
% 
% > 本章目标与证明范围：射影比较作为实质性定理模块。独立复几何路线在本章重建复数域上所需的代数工具；已读第五卷者按第20.1节的准确位置调用。解析凝聚性、高次扭转生成和上同调比较均有证明位置；更一般固有版本单列外引。[R5]
% 
% 纲要内容结束

\chapter{Chow、解析化与射影 GAGA}
\label{g4:ch:20}
\BookChapterNumber{g4:ch:20}

\begin{chapterabstract}
待撰写
\end{chapterabstract}

\begin{learninggoals}
待撰写
\end{learninggoals}

待撰写

% !TeX root = ../../Book4.tex
% 纲要标题：### 20.1 射影比较的代数输入与独立入口
% 纲要位置：几何学计划纲要.md，第 4010 行。
% 正文状态：未撰写；本文件目前是框架。

% 纲要内容开始
% 
% 独立入口只重建本章所需的复数域射影版本，不要求先读第五卷全部交换代数、形式函数或模空间理论。
% 
% * 复射影代数对象、概形和结构层的必要构造
% 
%   概形与层对应 V.10.1—10.3、V.13.1—13.3。
% 
% * 代数凝聚层、可逆层、张量与有限呈示
% * 相对仿射构造、Proj 和扭转线丛
% 
%   构造对应 V.13.5、V.14.1—14.2；相对 Spec 不依赖形式函数。
% 
% * 仿射无上同调与有限仿射覆盖的 Čech 模型
% 
%   对应 V.18.1—18.3；不用光滑单位分解替代代数层的无上同调。
% 
% * 射影空间扭转层的上同调及 Serre 生成、消失
% 
%   对应 V.19.1—19.2；解析侧凝聚性取 IV.14.4，高次扭转工具在 IV.20.3 建立。
% 
% * 解析化、局部环映射与态射方向
% 
%   解析化正合性所需的局部环性质在 IV.20.4 就地证明；不得从张量右正合直接推出正合。
% 
% 纲要内容结束

\section{射影比较的代数输入与独立入口}
\label{g4:sec:2001}

待撰写

\subsection{复射影代数对象、概形和结构层的必要构造}
\label{g4:subsec:2001:01}
待撰写

\subsection{代数凝聚层、可逆层、张量与有限呈示}
\label{g4:subsec:2001:02}
待撰写

\subsection{相对仿射构造、Proj 和扭转线丛}
\label{g4:subsec:2001:03}
待撰写

\subsection{仿射无上同调与有限仿射覆盖的 Čech 模型}
\label{g4:subsec:2001:04}
待撰写

\subsection{射影空间扭转层的上同调及 Serre 生成、消失}
\label{g4:subsec:2001:05}
待撰写

\subsection{解析化、局部环映射与态射方向}
\label{g4:subsec:2001:06}
待撰写
% 纲要条目排为编号小节标题，并进入目录与书签；说明段仍保留于纲要注释中。

% !TeX root = ../../Book4.tex
% 纲要标题：### 20.2 Chow 定理
% 纲要位置：几何学计划纲要.md，第 4035 行。
% 正文状态：未撰写；本文件目前是框架。

% 纲要内容开始
% 
% * 射影解析子集与齐次锥
% * 延拓及局部有限性引理
% * 代数方程的构造
% * 图像法与态射代数化
% 
% 纲要内容结束

\section{Chow 定理}
\label{g4:sec:2002}

待撰写

\subsection{射影解析子集与齐次锥}
\label{g4:subsec:2002:01}
待撰写

\subsection{延拓及局部有限性引理}
\label{g4:subsec:2002:02}
待撰写

\subsection{代数方程的构造}
\label{g4:subsec:2002:03}
待撰写

\subsection{图像法与态射代数化}
\label{g4:subsec:2002:04}
待撰写
% 纲要条目排为编号小节标题，并进入目录与书签；说明段仍保留于纲要注释中。

% !TeX root = ../../Book4.tex
% 纲要标题：### 20.3 解析凝聚层的射影工具
% 纲要位置：几何学计划纲要.md，第 4306 行。
% 正文状态：未撰写；本文件目前是框架。

% 纲要内容开始
% 
% * 高次扭转后的生成
% * 有限呈示与核的凝聚性
% * 解析上同调有限性
% * 与代数 Serre 工具的比较
% 
% 纲要内容结束

\section{解析凝聚层的射影工具}
\label{g4:sec:2003}

待撰写

\subsection{高次扭转后的生成}
\label{g4:subsec:2003:01}
待撰写

\subsection{有限呈示与核的凝聚性}
\label{g4:subsec:2003:02}
待撰写

\subsection{解析上同调有限性}
\label{g4:subsec:2003:03}
待撰写

\subsection{与代数 Serre 工具的比较}
\label{g4:subsec:2003:04}
待撰写
% 纲要条目排为编号小节标题，并进入目录与书签；说明段仍保留于纲要注释中。

% !TeX root = ../../Book4.tex
% 纲要标题：### 20.4 射影 GAGA 的连续证明
% 纲要位置：几何学计划纲要.md，第 4313 行。
% 正文状态：未撰写；本文件目前是框架。

% 纲要内容开始
% 
% * 解析化的正合性所需局部环性质
% * 扭转线丛的上同调比较
% * 有限呈示和长正合列归纳
% * 一般凝聚层的上同调比较
% * Hom 比较及全忠实性
% * 解析凝聚层的代数化构造
% * 代数化唯一性及张量、对偶相容
% 
% 纲要内容结束

\section{射影 GAGA 的连续证明}
\label{g4:sec:2004}

待撰写

\subsection{解析化的正合性所需局部环性质}
\label{g4:subsec:2004:01}
待撰写

\subsection{扭转线丛的上同调比较}
\label{g4:subsec:2004:02}
待撰写

\subsection{有限呈示和长正合列归纳}
\label{g4:subsec:2004:03}
待撰写

\subsection{一般凝聚层的上同调比较}
\label{g4:subsec:2004:04}
待撰写

\subsection{Hom 比较及全忠实性}
\label{g4:subsec:2004:05}
待撰写

\subsection{解析凝聚层的代数化构造}
\label{g4:subsec:2004:06}
待撰写

\subsection{代数化唯一性及张量、对偶相容}
\label{g4:subsec:2004:07}
待撰写
% 纲要条目排为编号小节标题，并进入目录与书签；说明段仍保留于纲要注释中。

% !TeX root = ../../Book4.tex
% 纲要标题：### 20.5 比较定理的应用与边界
% 纲要位置：几何学计划纲要.md，第 4059 行。
% 正文状态：未撰写；本文件目前是框架。

% 纲要内容开始
% 
% * 线丛、除子及扩张类
% * 复曲线解析与代数 Riemann–Roch
% * 非固有情形的失效
% * 一般紧非 Kähler 流形的限制
% 
% ---
% 
% 纲要内容结束

\section{比较定理的应用与边界}
\label{g4:sec:2005}

待撰写

\subsection{线丛、除子及扩张类}
\label{g4:subsec:2005:01}
待撰写

\subsection{复曲线解析与代数 Riemann–Roch}
\label{g4:subsec:2005:02}
待撰写

\subsection{非固有情形的失效}
\label{g4:subsec:2005:03}
待撰写

\subsection{一般紧非 Kähler 流形的限制}
\label{g4:subsec:2005:04}
待撰写
% 纲要条目排为编号小节标题，并进入目录与书签；说明段仍保留于纲要注释中。


\begin{chaptersummary}
待撰写
\end{chaptersummary}

\BookExercises
待撰写

% !TeX root = ../../Book4.tex
% 纲要标题：## 第21章　复环面、线丛与阿贝尔簇
% 纲要位置：几何学计划纲要.md，第 4332 行。
% 正文状态：未撰写；本文件目前是框架。

% 纲要内容开始
% 
% > 本章目标与证明范围：将 Appell–Humbert、Riemann 型条件和 theta 截面放入完整主线。一般复环面、可极化复环面和指定极化的阿贝尔簇是不同对象。
% 
% 纲要内容结束

\chapter{复环面、线丛与阿贝尔簇}
\label{g4:ch:21}
\BookChapterNumber{g4:ch:21}

\begin{chapterabstract}
待撰写
\end{chapterabstract}

\begin{learninggoals}
待撰写
\end{learninggoals}

待撰写

% !TeX root = ../../Book4.tex
% 纲要标题：### 21.1 复环面与格结构
% 纲要位置：几何学计划纲要.md，第 4072 行。
% 正文状态：未撰写；本文件目前是框架。

% 纲要内容开始
% 
% * 复向量空间模格
% * 同态、同源与对偶环面
% * 周期矩阵及换基
% * 全纯函数和亚纯函数的区别
% 
% 纲要内容结束

\section{复环面与格结构}
\label{g4:sec:2101}

待撰写

\subsection{复向量空间模格}
\label{g4:subsec:2101:01}
待撰写

\subsection{同态、同源与对偶环面}
\label{g4:subsec:2101:02}
待撰写

\subsection{周期矩阵及换基}
\label{g4:subsec:2101:03}
待撰写

\subsection{全纯函数和亚纯函数的区别}
\label{g4:subsec:2101:04}
待撰写
% 纲要条目排为编号小节标题，并进入目录与书签；说明段仍保留于纲要注释中。

% !TeX root = ../../Book4.tex
% 纲要标题：### 21.2 线丛的 Appell–Humbert 描述
% 纲要位置：几何学计划纲要.md，第 4343 行。
% 正文状态：未撰写；本文件目前是框架。

% 纲要内容开始
% 
% * 自同构因子与余循环
% * Hermitian 型及整交替部分
% * 半特征和等价关系
% * 线丛分类定理的证明
% * 第一陈类与 Pic⁰
% 
% 纲要内容结束

\section{线丛的 Appell–Humbert 描述}
\label{g4:sec:2102}

待撰写

\subsection{自同构因子与余循环}
\label{g4:subsec:2102:01}
待撰写

\subsection{Hermitian 型及整交替部分}
\label{g4:subsec:2102:02}
待撰写

\subsection{半特征和等价关系}
\label{g4:subsec:2102:03}
待撰写

\subsection{线丛分类定理的证明}
\label{g4:subsec:2102:04}
待撰写

\subsection{\texorpdfstring{第一陈类与 \(\operatorname{Pic}^0\)}{第一陈类与 Pic⁰}}
\label{g4:subsec:2102:05}
待撰写
% 纲要条目排为编号小节标题，并进入目录与书签；说明段仍保留于纲要注释中。

% !TeX root = ../../Book4.tex
% 纲要标题：### 21.3 极化及射影性
% 纲要位置：几何学计划纲要.md，第 4351 行。
% 正文状态：未撰写；本文件目前是框架。

% 纲要内容开始
% 
% * 正定 Riemann 型条件
% * 正线丛和 Kodaira 嵌入
% * 极化类型与对偶同源
% * 非可极化环面的模型
% 
% 纲要内容结束

\section{极化及射影性}
\label{g4:sec:2103}

待撰写

\subsection{正定 Riemann 型条件}
\label{g4:subsec:2103:01}
待撰写

\subsection{正线丛和 Kodaira 嵌入}
\label{g4:subsec:2103:02}
待撰写

\subsection{极化类型与对偶同源}
\label{g4:subsec:2103:03}
待撰写

\subsection{非可极化环面的模型}
\label{g4:subsec:2103:04}
待撰写
% 纲要条目排为编号小节标题，并进入目录与书签；说明段仍保留于纲要注释中。

% !TeX root = ../../Book4.tex
% 纲要标题：### 21.4 theta 函数和 Jacobi 簇
% 纲要位置：几何学计划纲要.md，第 4094 行。
% 正文状态：未撰写；本文件目前是框架。

% 纲要内容开始
% 
% * theta 级数的收敛与变换律
% * 截面空间及维数
% * 主极化和 theta 除子
% * Jacobi 模型及族中的周期
% 
% ---
% 
% 纲要内容结束

\section{theta 函数和 Jacobi 簇}
\label{g4:sec:2104}

待撰写

\subsection{theta 级数的收敛与变换律}
\label{g4:subsec:2104:01}
待撰写

\subsection{截面空间及维数}
\label{g4:subsec:2104:02}
待撰写

\subsection{主极化和 theta 除子}
\label{g4:subsec:2104:03}
待撰写

\subsection{Jacobi 模型及族中的周期}
\label{g4:subsec:2104:04}
待撰写
% 纲要条目排为编号小节标题，并进入目录与书签；说明段仍保留于纲要注释中。


\begin{chaptersummary}
待撰写
\end{chaptersummary}

\BookExercises
待撰写


% !TeX root = ../Book4.tex
% 纲要标题：# 第五编　伪全纯分析、紧化与 Floer 理论
% 纲要位置：几何学计划纲要.md，第 4367 行。
% 正文状态：未撰写；本文件目前是框架。

% 纲要内容开始
% 
% 纲要内容结束

\BookPart{伪全纯分析、紧化与 Floer 理论}{g4:part:05}
\BookPartGuide
待撰写

% !TeX root = ../../Book4.tex
% 纲要标题：## 第22章*　伪全纯方程与线性 Fredholm 理论
% 纲要位置：几何学计划纲要.md，第 4369 行。
% 正文状态：未撰写；本文件目前是框架。

% 纲要内容开始
% 
% > 本章目标与证明范围：将局部正则性和线性分析设为独立章。闭曲面与全实边界的指数分别处理；非线性模空间和紧化在后章，不在此提前宣布计数不变量存在。[R2] 非线性输入明确取 III.13.5—13.6；相关线性估计按 III.11—14 的对应节调用。
% 
% 纲要内容结束

\OptionalChapter{伪全纯方程与线性 Fredholm 理论}{g4:ch:22}
\BookChapterNumber{g4:ch:22}

\begin{chapterabstract}
待撰写
\end{chapterabstract}

\begin{learninggoals}
待撰写
\end{learninggoals}

待撰写

% !TeX root = ../../Book4.tex
% 纲要标题：### 22.1 Cauchy–Riemann 型方程
% 纲要位置：几何学计划纲要.md，第 4109 行。
% 正文状态：未撰写；本文件目前是框架。

% 纲要内容开始
% 
% * 相容或驯服概复结构
% * 非线性方程和能量—面积关系
% * 局部坐标表达及主部
% * 线性化与联络选择
% 
% 纲要内容结束

\section{Cauchy–Riemann 型方程}
\label{g4:sec:2201}

待撰写

\subsection{相容或驯服概复结构}
\label{g4:subsec:2201:01}
待撰写

\subsection{非线性方程和能量—面积关系}
\label{g4:subsec:2201:02}
待撰写

\subsection{局部坐标表达及主部}
\label{g4:subsec:2201:03}
待撰写

\subsection{线性化与联络选择}
\label{g4:subsec:2201:04}
待撰写
% 纲要条目排为编号小节标题，并进入目录与书签；说明段仍保留于纲要注释中。

% !TeX root = ../../Book4.tex
% 纲要标题：### 22.2 局部正则性和唯一延拓
% 纲要位置：几何学计划纲要.md，第 4380 行。
% 正文状态：未撰写；本文件目前是框架。

% 纲要内容开始
% 
% * W¹ᵖ 解的椭圆自举
% * 相似性原理的所需形式
% * 临界点、分支点与唯一延拓
% * 简单曲线与多重覆盖
% 
% 纲要内容结束

\section{局部正则性和唯一延拓}
\label{g4:sec:2202}

待撰写

\subsection{\texorpdfstring{\(W^{1,p}\) 解的椭圆自举}{W¹ᵖ 解的椭圆自举}}
\label{g4:subsec:2202:01}
待撰写

\subsection{相似性原理的所需形式}
\label{g4:subsec:2202:02}
待撰写

\subsection{临界点、分支点与唯一延拓}
\label{g4:subsec:2202:03}
待撰写

\subsection{简单曲线与多重覆盖}
\label{g4:subsec:2202:04}
待撰写
% 纲要条目排为编号小节标题，并进入目录与书签；说明段仍保留于纲要注释中。

% !TeX root = ../../Book4.tex
% 纲要标题：### 22.3 Banach 流形及 Fredholm
% 纲要位置：几何学计划纲要.md，第 4387 行。
% 正文状态：未撰写；本文件目前是框架。

% 纲要内容开始
% 
% * Sobolev 映射空间的图册
% * 截面方程的光滑性
% * 线性椭圆算子的 Fredholm 性
% * 核、余核和行列式线
% 
% 纲要内容结束

\section{Banach 流形及 Fredholm}
\label{g4:sec:2203}

待撰写

\subsection{Sobolev 映射空间的图册}
\label{g4:subsec:2203:01}
待撰写

\subsection{截面方程的光滑性}
\label{g4:subsec:2203:02}
待撰写

\subsection{线性椭圆算子的 Fredholm 性}
\label{g4:subsec:2203:03}
待撰写

\subsection{核、余核和行列式线}
\label{g4:subsec:2203:04}
待撰写
% 纲要条目排为编号小节标题，并进入目录与书签；说明段仍保留于纲要注释中。

% !TeX root = ../../Book4.tex
% 纲要标题：### 22.4 指数与边界条件
% 纲要位置：几何学计划纲要.md，第 4394 行。
% 正文状态：未撰写；本文件目前是框架。

% 纲要内容开始
% 
% * 闭曲面上的 Riemann–Roch 指数
% * 标记点和参数域自同构
% * 全实边界与 Maslov 指数
% * 指数规范及具体球面、圆盘算例
% 
% ---
% 
% 纲要内容结束

\section{指数与边界条件}
\label{g4:sec:2204}

待撰写

\subsection{闭曲面上的 Riemann–Roch 指数}
\label{g4:subsec:2204:01}
待撰写

\subsection{标记点和参数域自同构}
\label{g4:subsec:2204:02}
待撰写

\subsection{全实边界与 Maslov 指数}
\label{g4:subsec:2204:03}
待撰写

\subsection{指数规范及具体球面、圆盘算例}
\label{g4:subsec:2204:04}
待撰写
% 纲要条目排为编号小节标题，并进入目录与书签；说明段仍保留于纲要注释中。


\begin{chaptersummary}
待撰写
\end{chaptersummary}

\BookExercises
待撰写

% !TeX root = ../../Book4.tex
% 纲要标题：## 第23章*　曲线模空间、横截性与 Gromov 紧化
% 纲要位置：几何学计划纲要.md，第 4139 行。
% 正文状态：未撰写；本文件目前是框架。

% 纲要内容开始
% 
% > 本章目标与证明范围：把局部模空间、泛型性和全局紧性分别完成。主范围固定紧目标、闭曲面和能量控制；多重覆盖的横截困难明确保留，不用“取一般 J”一语越过。[R2]
% 
% 纲要内容结束

\OptionalChapter{曲线模空间、横截性与 Gromov 紧化}{g4:ch:23}
\BookChapterNumber{g4:ch:23}

\begin{chapterabstract}
待撰写
\end{chapterabstract}

\begin{learninggoals}
待撰写
\end{learninggoals}

待撰写

% !TeX root = ../../Book4.tex
% 纲要标题：### 23.1 模空间的局部构造
% 纲要位置：几何学计划纲要.md，第 4143 行。
% 正文状态：未撰写；本文件目前是框架。

% 纲要内容开始
% 
% * 参数化映射与自同构商
% * 正则点的隐函数定理
% * 标记点及求值映射
% * 稳定性与有限自同构
% 
% 纲要内容结束

\section{模空间的局部构造}
\label{g4:sec:2301}

待撰写

\subsection{参数化映射与自同构商}
\label{g4:subsec:2301:01}
待撰写

\subsection{正则点的隐函数定理}
\label{g4:subsec:2301:02}
待撰写

\subsection{标记点及求值映射}
\label{g4:subsec:2301:03}
待撰写

\subsection{稳定性与有限自同构}
\label{g4:subsec:2301:04}
待撰写
% 纲要条目排为编号小节标题，并进入目录与书签；说明段仍保留于纲要注释中。

% !TeX root = ../../Book4.tex
% 纲要标题：### 23.2 Fredholm 横截的泛函分析
% 纲要位置：几何学计划纲要.md，第 4150 行。
% 正文状态：未撰写；本文件目前是框架。

% 纲要内容开始
% 
% * 可分 Banach 流形和可数图册
% * 有限维核与补空间
% * 局部有限维约化
% * 所需正则性阶数
% 
% 纲要内容结束

\section{Fredholm 横截的泛函分析}
\label{g4:sec:2302}

待撰写

\subsection{可分 Banach 流形和可数图册}
\label{g4:subsec:2302:01}
待撰写

\subsection{有限维核与补空间}
\label{g4:subsec:2302:02}
待撰写

\subsection{局部有限维约化}
\label{g4:subsec:2302:03}
待撰写

\subsection{所需正则性阶数}
\label{g4:subsec:2302:04}
待撰写
% 纲要条目排为编号小节标题，并进入目录与书签；说明段仍保留于纲要注释中。

% !TeX root = ../../Book4.tex
% 纲要标题：### 23.3 Sard–Smale 定理的所选证明
% 纲要位置：几何学计划纲要.md，第 4157 行。
% 正文状态：未撰写；本文件目前是框架。

% 纲要内容开始
% 
% * 有限维 Sard 在局部约化中的使用
% * 正则值的局部剩余性
% * 可数覆盖和 Baire 论证
% * Fredholm 指数与正则性假设
% * 有参数边界的使用范围
% 
% 纲要内容结束

\section{Sard–Smale 定理的所选证明}
\label{g4:sec:2303}

待撰写

\subsection{有限维 Sard 在局部约化中的使用}
\label{g4:subsec:2303:01}
待撰写

\subsection{正则值的局部剩余性}
\label{g4:subsec:2303:02}
待撰写

\subsection{可数覆盖和 Baire 论证}
\label{g4:subsec:2303:03}
待撰写

\subsection{Fredholm 指数与正则性假设}
\label{g4:subsec:2303:04}
待撰写

\subsection{有参数边界的使用范围}
\label{g4:subsec:2303:05}
待撰写
% 纲要条目排为编号小节标题，并进入目录与书签；说明段仍保留于纲要注释中。

% !TeX root = ../../Book4.tex
% 纲要标题：### 23.4 简单曲线的泛型正则性
% 纲要位置：几何学计划纲要.md，第 4165 行。
% 正文状态：未撰写；本文件目前是框架。

% 纲要内容开始
% 
% * 通用 J-模空间
% * 单点注入附近的扰动
% * 唯一延拓与满射性
% * 求值映射的横截性
% * 多重覆盖的独立障碍
% 
% 纲要内容结束

\section{简单曲线的泛型正则性}
\label{g4:sec:2304}

待撰写

\subsection{通用 J-模空间}
\label{g4:subsec:2304:01}
待撰写

\subsection{单点注入附近的扰动}
\label{g4:subsec:2304:02}
待撰写

\subsection{唯一延拓与满射性}
\label{g4:subsec:2304:03}
待撰写

\subsection{求值映射的横截性}
\label{g4:subsec:2304:04}
待撰写

\subsection{多重覆盖的独立障碍}
\label{g4:subsec:2304:05}
待撰写
% 纲要条目排为编号小节标题，并进入目录与书签；说明段仍保留于纲要注释中。

% !TeX root = ../../Book4.tex
% 纲要标题：### 23.5 能量控制和去奇点
% 纲要位置：几何学计划纲要.md，第 4437 行。
% 正文状态：未撰写；本文件目前是框架。

% 纲要内容开始
% 
% * 均值不等式及梯度控制
% * 能量量子化
% * 有限能量去奇点定理
% * 重缩放和非恒定气泡
% 
% 纲要内容结束

\section{能量控制和去奇点}
\label{g4:sec:2305}

待撰写

\subsection{均值不等式及梯度控制}
\label{g4:subsec:2305:01}
待撰写

\subsection{能量量子化}
\label{g4:subsec:2305:02}
待撰写

\subsection{有限能量去奇点定理}
\label{g4:subsec:2305:03}
待撰写

\subsection{重缩放和非恒定气泡}
\label{g4:subsec:2305:04}
待撰写
% 纲要条目排为编号小节标题，并进入目录与书签；说明段仍保留于纲要注释中。

% !TeX root = ../../Book4.tex
% 纲要标题：### 23.6 Gromov 紧性的证明
% 纲要位置：几何学计划纲要.md，第 4180 行。
% 正文状态：未撰写；本文件目前是框架。

% 纲要内容开始
% 
% * 离开集中点的光滑子列
% * 气泡树的逐层抽取
% * 能量恒等式及颈部控制
% * 定义域稳定化
% * 稳定映射的收敛与紧性结论
% 
% 纲要内容结束

\section{Gromov 紧性的证明}
\label{g4:sec:2306}

待撰写

\subsection{离开集中点的光滑子列}
\label{g4:subsec:2306:01}
待撰写

\subsection{气泡树的逐层抽取}
\label{g4:subsec:2306:02}
待撰写

\subsection{能量恒等式及颈部控制}
\label{g4:subsec:2306:03}
待撰写

\subsection{定义域稳定化}
\label{g4:subsec:2306:04}
待撰写

\subsection{稳定映射的收敛与紧性结论}
\label{g4:subsec:2306:05}
待撰写
% 纲要条目排为编号小节标题，并进入目录与书签；说明段仍保留于纲要注释中。

% !TeX root = ../../Book4.tex
% 纲要标题：### 23.7* 四维附加几何
% 纲要位置：几何学计划纲要.md，第 4452 行。
% 正文状态：未撰写；本文件目前是框架。

% 纲要内容开始
% 
% * 交点正性
% * 伴随公式的代表版本
% * 自动横截性的条件
% * 低维性质不能直接搬到任意维数
% 
% ---
% 
% 纲要内容结束

\OptionalSection{四维附加几何}{g4:sec:2307}

待撰写

\subsection{交点正性}
\label{g4:subsec:2307:01}
待撰写

\subsection{伴随公式的代表版本}
\label{g4:subsec:2307:02}
待撰写

\subsection{自动横截性的条件}
\label{g4:subsec:2307:03}
待撰写

\subsection{低维性质不能直接搬到任意维数}
\label{g4:subsec:2307:04}
待撰写
% 纲要条目排为编号小节标题，并进入目录与书签；说明段仍保留于纲要注释中。


\begin{chaptersummary}
待撰写
\end{chaptersummary}

\BookExercises
待撰写

% !TeX root = ../../Book4.tex
% 纲要标题：## 第24章*　球—柱非挤压与辛容量
% 纲要位置：几何学计划纲要.md，第 4461 行。
% 正文状态：未撰写；本文件目前是框架。

% 纲要内容开始
% 
% > 本章目标与证明范围：标准非挤压给出连续完整证明，而非应用目录。固定 S²×环面的球面类和一点求值映射路线，使紧性、度与局部面积估计的使用能够逐项核验。 n=1 用面积直接处理，曲线模空间证明从 n≥2 开始。
% 
% 纲要内容结束

\OptionalChapter{球—柱非挤压与辛容量}{g4:ch:24}
\BookChapterNumber{g4:ch:24}

\begin{chapterabstract}
待撰写
\end{chapterabstract}

\begin{learninggoals}
待撰写
\end{learninggoals}

待撰写

% !TeX root = ../../Book4.tex
% 纲要标题：### 24.1 嵌入问题与紧化模型
% 纲要位置：几何学计划纲要.md，第 4465 行。
% 正文状态：未撰写；本文件目前是框架。

% 纲要内容开始
% 
% * 球、柱和面积的归一化
% * 把柱的相关紧区域放入 S²×T²ⁿ⁻²
% * 球面因子面积的任意小增量
% * 嵌入球上适应标准复结构的 J
% 
% 纲要内容结束

\section{嵌入问题与紧化模型}
\label{g4:sec:2401}

待撰写

\subsection{球、柱和面积的归一化}
\label{g4:subsec:2401:01}
待撰写

\subsection{\texorpdfstring{把柱的相关紧区域放入 \(S^2\times T^{2n-2}\)}{把柱的相关紧区域放入 S²×T²ⁿ⁻²}}
\label{g4:subsec:2401:02}
待撰写

\subsection{球面因子面积的任意小增量}
\label{g4:subsec:2401:03}
待撰写

\subsection{嵌入球上适应标准复结构的 J}
\label{g4:subsec:2401:04}
待撰写
% 纲要条目排为编号小节标题，并进入目录与书签；说明段仍保留于纲要注释中。

% !TeX root = ../../Book4.tex
% 纲要标题：### 24.2 指定球面类的紧性
% 纲要位置：几何学计划纲要.md，第 4472 行。
% 正文状态：未撰写；本文件目前是框架。

% 纲要内容开始
% 
% * 球面因子类及其辛面积
% * 该类的不可多重覆盖性
% * 球泡分裂的面积排除
% * 一点标记模空间的紧化
% 
% 纲要内容结束

\section{指定球面类的紧性}
\label{g4:sec:2402}

待撰写

\subsection{球面因子类及其辛面积}
\label{g4:subsec:2402:01}
待撰写

\subsection{该类的不可多重覆盖性}
\label{g4:subsec:2402:02}
待撰写

\subsection{球泡分裂的面积排除}
\label{g4:subsec:2402:03}
待撰写

\subsection{一点标记模空间的紧化}
\label{g4:subsec:2402:04}
待撰写
% 纲要条目排为编号小节标题，并进入目录与书签；说明段仍保留于纲要注释中。

% !TeX root = ../../Book4.tex
% 纲要标题：### 24.3 求值映射的度
% 纲要位置：几何学计划纲要.md，第 4479 行。
% 正文状态：未撰写；本文件目前是框架。

% 纲要内容开始
% 
% * 乘积复结构下的唯一球面纤维
% * 保持球内 J 的相对横截性
% 
%   正面积球面不能完全落在辛形式恰当的球内；结合简单曲线的单点注入性质，在球外扰动 J，保留球内的标准结构。
% 
% * 正则参数族和配边
% 
%   对一标记、模掉重参数化的球面模空间核验实维数为 2n；在固定类不分裂的紧性条件下，以一参数配边保持求值度。
% 
% * 求值映射的模2度为一
% * 对指定点的曲线存在
% 
%   先对正则数据由度为一取得满射，再以紧性传到需要的极限数据；不能仅以泛型点替代指定球心。
% 
% 纲要内容结束

\section{求值映射的度}
\label{g4:sec:2403}

待撰写

\subsection{乘积复结构下的唯一球面纤维}
\label{g4:subsec:2403:01}
待撰写

\subsection{保持球内 J 的相对横截性}
\label{g4:subsec:2403:02}
待撰写

\subsection{正则参数族和配边}
\label{g4:subsec:2403:03}
待撰写

\subsection{求值映射的模2度为一}
\label{g4:subsec:2403:04}
待撰写

\subsection{对指定点的曲线存在}
\label{g4:subsec:2403:05}
待撰写
% 纲要条目排为编号小节标题，并进入目录与书签；说明段仍保留于纲要注释中。

% !TeX root = ../../Book4.tex
% 纲要标题：### 24.4 局部面积与非挤压
% 纲要位置：几何学计划纲要.md，第 4495 行。
% 正文状态：未撰写；本文件目前是框架。

% 纲要内容开始
% 
% * 嵌入球中标准全纯曲线的单调性
% * 经过中心的面积下界
% * 整体球面面积上界
% * 缩小球和面积增量趋零
% * 得到 R≤r 的标准结论
% 
% 纲要内容结束

\section{局部面积与非挤压}
\label{g4:sec:2404}

待撰写

\subsection{嵌入球中标准全纯曲线的单调性}
\label{g4:subsec:2404:01}
待撰写

\subsection{经过中心的面积下界}
\label{g4:subsec:2404:02}
待撰写

\subsection{整体球面面积上界}
\label{g4:subsec:2404:03}
待撰写

\subsection{缩小球和面积增量趋零}
\label{g4:subsec:2404:04}
待撰写

\subsection{\texorpdfstring{得到 \(R\le r\) 的标准结论}{得到 R≤r 的标准结论}}
\label{g4:subsec:2404:05}
待撰写
% 纲要条目排为编号小节标题，并进入目录与书签；说明段仍保留于纲要注释中。

% !TeX root = ../../Book4.tex
% 纲要标题：### 24.5 容量及其他整体障碍
% 纲要位置：几何学计划纲要.md，第 4503 行。
% 正文状态：未撰写；本文件目前是框架。

% 纲要内容开始
% 
% * 容量的单调性、共形性和归一化
% * Gromov 宽度及基本计算
% * 与体积障碍的严格区别
% * 其他容量构造需要的附加理论
% 
% ---
% 
% 纲要内容结束

\section{容量及其他整体障碍}
\label{g4:sec:2405}

待撰写

\subsection{容量的单调性、共形性和归一化}
\label{g4:subsec:2405:01}
待撰写

\subsection{Gromov 宽度及基本计算}
\label{g4:subsec:2405:02}
待撰写

\subsection{与体积障碍的严格区别}
\label{g4:subsec:2405:03}
待撰写

\subsection{其他容量构造需要的附加理论}
\label{g4:subsec:2405:04}
待撰写
% 纲要条目排为编号小节标题，并进入目录与书签；说明段仍保留于纲要注释中。


\begin{chaptersummary}
待撰写
\end{chaptersummary}

\BookExercises
待撰写

% !TeX root = ../../Book4.tex
% 纲要标题：## 第25章*　Hamilton Floer 的作用、指数与渐近分析
% 纲要位置：几何学计划纲要.md，第 4248 行。
% 正文状态：未撰写；本文件目前是框架。

% 纲要内容开始
% 
% > 本章目标与证明范围：固定闭辛流形、[ω]与c₁在π₂上均消失、非退化 Hamilton、可缩一周期轨道及F₂系数。先建立作用和渐近算子，再谈模空间；不把能量有限自动等同于已证明的指数衰减。
% 
% 纲要内容结束

\OptionalChapter{Hamilton Floer 的作用、指数与渐近分析}{g4:ch:25}
\BookChapterNumber{g4:ch:25}

\begin{chapterabstract}
待撰写
\end{chapterabstract}

\begin{learninggoals}
待撰写
\end{learninggoals}

待撰写

% !TeX root = ../../Book4.tex
% 纲要标题：### 25.1 轨道、封口与作用
% 纲要位置：几何学计划纲要.md，第 4252 行。
% 正文状态：未撰写；本文件目前是框架。

% 纲要内容开始
% 
% * 周期轨道和封口等价
% * 作用的封口独立性
% * 统一取 A_H=−∫u*ω−∫Hdt
% * 变分得到 ∂su+J(∂tu−X_H)=0
% * 作用下降及能量恒等式
% 
% 纲要内容结束

\section{轨道、封口与作用}
\label{g4:sec:2501}

待撰写

\subsection{周期轨道和封口等价}
\label{g4:subsec:2501:01}
待撰写

\subsection{作用的封口独立性}
\label{g4:subsec:2501:02}
待撰写

\subsection{\texorpdfstring{统一取 \(A_H=-\int u^*\omega-\int H\,\mathrm{d}t\)}{统一取 A\_H=−∫u*\GeometryPdfSmallOmega{}−∫Hdt}}
\label{g4:subsec:2501:03}
待撰写

\subsection{\texorpdfstring{变分得到 \(\partial_su+J(\partial_tu-X_H)=0\)}{变分得到 ∂su+J(∂tu−X\_H)=0}}
\label{g4:subsec:2501:04}
待撰写

\subsection{作用下降及能量恒等式}
\label{g4:subsec:2501:05}
待撰写
% 纲要条目排为编号小节标题，并进入目录与书签；说明段仍保留于纲要注释中。

% !TeX root = ../../Book4.tex
% 纲要标题：### 25.2 渐近算子与有限能量极限
% 纲要位置：几何学计划纲要.md，第 4260 行。
% 正文状态：未撰写；本文件目前是框架。

% 纲要内容开始
% 
% * 非退化轨道的渐近自伴算子
% * 有限能量解的轨道极限
% * 谱隙与 L² 衰减不等式
% * 高阶指数衰减
% * 平移和加权范数
% 
% 纲要内容结束

\section{渐近算子与有限能量极限}
\label{g4:sec:2502}

待撰写

\subsection{非退化轨道的渐近自伴算子}
\label{g4:subsec:2502:01}
待撰写

\subsection{有限能量解的轨道极限}
\label{g4:subsec:2502:02}
待撰写

\subsection{\texorpdfstring{谱隙与 \(L^2\) 衰减不等式}{谱隙与 L² 衰减不等式}}
\label{g4:subsec:2502:03}
待撰写

\subsection{高阶指数衰减}
\label{g4:subsec:2502:04}
待撰写

\subsection{平移和加权范数}
\label{g4:subsec:2502:05}
待撰写
% 纲要条目排为编号小节标题，并进入目录与书签；说明段仍保留于纲要注释中。

% !TeX root = ../../Book4.tex
% 纲要标题：### 25.3 辛路径与轨道分次
% 纲要位置：几何学计划纲要.md，第 4532 行。
% 正文状态：未撰写；本文件目前是框架。

% 纲要内容开始
% 
% * 平凡化和封口的选择
% * Maslov 与 Conley–Zehnder 指数
% * 正向小旋转的规范
% * c₁ 消失与封口独立性
% * 与 Morse(−H) 对应的分次
% 
% 纲要内容结束

\section{辛路径与轨道分次}
\label{g4:sec:2503}

待撰写

\subsection{平凡化和封口的选择}
\label{g4:subsec:2503:01}
待撰写

\subsection{Maslov 与 Conley–Zehnder 指数}
\label{g4:subsec:2503:02}
待撰写

\subsection{正向小旋转的规范}
\label{g4:subsec:2503:03}
待撰写

\subsection{\texorpdfstring{\(c_1\) 消失与封口独立性}{c₁ 消失与封口独立性}}
\label{g4:subsec:2503:04}
待撰写

\subsection{\texorpdfstring{与 \(\operatorname{Morse}(-H)\) 对应的分次}{与 Morse(−H) 对应的分次}}
\label{g4:subsec:2503:05}
待撰写
% 纲要条目排为编号小节标题，并进入目录与书签；说明段仍保留于纲要注释中。

% !TeX root = ../../Book4.tex
% 纲要标题：### 25.4 轨线算子的 Fredholm 指数
% 纲要位置：几何学计划纲要.md，第 4276 行。
% 正文状态：未撰写；本文件目前是框架。

% 纲要内容开始
% 
% * 带权空间上的闭值域
% * 渐近算子的可逆性
% * 谱流和 Fredholm 指数定理
% * 核余核及平移方向
% * 指数差与预期维数
% 
% 纲要内容结束

\section{轨线算子的 Fredholm 指数}
\label{g4:sec:2504}

待撰写

\subsection{带权空间上的闭值域}
\label{g4:subsec:2504:01}
待撰写

\subsection{渐近算子的可逆性}
\label{g4:subsec:2504:02}
待撰写

\subsection{谱流和 Fredholm 指数定理}
\label{g4:subsec:2504:03}
待撰写

\subsection{核余核及平移方向}
\label{g4:subsec:2504:04}
待撰写

\subsection{指数差与预期维数}
\label{g4:subsec:2504:05}
待撰写
% 纲要条目排为编号小节标题，并进入目录与书签；说明段仍保留于纲要注释中。

% !TeX root = ../../Book4.tex
% 纲要标题：### 25.5 正则轨线模空间
% 纲要位置：几何学计划纲要.md，第 4548 行。
% 正文状态：未撰写；本文件目前是框架。

% 纲要内容开始
% 
% * H、J 的通用扰动
% * 线性化满射及泛型性
% * 平移作用和维数
% * 零能量轨线与常值柱面
% 
% ---
% 
% 纲要内容结束

\section{正则轨线模空间}
\label{g4:sec:2505}

待撰写

\subsection{H、J 的通用扰动}
\label{g4:subsec:2505:01}
待撰写

\subsection{线性化满射及泛型性}
\label{g4:subsec:2505:02}
待撰写

\subsection{平移作用和维数}
\label{g4:subsec:2505:03}
待撰写

\subsection{零能量轨线与常值柱面}
\label{g4:subsec:2505:04}
待撰写
% 纲要条目排为编号小节标题，并进入目录与书签；说明段仍保留于纲要注释中。


\begin{chaptersummary}
待撰写
\end{chaptersummary}

\BookExercises
待撰写

% !TeX root = ../../Book4.tex
% 纲要标题：## 第26章*　Floer 紧化、粘合与链复形
% 纲要位置：几何学计划纲要.md，第 4293 行。
% 正文状态：未撰写；本文件目前是框架。

% 纲要内容开始
% 
% > 本章目标与证明范围：这是构造 Floer 微分不可压缩的分析核心。分别证明排泡、断裂紧性、统一右逆、非线性修正和紧化边界结构；F₂ 只免除取向符号，不免除这些证明。
% 
% 纲要内容结束

\OptionalChapter{Floer 紧化、粘合与链复形}{g4:ch:26}
\BookChapterNumber{g4:ch:26}

\begin{chapterabstract}
待撰写
\end{chapterabstract}

\begin{learninggoals}
待撰写
\end{learninggoals}

待撰写

% !TeX root = ../../Book4.tex
% 纲要标题：### 26.1 能量界和断裂紧性
% 纲要位置：几何学计划纲要.md，第 4297 行。
% 正文状态：未撰写；本文件目前是框架。

% 纲要内容开始
% 
% * 作用窗与统一能量界
% * 梯度爆破导致球泡的抽取
% * [ω]在π₂上消失排除非恒定球泡
% * 时间平移后的收敛
% * 有限断裂及轨道序列控制
% 
% 纲要内容结束

\section{能量界和断裂紧性}
\label{g4:sec:2601}

待撰写

\subsection{作用窗与统一能量界}
\label{g4:subsec:2601:01}
待撰写

\subsection{梯度爆破导致球泡的抽取}
\label{g4:subsec:2601:02}
待撰写

\subsection{\texorpdfstring{\([\omega]\)在\(\pi_2\)上消失排除非恒定球泡}{[\GeometryPdfSmallOmega{}]在\GeometryPdfPi{}₂上消失排除非恒定球泡}}
\label{g4:subsec:2601:03}
待撰写

\subsection{时间平移后的收敛}
\label{g4:subsec:2601:04}
待撰写

\subsection{有限断裂及轨道序列控制}
\label{g4:subsec:2601:05}
待撰写
% 纲要条目排为编号小节标题，并进入目录与书签；说明段仍保留于纲要注释中。

% !TeX root = ../../Book4.tex
% 纲要标题：### 26.2 预粘合与误差
% 纲要位置：几何学计划纲要.md，第 4569 行。
% 正文状态：未撰写；本文件目前是框架。

% 纲要内容开始
% 
% * 断裂轨线的颈长参数
% * 截断函数和指数尾部
% * 预粘合方程误差估计
% * 平移规范和核方向的处理
% 
% 纲要内容结束

\section{预粘合与误差}
\label{g4:sec:2602}

待撰写

\subsection{断裂轨线的颈长参数}
\label{g4:subsec:2602:01}
待撰写

\subsection{截断函数和指数尾部}
\label{g4:subsec:2602:02}
待撰写

\subsection{预粘合方程误差估计}
\label{g4:subsec:2602:03}
待撰写

\subsection{平移规范和核方向的处理}
\label{g4:subsec:2602:04}
待撰写
% 纲要条目排为编号小节标题，并进入目录与书签；说明段仍保留于纲要注释中。

% !TeX root = ../../Book4.tex
% 纲要标题：### 26.3 统一右逆
% 纲要位置：几何学计划纲要.md，第 4576 行。
% 正文状态：未撰写；本文件目前是框架。

% 纲要内容开始
% 
% * 极限线性算子的分裂
% * 拼接右逆及交换子误差
% * 长颈下的统一算子范数
% * Neumann 级数修正
% * 权和范数选择的相容性
% 
% 纲要内容结束

\section{统一右逆}
\label{g4:sec:2603}

待撰写

\subsection{极限线性算子的分裂}
\label{g4:subsec:2603:01}
待撰写

\subsection{拼接右逆及交换子误差}
\label{g4:subsec:2603:02}
待撰写

\subsection{长颈下的统一算子范数}
\label{g4:subsec:2603:03}
待撰写

\subsection{Neumann 级数修正}
\label{g4:subsec:2603:04}
待撰写

\subsection{权和范数选择的相容性}
\label{g4:subsec:2603:05}
待撰写
% 纲要条目排为编号小节标题，并进入目录与书签；说明段仍保留于纲要注释中。

% !TeX root = ../../Book4.tex
% 纲要标题：### 26.4 非线性修正与边界图册
% 纲要位置：几何学计划纲要.md，第 4320 行。
% 正文状态：未撰写；本文件目前是框架。

% 纲要内容开始
% 
% * 二次余项的 Lipschitz 估计
% * 压缩映射或定量隐函数定理
% * 粘合解的唯一性
% * 所有邻近轨线的反向识别
% * 一维紧化的边界结构
% * 颈长坐标的光滑性
% 
%   证明粘合参数给出真正的端部图册，而不只是一组存在解。
% 
% * 边界覆盖的完备性
% 
%   利用反粘合排除额外端点；这是把计数等式转为 d²=0 的最后分析前提。
% 
% 纲要内容结束

\section{非线性修正与边界图册}
\label{g4:sec:2604}

待撰写

\subsection{二次余项的 Lipschitz 估计}
\label{g4:subsec:2604:01}
待撰写

\subsection{压缩映射或定量隐函数定理}
\label{g4:subsec:2604:02}
待撰写

\subsection{粘合解的唯一性}
\label{g4:subsec:2604:03}
待撰写

\subsection{所有邻近轨线的反向识别}
\label{g4:subsec:2604:04}
待撰写

\subsection{一维紧化的边界结构}
\label{g4:subsec:2604:05}
待撰写

\subsection{颈长坐标的光滑性}
\label{g4:subsec:2604:06}
待撰写

\subsection{边界覆盖的完备性}
\label{g4:subsec:2604:07}
待撰写
% 纲要条目排为编号小节标题，并进入目录与书签；说明段仍保留于纲要注释中。

% !TeX root = ../../Book4.tex
% 纲要标题：### 26.5 Floer 微分
% 纲要位置：几何学计划纲要.md，第 4335 行。
% 正文状态：未撰写；本文件目前是框架。

% 纲要内容开始
% 
% * 零维模空间的有限计数
% * 一维边界与两段轨线
% * d²=0 的证明
% * 链群、分次和系数的完整定义
% 
% ---
% 
% 纲要内容结束

\section{Floer 微分}
\label{g4:sec:2605}

待撰写

\subsection{零维模空间的有限计数}
\label{g4:subsec:2605:01}
待撰写

\subsection{一维边界与两段轨线}
\label{g4:subsec:2605:02}
待撰写

\subsection{\texorpdfstring{\(\mathrm{d}^2=0\) 的证明}{d²=0 的证明}}
\label{g4:subsec:2605:03}
待撰写

\subsection{链群、分次和系数的完整定义}
\label{g4:subsec:2605:04}
待撰写
% 纲要条目排为编号小节标题，并进入目录与书签；说明段仍保留于纲要注释中。


\begin{chaptersummary}
待撰写
\end{chaptersummary}

\BookExercises
待撰写

% !TeX root = ../../Book4.tex
% 纲要标题：## 第27章*　延拓、Morse 比较与 Arnold 不等式
% 纲要位置：几何学计划纲要.md，第 4344 行。
% 正文状态：未撰写；本文件目前是框架。

% 纲要内容开始
% 
% > 本章目标与证明范围：完成受限 Hamilton 理论的不变性和同调识别。比较采用明确的混合 Morse–Floer 构造；不以“小 Hamilton 的所有轨线显然与时间无关”替代证明。
% 
% 纲要内容结束

\OptionalChapter{延拓、Morse 比较与 Arnold 不等式}{g4:ch:27}
\BookChapterNumber{g4:ch:27}

\begin{chapterabstract}
待撰写
\end{chapterabstract}

\begin{learninggoals}
待撰写
\end{learninggoals}

待撰写

% !TeX root = ../../Book4.tex
% 纲要标题：### 27.1 延拓映射
% 纲要位置：几何学计划纲要.md，第 4612 行。
% 正文状态：未撰写；本文件目前是框架。

% 纲要内容开始
% 
% * 变化数据的能量恒等式
% * 延拓轨线的紧性和横截性
% * 链映射方程
% * 复合的长颈极限
% * 参数族链同伦及不变性
% 
% 纲要内容结束

\section{延拓映射}
\label{g4:sec:2701}

待撰写

\subsection{变化数据的能量恒等式}
\label{g4:subsec:2701:01}
待撰写

\subsection{延拓轨线的紧性和横截性}
\label{g4:subsec:2701:02}
待撰写

\subsection{链映射方程}
\label{g4:subsec:2701:03}
待撰写

\subsection{复合的长颈极限}
\label{g4:subsec:2701:04}
待撰写

\subsection{参数族链同伦及不变性}
\label{g4:subsec:2701:05}
待撰写
% 纲要条目排为编号小节标题，并进入目录与书签；说明段仍保留于纲要注释中。

% !TeX root = ../../Book4.tex
% 纲要标题：### 27.2 Morse–Floer 混合模空间
% 纲要位置：几何学计划纲要.md，第 4356 行。
% 正文状态：未撰写；本文件目前是框架。

% 纲要内容开始
% 
% * 半圆柱和梯度半轨线
% * 匹配条件及指数
% * 紧化和计数
% * 与第二卷 Morse 复形的链映射
% 
% 纲要内容结束

\section{Morse–Floer 混合模空间}
\label{g4:sec:2702}

待撰写

\subsection{半圆柱和梯度半轨线}
\label{g4:subsec:2702:01}
待撰写

\subsection{匹配条件及指数}
\label{g4:subsec:2702:02}
待撰写

\subsection{紧化和计数}
\label{g4:subsec:2702:03}
待撰写

\subsection{与第二卷 Morse 复形的链映射}
\label{g4:subsec:2702:04}
待撰写
% 纲要条目排为编号小节标题，并进入目录与书签；说明段仍保留于纲要注释中。

% !TeX root = ../../Book4.tex
% 纲要标题：### 27.3 比较映射互逆的证明
% 纲要位置：几何学计划纲要.md，第 4627 行。
% 正文状态：未撰写；本文件目前是框架。

% 纲要内容开始
% 
% * 反向混合构造
% * 颈长趋无穷的复合分解
% * 颈长趋零与恒等模型
% * 参数边界和链同伦
% * HF 与奇异同调的归一化同构
% 
% 纲要内容结束

\section{比较映射互逆的证明}
\label{g4:sec:2703}

待撰写

\subsection{反向混合构造}
\label{g4:subsec:2703:01}
待撰写

\subsection{颈长趋无穷的复合分解}
\label{g4:subsec:2703:02}
待撰写

\subsection{颈长趋零与恒等模型}
\label{g4:subsec:2703:03}
待撰写

\subsection{参数边界和链同伦}
\label{g4:subsec:2703:04}
待撰写

\subsection{HF 与奇异同调的归一化同构}
\label{g4:subsec:2703:05}
待撰写
% 纲要条目排为编号小节标题，并进入目录与书签；说明段仍保留于纲要注释中。

% !TeX root = ../../Book4.tex
% 纲要标题：### 27.4 Arnold 型结论
% 纲要位置：几何学计划纲要.md，第 4635 行。
% 正文状态：未撰写；本文件目前是框架。

% 纲要内容开始
% 
% * 非退化轨道数与链群维数
% * Betti 数和下界
% * Hamilton 不动点的解释
% * 退化、非可缩及一般球面类的额外问题
% 
% 纲要内容结束

\section{Arnold 型结论}
\label{g4:sec:2704}

待撰写

\subsection{非退化轨道数与链群维数}
\label{g4:subsec:2704:01}
待撰写

\subsection{Betti 数和下界}
\label{g4:subsec:2704:02}
待撰写

\subsection{Hamilton 不动点的解释}
\label{g4:subsec:2704:03}
待撰写

\subsection{退化、非可缩及一般球面类的额外问题}
\label{g4:subsec:2704:04}
待撰写
% 纲要条目排为编号小节标题，并进入目录与书签；说明段仍保留于纲要注释中。

% !TeX root = ../../Book4.tex
% 纲要标题：### 27.5 作用滤过与谱不变量
% 纲要位置：几何学计划纲要.md，第 4642 行。
% 正文状态：未撰写；本文件目前是框架。

% 纲要内容开始
% 
% * 过滤子复形和延拓估计
% * 同调类的谱数
% * 谱性及连续性所需有限性
% * Hamilton 距离与基本下界
% * 与持续结构附录的接口
% 
% ---
% 
% 纲要内容结束

\section{作用滤过与谱不变量}
\label{g4:sec:2705}

待撰写

\subsection{过滤子复形和延拓估计}
\label{g4:subsec:2705:01}
待撰写

\subsection{同调类的谱数}
\label{g4:subsec:2705:02}
待撰写

\subsection{谱性及连续性所需有限性}
\label{g4:subsec:2705:03}
待撰写

\subsection{Hamilton 距离与基本下界}
\label{g4:subsec:2705:04}
待撰写

\subsection{与持续结构附录的接口}
\label{g4:subsec:2705:05}
待撰写
% 纲要条目排为编号小节标题，并进入目录与书签；说明段仍保留于纲要注释中。


\begin{chaptersummary}
待撰写
\end{chaptersummary}

\BookExercises
待撰写

% !TeX root = ../../Book4.tex
% 纲要标题：## 第28章*　Lagrange Floer、量子乘积与曲线计数
% 纲要位置：几何学计划纲要.md，第 4388 行。
% 正文状态：未撰写；本文件目前是框架。

% 纲要内容开始
% 
% > 本章目标与证明范围：在已有完整 Hamilton 模块上深入说明新障碍，并完成一个精确无泡的 Lagrange 代表模型。一般虚类和 Fukaya 构造保留为明确外部框架，而不是声称本章已覆盖全部理论。
% 
% 纲要内容结束

\OptionalChapter{Lagrange Floer、量子乘积与曲线计数}{g4:ch:28}
\BookChapterNumber{g4:ch:28}

\begin{chapterabstract}
待撰写
\end{chapterabstract}

\begin{learninggoals}
待撰写
\end{learninggoals}

待撰写

% !TeX root = ../../Book4.tex
% 纲要标题：### 28.1 全实边界的分析
% 纲要位置：几何学计划纲要.md，第 4656 行。
% 正文状态：未撰写；本文件目前是框架。

% 纲要内容开始
% 
% * 条带方程及交点渐近
% * 相对作用和 Maslov 分次
% * 边界正则性
% * 圆盘冒泡与 m₀ 障碍
% 
% 纲要内容结束

\section{全实边界的分析}
\label{g4:sec:2801}

待撰写

\subsection{条带方程及交点渐近}
\label{g4:subsec:2801:01}
待撰写

\subsection{相对作用和 Maslov 分次}
\label{g4:subsec:2801:02}
待撰写

\subsection{边界正则性}
\label{g4:subsec:2801:03}
待撰写

\subsection{\texorpdfstring{圆盘冒泡与 \(m_0\) 障碍}{圆盘冒泡与 m₀ 障碍}}
\label{g4:subsec:2801:04}
待撰写
% 纲要条目排为编号小节标题，并进入目录与书签；说明段仍保留于纲要注释中。

% !TeX root = ../../Book4.tex
% 纲要标题：### 28.2 精确 Lagrange 的代表模块
% 纲要位置：几何学计划纲要.md，第 4663 行。
% 正文状态：未撰写；本文件目前是框架。

% 纲要内容开始
% 
% * 余切丛中的零截面和精确图
% * 紧支或无穷远凸性条件
% * 无球泡、无圆盘泡的能量论证
% * 交点 Floer 与 Morse 比较
% * 非紧目标的逃逸控制单列证明
% 
% 纲要内容结束

\section{精确 Lagrange 的代表模块}
\label{g4:sec:2802}

待撰写

\subsection{余切丛中的零截面和精确图}
\label{g4:subsec:2802:01}
待撰写

\subsection{紧支或无穷远凸性条件}
\label{g4:subsec:2802:02}
待撰写

\subsection{无球泡、无圆盘泡的能量论证}
\label{g4:subsec:2802:03}
待撰写

\subsection{交点 Floer 与 Morse 比较}
\label{g4:subsec:2802:04}
待撰写

\subsection{非紧目标的逃逸控制单列证明}
\label{g4:subsec:2802:05}
待撰写
% 纲要条目排为编号小节标题，并进入目录与书签；说明段仍保留于纲要注释中。

% !TeX root = ../../Book4.tex
% 纲要标题：### 28.3 乘积及同伦代数
% 纲要位置：几何学计划纲要.md，第 4407 行。
% 正文状态：未撰写；本文件目前是框架。

% 纲要内容开始
% 
% * 三端曲面和配对模空间
% * 紧化边界产生的结合关系
% * 量子修正与球面计数
% * A∞ 关系的几何来源
% 
% 纲要内容结束

\section{乘积及同伦代数}
\label{g4:sec:2803}

待撰写

\subsection{三端曲面和配对模空间}
\label{g4:subsec:2803:01}
待撰写

\subsection{紧化边界产生的结合关系}
\label{g4:subsec:2803:02}
待撰写

\subsection{量子修正与球面计数}
\label{g4:subsec:2803:03}
待撰写

\subsection{\texorpdfstring{\(A_\infty\) 关系的几何来源}{A∞ 关系的几何来源}}
\label{g4:subsec:2803:04}
待撰写
% 纲要条目排为编号小节标题，并进入目录与书签；说明段仍保留于纲要注释中。

% !TeX root = ../../Book4.tex
% 纲要标题：### 28.4 一般计数的额外结构
% 纲要位置：几何学计划纲要.md，第 4414 行。
% 正文状态：未撰写；本文件目前是框架。

% 纲要内容开始
% 
% * 稳定映射和虚维数
% * 多重覆盖、取向及虚扰动
% * Gromov–Witten 与 Fukaya 的不同对象
% * 只在已指定的几何条件下使用相应定理
% 
% ---
% 
% 纲要内容结束

\section{一般计数的额外结构}
\label{g4:sec:2804}

待撰写

\subsection{稳定映射和虚维数}
\label{g4:subsec:2804:01}
待撰写

\subsection{多重覆盖、取向及虚扰动}
\label{g4:subsec:2804:02}
待撰写

\subsection{Gromov–Witten 与 Fukaya 的不同对象}
\label{g4:subsec:2804:03}
待撰写

\subsection{只在已指定的几何条件下使用相应定理}
\label{g4:subsec:2804:04}
待撰写
% 纲要条目排为编号小节标题，并进入目录与书签；说明段仍保留于纲要注释中。


\begin{chaptersummary}
待撰写
\end{chaptersummary}

\BookExercises
待撰写


% !TeX root = ../Book4.tex
% 纲要标题：# 第六编　复形变、特殊结构与周期
% 纲要位置：几何学计划纲要.md，第 4423 行。
% 正文状态：未撰写；本文件目前是框架。

% 纲要内容开始
% 
% 纲要内容结束

\BookPart{复形变、特殊结构与周期}{g4:part:06}
\BookPartGuide
待撰写

% !TeX root = ../../Book4.tex
% 纲要标题：## 第29章*　Kodaira–Spencer 与 Kuranishi 形变
% 纲要位置：几何学计划纲要.md，第 4697 行。
% 正文状态：未撰写；本文件目前是框架。

% 纲要内容开始
% 
% > 本章目标与证明范围：从实际光滑族到无穷小类，再以非线性可积方程构造局部障碍模型。Teichmüller 已在曲线章处理，本章不再用曲线参数替代一般复流形的障碍理论。 非线性输入明确取 III.13.5—13.6；相关线性估计按 III.11—14 的对应节调用。
% 
% 纲要内容结束

\OptionalChapter{Kodaira–Spencer 与 Kuranishi 形变}{g4:ch:29}
\BookChapterNumber{g4:ch:29}

\begin{chapterabstract}
待撰写
\end{chapterabstract}

\begin{learninggoals}
待撰写
\end{learninggoals}

待撰写

% !TeX root = ../../Book4.tex
% 纲要标题：### 29.1 光滑族与无穷小变化
% 纲要位置：几何学计划纲要.md，第 4701 行。
% 正文状态：未撰写；本文件目前是框架。

% 纲要内容开始
% 
% * 固有淹没的 Ehresmann 平凡化
% * 复结构随参数的变化
% * Kodaira–Spencer 映射
% * 平凡变化与向量场作用
% 
% 纲要内容结束

\section{光滑族与无穷小变化}
\label{g4:sec:2901}

待撰写

\subsection{固有淹没的 Ehresmann 平凡化}
\label{g4:subsec:2901:01}
待撰写

\subsection{复结构随参数的变化}
\label{g4:subsec:2901:02}
待撰写

\subsection{Kodaira–Spencer 映射}
\label{g4:subsec:2901:03}
待撰写

\subsection{平凡变化与向量场作用}
\label{g4:subsec:2901:04}
待撰写
% 纲要条目排为编号小节标题，并进入目录与书签；说明段仍保留于纲要注释中。

% !TeX root = ../../Book4.tex
% 纲要标题：### 29.2 形变方程和上同调
% 纲要位置：几何学计划纲要.md，第 4436 行。
% 正文状态：未撰写；本文件目前是框架。

% 纲要内容开始
% 
% * Beltrami 型丛值(0,1)形式
% * Maurer–Cartan 可积方程
% * 一阶空间及 H¹(T)
% * 二阶障碍和 H²(T)
% * 障碍空间不等于每个类都实际出现
% 
% 纲要内容结束

\section{形变方程和上同调}
\label{g4:sec:2902}

待撰写

\subsection{Beltrami 型丛值(0,1)形式}
\label{g4:subsec:2902:01}
待撰写

\subsection{Maurer–Cartan 可积方程}
\label{g4:subsec:2902:02}
待撰写

\subsection{\texorpdfstring{一阶空间及 \(H^1(\mathrm{T})\)}{一阶空间及 H¹(T)}}
\label{g4:subsec:2902:03}
待撰写

\subsection{\texorpdfstring{二阶障碍和 \(H^2(\mathrm{T})\)}{二阶障碍和 H²(T)}}
\label{g4:subsec:2902:04}
待撰写

\subsection{障碍空间不等于每个类都实际出现}
\label{g4:subsec:2902:05}
待撰写
% 纲要条目排为编号小节标题，并进入目录与书签；说明段仍保留于纲要注释中。

% !TeX root = ../../Book4.tex
% 纲要标题：### 29.3 Kuranishi 有限维约化
% 纲要位置：几何学计划纲要.md，第 4444 行。
% 正文状态：未撰写；本文件目前是框架。

% 纲要内容开始
% 
% * 规范固定及 Hodge 分解
% * Green 算子和非线性不动点
% * 障碍映射的构造
% * 小解的完备性和半通用性
% * 自同构与通用族的区别
% 
% 纲要内容结束

\section{Kuranishi 有限维约化}
\label{g4:sec:2903}

待撰写

\subsection{规范固定及 Hodge 分解}
\label{g4:subsec:2903:01}
待撰写

\subsection{Green 算子和非线性不动点}
\label{g4:subsec:2903:02}
待撰写

\subsection{障碍映射的构造}
\label{g4:subsec:2903:03}
待撰写

\subsection{小解的完备性和半通用性}
\label{g4:subsec:2903:04}
待撰写

\subsection{自同构与通用族的区别}
\label{g4:subsec:2903:05}
待撰写
% 纲要条目排为编号小节标题，并进入目录与书签；说明段仍保留于纲要注释中。

% !TeX root = ../../Book4.tex
% 纲要标题：### 29.4 实际模型与跳跃
% 纲要位置：几何学计划纲要.md，第 4452 行。
% 正文状态：未撰写；本文件目前是框架。

% 纲要内容开始
% 
% * 复环面和曲线的无障碍情形
% * 丛和嵌入的形变
% * 上同调维数跳跃
% * 全局模空间所需的新工具
% 
% ---
% 
% 纲要内容结束

\section{实际模型与跳跃}
\label{g4:sec:2904}

待撰写

\subsection{复环面和曲线的无障碍情形}
\label{g4:subsec:2904:01}
待撰写

\subsection{丛和嵌入的形变}
\label{g4:subsec:2904:02}
待撰写

\subsection{上同调维数跳跃}
\label{g4:subsec:2904:03}
待撰写

\subsection{全局模空间所需的新工具}
\label{g4:subsec:2904:04}
待撰写
% 纲要条目排为编号小节标题，并进入目录与书签；说明段仍保留于纲要注释中。


\begin{chaptersummary}
待撰写
\end{chaptersummary}

\BookExercises
待撰写

% !TeX root = ../../Book4.tex
% 纲要标题：## 第30章*　复曲面、吹起与双有理结构
% 纲要位置：几何学计划纲要.md，第 4461 行。
% 正文状态：未撰写；本文件目前是框架。

% 纲要内容开始
% 
% > 本章目标与证明范围：以可计算曲面不变量和修改为主体，分类作为具有明确输入的结果。把复二维曲面与黎曼曲面、代数曲面的极小模型清楚区分。
% 
% 纲要内容结束

\OptionalChapter{复曲面、吹起与双有理结构}{g4:ch:30}
\BookChapterNumber{g4:ch:30}

\begin{chapterabstract}
待撰写
\end{chapterabstract}

\begin{learninggoals}
待撰写
\end{learninggoals}

待撰写

% !TeX root = ../../Book4.tex
% 纲要标题：### 30.1 曲面不变量
% 纲要位置：几何学计划纲要.md，第 4465 行。
% 正文状态：未撰写；本文件目前是框架。

% 纲要内容开始
% 
% * Hodge 数、Euler 数及 signature
% * 典范类和交形式
% * Kodaira 维数及多典范映射
% 
% 纲要内容结束

\section{曲面不变量}
\label{g4:sec:3001}

待撰写

\subsection{Hodge 数、Euler 数及 signature}
\label{g4:subsec:3001:01}
待撰写

\subsection{典范类和交形式}
\label{g4:subsec:3001:02}
待撰写

\subsection{Kodaira 维数及多典范映射}
\label{g4:subsec:3001:03}
待撰写
% 纲要条目排为编号小节标题，并进入目录与书签；说明段仍保留于纲要注释中。

% !TeX root = ../../Book4.tex
% 纲要标题：### 30.2 吹起和例外曲线
% 纲要位置：几何学计划纲要.md，第 4743 行。
% 正文状态：未撰写；本文件目前是框架。

% 纲要内容开始
% 
% * 复解析坐标中的吹起
% * 例外曲线和自交
% * 典范类及上同调的变化
% * 收缩的条件和所选模型
% 
% 纲要内容结束

\section{吹起和例外曲线}
\label{g4:sec:3002}

待撰写

\subsection{复解析坐标中的吹起}
\label{g4:subsec:3002:01}
待撰写

\subsection{例外曲线和自交}
\label{g4:subsec:3002:02}
待撰写

\subsection{典范类及上同调的变化}
\label{g4:subsec:3002:03}
待撰写

\subsection{收缩的条件和所选模型}
\label{g4:subsec:3002:04}
待撰写
% 纲要条目排为编号小节标题，并进入目录与书签；说明段仍保留于纲要注释中。

% !TeX root = ../../Book4.tex
% 纲要标题：### 30.3 典型曲面的几何
% 纲要位置：几何学计划纲要.md，第 4478 行。
% 正文状态：未撰写；本文件目前是框架。

% 纲要内容开始
% 
% * 有理与直纹曲面
% * K3、复环面和椭圆曲面
% * 非 Kähler 曲面
% * 分类定理各部分的准确外部输入
% 
% 纲要内容结束

\section{典型曲面的几何}
\label{g4:sec:3003}

待撰写

\subsection{有理与直纹曲面}
\label{g4:subsec:3003:01}
待撰写

\subsection{K3、复环面和椭圆曲面}
\label{g4:subsec:3003:02}
待撰写

\subsection{非 Kähler 曲面}
\label{g4:subsec:3003:03}
待撰写

\subsection{分类定理各部分的准确外部输入}
\label{g4:subsec:3003:04}
待撰写
% 纲要条目排为编号小节标题，并进入目录与书签；说明段仍保留于纲要注释中。

% !TeX root = ../../Book4.tex
% 纲要标题：### 30.4 解析与代数的比较
% 纲要位置：几何学计划纲要.md，第 4757 行。
% 正文状态：未撰写；本文件目前是框架。

% 纲要内容开始
% 
% * 射影性及极化
% * 同一拓扑类型的不同复结构
% * 与第五卷代数曲面分类的接口
% 
% ---
% 
% 纲要内容结束

\section{解析与代数的比较}
\label{g4:sec:3004}

待撰写

\subsection{射影性及极化}
\label{g4:subsec:3004:01}
待撰写

\subsection{同一拓扑类型的不同复结构}
\label{g4:subsec:3004:02}
待撰写

\subsection{与第五卷代数曲面分类的接口}
\label{g4:subsec:3004:03}
待撰写
% 纲要条目排为编号小节标题，并进入目录与书签；说明段仍保留于纲要注释中。


\begin{chaptersummary}
待撰写
\end{chaptersummary}

\BookExercises
待撰写

% !TeX root = ../../Book4.tex
% 纲要标题：## 第31章*　复 Monge–Ampère、Calabi–Yau 与特殊度量
% 纲要位置：几何学计划纲要.md，第 4765 行。
% 正文状态：未撰写；本文件目前是框架。

% 纲要内容开始
% 
% > 本章目标与证明范围：比“特殊度量简介”深入到方程、线性化和先验估计的实际结构。Calabi–Yau 存在定理作为明确深层输入；正文完成唯一性、连续方法框架及所选估计，不宣称已在一章内给出全部高阶分析。
% 
% 纲要内容结束

\OptionalChapter{复 Monge–Ampère、Calabi–Yau 与特殊度量}{g4:ch:31}
\BookChapterNumber{g4:ch:31}

\begin{chapterabstract}
待撰写
\end{chapterabstract}

\begin{learninggoals}
待撰写
\end{learninggoals}

待撰写

% !TeX root = ../../Book4.tex
% 纲要标题：### 31.1 Ricci 形式与典范类
% 纲要位置：几何学计划纲要.md，第 4769 行。
% 正文状态：未撰写；本文件目前是框架。

% 纲要内容开始
% 
% * Ricci 形式的局部表达
% * 第一陈类及体积形式
% * Kähler 类中的势函数
% * 几何问题到复 Monge–Ampère 方程
% 
% 纲要内容结束

\section{Ricci 形式与典范类}
\label{g4:sec:3101}

待撰写

\subsection{Ricci 形式的局部表达}
\label{g4:subsec:3101:01}
待撰写

\subsection{第一陈类及体积形式}
\label{g4:subsec:3101:02}
待撰写

\subsection{Kähler 类中的势函数}
\label{g4:subsec:3101:03}
待撰写

\subsection{几何问题到复 Monge–Ampère 方程}
\label{g4:subsec:3101:04}
待撰写
% 纲要条目排为编号小节标题，并进入目录与书签；说明段仍保留于纲要注释中。

% !TeX root = ../../Book4.tex
% 纲要标题：### 31.2 线性化与唯一性
% 纲要位置：几何学计划纲要.md，第 4776 行。
% 正文状态：未撰写；本文件目前是框架。

% 纲要内容开始
% 
% * 方程线性化的 Laplace 算子
% * 积分约束和规范化
% * 最大值或积分法的唯一性
% * 连续方法的开放性
% 
% 纲要内容结束

\section{线性化与唯一性}
\label{g4:sec:3102}

待撰写

\subsection{方程线性化的 Laplace 算子}
\label{g4:subsec:3102:01}
待撰写

\subsection{积分约束和规范化}
\label{g4:subsec:3102:02}
待撰写

\subsection{最大值或积分法的唯一性}
\label{g4:subsec:3102:03}
待撰写

\subsection{连续方法的开放性}
\label{g4:subsec:3102:04}
待撰写
% 纲要条目排为编号小节标题，并进入目录与书签；说明段仍保留于纲要注释中。

% !TeX root = ../../Book4.tex
% 纲要标题：### 31.3 闭性所需的估计链
% 纲要位置：几何学计划纲要.md，第 4511 行。
% 正文状态：未撰写；本文件目前是框架。

% 纲要内容开始
% 
% * C⁰ 控制及其额外输入
% * 二阶迹估计的主要微分不等式
% * 高阶正则性与自举
% * 非塌缩和退化边界
% * 各估计的证明或外引位置逐项列明
% 
% 纲要内容结束

\section{闭性所需的估计链}
\label{g4:sec:3103}

待撰写

\subsection{\texorpdfstring{\(C^0\) 控制及其额外输入}{C⁰ 控制及其额外输入}}
\label{g4:subsec:3103:01}
待撰写

\subsection{二阶迹估计的主要微分不等式}
\label{g4:subsec:3103:02}
待撰写

\subsection{高阶正则性与自举}
\label{g4:subsec:3103:03}
待撰写

\subsection{非塌缩和退化边界}
\label{g4:subsec:3103:04}
待撰写

\subsection{各估计的证明或外引位置逐项列明}
\label{g4:subsec:3103:05}
待撰写
% 纲要条目排为编号小节标题，并进入目录与书签；说明段仍保留于纲要注释中。

% !TeX root = ../../Book4.tex
% 纲要标题：### 31.4 存在定理的几何后果
% 纲要位置：几何学计划纲要.md，第 4791 行。
% 正文状态：未撰写；本文件目前是框架。

% 纲要内容开始
% 
% * Calabi–Yau 定理的精确紧 Kähler 版本
% * Ricci 平坦与完整群约束
% * 超 Kähler 和 K3 模型
% * Calabi–Yau 一词的不同定义约定
% 
% 纲要内容结束

\section{存在定理的几何后果}
\label{g4:sec:3104}

待撰写

\subsection{Calabi–Yau 定理的精确紧 Kähler 版本}
\label{g4:subsec:3104:01}
待撰写

\subsection{Ricci 平坦与完整群约束}
\label{g4:subsec:3104:02}
待撰写

\subsection{超 Kähler 和 K3 模型}
\label{g4:subsec:3104:03}
待撰写

\subsection{Calabi–Yau 一词的不同定义约定}
\label{g4:subsec:3104:04}
待撰写
% 纲要条目排为编号小节标题，并进入目录与书签；说明段仍保留于纲要注释中。

% !TeX root = ../../Book4.tex
% 纲要标题：### 31.5 其他典范度量问题
% 纲要位置：几何学计划纲要.md，第 4798 行。
% 正文状态：未撰写；本文件目前是框架。

% 纲要内容开始
% 
% * Kähler–Einstein 方程的符号差异
% * 自同构和 Futaki 型障碍
% * 稳定性与存在性定理的准确外部接口
% 
% ---
% 
% 纲要内容结束

\section{其他典范度量问题}
\label{g4:sec:3105}

待撰写

\subsection{Kähler–Einstein 方程的符号差异}
\label{g4:subsec:3105:01}
待撰写

\subsection{自同构和 Futaki 型障碍}
\label{g4:subsec:3105:02}
待撰写

\subsection{稳定性与存在性定理的准确外部接口}
\label{g4:subsec:3105:03}
待撰写
% 纲要条目排为编号小节标题，并进入目录与书签；说明段仍保留于纲要注释中。

% !TeX root = ../../Book4.tex
% 纲要标题：### 31.6*　Kähler–Ricci 流与势方程
% 纲要位置：几何学计划纲要.md，第 4806 行。
% 正文状态：未撰写；本文件目前是框架。
% 写作范围、先修与边界：先修为 III.27 的 Ricci 流、III.13.5—13.6 的非线性共同工具，以及 IV.31.1—31.3 的 Ricci 形式、势和复 Monge–Ampère 方程。本节固定紧 Kähler 流形，核对 Ricci 形式和第一陈类的归一化，推导未归一化流的上同调类演化与抛物势方程，并计算 Kähler–Einstein 初始度量的缩放解；正性给出的时间约束与最大光滑存在时间定理须分开。最大存在时间、特定长期收敛和奇点结论只准确外引，不追加完整抛物估计链、一般奇点分类或与 MMP 等价的承诺；既有 IV.31.5 保持不变。

% 纲要内容开始
% 
% > 写作范围与外部输入：先修为 III.27 的 Ricci 流、III.13.5—13.6 的非线性共同工具，以及 IV.31.1—31.3 的 Ricci 形式、势和复 Monge–Ampère 方程。本节固定紧 Kähler 流形，核对 Ricci 形式和第一陈类的归一化，推导未归一化流的上同调类演化与抛物势方程，并计算 Kähler–Einstein 初始度量的缩放解；正性给出的时间约束与最大光滑存在时间定理须分开。最大存在时间、特定长期收敛和奇点结论只准确外引，不追加完整抛物估计链、一般奇点分类或与 MMP 等价的承诺；既有 IV.31.5 保持不变。
% 
% * Kähler–Ricci 流的形式方程与符号
% * Kähler 类的演化与正性时间范围
% * 随时间变化的背景形式与势方程
% * Kähler–Einstein 度量的缩放模型
% * 长期行为与奇点结论的适用范围
% 
% ---
% 
% 纲要内容结束

\OptionalSection{Kähler–Ricci 流与势方程}{g4:sec:3106}

待撰写

\subsection{Kähler–Ricci 流的形式方程与符号}
\label{g4:subsec:3106:01}
待撰写

\subsection{Kähler 类的演化与正性时间范围}
\label{g4:subsec:3106:02}
待撰写

\subsection{随时间变化的背景形式与势方程}
\label{g4:subsec:3106:03}
待撰写

\subsection{Kähler–Einstein 度量的缩放模型}
\label{g4:subsec:3106:04}
待撰写

\subsection{长期行为与奇点结论的适用范围}
\label{g4:subsec:3106:05}
待撰写

% 纲要条目排为编号小节标题，并进入目录与书签；说明段仍保留于纲要注释中。


\begin{chaptersummary}
待撰写
\end{chaptersummary}

\BookExercises
待撰写

% !TeX root = ../../Book4.tex
% 纲要标题：## 第32章*　Hodge 变分、Gauss–Manin 与退化
% 纲要位置：几何学计划纲要.md，第 4534 行。
% 正文状态：未撰写；本文件目前是框架。

% 纲要内容开始
% 
% > 本章目标与证明范围：由上同调局部系统构造平坦联络，再研究其与全纯 Hodge 滤过的张力。周期映射、横截性和退化各有实际计算，不止列出三个名称。
% 
% 纲要内容结束

\OptionalChapter{Hodge 变分、Gauss–Manin 与退化}{g4:ch:32}
\BookChapterNumber{g4:ch:32}

\begin{chapterabstract}
待撰写
\end{chapterabstract}

\begin{learninggoals}
待撰写
\end{learninggoals}

待撰写

% !TeX root = ../../Book4.tex
% 纲要标题：### 32.1 整 Hodge 结构与极化
% 纲要位置：几何学计划纲要.md，第 4822 行。
% 正文状态：未撰写；本文件目前是框架。

% 纲要内容开始
% 
% * 整格、权和型分解
% * Hodge 滤过的等价表述
% * 极化及双线性关系
% * 纯结构与混合结构的区别
% 
% 纲要内容结束

\section{整 Hodge 结构与极化}
\label{g4:sec:3201}

待撰写

\subsection{整格、权和型分解}
\label{g4:subsec:3201:01}
待撰写

\subsection{Hodge 滤过的等价表述}
\label{g4:subsec:3201:02}
待撰写

\subsection{极化及双线性关系}
\label{g4:subsec:3201:03}
待撰写

\subsection{纯结构与混合结构的区别}
\label{g4:subsec:3201:04}
待撰写
% 纲要条目排为编号小节标题，并进入目录与书签；说明段仍保留于纲要注释中。

% !TeX root = ../../Book4.tex
% 纲要标题：### 32.2 Gauss–Manin 联络
% 纲要位置：几何学计划纲要.md，第 4545 行。
% 正文状态：未撰写；本文件目前是框架。

% 纲要内容开始
% 
% * 光滑固有族的局部系统
% * 相对 de Rham 复形
% * 由微分导出联络
% * 平坦性和单值
% * 基变换与局部平凡化选择
% 
% 纲要内容结束

\section{Gauss–Manin 联络}
\label{g4:sec:3202}

待撰写

\subsection{光滑固有族的局部系统}
\label{g4:subsec:3202:01}
待撰写

\subsection{相对 de Rham 复形}
\label{g4:subsec:3202:02}
待撰写

\subsection{由微分导出联络}
\label{g4:subsec:3202:03}
待撰写

\subsection{平坦性和单值}
\label{g4:subsec:3202:04}
待撰写

\subsection{基变换与局部平凡化选择}
\label{g4:subsec:3202:05}
待撰写
% 纲要条目排为编号小节标题，并进入目录与书签；说明段仍保留于纲要注释中。

% !TeX root = ../../Book4.tex
% 纲要标题：### 32.3 周期映射与横截性
% 纲要位置：几何学计划纲要.md，第 4837 行。
% 正文状态：未撰写；本文件目前是框架。

% 纲要内容开始
% 
% * 周期域及滤过参数
% * Griffiths 横截性的推导
% * Kodaira–Spencer 与周期微分
% * 局部 Torelli 的所选模型
% 
% 纲要内容结束

\section{周期映射与横截性}
\label{g4:sec:3203}

待撰写

\subsection{周期域及滤过参数}
\label{g4:subsec:3203:01}
待撰写

\subsection{Griffiths 横截性的推导}
\label{g4:subsec:3203:02}
待撰写

\subsection{Kodaira–Spencer 与周期微分}
\label{g4:subsec:3203:03}
待撰写

\subsection{局部 Torelli 的所选模型}
\label{g4:subsec:3203:04}
待撰写
% 纲要条目排为编号小节标题，并进入目录与书签；说明段仍保留于纲要注释中。

% !TeX root = ../../Book4.tex
% 纲要标题：### 32.4 退化与极限结构
% 纲要位置：几何学计划纲要.md，第 4560 行。
% 正文状态：未撰写；本文件目前是框架。

% 纲要内容开始
% 
% * 穿孔圆盘族和单值
% * 准幺幂单值的定理范围
% * 对数、幂零算子及权滤过
% * 极限混合 Hodge 结构的额外输入
% * 结点椭圆曲线的显式计算
% 
% ---
% 
% 纲要内容结束

\section{退化与极限结构}
\label{g4:sec:3204}

待撰写

\subsection{穿孔圆盘族和单值}
\label{g4:subsec:3204:01}
待撰写

\subsection{准幺幂单值的定理范围}
\label{g4:subsec:3204:02}
待撰写

\subsection{对数、幂零算子及权滤过}
\label{g4:subsec:3204:03}
待撰写

\subsection{极限混合 Hodge 结构的额外输入}
\label{g4:subsec:3204:04}
待撰写

\subsection{结点椭圆曲线的显式计算}
\label{g4:subsec:3204:05}
待撰写
% 纲要条目排为编号小节标题，并进入目录与书签；说明段仍保留于纲要注释中。


\begin{chaptersummary}
待撰写
\end{chaptersummary}

\BookExercises
待撰写

% !TeX root = ../../Book4.tex
% 纲要标题：## 第33章*　镜像对称的可计算模型与研究接口
% 纲要位置：几何学计划纲要.md，第 4854 行。
% 正文状态：未撰写；本文件目前是框架。

% 纲要内容开始
% 
% > 本章目标与证明范围：以椭圆曲线和二维环面形成一个具体模型专题，再区分枚举、周期和范畴三个不同镜像命题。一般镜像不是统一适用于全部对象的已证定理。
% 
% 纲要内容结束

\OptionalChapter{镜像对称的可计算模型与研究接口}{g4:ch:33}
\BookChapterNumber{g4:ch:33}

\begin{chapterabstract}
待撰写
\end{chapterabstract}

\begin{learninggoals}
待撰写
\end{learninggoals}

待撰写

% !TeX root = ../../Book4.tex
% 纲要标题：### 33.1 二维环面的两种参数
% 纲要位置：几何学计划纲要.md，第 4574 行。
% 正文状态：未撰写；本文件目前是框架。

% 纲要内容开始
% 
% * 复模参数和辛面积
% * B 场与复化 Kähler 参数
% * 格对偶和基本变换
% 
% 纲要内容结束

\section{二维环面的两种参数}
\label{g4:sec:3301}

待撰写

\subsection{复模参数和辛面积}
\label{g4:subsec:3301:01}
待撰写

\subsection{B 场与复化 Kähler 参数}
\label{g4:subsec:3301:02}
待撰写

\subsection{格对偶和基本变换}
\label{g4:subsec:3301:03}
待撰写
% 纲要条目排为编号小节标题，并进入目录与书签；说明段仍保留于纲要注释中。

% !TeX root = ../../Book4.tex
% 纲要标题：### 33.2 直线 Lagrange 与全纯线丛
% 纲要位置：几何学计划纲要.md，第 4580 行。
% 正文状态：未撰写；本文件目前是框架。

% 纲要内容开始
% 
% * 斜率、交点和 Floer 的有限模型
% * 线丛次数及截面
% * theta 函数和三角形计数的代表计算
% * 选择、收敛及系数条件
% 
% 纲要内容结束

\section{直线 Lagrange 与全纯线丛}
\label{g4:sec:3302}

待撰写

\subsection{斜率、交点和 Floer 的有限模型}
\label{g4:subsec:3302:01}
待撰写

\subsection{线丛次数及截面}
\label{g4:subsec:3302:02}
待撰写

\subsection{theta 函数和三角形计数的代表计算}
\label{g4:subsec:3302:03}
待撰写

\subsection{选择、收敛及系数条件}
\label{g4:subsec:3302:04}
待撰写
% 纲要条目排为编号小节标题，并进入目录与书签；说明段仍保留于纲要注释中。

% !TeX root = ../../Book4.tex
% 纲要标题：### 33.3 从模型到不同镜像问题
% 纲要位置：几何学计划纲要.md，第 4871 行。
% 正文状态：未撰写；本文件目前是框架。

% 纲要内容开始
% 
% * 周期方程与特殊函数
% * Gromov–Witten 势与枚举
% * 导出范畴和 Fukaya 范畴的额外构造
% * 定理、猜想和例证逐项标明
% 
% ---
% 
% 纲要内容结束

\section{从模型到不同镜像问题}
\label{g4:sec:3303}

待撰写

\subsection{周期方程与特殊函数}
\label{g4:subsec:3303:01}
待撰写

\subsection{Gromov–Witten 势与枚举}
\label{g4:subsec:3303:02}
待撰写

\subsection{导出范畴和 Fukaya 范畴的额外构造}
\label{g4:subsec:3303:03}
待撰写

\subsection{定理、猜想和例证逐项标明}
\label{g4:subsec:3303:04}
待撰写
% 纲要条目排为编号小节标题，并进入目录与书签；说明段仍保留于纲要注释中。


\begin{chaptersummary}
待撰写
\end{chaptersummary}

\BookExercises
待撰写


% !TeX root = ../Book4.tex
% 纲要标题：# 第四卷附录
% 纲要位置：几何学计划纲要.md，第 4880 行。
% 正文状态：未撰写；本文件目前是框架。

% 纲要内容开始
% 
% 纲要内容结束

\appendix
\renewcommand{\thechapter}{\Alph{chapter}}
\ctexset{chapter={number=\Alph{chapter}}}
\BookAppendixPart{g4:part:appendices}

% !TeX root = ../../Book4.tex
% 纲要标题：## 附录 A　Stein 空间与解析层
% 纲要位置：几何学计划纲要.md，第 4598 行。
% 正文状态：未撰写；本文件目前是框架。

% 纲要内容开始
% 
% > 局部多复变之后的独立整体路线。正文需要的凝聚性不留作附录黑箱。
% 
% * **Stein 条件与模型**
%   * 全纯凸性和可分性
%   * 严格 psh 穷竭
%   * 仿射簇、域及非 Stein 模型
% 
% * **Cartan 方法**
%   * 定理 A、B 的准确对象
%   * Cousin 问题与上同调消失
%   * Oka 原理的所选陈述
% 
% * **∂̄ 与整体嵌入**
%   * 带权估计和穷竭
%   * 函数分离及嵌入问题
%   * 奇异 Stein 空间的额外工具
% 
% ---
% 
% 纲要内容结束

\chapter{Stein 空间与解析层}
\label{g4:app:A}
\BookChapterNumber{g4:app:A}

待撰写

\section{Stein 条件与模型}
\label{g4:sec:A01}
待撰写

\subsection{全纯凸性和可分性}
\label{g4:subsec:A01:01}
待撰写

\subsection{严格 psh 穷竭}
\label{g4:subsec:A01:02}
待撰写

\subsection{仿射簇、域及非 Stein 模型}
\label{g4:subsec:A01:03}
待撰写

\section{Cartan 方法}
\label{g4:sec:A02}
待撰写

\subsection{定理 A、B 的准确对象}
\label{g4:subsec:A02:01}
待撰写

\subsection{Cousin 问题与上同调消失}
\label{g4:subsec:A02:02}
待撰写

\subsection{Oka 原理的所选陈述}
\label{g4:subsec:A02:03}
待撰写

\section{\texorpdfstring{\(\bar{\partial}\) 与整体嵌入}{\={∂} 与整体嵌入}}
\label{g4:sec:A03}
待撰写

\subsection{带权估计和穷竭}
\label{g4:subsec:A03:01}
待撰写

\subsection{函数分离及嵌入问题}
\label{g4:subsec:A03:02}
待撰写

\subsection{奇异 Stein 空间的额外工具}
\label{g4:subsec:A03:03}
待撰写

% !TeX root = ../../Book4.tex
% 纲要标题：## 附录 B　复奇点、Milnor 纤维与混合结构
% 纲要位置：几何学计划纲要.md，第 4619 行。
% 正文状态：未撰写；本文件目前是框架。

% 纲要内容开始
% 
% > 从孤立超曲面奇点进入拓扑和单值，选定 Milnor 纤维化及结点、尖点算例深入说明。
% 
% * **Milnor 纤维化**
%   * 小球、小球面及横截性
%   * 穿孔邻域到圆周的映射
%   * 纤维、链接与选择无关性
%   * 孤立奇点条件的作用
% 
% * **消失圈和单值**
%   * Morse 型平滑化
%   * 交形式及 Picard–Lefschetz 变换
%   * 结点与尖点的计算
%   * Milnor 数及 Jacobian 代数
% 
% * **混合结构**
%   * 局部与全局退化的关系
%   * 权与 Hodge 滤过
%   * 一般混合 Hodge 定理的引用范围
% 
% ---
% 
% 纲要内容结束

\chapter{复奇点、Milnor 纤维与混合结构}
\label{g4:app:B}
\BookChapterNumber{g4:app:B}

待撰写

\section{Milnor 纤维化}
\label{g4:sec:B01}
待撰写

\subsection{小球、小球面及横截性}
\label{g4:subsec:B01:01}
待撰写

\subsection{穿孔邻域到圆周的映射}
\label{g4:subsec:B01:02}
待撰写

\subsection{纤维、链接与选择无关性}
\label{g4:subsec:B01:03}
待撰写

\subsection{孤立奇点条件的作用}
\label{g4:subsec:B01:04}
待撰写

\section{消失圈和单值}
\label{g4:sec:B02}
待撰写

\subsection{Morse 型平滑化}
\label{g4:subsec:B02:01}
待撰写

\subsection{交形式及 Picard–Lefschetz 变换}
\label{g4:subsec:B02:02}
待撰写

\subsection{结点与尖点的计算}
\label{g4:subsec:B02:03}
待撰写

\subsection{Milnor 数及 Jacobian 代数}
\label{g4:subsec:B02:04}
待撰写

\section{混合结构}
\label{g4:sec:B03}
待撰写

\subsection{局部与全局退化的关系}
\label{g4:subsec:B03:01}
待撰写

\subsection{权与 Hodge 滤过}
\label{g4:subsec:B03:02}
待撰写

\subsection{一般混合 Hodge 定理的引用范围}
\label{g4:subsec:B03:03}
待撰写

% !TeX root = ../../Book4.tex
% 纲要标题：## 附录 C　Poisson 几何与量子化
% 纲要位置：几何学计划纲要.md，第 4926 行。
% 正文状态：未撰写；本文件目前是框架。

% 纲要内容开始
% 
% > 与正文退化 Poisson 模块衔接，给出整性和一个低阶形变的完整模型。
% 
% * **辛叶与横向结构**
%   * Lie–Poisson 叶
%   * Casimir 函数
%   * 秩跳跃及局部分裂定理的范围
% 
% * **预量子线丛**
%   * 曲率和第一 Chern 类
%   * 整性充要条件的证明
%   * 环面、球面和余伴随轨道
% 
% * **两种量子化**
%   * 极化与截面空间
%   * 星积、结合性和一阶 Poisson 括号
%   * Moyal 乘积的核验
%   * 一般存在定理的额外输入
% 
% ---
% 
% 纲要内容结束

\chapter{Poisson 几何与量子化}
\label{g4:app:C}
\BookChapterNumber{g4:app:C}

待撰写

\section{辛叶与横向结构}
\label{g4:sec:C01}
待撰写

\subsection{Lie–Poisson 叶}
\label{g4:subsec:C01:01}
待撰写

\subsection{Casimir 函数}
\label{g4:subsec:C01:02}
待撰写

\subsection{秩跳跃及局部分裂定理的范围}
\label{g4:subsec:C01:03}
待撰写

\section{预量子线丛}
\label{g4:sec:C02}
待撰写

\subsection{曲率和第一 Chern 类}
\label{g4:subsec:C02:01}
待撰写

\subsection{整性充要条件的证明}
\label{g4:subsec:C02:02}
待撰写

\subsection{环面、球面和余伴随轨道}
\label{g4:subsec:C02:03}
待撰写

\section{两种量子化}
\label{g4:sec:C03}
待撰写

\subsection{极化与截面空间}
\label{g4:subsec:C03:01}
待撰写

\subsection{星积、结合性和一阶 Poisson 括号}
\label{g4:subsec:C03:02}
待撰写

\subsection{Moyal 乘积的核验}
\label{g4:subsec:C03:03}
待撰写

\subsection{一般存在定理的额外输入}
\label{g4:subsec:C03:04}
待撰写

% !TeX root = ../../Book4.tex
% 纲要标题：## 附录 D　接触拓扑与填充障碍
% 纲要位置：几何学计划纲要.md，第 4664 行。
% 正文状态：未撰写；本文件目前是框架。

% 纲要内容开始
% 
% > 保留三维及高维两条支线，各自说明维数、紧性和可用不变量。
% 
% * **Legendrian 的整体不变量**
%   * 前图及尖点
%   * Thurston–Bennequin 与旋转数
%   * 稳定化和同位的区别
% 
% * **填充和配边**
%   * 强弱精确填充
%   * 典型不可填充模型及引用定理
%   * 伪全纯曲线产生的障碍
% 
% * **高维柔性与刚性**
%   * 形式解与实际解
%   * loose Legendrian 的准确对象
%   * h-原理的条件和适用界线
% 
% ---
% 
% 纲要内容结束

\chapter{接触拓扑与填充障碍}
\label{g4:app:D}
\BookChapterNumber{g4:app:D}

待撰写

\section{Legendrian 的整体不变量}
\label{g4:sec:D01}
待撰写

\subsection{前图及尖点}
\label{g4:subsec:D01:01}
待撰写

\subsection{Thurston–Bennequin 与旋转数}
\label{g4:subsec:D01:02}
待撰写

\subsection{稳定化和同位的区别}
\label{g4:subsec:D01:03}
待撰写

\section{填充和配边}
\label{g4:sec:D02}
待撰写

\subsection{强弱精确填充}
\label{g4:subsec:D02:01}
待撰写

\subsection{典型不可填充模型及引用定理}
\label{g4:subsec:D02:02}
待撰写

\subsection{伪全纯曲线产生的障碍}
\label{g4:subsec:D02:03}
待撰写

\section{高维柔性与刚性}
\label{g4:sec:D03}
待撰写

\subsection{形式解与实际解}
\label{g4:subsec:D03:01}
待撰写

\subsection{loose Legendrian 的准确对象}
\label{g4:subsec:D03:02}
待撰写

\subsection{h-原理的条件和适用界线}
\label{g4:subsec:D03:03}
待撰写

% !TeX root = ../../Book4.tex
% 纲要标题：## 附录 E　作用滤过与持续 Floer 结构
% 纲要位置：几何学计划纲要.md，第 4685 行。
% 正文状态：未撰写；本文件目前是框架。

% 纲要内容开始
% 
% > 建立有限型过滤复形的代数理论，再与正文受限 Hamilton 模块连接。
% 
% * **过滤链复形**
%   * 作用值、子复形与商复形
%   * 延拓造成的过滤移位
%   * 谱数的定义和存在
% 
% * **持续模与稳定性**
%   * 有限型区间分解
%   * 交错和瓶颈距离
%   * 条形码稳定性所需条件
%   * 无限维与 Novikov 情形的修正
% 
% * **Hamilton 距离**
%   * Hofer 长度及范数
%   * 谱不变量给出的估计
%   * 具体平移和支撑模型
% 
% ---
% 
% 纲要内容结束

\chapter{作用滤过与持续 Floer 结构}
\label{g4:app:E}
\BookChapterNumber{g4:app:E}

待撰写

\section{过滤链复形}
\label{g4:sec:E01}
待撰写

\subsection{作用值、子复形与商复形}
\label{g4:subsec:E01:01}
待撰写

\subsection{延拓造成的过滤移位}
\label{g4:subsec:E01:02}
待撰写

\subsection{谱数的定义和存在}
\label{g4:subsec:E01:03}
待撰写

\section{持续模与稳定性}
\label{g4:sec:E02}
待撰写

\subsection{有限型区间分解}
\label{g4:subsec:E02:01}
待撰写

\subsection{交错和瓶颈距离}
\label{g4:subsec:E02:02}
待撰写

\subsection{条形码稳定性所需条件}
\label{g4:subsec:E02:03}
待撰写

\subsection{无限维与 Novikov 情形的修正}
\label{g4:subsec:E02:04}
待撰写

\section{Hamilton 距离}
\label{g4:sec:E03}
待撰写

\subsection{Hofer 长度及范数}
\label{g4:subsec:E03:01}
待撰写

\subsection{谱不变量给出的估计}
\label{g4:subsec:E03:02}
待撰写

\subsection{具体平移和支撑模型}
\label{g4:subsec:E03:03}
待撰写

% !TeX root = ../../Book4.tex
% 纲要标题：## 附录 F　复几何、算术与表示
% 纲要位置：几何学计划纲要.md，第 4707 行。
% 正文状态：未撰写；本文件目前是框架。

% 纲要内容开始
% 
% > 保持与数论和表示论的联系，但不把相应整套理论追加为本卷必修。
% 
% * **模曲线与阿贝尔簇**
%   * 格、复乘及代数模型
%   * 极化和模参数
%   * 算术下降所需的额外数据
% 
% * **旗流形与全纯表示**
%   * 旗簇和齐性线丛
%   * 截面表示的低秩模型
%   * Borel–Weil 的精确陈述
% 
% * **周期与几何化**
%   * 周期积分的微分方程
%   * Gauss–Manin 表示
%   * 几何表示论需要的新范畴
% 
% ---
% 
% 纲要内容结束

\chapter{复几何、算术与表示}
\label{g4:app:F}
\BookChapterNumber{g4:app:F}

待撰写

\section{模曲线与阿贝尔簇}
\label{g4:sec:F01}
待撰写

\subsection{格、复乘及代数模型}
\label{g4:subsec:F01:01}
待撰写

\subsection{极化和模参数}
\label{g4:subsec:F01:02}
待撰写

\subsection{算术下降所需的额外数据}
\label{g4:subsec:F01:03}
待撰写

\section{旗流形与全纯表示}
\label{g4:sec:F02}
待撰写

\subsection{旗簇和齐性线丛}
\label{g4:subsec:F02:01}
待撰写

\subsection{截面表示的低秩模型}
\label{g4:subsec:F02:02}
待撰写

\subsection{Borel–Weil 的精确陈述}
\label{g4:subsec:F02:03}
待撰写

\section{周期与几何化}
\label{g4:sec:F03}
待撰写

\subsection{周期积分的微分方程}
\label{g4:subsec:F03:01}
待撰写

\subsection{Gauss–Manin 表示}
\label{g4:subsec:F03:02}
待撰写

\subsection{几何表示论需要的新范畴}
\label{g4:subsec:F03:03}
待撰写



\backmatter
\cleardoublepage
\pagenumbering{roman}
\bookmarksetup{startatroot}
% !TeX root = ../Book4.tex
\cleardoublepage
% !TeX root = ../../Book4.tex
\chapter*{后记}
\addcontentsline{toc}{chapter}{后记}
待撰写

\cleardoublepage
% !TeX root = ../../Book4.tex
\PrintBookBibliography
% 仅在核读并引用真实资料后补入书目。

\cleardoublepage
% !TeX root = ../../Book4.tex
\PrintBookIndex{chinese}{中文索引}

\cleardoublepage
% !TeX root = ../../Book4.tex
\PrintBookIndex{foreign}{外文索引}

\cleardoublepage
% !TeX root = ../../Book4.tex
\PrintBookIndex{symbols}{符号索引}


\end{document}
